\documentclass[11pt,a4paper,openright,twoside]{memoir}

% === Preamble for Introduction to Doxonomics ===

% --- Encoding and fonts ---
\usepackage[utf8]{inputenc}
\usepackage[T1]{fontenc}
\usepackage{lmodern}

% --- Mathematics ---
\usepackage{amsmath}
\usepackage{amssymb}
\usepackage{amsthm}
\usepackage{mathtools}
\usepackage{bm}              % bold math symbols
\usepackage{mathrsfs}         % \mathscr
\usepackage{stmaryrd}         % additional math symbols

% --- Graphics and diagrams ---
\usepackage{tikz}
\usetikzlibrary{
  arrows.meta,
  calc,
  positioning,
  shapes.geometric,
  shapes.misc,
  fit,
  backgrounds,
  decorations.markings,
  decorations.pathreplacing,
  matrix,
  chains,
  graphs,
  quotes,
}
\usepackage{pgfplots}
\pgfplotsset{compat=1.18}

% --- Tables ---
\usepackage{booktabs}
\usepackage{array}
\usepackage{longtable}

% --- Colors ---
\usepackage{xcolor}
\definecolor{doxoblue}{RGB}{31, 78, 121}
\definecolor{doxogray}{RGB}{88, 88, 88}
\definecolor{doxored}{RGB}{153, 42, 42}
\definecolor{doxogreen}{RGB}{42, 111, 63}
\definecolor{doxolime}{RGB}{60, 179, 113}

% --- Bibliography ---
\usepackage[
  backend=biber,
  style=authoryear,
  sorting=nyt,
  maxbibnames=99,
  maxcitenames=2,
  giveninits=true,
  uniquename=init,
  dashed=false,
]{biblatex}
\addbibresource{bibliography/references.bib}

% --- Cross-references and links ---
\usepackage{hyperref}
\hypersetup{
  colorlinks=true,
  linkcolor=doxoblue,
  citecolor=doxogreen,
  urlcolor=doxored,
}
\usepackage[capitalize,noabbrev]{cleveref}

% --- Indexing ---
\usepackage{makeidx}
\makeindex

% --- Epigraphs ---
\usepackage{epigraph}
\setlength{\epigraphwidth}{0.7\textwidth}

% --- Miscellaneous ---
\usepackage{enumitem}
\usepackage{microtype}

% === Theorem-like environments ===
\theoremstyle{plain}
\newtheorem{theorem}{Theorem}[chapter]
\newtheorem{proposition}[theorem]{Proposition}
\newtheorem{lemma}[theorem]{Lemma}
\newtheorem{corollary}[theorem]{Corollary}
\newtheorem{conjecture}[theorem]{Conjecture}

\theoremstyle{definition}
\newtheorem{definition}[theorem]{Definition}
\newtheorem{axiom}[theorem]{Axiom}
\newtheorem{principle}[theorem]{Principle}
\newtheorem{model}[theorem]{Model}

\theoremstyle{remark}
\newtheorem{remark}[theorem]{Remark}
\newtheorem{example}[theorem]{Example}
\newtheorem{exercise}{Exercise}[chapter]

% === Part interludes ===
\newenvironment{interlude}{%
  \cleardoublepage
  \thispagestyle{empty}
  \vspace*{3cm}
  \begin{center}\rule{4cm}{0.4pt}\end{center}
  \vspace{0.5cm}
  \begin{quote}\itshape
}{%
  \end{quote}
  \vspace{1cm}
  \begin{center}\rule{4cm}{0.4pt}\end{center}
  \clearpage
}

% === Doxonomics-specific notation ===

% Belief strength
\newcommand{\belief}[1]{B_{#1}}
\newcommand{\beliefvec}{\bm{B}}

% Funding flow
\newcommand{\funding}[1]{F_{#1}}
\newcommand{\fundingvec}{\bm{F}}

% Channel credibility
\newcommand{\credibility}[1]{C_{#1}}

% Population reach
\newcommand{\reach}[1]{P_{#1}}

% Temporal persistence
\newcommand{\persistence}[1]{T_{#1}}

% Relevance
\newcommand{\relevance}[2]{R_{#1,#2}}

% Countervailing force
\newcommand{\counter}[1]{X_{#1}}

% Epistemic capital
\newcommand{\ecap}{\mathcal{E}}

% Narrative infrastructure
\newcommand{\narr}{\mathcal{N}}

% Doxonomic system
\newcommand{\doxsys}{\mathfrak{D}}

% Belief production chain
\newcommand{\bpc}{\mathcal{B}}

% Institutional ecology
\newcommand{\insteco}{\mathcal{I}}

% Epistemic concentration (Herfindahl-like)
\newcommand{\ehhi}{\mathrm{EHI}}

% Ideological ROI
\newcommand{\iroi}{\mathrm{IROI}}

% Doxa operator (maps resources to beliefs)
\newcommand{\doxop}{\mathbf{D}}

% Sheaf of local beliefs
\newcommand{\sheaf}{\mathscr{F}}

% Category of institutions
\newcommand{\instcat}{\mathbf{Inst}}

% Spaces
\newcommand{\beliefspace}{\mathcal{B}}
\newcommand{\narspace}{\mathcal{N}}
\newcommand{\agentspace}{\mathcal{A}}

% Operators
\newcommand{\E}{\mathbb{E}}     % expectation
\newcommand{\R}{\mathbb{R}}     % reals
\newcommand{\N}{\mathbb{N}}     % naturals
\newcommand{\Prob}{\mathbb{P}}  % probability

% === Memoir adjustments ===

% Chapter style
\chapterstyle{veelo}

% Page headers
\pagestyle{headings}

% Part page style
\renewcommand{\partnamefont}{\normalfont\Large\scshape}
\renewcommand{\partnumfont}{\normalfont\Large\scshape}
\renewcommand{\parttitlefont}{\normalfont\Huge\bfseries}

% Section depths
\setsecnumdepth{subsection}
\settocdepth{subsection}


\begin{document}

% === Front Matter ===
\frontmatter

% Full-black cover page matching Piatra . Institute aesthetic
\thispagestyle{empty}
\begin{tikzpicture}[remember picture, overlay]
  % Black background filling entire page
  \fill[black] (current page.south west) rectangle (current page.north east);

  % Piatra . Institute name at top
  \node[white, font=\large\scshape] at ([yshift=-4cm]current page.north) {%
    Piatra {\color{doxolime}\,.}\, Institute%
  };

  % Main title
  \node[white, font=\fontsize{36}{42}\selectfont\bfseries, text width=12cm, align=center]
    at ([yshift=2cm]current page.center) {%
    Introduction to\\[6pt]Doxonomics%
  };

  % Subtitle
  \node[doxolime, font=\Large\itshape, text width=12cm, align=center]
    at ([yshift=-1.5cm]current page.center) {%
    Toward a Science of\\[2pt]How Power Shapes Belief%
  };

  % Decorative line
  \draw[doxolime, line width=0.8pt] ([yshift=-3.5cm, xshift=-4cm]current page.center)
    -- ([yshift=-3.5cm, xshift=4cm]current page.center);

  % Year at bottom
  \node[white, font=\large] at ([yshift=2.5cm]current page.south) {%
    \the\year%
  };
\end{tikzpicture}
\clearpage


\begin{titlingpage}
\calccentering{\unitlength}
\begin{adjustwidth*}{\unitlength}{-\unitlength}
\centering

\vspace*{2cm}

{\large\scshape Piatra \,.\, Institute}

\vspace{3cm}

{\Huge\bfseries Introduction to Doxonomics}

\vspace{0.8cm}

{\Large\itshape Toward a Science of How Power Shapes Belief}

\vspace{4cm}

\epigraph{The ideas of the ruling class are in every epoch the ruling ideas: i.e., the class which is the ruling material force of society is at the same time its ruling intellectual force.}{Karl Marx \& Friedrich Engels, \textit{The German Ideology} (1846)}

\vfill

{\large\the\year}

\end{adjustwidth*}
\end{titlingpage}


\cleardoublepage
\tableofcontents

\cleardoublepage
\chapter*{Preface}
\addcontentsline{toc}{chapter}{Preface}

\noindent
Every society produces goods, laws, and infrastructures.
Every society also produces beliefs---about what people are, what they deserve, what is possible, and what is natural.
The first kind of production is studied by economics, engineering, and political science.
The second kind has no dedicated science.

This book proposes one.

\textbf{Doxonomics}---from the Greek \textit{doxa} (opinion, common belief, public orthodoxy) and \textit{-nomics} (systematic study of allocation and structure)---is the science of how material power shapes public belief.
Its central question is not whether a belief is true, but what it cost to make it appear self-evident: who financed it, which institutions processed it, through what channels it circulated, and what social effects followed from its dominance.

The need for such a field is not difficult to motivate.
Global advertising exceeded one trillion dollars in 2024.
Lobbying, think-tank funding, public relations, and curriculum sponsorship add hundreds of billions more.
Existing disciplines study parts of this process: sociology of knowledge examines the social formation of ideas, political communication traces agenda-setting, media studies analyzes representation, propaganda studies maps persuasion techniques, and social epistemology investigates the conditions of collective knowledge.
What no existing field does is integrate these strands around \emph{explicit accounting of material support, institutional routing, and measurable belief effects}.
Political economy studies the distribution of material power but not its conversion into common sense.
Sociology of knowledge studies the social formation of ideas but rarely performs financial accounting on them.
Media studies analyzes representation but seldom asks which representations are continuously subsidized and to what measurable effect.
Doxonomics sits at the intersection of all these fields and adds what they individually lack: a unified framework for auditing the industrial production of belief.

This textbook is mathematically serious.
We take the position that critique without formal tools remains vulnerable to the same vagueness it diagnoses.
The field needs stock-and-flow models, causal inference, network analysis, information-theoretic measures, and---where the structure demands it---sheaf-theoretic coherence conditions.
Many of the models presented here are heuristic or conjectural rather than empirically calibrated; we mark them as such.
Their purpose is to clarify structure and generate testable predictions, not to simulate maturity the field has not yet earned.
Where a result is a proven theorem, we say so.
Where it is a stylized model, a conjecture, or a candidate formalization, we say that too.

Throughout the text, we reference interactive \emph{playgrounds}: computational environments where the reader can manipulate funding parameters, observe belief diffusion, simulate narrative cascades, and test the sensitivity of doxonomic models to their assumptions.
These are available at the Piatra Institute's companion site.

The book proceeds in five parts.
Part~I lays the conceptual and historical foundations.
Part~II develops the core formal apparatus---epistemic capital, narrative infrastructure, common-sense capture, doxonomic power, and sheaf-like coherence.
Part~III introduces the methods: measurement, accounting, network analysis, causal inference, computational text analysis, and experimental design.
Part~IV applies the machinery to nine domains, from the managed narrative of selfish human nature to crisis doxonomics.
Part~V turns to resistance, repair, and the ethics of studying belief production without becoming belief engineers.

No single author could write this book.
It emerges from a research program at the Piatra Institute, and it will remain open to revision as the field develops.
The reader is invited not only to learn doxonomics but to extend it.

\vspace{1em}
\noindent\textit{Piatra Institute, \the\year}


\cleardoublepage
\chapter*{Notation}
\addcontentsline{toc}{chapter}{Notation}

\noindent
The following symbols and conventions are used throughout the text.
Page numbers indicate the first substantive appearance of each symbol.

\section*{Belief Variables}

\begin{longtable}{>{\raggedright}p{3cm} p{4cm} p{6.5cm}}
\toprule
\textbf{Symbol} & \textbf{Name} & \textbf{Description} \\
\midrule
\endhead
$\belief{i}$ & Belief strength & Social strength of belief $i$ in the population \\
$\belief{i}(t)$ & Belief at time $t$ & Time-indexed belief strength \\
$\beliefvec$ & Belief vector & Vector of all belief strengths $(\belief{1}, \ldots, \belief{n})$ \\
$\beliefspace$ & Belief space & The space of all possible belief configurations \\
$\counter{i}$ & Countervailing force & Aggregate pressure opposing belief $i$ \\
\bottomrule
\end{longtable}

\section*{Funding and Channels}

\begin{longtable}{>{\raggedright}p{3cm} p{4cm} p{6.5cm}}
\toprule
\textbf{Symbol} & \textbf{Name} & \textbf{Description} \\
\midrule
\endhead
$\funding{k}$ & Funding flow & Resources flowing through channel $k$ \\
$\fundingvec$ & Funding vector & Vector of all funding flows $(\funding{1}, \ldots, \funding{m})$ \\
$\credibility{k}$ & Credibility & Credibility weight of channel $k$ \\
$\reach{k}$ & Population reach & Fraction of population exposed via channel $k$ \\
$\persistence{k}$ & Temporal persistence & Duration and repetition factor of channel $k$ \\
$\relevance{k}{i}$ & Relevance & Relevance of channel $k$ to belief $i$ \\
\bottomrule
\end{longtable}

\section*{Structural Objects}

\begin{longtable}{>{\raggedright}p{3cm} p{4cm} p{6.5cm}}
\toprule
\textbf{Symbol} & \textbf{Name} & \textbf{Description} \\
\midrule
\endhead
$\doxsys$ & Doxonomic system & A structured ensemble of actors, resources, institutions, and channels producing belief \\
$\bpc$ & Belief production chain & The full sequence from funding to behavioral consequence \\
$\narr$ & Narrative infrastructure & The durable channels through which a society reproduces its common sense \\
$\insteco$ & Institutional ecology & The ensemble of institutions stabilizing a belief regime \\
$\ecap$ & Epistemic capital & Resources convertible to credibility, attention, or authority \\
$\doxop$ & Doxa operator & The operator mapping resource allocations to belief distributions \\
\bottomrule
\end{longtable}

\section*{Spaces and Categories}

\begin{longtable}{>{\raggedright}p{3cm} p{4cm} p{6.5cm}}
\toprule
\textbf{Symbol} & \textbf{Name} & \textbf{Description} \\
\midrule
\endhead
$\narspace$ & Narrative space & The space of available narratives \\
$\agentspace$ & Agent space & The space of doxonomic agents \\
$\instcat$ & Category of institutions & Category whose objects are institutions and whose morphisms are influence relations \\
$\sheaf$ & Sheaf of local beliefs & Sheaf assigning local belief systems to institutional contexts \\
$\R$ & Real numbers & \\
$\N$ & Natural numbers & \\
\bottomrule
\end{longtable}

\section*{Operators and Indices}

\begin{longtable}{>{\raggedright}p{3cm} p{4cm} p{6.5cm}}
\toprule
\textbf{Symbol} & \textbf{Name} & \textbf{Description} \\
\midrule
\endhead
$\ehhi$ & Epistemic Herfindahl Index & Concentration score for belief-production market share \\
$\iroi$ & Ideological ROI & Return on ideological investment: belief change per unit funding \\
$\E[\cdot]$ & Expectation & Expected value operator \\
$\Prob(\cdot)$ & Probability & Probability measure \\
\bottomrule
\end{longtable}

\section*{Conventions}

\begin{itemize}[nosep]
  \item Subscript $i$ indexes beliefs; subscript $k$ indexes channels or institutions.
  \item Time dependence is written $f(t)$; discrete time steps use $f_t$.
  \item Bold lowercase ($\bm{x}$) denotes vectors; bold uppercase ($\bm{A}$) denotes matrices.
  \item Calligraphic letters ($\mathcal{B}, \mathcal{N}$) denote spaces or structured sets.
  \item Fraktur ($\mathfrak{D}$) denotes systems or large composite structures.
  \item Script ($\mathscr{F}$) denotes sheaves.
  \item Sans-serif ($\mathsf{Cat}$) denotes categories when disambiguation is needed.
\end{itemize}


% === Main Matter ===
\mainmatter

% --- Part I: Foundations ---
\part{Foundations}
\label{part:foundations}

\chapter{What Is Doxonomics?}
\label{ch:what-is-doxonomics}

\epigraph{The most potent weapon in the hands of the oppressor is the mind of the oppressed.}{Steve Biko}

\noindent
Human beings do not merely inhabit economies and legal systems.
They inhabit interpretive worlds: shared assumptions about what people are, what is possible, what is deserved, and what cannot be changed.
These assumptions are rarely spontaneous.
They are funded, repeated, institutionalized, and defended.

\textbf{Doxonomics}\index{doxonomics} is the study of how material power shapes public belief.
The term derives from the Greek \textit{doxa} (opinion, common belief, public orthodoxy) and denotes the systematic investigation of how beliefs become socially normal.
Its central question is not whether a belief is true, but:

\begin{quote}
\textit{What did it cost to make this belief appear self-evident?}
\end{quote}

This question is not rhetorical.
It is a research question with empirical content.
Global advertising expenditure surpassed one trillion dollars in 2024.
U.S.\ lobbying and political spending exceed tens of billions annually.
Philanthropic grants to policy-oriented think tanks run into hundreds of millions per year, channeled through networks of foundations whose 990 filings are public record.
University curricula shape the foundational assumptions of millions of students per decade.
Platform algorithms determine which narratives reach which audiences at what frequency.
Each of these is a measurable flow of resources into the production of belief, and none of them is systematically tracked by any existing discipline as a \emph{unified} phenomenon.

Doxonomics proposes to track them.


\section{The Object of Study}
\label{sec:ch1-object}

Modern societies devote enormous resources to the production of goods, laws, and infrastructures.
They also devote immense resources to the production of legitimacy, resignation, aspiration, and moral common sense.
A society does not only distribute wealth; it distributes interpretations of the world.
Some interpretations are amplified until they appear obvious.
Others are starved of attention until they appear unrealistic or unintelligible.

Doxonomics studies this process.
It asks how money becomes belief, how institutions become credibility, how repetition becomes reality, and how narratives become embedded in everyday life \autocite{bourdieu1991, herman1988}.

Consider a concrete example.
The proposition ``human nature is selfish'' is not merely a philosophical claim.
It is an anthropological assumption embedded in economics curricula (taught to over ten million students per decade in the United States alone), management theory (principal-agent models presuppose worker self-interest), policy design (incentive-based regulation assumes that agents will not cooperate without material reward), advertising (which addresses the consumer as an isolated desire-maximizer), and popular culture (which narrates life as competition).
Each of these institutional sites has a budget.
Each produces narrative output.
Each reaches an audience.
The combined effect is an interpretive environment in which ``selfishness is natural'' operates as background reality rather than as a contestable claim.

Is this belief \emph{true}?
That is an empirical question, and the experimental evidence is considerably more mixed than the institutional consensus suggests (Chapter~\ref{ch:human-nature}).
The doxonomic question is different: \emph{what institutional apparatus sustains this belief, at what cost, through which channels, and with what social effects?}

\begin{definition}[Doxonomic System]
\label{def:doxonomic-system}
\index{doxonomic system}
A \emph{doxonomic system} is a tuple
\[
  \doxsys = (\agentspace, \insteco, \narr, \beliefspace, \Phi, \Psi)
\]
where:
\begin{itemize}[nosep]
  \item $\agentspace$ is a set of agents (funders, intermediaries, audiences),
  \item $\insteco$ is an institutional ecology (the ensemble of belief-processing institutions),
  \item $\narr$ is a narrative space (the set of available narratives and framings),
  \item $\beliefspace$ is a belief space (the space of possible population-level belief configurations),
  \item $\Phi: \R^m_{\geq 0} \times \insteco \to \narr$ is the \emph{production function} mapping funding and institutions to narrative outputs,
  \item $\Psi: \narr \times \beliefspace \to \beliefspace$ is the \emph{internalization function} mapping narrative exposure and prior beliefs to updated beliefs.
\end{itemize}
\end{definition}

\begin{remark}
This definition is deliberately abstract.
It does not specify the internal structure of agents, the topology of the institutional ecology, or the dimensionality of the belief space.
These are specified in later chapters as needed.
The purpose of the tuple is to establish that doxonomics has a \emph{formal object}: a mathematical structure that can be analyzed, compared, and measured.
\end{remark}

The aim is not conspiracy thinking.
A doxonomic system need not have a single controller.
Coordination may be deliberate, distributed, emergent, or structural.
A belief can dominate because many institutions, each acting locally and rationally, produce the same directional effect.
\textcite{herman1988} demonstrated this for media: five structural filters (ownership, advertising dependence, sourcing, flak, and ideological convergence) produce systematic bias without requiring any explicit coordination.
Doxonomics generalizes this insight beyond media to the full institutional landscape of belief production.

In this respect, doxonomics is closer to epidemiology than to melodrama.
Epidemiology does not require a single villain who spreads disease.
It maps transmission pathways, identifies risk factors, measures incidence, and designs interventions.
Doxonomics does the same for belief.


\section{The Core Conversion Problem}
\label{sec:ch1-conversion}

The foundational problem of doxonomics is the \textbf{conversion problem}\index{conversion problem}:

\begin{quote}
\textit{How are resources converted into durable public belief?}
\end{quote}

This conversion is never one-to-one.
One million dollars spent on direct advertising is not equivalent to one million dollars spent on textbooks, academic chairs, or platform placement.
Different channels have different reach, credibility, persistence, and downstream effects.
A think-tank report that shapes a government white paper that shapes a curriculum guideline that shapes twenty years of classroom instruction has a different doxonomic trajectory than a viral social media post that is forgotten in a week.

\begin{model}[First-Order Belief Equation]
\label{mod:first-order}
\index{belief equation}
As a first approximation, the social strength of belief $i$ at time $t$ is
\begin{equation}
\label{eq:first-order}
  \belief{i}(t) = \sum_{k=1}^{m} \funding{k}(t) \cdot \credibility{k} \cdot \relevance{k}{i} \cdot \reach{k} \cdot \persistence{k} - \counter{i}(t)
\end{equation}
where $\funding{k}(t)$ is funding through channel $k$, $\credibility{k}$ is the credibility weight of that channel, $\relevance{k}{i}$ is the relevance of channel $k$ to belief $i$, $\reach{k}$ is the population fraction reached, $\persistence{k}$ is the temporal persistence factor, and $\counter{i}(t)$ is the aggregate countervailing pressure against belief $i$.
\end{model}

This equation is a conceptual scaffold, not a mature model.
It is linear, static, and treats parameters as exogenous constants when in reality they are endogenous, time-varying, and interdependent.
Its purpose is to formalize a basic claim: beliefs are not just argued; they are \emph{financed exposures mediated by institutions}.
Later chapters will replace this linear approximation with dynamical models (Chapter~\ref{ch:formal-models}), threshold effects (Chapter~\ref{ch:common-sense-capture}), and network diffusion (Chapter~\ref{ch:network-analysis}).

Even in its crude form, the equation reveals something important.
Consider two beliefs:
\begin{itemize}[nosep]
  \item $p_1$: ``human nature is selfish,'' promoted through economics curricula ($\credibility{} \approx 0.9$, $\persistence{} \approx 30$ years), management training, advertising, and entertainment.
  \item $p_2$: ``humans are conditional cooperators,'' supported by experimental evidence but with minimal institutional infrastructure ($\funding{} \approx 0$, $\persistence{} \approx 0$).
\end{itemize}
The first-order model predicts $\belief{1} \gg \belief{2}$ regardless of the evidential merits of $p_1$ versus $p_2$.
This is the asymmetry that doxonomics studies.

\begin{remark}
The equation should not be read as claiming that \emph{only} money matters.
Evidence, lived experience, peer discussion, and individual reasoning all contribute to belief formation.
But the equation identifies a structural advantage: beliefs with institutional support have a baseline ``subsidy'' that beliefs without institutional support lack.
The playing field of ideas is not level; doxonomics measures the tilt.
\end{remark}


\section{The Belief Production Chain}
\label{sec:ch1-chain}

The full process from expenditure to social consequence is the \textbf{belief production chain}\index{belief production chain}.
It consists of seven stages, each with its own logic, costs, and failure modes:

\begin{enumerate}[nosep]
  \item \textbf{Funding}: resources are allocated to belief-production channels.
  \item \textbf{Institutional processing}: institutions (universities, think tanks, media, PR firms) transform funding into credentialed narrative.
  \item \textbf{Narrative output}: the processed material takes the form of research papers, op-eds, textbooks, films, curricula, advertisements, policy briefs.
  \item \textbf{Circulation}: narratives reach audiences through media, education, social networks, and platform algorithms.
  \item \textbf{Internalization}: audiences update their beliefs, norms, and practical dispositions in response to narrative exposure.
  \item \textbf{Behavioral consequence}: internalized beliefs alter behavior: cooperation, voting, consumption, labor compliance, tolerance of inequality.
  \item \textbf{Reproduction}: behavioral changes generate new institutional demands, electoral outcomes, and philanthropic priorities that feed back into the funding stage.
\end{enumerate}

\begin{center}
\begin{tikzpicture}[
  node distance=1.8cm,
  every node/.style={draw, rounded corners, minimum height=1cm, minimum width=2.2cm, align=center, font=\small},
  arrow/.style={-Stealth, thick}
]
  \node (F) {Funding};
  \node (I) [right=of F] {Institutional\\Processing};
  \node (N) [right=of I] {Narrative\\Output};
  \node (C) [right=of N] {Circulation};
  \node (Int) [below=1.2cm of C] {Internalization};
  \node (B) [left=of Int] {Behavioral\\Consequence};
  \node (R) [left=of B] {Reproduction};

  \draw[arrow] (F) -- (I);
  \draw[arrow] (I) -- (N);
  \draw[arrow] (N) -- (C);
  \draw[arrow] (C) -- (Int);
  \draw[arrow] (Int) -- (B);
  \draw[arrow] (B) -- (R);
  \draw[arrow, dashed] (R) -- (F) node[midway, draw=none, minimum height=0pt, minimum width=0pt, above, font=\footnotesize\itshape] {feedback};
\end{tikzpicture}
\end{center}

The dashed feedback arrow is critical.
The chain closes when the beliefs produced by the system generate exactly the behaviors and institutions that sustain those beliefs.
At that point the system has reached a \emph{doxonomic equilibrium}.
At equilibrium, the belief appears to reproduce itself spontaneously, without visible effort.
This is the condition that Bourdieu called \textit{doxa}: a belief so deeply embedded that it is no longer experienced as a belief at all.

Chapter~\ref{ch:belief-production-chain} formalizes each stage as a mathematical operator and the full chain as their composition.

\begin{example}[The Neoliberal Belief Production Chain]
\label{ex:neoliberal-chain}
The belief ``markets are the most efficient allocators of resources'' can be traced through a concrete chain:
\begin{enumerate}[nosep]
  \item \textbf{Funding}: Koch, Scaife, Olin, and Bradley foundations directed hundreds of millions of dollars into free-market intellectual infrastructure from the 1970s onward \autocite{mayer2016}.
  \item \textbf{Institutional processing}: Heritage Foundation, Cato Institute, American Enterprise Institute, and Mont P\`{e}lerin Society produced research, policy briefs, and expert testimony \autocite{mirowski2009}.
  \item \textbf{Narrative output}: op-eds, books (\textit{Free to Choose}, \textit{The Road to Serfdom}), textbooks (Mankiw's \textit{Principles of Economics}), and television programs.
  \item \textbf{Circulation}: editorial pages of \textit{Wall Street Journal}, \textit{Financial Times}; economics curricula at thousands of universities; cable news; social media.
  \item \textbf{Internalization}: generations of students, policymakers, journalists, and voters absorbed market-efficiency assumptions as common sense.
  \item \textbf{Behavioral consequence}: deregulation, privatization, austerity, reduced union membership, increased tolerance for inequality.
  \item \textbf{Reproduction}: beneficiaries of deregulation and low taxation reinvested in the same think tanks and foundations, closing the loop.
\end{enumerate}
This is not a conspiracy theory.
It is a documented institutional history with named actors, traceable funding, and measurable outcomes.
The doxonomic contribution is to treat it as a \emph{system} with formal structure, not merely as a sequence of political events.
\end{example}


\section{What Doxonomics Is Not}
\label{sec:ch1-is-not}

Several misreadings should be preempted.

\textbf{Doxonomics is not a claim that all beliefs are manufactured.}
Many beliefs arise from direct experience, logical inference, or scientific investigation.
Doxonomics studies the \emph{differential amplification} of beliefs: the fact that some beliefs receive vastly more institutional support than others.
It does not claim that evidence is irrelevant.

\textbf{Doxonomics is not a claim that truth is reducible to power.}
A belief can be both true and heavily funded, or false and poorly funded, or any combination.
The doxonomic audit reveals the material scaffolding of belief; the epistemic evaluation remains a separate task.
Chapter~\ref{ch:ontology-belief} develops this distinction in detail.

\textbf{Doxonomics is not a claim that people are passive.}
Audiences interpret, resist, reinterpret, ignore, and subvert the narratives they encounter.
A doxonomic system that spends billions may still fail to achieve capture.
Understanding why it fails is as important as understanding why it succeeds.
The field studies production, not automatic absorption.

\textbf{Doxonomics is not conspiracy theory.}
Conspiracy theories posit secret coordination by a small group.
Doxonomic systems can operate through open, legal, well-documented institutional processes.
The funding is often on public record (IRS 990 filings, lobbying disclosures, advertising expenditure reports).
The coordination is often structural (shared incentives, overlapping networks, common training) rather than secret.
The analogy is market competition: firms in the same industry produce similar products not because they conspire, but because they face similar incentives.


\section{Scope and Limits}
\label{sec:ch1-scope}

\begin{principle}[Scope of Doxonomics]
\label{prin:scope}
Doxonomics studies the \emph{differential amplification} of beliefs: the mechanisms by which some beliefs receive vastly more institutional support, repetition, and prestige than others, and the social consequences of this asymmetry.
\end{principle}

The scope is deliberately limited.
Doxonomics does not attempt to explain all belief formation.
It does not replace psychology (which studies individual cognition), epistemology (which studies the conditions for justified belief), or political theory (which studies the legitimacy of institutional arrangements).
It adds something that none of these fields provides: a systematic accounting of the material infrastructure of public belief.

A society that can estimate inflation, unemployment, and carbon intensity, but cannot estimate the industrial production of its own common sense, remains partially blind to itself.
Doxonomics seeks to reduce that blindness.


\section{Neighboring Fields and the Distinctive Contribution}
\label{sec:ch1-neighbors}

Doxonomics draws on and extends several existing disciplines, but is identical to none of them.

\begin{center}
\begin{tikzpicture}[
  field/.style={draw, rounded corners, minimum width=3.5cm, minimum height=0.7cm, align=center, font=\small},
  note/.style={font=\small, text width=6cm, align=left},
]
\matrix[row sep=0.5cm, column sep=0.4cm] {
  \node[field] {Political Economy}; &
  \node[note] {Studies material power but not its conversion into common sense.}; \\
  \node[field] {Sociology of Knowledge}; &
  \node[note] {Studies social formation of ideas but rarely performs financial accounting.}; \\
  \node[field] {Media Studies}; &
  \node[note] {Analyzes representation but not which representations are subsidized.}; \\
  \node[field] {Propaganda Studies}; &
  \node[note] {Maps persuasion techniques but often lacks formal models.}; \\
  \node[field] {Social Epistemology}; &
  \node[note] {Studies collective knowledge but not its material infrastructure.}; \\
  \node[field] {Behavioral Economics}; &
  \node[note] {Studies cognitive bias but not the institutional environment that exploits it.}; \\
  \node[field, fill=doxoblue!15] {\textbf{Doxonomics}}; &
  \node[note] {\textbf{Integrates all of the above around explicit accounting of material support, institutional routing, and measurable belief effects.}}; \\
};
\end{tikzpicture}
\end{center}

The distinctive contribution is the insistence on \emph{accounting}: tracing specific dollar amounts through specific institutions to specific narrative outputs with measurable effects on specific populations.
This is what makes doxonomics a potential science rather than a critical commentary.


\section{The Structure of This Book}
\label{sec:ch1-structure}

The book proceeds in five parts.

\textbf{Part~I (Foundations)} establishes the conceptual base: the ontology of belief, the intellectual prehistory, the core units of analysis, and the belief production chain.

\textbf{Part~II (Concepts and Models)} develops the formal apparatus: epistemic capital, narrative infrastructure, common-sense capture, doxonomic power, the economics of belief production, dynamical models, and sheaf-theoretic coherence.

\textbf{Part~III (Methods)} introduces the empirical toolkit: doxometric measurement, belief-flow accounting, network analysis, causal inference, text analysis, historical methods, experimental design, and model validation.

\textbf{Part~IV (Applications)} applies the machinery to specific domains: selfish-human-nature narratives, meritocracy, consumerism, inequality, race and coloniality, austerity, work discipline, national myth, technology, and crisis.

\textbf{Part~V (Resistance, Repair, and Reconstruction)} turns to counter-strategies, epistemic democracy, belief ecology design, ethics, and open problems.

Throughout, interactive \emph{playgrounds} allow the reader to manipulate funding parameters, simulate diffusion, and test model sensitivity.
These are available at the Piatra Institute's companion site.


\section{Notes and References}
\label{sec:ch1-notes}

The concept of belief as a social product has deep roots.
\textcite{marx1846} argued that the ruling ideas of any epoch are the ideas of the ruling class.
\textcite{gramsci1971} developed this into a theory of hegemony: the production of consent through civil society rather than through coercion alone.
\textcite{bourdieu1984} introduced the concept of \textit{doxa} as the unquestioned background of social life, extended in \textcite{bourdieu1991} to the analysis of language and symbolic power.
\textcite{berger1966} provided the foundational account of reality as a social construction.
\textcite{mannheim1936} developed the sociology of knowledge as a systematic discipline.
\textcite{althusser1970} formalized the concept of ideological state apparatuses.

Turning to media specifically, \textcite{lippmann1922} pioneered the analysis of how mass media manufacture public opinion.
\textcite{herman1988} provided the most systematic media-level account with the propaganda model.
\textcite{ellul1965} analyzed propaganda as a total social phenomenon, not merely a state tool.
The conversion problem connects to \textcite{bernays1928} on the engineering of consent.
For the economic logic of belief subsidies, see \textcite{coase1974} and \textcite{stigler1961}.

On the institutional history of specific belief production systems, \textcite{mirowski2009} traces the Mont P\`{e}lerin Society and the construction of neoliberal common sense.
\textcite{mayer2016} documents the dark-money networks behind American conservative ideology.
\textcite{oreskes2010} details the manufacture of doubt about climate science and tobacco harms.
\textcite{slobodian2018} provides a complementary genealogy of neoliberal internationalism.


\section{Exercises}

\begin{exercise}
\label{ex:ch1-identify}
Identify three widely held beliefs in your society that present themselves as descriptions of human nature or social reality (e.g., ``competition drives excellence,'' ``there is no alternative to the market'').
For each, list two institutions that regularly repeat them and estimate (order of magnitude) the annual budget of those institutions.
\end{exercise}

\begin{exercise}
\label{ex:ch1-equation}
Using the first-order belief equation~\eqref{eq:first-order}, argue qualitatively why a belief can persist even when the underlying evidence shifts against it, provided $\funding{k}$ and $\persistence{k}$ remain high.
Give a real-world example.
\end{exercise}

\begin{exercise}
\label{ex:ch1-chain}
Consider the belief ``markets are self-correcting.''
Map it onto the belief production chain, identifying at least one concrete actor or institution at each stage.
Where is the chain strongest?
Where is it most vulnerable to disruption?
\end{exercise}

\begin{exercise}
\label{ex:ch1-tuple}
Extend Definition~\ref{def:doxonomic-system} to include a \emph{resistance function} $\Omega: \beliefspace \to \R^m_{\geq 0}$ that models how population beliefs generate counter-pressure on funding channels.
Write down a fixed-point condition for equilibrium and discuss what it would mean for a society to be at such a fixed point.
\end{exercise}

\begin{exercise}
\label{ex:ch1-comparison}
Choose a belief that has \emph{lost} institutional support over the past century (e.g., the divine right of kings, explicit racial hierarchy as scientific consensus, the gold standard as natural law).
Trace its decline through the belief production chain.
At which stage did the chain break first?
\end{exercise}

\begin{exercise}
\label{ex:ch1-nonmonetary}
Identify a doxonomic system that operates primarily through non-monetary channels, such as family socialization, religious ritual, or architectural design.
How would you modify Definition~\ref{def:doxonomic-system} to accommodate non-monetary resource flows?
What does this reveal about the limits of a funding-centered framework?
\end{exercise}

\begin{exercise}
\label{ex:ch1-epidemiology}
The text compares doxonomics to epidemiology.
Push the analogy further: what would the doxonomic equivalents of ``pathogen,'' ``transmission vector,'' ``incubation period,'' ``herd immunity,'' and ``vaccination'' be?
Where does the analogy break down?
\end{exercise}

\chapter{The Ontology of Public Belief}
\label{ch:ontology-belief}

\epigraph{What is most difficult to perceive is what is taken for granted.}{Pierre Bourdieu}

\noindent
Before we can study how beliefs are produced, we must clarify what a belief \emph{is}, and what distinguishes a \emph{public} belief from a private mental state.
Doxonomics is not psychology.
Its object is not the individual mind but the social distribution of interpretive commitments: what a population takes to be true, normal, natural, or inevitable, and how that distribution is shaped by institutional forces.

This chapter provides the ontological ground floor.
It distinguishes opinion from belief from ideology from common sense from doxa, formalizes each, addresses the gap between private thought and public expression, introduces anthropological regimes, develops a layered ontology that goes beyond propositions, and confronts the relationship between truth and power that any such project must face.


\section{Opinion, Belief, Ideology, Common Sense, Doxa}
\label{sec:ch2-taxonomy}

Ordinary language conflates several distinct objects.
A newspaper editorial is an ``opinion''; a religious conviction is a ``belief''; a political program is an ``ideology''; the assumption that growth is good is ``common sense''; and the taken-for-granted nature of wage labor is something so deep it barely has a name.
Doxonomics needs sharp distinctions among these because each has a different production cost, institutional profile, and vulnerability to disruption.

\begin{definition}[Belief]
\label{def:belief}
\index{belief}
A \emph{belief} is a proposition $p$ together with a subjective probability $\pi(p) \in [0,1]$ that an agent assigns to its truth.
A \emph{public belief} is a measurable aggregate: a distribution $\mu$ over the population such that $\mu(p) = \E_{j \in \text{Pop}}[\pi_j(p)]$ gives the population-average credence in $p$.
\end{definition}

\begin{remark}
\label{rem:belief-limitations}
This definition is precise but deliberately narrow.
It models belief as propositional and probabilistic, suitable for claims like ``most people are selfish'' or ``markets are efficient,'' which can be assigned truth values.
Much of what doxonomics studies, however, does not fit neatly into this format.
We address this limitation in Section~\ref{sec:ch2-layers} below.
The propositional model is a workable first approximation, not a claim about the deep structure of cognition.
\end{remark}

An \emph{opinion}\index{opinion} is a belief that the holder recognizes as contestable, one among several possible views.
The distinguishing feature is reflexive awareness: the opinion-holder knows that others disagree and does not experience the disagreement as irrational.
Opinions are relatively cheap to produce and easy to change.

\begin{definition}[Ideology]
\label{def:ideology}
\index{ideology}
An \emph{ideology} is a structured set of propositions $\mathcal{P} = \{p_1, \ldots, p_n\}$ equipped with an internal logic (deductive, narrative, or affective) such that acceptance of a subset tends to raise credence in the others.
Formally, for $S \subset \mathcal{P}$ and $p_j \notin S$:
\[
  \E[\pi(p_j) \mid \pi(p_i) > \tau \text{ for all } p_i \in S] > \E[\pi(p_j)]
\]
The \emph{coherence strength} of the ideology is
\[
  \kappa(\mathcal{P}) = \frac{1}{|\mathcal{P}|^2} \sum_{i \neq j} \mathrm{Corr}(\pi(p_i), \pi(p_j))
\]
measuring how tightly the propositions are bundled in the population's belief structure.
\end{definition}

\begin{example}
\label{ex:neoliberal-ideology}
The neoliberal ideology bundles: $p_1$ = ``markets are efficient,'' $p_2$ = ``government intervention distorts,'' $p_3$ = ``individuals are responsible for their outcomes,'' $p_4$ = ``deregulation promotes growth,'' $p_5$ = ``inequality reflects differential contribution.''
Accepting $p_1$ raises credence in $p_2$--$p_5$ not by logical entailment but by narrative coherence: they ``fit together'' as a worldview.
A person exposed to $p_1$ through economics education is more likely to accept $p_3$ through meritocratic framing, even though the two claims are logically independent.
The doxonomic significance of ideology is that bundling reduces the per-proposition production cost: investing in one node of the bundle raises credence across the whole set.
\end{example}

\begin{definition}[Common Sense]
\label{def:common-sense}
\index{common sense}
A proposition $p$ is \emph{common sense} in a population if:
\begin{enumerate}[nosep]
  \item $\mu(p) > \theta$ for a high threshold $\theta$ (widespread adoption),
  \item $\mathrm{Var}(\pi_j(p))$ is low (tight concentration: near-universal belief, not merely high average),
  \item $p$ is rarely challenged in mainstream institutional discourse.
\end{enumerate}
Common sense is not merely widespread belief; it is belief whose distribution is tightly concentrated around high credence and whose contestation is socially costly.
\end{definition}

\textcite{bourdieu1984} introduced \textit{doxa}\index{doxa} as the deepest layer: propositions so thoroughly embedded that they are not perceived as beliefs at all.
They are experienced as the texture of reality rather than as assertions.
A person operating within doxa does not \emph{agree} that wage labor is natural; they simply cannot imagine economic life organized otherwise.

\begin{definition}[Doxa]
\label{def:doxa}
\index{doxa!formal definition}
A proposition $p$ is \emph{doxic} in a population if it satisfies the conditions for common sense and, additionally:
\begin{enumerate}[nosep]
  \setcounter{enumi}{3}
  \item $p$ is not represented in public discourse as a contestable claim. It is invisible as a proposition.
\end{enumerate}
The transition from common sense to doxa is the transition from ``everyone agrees'' to ``there is nothing to agree or disagree about.''
This is the deepest form of doxonomic achievement: the belief has been absorbed into the infrastructure of perception.
\end{definition}

\begin{remark}
The invisibility condition (4) is the hardest to operationalize.
One proxy: the frequency with which $p$ appears as the \emph{conclusion} of an argument versus the frequency with which it appears as an unstated \emph{premise}.
A doxic proposition is almost never argued for; it is argued \emph{from}.
Corpus analysis (Chapter~\ref{ch:text-media}) can detect this asymmetry.
\end{remark}

The doxonomic production cost increases as we move from opinion to doxa.
Changing an opinion requires a persuasive argument.
Changing common sense requires sustained institutional effort.
Changing doxa requires making the invisible visible, which is the most difficult and most important task of doxonomic analysis.

\begin{center}
\begin{tikzpicture}[
  every node/.style={font=\small},
  level/.style={draw, thick, rounded corners, minimum width=4cm, minimum height=0.7cm, align=center},
]
\matrix[row sep=0.4cm] {
  \node[level, fill=doxoblue!5] {Opinion}; &
  \node[font=\small, text width=7cm] {Recognized as contestable. Cheap to produce. Easy to change.}; \\
  \node[level, fill=doxoblue!15] {Belief}; &
  \node[font=\small, text width=7cm] {Held with credence. May or may not be reflexively examined.}; \\
  \node[level, fill=doxoblue!25] {Ideology}; &
  \node[font=\small, text width=7cm] {Bundled propositions. Internal coherence reinforces adoption.}; \\
  \node[level, fill=doxoblue!40] {Common Sense}; &
  \node[font=\small, text width=7cm] {Near-universal, low variance, rarely challenged. Costly to contest.}; \\
  \node[level, fill=doxoblue!60, text=white] {Doxa}; &
  \node[font=\small, text width=7cm] {Invisible as a claim. Experienced as reality. Most durable.}; \\
};
\end{tikzpicture}
\end{center}


\section{Private Belief and Public Expression}
\label{sec:ch2-private-public}

The distinction between what people privately think and what they publicly act on is not merely philosophical.
It has precise game-theoretic structure, and it is doxonomically consequential, because a doxonomic system need not convince people in order to govern them.
It suffices to maintain the \emph{appearance} of consensus.

\begin{definition}[Pluralistic Ignorance]
\label{def:pluralistic-ignorance}
\index{pluralistic ignorance}
A population exhibits \emph{pluralistic ignorance} with respect to $p$ if:
\begin{enumerate}[nosep]
  \item $\mu_{\text{private}}(p) < \tau$: most individuals privately doubt $p$,
  \item $\mu_{\text{expressed}}(p) > \tau$: most individuals publicly affirm $p$,
  \item this gap is sustained by the false belief that most others sincerely hold $p$.
\end{enumerate}
\end{definition}

\begin{example}
\label{ex:pluralistic-ignorance}
In many corporate environments, employees privately doubt that constant growth is sustainable or that competition among colleagues improves performance.
But each employee, observing that everyone else appears to endorse these claims (in meetings, in performance reviews, in public communication), concludes that their own doubt is unusual.
The result: a collective performance of belief that no individual sincerely holds.
The doxonomic system need only maintain the performance, not the conviction.
\end{example}

\textcite{kuran1995} formalized this as \emph{preference falsification}: individuals misrepresent their private beliefs to conform to perceived public norms, and the aggregate misrepresentation creates the very norms it appears to obey.
The cost of maintaining false consensus is typically lower than the cost of genuine persuasion, which is one reason pluralistic ignorance is a recurrent feature of doxonomic regimes.

\begin{model}[Preference Falsification Game]
\label{mod:falsification}
\index{preference falsification}
Let $n$ agents each have a private belief $\pi_j \in [0,1]$ and choose a public expression $e_j \in \{0, 1\}$ (affirm or dissent).
Each agent faces:
\begin{itemize}[nosep]
  \item a \emph{conformity payoff} $c \cdot |\bar{e} - e_j|$ penalizing deviation from the average expression,
  \item a \emph{sincerity cost} $s \cdot |e_j - \mathbf{1}[\pi_j > 0.5]|$ from expressing what one does not believe,
  \item a \emph{reputation payoff} $r(e_j)$ from institutional rewards for affirming $p$.
\end{itemize}
When $c + r > s$, agents affirm $p$ regardless of private belief.
The equilibrium expression $\bar{e}^* = 1$ is then supported even if $\mu_{\text{private}}(p) < 0.5$.
\end{model}

\begin{proposition}[Fragility of Pluralistic Ignorance]
\label{prop:fragility}
A pluralistic-ignorance equilibrium is fragile: if a sufficient number of agents simultaneously reveal their private beliefs, the equilibrium collapses.
Formally, if $k > k^*$ agents express dissent, the conformity payoff reverses sign and a cascade of honest expression follows.
The threshold $k^*$ depends on the distribution of private beliefs and the conformity parameter $c$.
This is the mechanism underlying the sudden collapse of apparently stable regimes, what \textcite{kuran1995} calls ``the element of surprise in revolutions.''
\end{proposition}

This has a direct doxonomic implication: some apparently dominant belief regimes are held together by expression management, not by genuine conviction.
Doxonomic systems that rely on pluralistic ignorance rather than deep internalization are cheaper to maintain but more vulnerable to sudden collapse.
Chapter~\ref{ch:how-belief-regimes-fracture} develops this further.


\section{Anthropological Regimes}
\label{sec:ch2-anthropological}

The most consequential beliefs are not about policy but about \emph{what human beings are}.
Policy disputes presuppose an anthropology: if humans are selfish, then institutions must use incentives; if they are cooperative, then institutions can rely on reciprocity; if they are plastic, then the institutional environment itself determines the outcome.

\begin{definition}[Anthropological Regime]
\label{def:anthropological-regime}
\index{anthropological regime}
An \emph{anthropological regime} is a dominant practical model $\alpha$ of human nature, specifying:
\begin{enumerate}[nosep]
  \item a \emph{motivational theory} (what drives people: self-interest, status, care, duty, fear, \ldots),
  \item a \emph{social ontology} (whether persons are fundamentally isolated or constitutively relational),
  \item a \emph{capacity claim} (what people can and cannot do collectively: cooperate without coercion? govern themselves? share resources sustainably?).
\end{enumerate}
Formally, $\alpha \in \mathcal{A}$ where $\mathcal{A}$ is the space of such models, and the population's effective anthropological regime is
\[
  \hat{\alpha} = \arg\max_{\alpha \in \mathcal{A}} \sum_j w_j \cdot \mathbf{1}[\alpha_j = \alpha]
\]
where $w_j$ are institutional weights reflecting not headcount but institutional embeddedness.
\end{definition}

The institutional weighting $w_j$ is essential.
An anthropological regime can dominate even if most individuals privately doubt it, provided it is embedded in the institutions that structure daily life: workplaces, schools, legal systems, media.
A worker who privately believes in cooperation but operates within a workplace designed around the assumption of shirking lives under a selfish-anthropology regime regardless of private conviction.

\begin{example}[Competing Anthropological Regimes]
\label{ex:competing-regimes}
\index{anthropological regime!examples}
Three historically significant regimes:
\begin{description}[nosep]
  \item[Hobbesian] ($\alpha_H$): humans are competitive by nature; order requires sovereignty, enforcement, and hierarchy.
    Institutional expression: criminal justice systems, competitive labor markets, principal-agent contracts.
  \item[Smithian-Neoclassical] ($\alpha_S$): humans are rational self-interest maximizers; social welfare emerges from individual optimization within markets.
    Institutional expression: economics curricula, cost-benefit analysis, incentive-based regulation.
  \item[Cooperative] ($\alpha_C$): humans are conditional cooperators; collective action is possible under appropriate institutional conditions.
    Institutional expression: commons governance, cooperatives, public goods provision, Scandinavian welfare states.
\end{description}
The doxonomic question is not which is ``true'' (all capture real tendencies) but which receives the most institutional support, and how that support shapes the evidence environment.
\end{example}

Anthropological regimes have a distinctive doxonomic property: they can be \emph{self-fulfilling}.
If institutions are designed around the assumption that people are selfish, they create environments that reward selfish behavior, which generates behavioral evidence for the selfish assumption, which justifies maintaining the institutions.
Chapter~\ref{ch:human-nature} analyzes this feedback loop in detail.


\section{Beyond Propositions: A Layered Ontology}
\label{sec:ch2-layers}

The propositional-credence framework of the previous sections is a useful formal tool, but it does not exhaust the phenomena doxonomics must study.
Many politically consequential ``beliefs'' are not well-modeled as explicit credences in propositions.

\begin{itemize}[nosep]
  \item \textbf{Practical schemas} (Bourdieu's \textit{habitus}\index{habitus}): embodied dispositions that organize perception and action without passing through conscious propositional evaluation.
  A person raised in a competitive school system does not decide that ``life is competition''; they \emph{perceive} the world through competitive categories, feel anxiety when not competing, and organize action around comparative performance, all without propositional mediation.

  \item \textbf{Social norms}\index{social norm}: shared expectations about behavior that agents follow without necessarily believing them to be true or good.
  A norm can govern behavior even when no one endorses it \autocite{cialdini2006}.

  \item \textbf{Identity commitments}\index{identity commitment}: affiliations (national, religious, professional, racial, gendered) that structure interpretation prior to evidence evaluation.
  Identity acts as a \emph{prior filter}: evidence consistent with identity is more readily absorbed; evidence threatening to identity is resisted or reinterpreted \autocite{kahan2012}.

  \item \textbf{Affective orientations}\index{affective orientation}: trust, fear, disgust, and solidarity toward social categories, felt rather than believed.
  Racial prejudice, for example, often operates through affective response rather than propositional endorsement: a person may sincerely disavow racist beliefs while exhibiting implicit bias in behavior.

  \item \textbf{Tacit classifications}\index{tacit classification}: the taken-for-granted categories through which experience is organized, such as ``worker,'' ``citizen,'' ``consumer,'' ``taxpayer,'' ``illegal.''
  These classifications are not beliefs \emph{about} the world; they are the \emph{grid} through which the world is apprehended.
  Changing a classification system is among the most difficult and consequential doxonomic interventions.
\end{itemize}

A person can \emph{enact} an anthropological regime they cannot state propositionally.
A worker who ``knows'' that competition is natural may not hold an explicit belief to that effect; the assumption is built into their behavioral repertoire, their emotional responses, and their sense of what is and is not possible.

\begin{principle}[Layered Ontology]
\label{prin:layered-ontology}
\index{layered ontology}
Doxonomic analysis distinguishes at least five layers of socially shaped cognition:
\begin{enumerate}[nosep]
  \item \textbf{Propositional belief}: explicit credence in a statement ($\pi_j(p)$).
  \item \textbf{Norm perception}: belief about what others believe or do ($\hat{\mu}_j$).
  \item \textbf{Practical schema}: embodied disposition shaping action without deliberation.
  \item \textbf{Identity commitment}: group affiliation structuring interpretation.
  \item \textbf{Affective orientation}: emotional stance toward categories and institutions.
\end{enumerate}
The probability representation used in most formal models of this book operates at layer~1.
This is a \emph{local convenience}, not a full ontology of ideological life.
Layers 2--5 are acknowledged as important and addressed qualitatively where needed; their full formalization is an open problem for the field.
\end{principle}

\begin{remark}
The layers differ in their doxonomic production pathways.
Layer~1 (propositional belief) is most responsive to evidence and argument, and therefore to traditional persuasion.
Layer~3 (practical schema) is produced primarily through prolonged institutional immersion (schooling, workplace discipline, military training) and is resistant to propositional counter-argument.
Layer~5 (affective orientation) is shaped by repetition, association, and social identity, often below the threshold of conscious processing.
A doxonomic system that operates at multiple layers simultaneously is more durable than one operating at layer~1 alone, because it cannot be dislodged by argument.
\end{remark}


\section{Truth, Power, and the Scope of Doxonomic Analysis}
\label{sec:ch2-truth-power}

Doxonomics studies how beliefs are produced, not whether they are true.
But this scope statement requires care, because the relationship between truth and power is not simple.
A naively stated doxonomics could slide into the claim that ``all socially successful beliefs are just power effects,'' which is both philosophically untenable and empirically false.

A belief may be:
\begin{enumerate}[nosep]
  \item funded \emph{because} it is true (scientific consensus attracts institutional support: the germ theory of disease, the efficacy of vaccines),
  \item made to \emph{appear} true by subsidy (repetition creates familiarity, which is mistaken for evidence: ``tax cuts pay for themselves''),
  \item partly true but selectively amplified (``markets create wealth'' is true in some senses, but the selective emphasis obscures distributional effects, externalities, and power asymmetries),
  \item normatively useful but descriptively false (``hard work always pays off'' motivates effort but misdescribes the economy),
  \item descriptively accurate only under institutions built around it (``people are selfish'' becomes true when institutions reward selfishness and punish cooperation: the self-fulfilling case).
\end{enumerate}

Each type requires a different doxonomic analysis.
Type~1 is not a doxonomic pathology at all; it is the normal operation of evidence-responsive institutions.
Type~2 is pure doxonomic distortion.
Types~3--5 are mixed cases that require careful disentangling.

\begin{principle}[Separation of Audit and Evaluation]
\label{prin:separation}
\index{audit and evaluation, separation of}
The doxonomic audit (tracing material scaffolding) and the epistemic evaluation (assessing truth status) are \emph{distinct operations}.
A heavily funded belief is not automatically suspect; a poorly funded belief is not automatically true.
The audit reveals the scaffolding.
The evaluation requires independent evidence and argument.
Conflating the two is a category error that doxonomics must avoid \autocite{mannheim1936, fricker2007}.
\end{principle}

\begin{remark}
This principle protects the field in both directions.
It prevents critics from dismissing doxonomics as ``everything you disagree with is ideology.''
And it prevents practitioners from assuming that revealing a belief's material support settles the question of its truth.
The most interesting doxonomic cases are precisely those where the truth status is mixed or uncertain, and where the material scaffolding is doing real epistemic work, for better or worse.
\end{remark}


\section{Belief as Probability Distribution}
\label{sec:ch2-probability}

With the layered ontology acknowledged, we proceed to the formal representation used in most of this book.
For mathematical purposes, we represent a population's propositional beliefs as a probability measure $\mu$ on a proposition space $(\Omega, \mathcal{F})$, where $\Omega$ is the set of possible world-models and $\mathcal{F}$ is a $\sigma$-algebra of propositions.

\begin{remark}
\label{rem:distribution-not-mean}
Population-level belief is better represented as a \emph{distribution} over individual credences than as a single mean.
A society with 50/50 polarization and a society with universal moderate uncertainty can have the same mean credence $\mu(p) = 0.5$, but these are radically different doxonomic situations.
Throughout this book, when we write $\belief{i}$, we refer to a summary statistic.
Full doxonomic analysis should attend to:
\begin{itemize}[nosep]
  \item the \emph{shape} of the credence distribution (unimodal, bimodal, skewed),
  \item its \emph{variance} and kurtosis,
  \item its \emph{clustering by subgroup} (class, race, education, media environment),
  \item the \emph{gap between private and expressed beliefs} (pluralistic ignorance),
  \item the \emph{elite-mass divergence} (do elites and populations hold the same beliefs?).
\end{itemize}
\end{remark}

\begin{definition}[Belief Distribution]
\label{def:belief-distribution}
\index{belief distribution}
The \emph{belief distribution} for proposition $p$ in a population of $N$ agents is the empirical measure
\[
  \hat{F}_p(x) = \frac{1}{N} \sum_{j=1}^{N} \mathbf{1}[\pi_j(p) \leq x]
\]
We define the following summary statistics:
\begin{itemize}[nosep]
  \item Mean belief: $\bar{\mu}(p) = \frac{1}{N}\sum_j \pi_j(p)$
  \item Belief variance: $\sigma^2(p) = \frac{1}{N}\sum_j (\pi_j(p) - \bar{\mu}(p))^2$
  \item Polarization index: $\mathrm{Pol}(p) = 4\bar{\mu}(p)(1 - \bar{\mu}(p)) \cdot (1 + \sigma^2(p)/\bar{\mu}(p)(1-\bar{\mu}(p)))$ when $\bar{\mu}(p) \in (0,1)$
\end{itemize}
A belief is \emph{polarized} when $\hat{F}_p$ is bimodal (clustered near 0 and 1) and $\mathrm{Pol}(p)$ is high.
\end{definition}

\begin{proposition}[Belief Updating Under Unequal Exposure]
\label{prop:unequal-updating}
Let agents update beliefs via Bayes' rule.
If agent $j$ observes evidence $e_j$ drawn from a distribution $\nu_j$ that depends on the channels $j$ is exposed to, then even rational agents will converge to different posteriors when their evidence streams differ systematically.
Formally, if $\nu_j \neq \nu_k$ for agents $j, k$ in different media environments, then generically $\pi_j(p \mid e_j) \neq \pi_k(p \mid e_k)$ even in the limit.
\end{proposition}

\begin{proof}[Proof sketch]
Under Bayes' rule, the posterior after $T$ signals is
\[
  \pi_j^{(T)}(p) = \frac{\pi_j^{(0)}(p) \prod_{t=1}^{T} L(s_t^{(j)} \mid p)}{\pi_j^{(0)}(p) \prod_{t=1}^{T} L(s_t^{(j)} \mid p) + (1 - \pi_j^{(0)}(p)) \prod_{t=1}^{T} L(s_t^{(j)} \mid \neg p)}
\]
When the signal distribution $\nu_j$ systematically favors $p$ (more signals consistent with $p$ than with $\neg p$), the likelihood ratio $\prod_t L(s_t \mid p)/L(s_t \mid \neg p) \to \infty$, driving $\pi_j^{(T)} \to 1$.
Two agents with the same prior but different signal distributions will diverge.
The divergence rate is bounded below by the KL divergence $D_{\mathrm{KL}}(\nu_j \| \nu_k)$.
\end{proof}

This is the Bayesian foundation of doxonomics: differential exposure to evidence, not irrationality, suffices to produce differential belief.
Rational agents in asymmetric information environments will hold different beliefs. The doxonomic question is: who controls the information environments?

\begin{remark}
The Bayesian model assumes agents are rational updaters.
In practice, agents also exhibit motivated reasoning (updating asymmetrically depending on whether evidence confirms or threatens identity), confirmation bias (seeking evidence consistent with prior beliefs), source credibility weighting (discounting evidence from distrusted sources regardless of content), and bounded attention (processing only a fraction of available information).
Each of these makes the doxonomic effect \emph{stronger} than the pure Bayesian model predicts: biased exposure combined with biased processing produces faster and more durable divergence.
\end{remark}


\section{Operationalizing the Invisible: How to Detect Doxa}
\label{sec:ch2-detecting-doxa}

The deepest doxonomic products are, by definition, hard to see.
If a proposition is truly doxic (invisible as a claim), how does the researcher identify it?

Several heuristic methods are available:

\begin{enumerate}[nosep]
  \item \textbf{Premise analysis}: search large corpora for propositions that appear overwhelmingly as unstated premises rather than as explicit conclusions.
  If ``economic growth is desirable'' appears in 10,000 texts as a background assumption but is explicitly argued for in only 50, it is behaving doxically.

  \item \textbf{Contrast with other societies}: propositions that are doxic in one society but explicitly debated in another are candidates for doxonomic analysis.
  The assumption that health care is a market commodity is doxic in the United States but explicitly contested in most European countries.

  \item \textbf{Historical discontinuity}: if a proposition was once explicitly argued for (with named advocates and institutional campaigns) but is now ``just how things are,'' the transition from argument to doxa is itself a doxonomic achievement worth tracing.

  \item \textbf{Anomaly response}: observe how institutions respond when someone challenges a doxic proposition.
  If the response is not counter-argument but \emph{confusion, dismissal, or ridicule} (``that's just how the world works''), the proposition is likely doxic.

  \item \textbf{Cross-domain consistency}: if the same anthropological assumption appears in economics, management, law, media, and everyday speech without any of these domains citing the others, it has achieved the kind of distributed invisibility characteristic of doxa.
\end{enumerate}

\begin{exercise}
\label{ex:ch2-detect-doxa}
Apply the five heuristic methods above to the proposition ``economic growth is the primary objective of government.''
For each method, describe what evidence you would look for and what a positive finding would indicate.
\end{exercise}


\section{Notes and References}
\label{sec:ch2-notes}

The taxonomy of belief, ideology, and common sense draws on \textcite{gramsci1971}, \textcite{bourdieu1984}, and \textcite{foucault1971}.
The social construction of reality as a general framework is developed in \textcite{berger1966}; the sociology of knowledge as a discipline originates with \textcite{mannheim1936}.
\textcite{althusser1970} provides a structural account of how institutions reproduce ideology through ideological state apparatuses.

Pluralistic ignorance was formalized by \textcite{kuran1995}.
The game-theoretic structure of public belief connects to coordination games in \textcite{schelling1978}.
On cognitive dissonance and belief maintenance under contradictory evidence, see \textcite{festinger1957}.
Cultural cognition and identity-protective reasoning are developed in \textcite{kahan2012} and \textcite{mercier2017}.
Social identity theory is foundational for understanding how group affiliation structures belief \autocite{tajfel1979}.

Bayesian updating under asymmetric information exposure is treated in \textcite{gentzkow2010}.
The concept of anthropological regimes extends \textcite{gintis2005}.
On epistemic injustice (how social power shapes who is heard and believed), see \textcite{fricker2007} and \textcite{medina2013}.
The layered ontology draws on Bourdieu's habitus, \textcite{hall1980} on encoding/decoding, and the pragmatist tradition of \textcite{dewey1927}.
On the relation between truth and power in social analysis, see \textcite{mannheim1936} and \textcite{jasanoff2004}.
On performativity (the way models and descriptions can remake the reality they claim to describe), see \textcite{mackenzie2006}.

The operationalization of doxa through corpus analysis connects to methods in Chapter~\ref{ch:text-media}.
For mass opinion formation as a multi-layered process, see \textcite{zaller1992}.
On propaganda as a total social phenomenon rather than merely a state tool, see \textcite{ellul1965}.


\section{Exercises}

\begin{exercise}
\label{ex:ch2-coordination}
Give a formal example of pluralistic ignorance using a coordination game with $n$ players.
Show that an equilibrium exists in which all players publicly affirm $p$ despite privately doubting it.
Compute the critical mass $k^*$ of dissenters needed to collapse the equilibrium as a function of the conformity parameter $c$.
\end{exercise}

\begin{exercise}
\label{ex:ch2-kl}
Consider two populations exposed to different evidence streams $\nu_A$ and $\nu_B$ about human cooperation.
Using Bayes' rule, show that after $T$ periods their posterior beliefs on ``human nature is selfish'' diverge by an amount proportional to the KL divergence $D_{\mathrm{KL}}(\nu_A \| \nu_B)$.
What happens when $T \to \infty$?
\end{exercise}

\begin{exercise}
\label{ex:ch2-classify}
Classify the following as opinion, belief, ideology, common sense, or doxa in contemporary Western societies:
(a)~``hard work leads to success,''
(b)~``the earth orbits the sun,''
(c)~``economic growth is good,''
(d)~``taxation is theft,''
(e)~``a person's worth is measured by their productivity,''
(f)~``democratic elections are the legitimate way to choose leaders.''
Justify each classification using the formal definitions above.
For any that you classify as common sense or doxa, describe what institutional apparatus maintains them.
\end{exercise}

\begin{exercise}
\label{ex:ch2-layers}
For the belief ``inequality is a natural consequence of differences in talent and effort,'' identify how this belief operates at each of the five layers of the layered ontology.
At which layer is it most deeply embedded?
At which layer would intervention be most effective?
\end{exercise}

\begin{exercise}
\label{ex:ch2-selffullfilling}
Formalize the self-fulfilling property of anthropological regimes.
Let institutions design rules based on regime $\alpha$, and let agents respond to those rules.
Under what conditions does the observed behavior confirm $\alpha$ regardless of the ``true'' human nature?
(Hint: model this as a game where the institutional designer chooses a mechanism based on $\alpha$, and agents best-respond.)
\end{exercise}

\begin{exercise}
\label{ex:ch2-coherence}
Compute the coherence strength $\kappa(\mathcal{P})$ for a hypothetical ideology with five propositions.
Suppose that beliefs in $p_1$ and $p_2$ are highly correlated ($r = 0.8$), $p_3$ and $p_4$ are moderately correlated ($r = 0.4$), and all other pairs are uncorrelated ($r = 0$).
What is $\kappa$?
Is this a ``strong'' ideology by your measure?
\end{exercise}

\begin{exercise}
\label{ex:ch2-doxa-detection}
Choose a proposition you suspect is doxic in your society.
Apply at least three of the five heuristic detection methods from Section~\ref{sec:ch2-detecting-doxa}.
Report what you find.
If the proposition turns out to be common sense rather than doxa, explain what evidence would distinguish the two.
\end{exercise}

\chapter{A Brief Prehistory of Doxonomics}
\label{ch:prehistory}

\epigraph{The philosophers have only interpreted the world, in various ways; the point is to change it.}{Karl Marx, \textit{Theses on Feuerbach} (1845)}

\noindent
Doxonomics is new as a named discipline.
But the insight that beliefs are shaped by material power is ancient.
This chapter traces the intellectual lineage, extracting from each predecessor the formal insight that doxonomics inherits.
The genealogy is not meant to be exhaustive. It is organized around the question: \emph{who saw the conversion problem, and what tools did they develop to analyze it?}


\section{Sophists and the Market for Persuasion}
\label{sec:ch3-sophists}

The Greek sophists were, in effect, the first consultants in belief production.
Protagoras, Gorgias, Thrasymachus, and their colleagues sold rhetorical skill, the capacity to make the weaker argument appear stronger.
They charged fees.
They competed for students and patrons.
They were, in modern terms, a market for persuasion services.

Plato attacked them for it, famously in the \textit{Gorgias} and the \textit{Republic}.
His objection was that rhetoric separated persuasion from truth: the sophist could make any proposition plausible, regardless of its accuracy, provided the audience was susceptible and the speaker was skilled.
This is already a proto-doxonomic critique.

\begin{remark}
The sophists reveal the first doxonomic principle: \emph{the ability to shape belief is an economic good with a market price}.
In modern notation, there exists a market-clearing price $p^*$ for persuasion services, which implies that wealthier actors can purchase more belief-shaping capacity, a precursor to the concept of epistemic capital ($\ecap$).
The Platonic response (that truth should govern belief) is the normative position that doxonomics, as a descriptive science, neither endorses nor rejects but takes as a benchmark against which to measure distortion.
\end{remark}

What the sophists lacked was a structural analysis.
They understood individual persuasion.
They did not theorize the institutional systems by which persuasion becomes common sense across an entire polity.


\section{Hobbes, Mandeville, and the Construction of Selfish Anthropology}
\label{sec:ch3-hobbes}

\textcite{hobbes1651} proposed that the natural condition of humanity is a war of all against all, requiring absolute sovereignty to impose order.
\textcite{mandeville1714} argued in \textit{The Fable of the Bees} that private vices produce public benefits. Selfishness is not only natural but socially productive, and attempts to enforce virtue would impoverish society.

These are not merely philosophical positions.
They are \emph{anthropological regimes} (Definition~\ref{def:anthropological-regime}) whose institutional installation required centuries of work: in political theory, in economics curricula, in management doctrine, in legal reasoning.
By the late twentieth century, the Hobbesian-Mandevillean anthropology had become operational common sense in most Western institutions, not because it won a philosophical argument, but because it was embedded in the design of markets, contracts, incentive systems, and governance structures.

\begin{remark}
\label{rem:hobbes-doxonomic}
The doxonomic question about Hobbes is not ``was he right?'' but ``what institutional investment was required to make his anthropology the default assumption of modern governance, and what alternative anthropologies were crowded out in the process?''
\textcite{bowles2011} and \textcite{ostrom1990} provide evidence that the crowded-out alternatives had substantial empirical support.
\end{remark}

\textcite{smith1776} occupies an ambiguous position.
The \textit{Wealth of Nations} is routinely cited as a founding text of the self-interest paradigm (``It is not from the benevolence of the butcher, the brewer, or the baker that we expect our dinner'').
But Smith also wrote \textit{The Theory of Moral Sentiments} (1759), which grounds human sociality in sympathy, fellow-feeling, and the desire for mutual recognition.
The selective amplification of Smith-as-self-interest-theorist, and the suppression of Smith-as-moral-philosopher, is itself a doxonomic achievement, one that \textcite{mcCloskey2006} has documented in detail.


\section{Marx: Ideology as Material Production}
\label{sec:ch3-marx}

\textcite{marx1846} proposed the most radical version of the thesis: the ruling ideas of any epoch are the ideas of the ruling class, produced not by the force of argument alone but by control over the means of mental production.
In the same way that the class which owns the factories controls the production of goods, it also controls, through ownership of newspapers, publishing houses, universities, and churches, the production of ideas.

\begin{principle}[Marxian Doxonomic Principle]
\label{prin:marx}
\index{Marx}
The class that controls the means of material production simultaneously controls the means of intellectual production.
In doxonomic terms: the funding vector $\fundingvec$ is not exogenous but is determined by the distribution of economic power.
Formally, let $W$ be the wealth distribution and $\fundingvec = g(W)$ for some function $g$ that maps material power to doxonomic investment.
If $g$ is monotonically increasing (wealthier actors invest more in belief production), then doxonomic output is structurally biased toward the interests of the wealthy.
\end{principle}

Marx's contribution to doxonomics is threefold:
\begin{enumerate}[nosep]
  \item the recognition that ideas have material conditions of production,
  \item the claim that the content of dominant ideas systematically reflects the interests of dominant groups,
  \item the concept of ideology as \emph{false consciousness}: beliefs that serve the interests of the powerful while being experienced by the powerless as their own convictions.
\end{enumerate}

His limitation, from a doxonomic standpoint, was insufficient formalization and insufficient attention to the mechanisms of belief transmission.
Marx knew \emph{that} ideas served class interests; he did not provide the accounting tools to trace \emph{how much} was spent, \emph{through which institutions}, with \emph{what measurable effects}.
He also underestimated the autonomy of institutional processing: ideas are not simply piped from the ruling class to the population but are transformed, filtered, and repackaged by intermediary institutions with their own logics.

\textcite{althusser1970} partially addressed this with the concept of \emph{ideological state apparatuses}\index{ideological state apparatus}: schools, churches, media, family, and legal systems that reproduce ideology not through direct coercion but through daily institutional practice.
In doxonomic terms, Althusser identified the institutional ecology $\insteco$ as a set of belief-processing functions that operate with relative autonomy from direct class instruction.


\section{Gramsci: Hegemony}
\label{sec:ch3-gramsci}

\textcite{gramsci1971} introduced \emph{hegemony}\index{hegemony}: the process by which a ruling group secures consent not through coercion but through cultural and intellectual leadership.
Hegemony operates through civil society (schools, churches, media, professional associations, cultural institutions) rather than through the state alone.

Gramsci's key innovation was the recognition that domination through belief is \emph{more effective and more durable} than domination through force.
A state that relies on police is fragile; a state whose citizens genuinely believe that the existing order is natural, just, or inevitable barely needs police at all.

In doxonomic terms, Gramsci's contribution is the recognition that $\insteco$ (the institutional ecology) is the primary site of belief production, and that durable power requires not just force but narrative infrastructure $\narr$.
He also introduced a crucial dynamic concept: hegemony is never total, always contested, and must be continuously \emph{renewed} through institutional work.
This anticipates the doxonomic concept of maintenance cost: even a captured belief requires ongoing investment to remain dominant.

\begin{remark}
Gramsci also theorized \emph{counter-hegemony}: the construction of alternative common sense by subordinate groups through their own institutional networks (workers' councils, alternative media, pedagogical movements).
This prefigures the counter-doxonomic strategies of Chapter~\ref{ch:counter-doxonomic}.
\end{remark}


\section{Lippmann and Dewey: The Public Opinion Debate}
\label{sec:ch3-lippmann-dewey}

In the 1920s, two American thinkers framed a debate that remains central to doxonomics.

\textcite{lippmann1922} argued in \textit{Public Opinion} that ordinary citizens are incapable of understanding the complexity of modern governance.
They perceive the world through ``stereotypes,'' simplified pictures in their heads, produced by media and institutional sources rather than by direct experience.
Public opinion is therefore not the spontaneous expression of citizen judgment but a \emph{manufactured product}, shaped by those who control the information environment.

Lippmann's doxonomic insight: the production function $\Phi$ is controlled by a relatively small number of institutional actors (editors, publishers, experts, officials), and the population's beliefs are the \emph{output} of this function, not its independent input.

\textcite{dewey1927} responded in \textit{The Public and Its Problems} that the solution was not technocratic management of opinion (Lippmann's implicit conclusion) but the \emph{democratization of communication}.
If the public is poorly informed, the remedy is not to hand belief production to experts but to create institutional conditions (free media, public education, open deliberation) under which citizens can form genuine judgments.

\begin{remark}
The Lippmann-Dewey debate maps directly onto a tension within doxonomics.
Lippmann's position leads to a \emph{diagnostic} doxonomics: mapping how belief is produced by elites and processed by publics.
Dewey's position leads to a \emph{constructive} doxonomics: designing institutional environments for democratic belief formation.
This book attempts both, but the tension is real and unresolved.
Chapter~\ref{ch:epistemic-democracy} returns to it.
\end{remark}


\section{Bernays: Engineering Consent}
\label{sec:ch3-bernays}

\textcite{bernays1928} made the conversion problem explicit.
Drawing on his uncle Freud's psychology, Bernays developed techniques for converting corporate funding into public opinion: the staged event, the third-party endorsement, the appeal to unconscious desire, the manipulation of group psychology.

His 1929 ``Torches of Freedom'' campaign (in which he hired women to smoke cigarettes in the New York Easter Parade, framing smoking as a feminist act) is a canonical example of doxonomic engineering.
The client was the American Tobacco Company.
The narrative output was ``smoking = women's liberation.''
The channel was news media (journalists covered the stunt as a news event, lending it third-party credibility).
The behavioral consequence was a measurable increase in women's smoking.

\begin{remark}
Bernays is the first explicit \emph{doxonomic engineer}: a practitioner who consciously understood and optimized the production function $\Phi: \R^m_{\geq 0} \times \insteco \to \narr$.
Public relations, as he conceived it, is applied doxonomics, the deliberate, funded transformation of institutional inputs into belief outputs.
His candor was remarkable.
The opening of \textit{Propaganda}: ``The conscious and intelligent manipulation of the organized habits and opinions of the masses is an important element in democratic society.
Those who manipulate this unseen mechanism of society constitute an invisible government which is the true ruling power of our country.''
\end{remark}

\textcite{ellul1965} extended this analysis.
Where Bernays focused on technique, Ellul argued that propaganda is not merely a tool used by states or corporations but a \emph{total social phenomenon}: modern societies produce propaganda as naturally as they produce consumer goods, because the complexity of modern life creates a demand for simplified interpretive frameworks that only organized narrative production can supply.
In doxonomic terms, Ellul's point is that the demand side of the conversion problem is endogenous: complex societies create the audience conditions under which doxonomic production becomes effective.


\section{The Frankfurt School: Culture Industry}
\label{sec:ch3-frankfurt}

\textcite{adorno1944} introduced the concept of the \emph{culture industry}\index{culture industry}: the industrial-scale production of cultural goods (film, music, radio, magazines) that standardize consciousness and integrate populations into the existing order.

The Frankfurt School's doxonomic insight was that \emph{entertainment is a channel of belief production}, not merely a diversion from it.
A Hollywood film that represents success as individual achievement, romance as consumption, and conflict as resolved through violence is producing an anthropological regime, not by arguing for it but by \emph{narrating} it as the natural shape of human life.

In doxonomic terms, the culture industry operates at layers 3--5 of the layered ontology (Principle~\ref{prin:layered-ontology}): practical schemas, identity commitments, and affective orientations.
It shapes not what people explicitly believe but what they desire, expect, fear, and find normal.
This makes it both more powerful and harder to trace than propositional propaganda.

The limitation of the Frankfurt School, from a doxonomic standpoint, was a tendency toward totalization: the culture industry was presented as so pervasive that resistance seemed impossible.
Later work (by Stuart Hall, the Birmingham cultural studies tradition, and audience reception theorists) showed that audiences are not passive consumers but active interpreters who negotiate, resist, and reinterpret cultural products \autocite{hall1980}.


\section{Foucault: Discourse, Power, and the Production of Truth}
\label{sec:ch3-foucault}

\textcite{foucault1971} shifted attention from who speaks to the \emph{rules of discourse}: the conditions under which a statement counts as knowledge, who is authorized to make it, and what institutional apparatus validates it.

In Foucault's analysis, power does not merely suppress truth; it \emph{produces} truth.
Each society has its ``regime of truth'': the types of discourse it accepts as true, the mechanisms that enable one to distinguish true from false statements, the techniques for acquiring truth, and the status of those who are charged with saying what counts as true.

In doxonomic terms, Foucault's contribution is the insight that $\credibility{k}$ is not a fixed property of a channel but is itself socially produced.
The rules determining what counts as expertise, evidence, or common sense are themselves objects of doxonomic investment.
This is the fourth dimension of doxonomic power (Definition~\ref{def:doxonomic-power} in Chapter~\ref{ch:doxonomic-power}): the power to shape the meta-rules of belief evaluation.

\begin{remark}
Foucault pushes doxonomics toward a second-order analysis.
The first-order model asks: how are beliefs produced?
The Foucauldian question asks: how are the \emph{rules for evaluating beliefs} produced?
Who decides what counts as evidence, what counts as expertise, what counts as a reasonable question?
These meta-rules are themselves doxonomic products, and among the most consequential ones.
A society that controls the definition of ``science,'' ``expert,'' and ``evidence'' controls the credibility weights $\credibility{k}$ for all channels simultaneously.
\end{remark}


\section{Bourdieu: Symbolic Power, Doxa, and Reproduction}
\label{sec:ch3-bourdieu}

\textcite{bourdieu1984, bourdieu1991} provided the richest pre-doxonomic framework.
His contributions map onto doxonomic concepts with unusual precision:

\begin{center}
\begin{tabular}{ll}
\toprule
\textbf{Bourdieu} & \textbf{Doxonomics} \\
\midrule
Symbolic capital & Epistemic capital $\ecap$ \\
Doxa & Common-sense capture (Definition~\ref{def:doxa}) \\
Field & Institutional ecology $\insteco$ \\
Habitus & Practical schema (layer 3) \\
Cultural reproduction & Belief production chain $\bpc$ \\
Misrecognition & Invisible as a claim (doxa condition 4) \\
\bottomrule
\end{tabular}
\end{center}

Bourdieu's analysis of the educational system as a reproduction mechanism is particularly relevant.
In \textit{Distinction}, he showed that the French educational system does not merely transmit knowledge; it converts inherited economic capital into credentialed cultural capital, laundering class privilege through the appearance of meritocratic achievement.
The school system is simultaneously an institution that produces knowledge and an institution that produces the \emph{belief} that the knowledge hierarchy reflects natural differences in ability.

What Bourdieu lacked was the accounting machinery.
He showed that symbolic power reproduces itself through institutions but did not build the input-output tables, trace the funding flows, or estimate the conversion rates.
Doxonomics inherits his sociology and adds the ledger.


\section{Herman and Chomsky: The Propaganda Model}
\label{sec:ch3-herman-chomsky}

\textcite{herman1988} provided the most systematic media-level doxonomic model.
The propaganda model identifies five ``filters'' through which media content passes before reaching the public:

\begin{enumerate}[nosep]
  \item \textbf{Ownership}: media corporations are owned by large conglomerates with interests in maintaining the status quo.
  \item \textbf{Advertising}: the primary revenue source for most media, giving advertisers structural influence over content.
  \item \textbf{Sourcing}: dependence on official sources (government, corporate PR) for raw material.
  \item \textbf{Flak}: organized campaigns to discredit or punish media that deviate from acceptable narratives.
  \item \textbf{Ideology}: shared assumptions (anticommunism in the original formulation; market fundamentalism in later versions) that frame what counts as newsworthy.
\end{enumerate}

In doxonomic terms, the five filters are structural parameters of the narrative infrastructure $\narr$ that systematically bias output without requiring any explicit conspiracy.
The propaganda model's great strength is that it is \emph{testable}: it generates predictions about which stories will be covered, how they will be framed, and which perspectives will be excluded, based on structural properties of the media system rather than on the intentions of individual journalists.

Its limitation is its focus on media.
Doxonomics extends the same structural logic to the full institutional landscape: education, think tanks, platforms, professional training, entertainment, and the built environment.


\section{Colonial Knowledge Production}
\label{sec:ch3-colonial}

The most large-scale doxonomic operations in history were not conducted by think tanks or advertising agencies.
They were conducted by colonial empires, which invested enormous resources in producing knowledge systems (racial taxonomies, civilizational hierarchies, linguistic classifications, legal codes) that justified extraction and governance.

\textcite{said1978} showed that ``Orientalism'' was not merely prejudice but an institutionally sustained body of knowledge: endowed chairs at universities, specialized journals, government translation offices, and training programs for colonial administrators, all dedicated to producing authoritative accounts of the colonized world.
The narrative infrastructure was vast: the British Indian Civil Service alone trained thousands of administrators in a specific interpretive framework for understanding ``the Oriental mind.''

\textcite{fanon1961} analyzed the internalization stage: how colonial education systems replaced indigenous knowledge with the colonizer's worldview, producing subjects who internalized their own subordination.
The colonial school taught not just the colonizer's language but the colonizer's anthropology: the belief that European civilization was the pinnacle of human achievement, that indigenous cultures were backward, and that colonial rule was a civilizing mission.

In doxonomic terms, colonialism built the most extensive narrative infrastructure $\narr$ in pre-digital history, reaching hundreds of millions through missionary schools, colonial administration, and restructured legal systems.
The persistence of these narratives after formal decolonization (what decolonial theorists call \emph{coloniality}\index{coloniality}) is a doxonomic phenomenon: the institutional infrastructure outlasts the political regime that built it.
Chapter~\ref{ch:race-coloniality-gender} develops this analysis.


\section{The Digital Turn: Platforms as Doxonomic Infrastructure}
\label{sec:ch3-digital}

The early twenty-first century introduced a qualitative shift in doxonomic infrastructure.
Social media platforms (Facebook, YouTube, Twitter/X, TikTok) became simultaneously:
\begin{itemize}[nosep]
  \item \emph{channels} of narrative circulation (reaching billions),
  \item \emph{institutions} of narrative processing (algorithmic curation decides what is amplified),
  \item \emph{data-collection systems} that enable micro-targeted doxonomic operations (personalized belief exposure based on behavioral profiles).
\end{itemize}

\textcite{zuboff2019} analyzed this as ``surveillance capitalism'': the extraction and commodification of behavioral data for predictive and manipulative purposes.
In doxonomic terms, platforms represent a new kind of narrative infrastructure in which the exposure matrix $E$ (who sees what) is not determined by editorial judgment or public policy but by engagement-maximizing algorithms whose objective function is revenue, not truth or public welfare.

This makes the conversion problem both easier and harder to study.
Easier, because platform advertising data provides unprecedented visibility into doxonomic spending (Meta's Ad Library, Google's Transparency Report).
Harder, because algorithmic curation operates at a speed and granularity that traditional doxonomic accounting cannot yet match.


\section{Synthesis: What Doxonomics Inherits}
\label{sec:ch3-synthesis}

Each predecessor contributed a piece.
Doxonomics assembles them:

\begin{center}
\begin{tabular}{lp{8cm}}
\toprule
\textbf{Predecessor} & \textbf{Doxonomic Inheritance} \\
\midrule
Sophists & Persuasion is a priced good; $\ecap$ has a market. \\
Hobbes/Mandeville & Anthropological regimes are installed, not discovered. \\
Marx & $\fundingvec$ is endogenous to the wealth distribution. \\
Althusser & $\insteco$ operates with relative autonomy from direct class instruction. \\
Gramsci & Durable power requires $\narr$, not just coercion; hegemony must be continuously renewed. \\
Lippmann & $\Phi$ is controlled by a small number of institutional actors. \\
Dewey & The remedy is democratized communication, not technocratic management. \\
Bernays & Applied doxonomics: optimizing $\Phi$ for commercial and political clients. \\
Ellul & Demand for simplified belief is endogenous to social complexity. \\
Frankfurt School & Entertainment is a doxonomic channel operating at layers 3--5. \\
Foucault & $\credibility{k}$ is itself socially produced (second-order doxonomics). \\
Bourdieu & Symbolic capital, doxa, reproduction: the richest pre-doxonomic vocabulary. \\
Herman/Chomsky & Structural filters produce systematic bias without conspiracy. \\
Said/Fanon & Colonial empires built the largest pre-digital $\narr$ in history. \\
Zuboff & Platforms as algorithmic doxonomic infrastructure with micro-targeting. \\
\bottomrule
\end{tabular}
\end{center}

What none of these predecessors provided, and what doxonomics aims to build, is a \emph{unified formal framework} that integrates accounting, network analysis, causal inference, measurement, and dynamical modeling into a single disciplinary structure.
The tools exist in scattered form across economics, sociology, media studies, and computer science.
The task of this book is to assemble them.


\section{A Timeline}

\begin{center}
\begin{tikzpicture}[
  every node/.style={font=\footnotesize},
  name/.style={font=\footnotesize\bfseries, rotate=45, anchor=north east},
]
  \draw[thick, -Stealth] (0,0) -- (15,0);

  \foreach \x/\t/\d in {
    0.3/Hobbes/1651,
    1.5/Mandeville/1714,
    2.7/Marx/1846,
    3.9/Lippmann/1922,
    5.0/Bernays/1928,
    6.1/Adorno/1944,
    7.3/Ellul/1962,
    8.5/Gramsci/1971,
    9.5/Foucault/1971,
    10.5/Said/1978,
    11.5/Bourdieu/1984,
    12.5/Herman/1988,
    13.7/Zuboff/2019
  } {
    \draw (\x, 0.15) -- (\x, -0.15);
    \node[name] at (\x, -0.25) {\t};
    \node[font=\tiny, anchor=south] at (\x, 0.2) {\d};
  }
\end{tikzpicture}
\end{center}
\vspace{1.5em}


\section{Notes and References}
\label{sec:ch3-notes}

On the sophists and rhetoric, see Plato's \textit{Gorgias} and \textit{Republic}.
\textcite{hobbes1651} and \textcite{mandeville1714} establish the selfish anthropology; on its selective amplification, see \textcite{mcCloskey2006}.
\textcite{marx1846} is the foundational text on ideology as material production; \textcite{althusser1970} develops the institutional apparatus.
\textcite{gramsci1971} develops hegemony and counter-hegemony.
\textcite{lippmann1922} and \textcite{dewey1927} frame the public-opinion debate.
\textcite{bernays1928} codifies propaganda technique; \textcite{ellul1965} extends it to a total social phenomenon.
\textcite{adorno1944} develops the culture industry concept.
\textcite{foucault1971} restructures the analysis of discourse and truth-production.
\textcite{bourdieu1984, bourdieu1991} provide the symbolic-power framework.
\textcite{herman1988} extends these ideas into a testable media model.
\textcite{said1978} and \textcite{fanon1961} analyze colonial knowledge production.
\textcite{zuboff2019} diagnoses platform-mediated surveillance capitalism.
On the institutional history of neoliberalism as a doxonomic project, see \textcite{mirowski2009}, \textcite{mayer2016}, and \textcite{slobodian2018}.
\textcite{stanley2015} provides a recent philosophical analysis of propaganda in democratic societies.
\textcite{hall1980} on encoding/decoding introduced the concept of negotiated readings that resist the dominant code.


\section{Exercises}

\begin{exercise}
For each thinker in this chapter, state the corresponding doxonomic concept from the formal framework of Chapter~\ref{ch:what-is-doxonomics} (e.g., Marx $\to$ endogenous $\fundingvec$; Bernays $\to$ applied optimization of $\Phi$).
\end{exercise}

\begin{exercise}
Bernays staged the ``Torches of Freedom'' campaign (1929) to associate women's smoking with feminist liberation.
Analyze this as a doxonomic operation: identify $\funding{k}$, $\credibility{k}$, $\relevance{k}{i}$, $\reach{k}$, and $\persistence{k}$.
Was the campaign operating at layer 1 (propositional belief) or at deeper layers?
\end{exercise}

\begin{exercise}
Compare Foucault's ``order of discourse'' with the concept of credibility weights $\credibility{k}$.
In what sense does Foucault's analysis go beyond the first-order model?
Write down a second-order model in which $\credibility{k}$ is itself produced by a meta-level doxonomic system, and discuss the implications for measurement.
\end{exercise}

\begin{exercise}
The propaganda model identifies five structural filters.
For each filter, write the corresponding formal constraint on the production function $\Phi$ or the narrative space $\narr$.
Which filters are most affected by the transition from legacy media to platform-based media?
\end{exercise}

\begin{exercise}
Choose a colonial education system (e.g., British India, French West Africa, Belgian Congo).
Reconstruct its belief production chain.
How did it differ from the neoliberal chain described in Example~\ref{ex:neoliberal-chain}?
What institutional features made colonial belief production particularly durable?
\end{exercise}

\begin{exercise}
The chapter presents the Lippmann-Dewey debate as mapping onto a tension between diagnostic and constructive doxonomics.
Is it possible to practice both simultaneously, or does the diagnostic stance (``belief is manufactured'') undermine the constructive stance (``democratic belief formation is possible'')?
Argue your position.
\end{exercise}

\begin{exercise}
Identify a doxonomic insight from a tradition \emph{not} covered in this chapter, for example from Islamic jurisprudence, Confucian statecraft, Buddhist epistemology, or indigenous knowledge systems.
State the insight in formal doxonomic language.
\end{exercise}

\chapter{Core Units of Analysis}
\label{ch:core-units}

\epigraph{All models are wrong, but some are useful.}{George E.\ P.\ Box}

\noindent
Doxonomics requires a precise vocabulary of parts.
Before we can model how beliefs are produced, we need formal definitions of the entities involved: who produces belief, through what institutional machinery, via which channels, with what resources, carrying what content, reaching which audiences, and feeding back into what structures.

This chapter defines seven core units (agents, institutions, channels, budgets, narratives, audiences, and feedback loops) and gives each a mathematical specification.
Together they constitute the component library from which doxonomic systems are assembled.


\section{Agents}
\label{sec:ch4-agents}

\begin{definition}[Doxonomic Agent]
\label{def:agent}
\index{agent}
A \emph{doxonomic agent} is an entity $a \in \agentspace$ characterized by:
\begin{itemize}[nosep]
  \item a \emph{type} $\tau(a) \in \{\text{funder}, \text{institution}, \text{intermediary}, \text{audience member}\}$,
  \item a \emph{resource endowment} $r(a) \in \R_{\geq 0}$,
  \item a \emph{belief state} $\pi_a \in \Delta(\Omega)$ (a probability distribution over propositions),
  \item a \emph{strategic objective} $u_a: \beliefspace \to \R$ (a utility function over population-level belief states).
\end{itemize}
\end{definition}

The four types play structurally different roles:

\textbf{Funders} allocate resources to channels.
They have preferences about what the population believes, preferences that may be ideological (a foundation promoting free-market ideas), commercial (a corporation promoting consumerism), or governmental (a state promoting national identity).
Their strategic problem is: given budget $r(a)$, how to allocate across channels to maximize $u_a(\beliefvec)$.

\textbf{Institutions} process funding into narrative outputs.
They are the factories of belief production: universities, think tanks, media organizations, PR firms, religious organizations, government agencies.
Their output depends not only on funding but on their internal norms, personnel, accumulated credibility, and organizational culture.

\textbf{Intermediaries} amplify, filter, translate, or repackage narratives as they move through the system.
A journalist who writes up a think-tank report, a social media influencer who popularizes an academic finding, a textbook author who translates research into pedagogical material: each is an intermediary who transforms the narrative in the process of transmitting it.

\textbf{Audience members} receive, interpret, and may resist or adopt narratives.
They are not passive.
They bring prior beliefs, identity commitments, lived experience, and cognitive resources to the encounter with any narrative.
The same narrative may persuade one audience segment and provoke backlash in another.

\begin{remark}
\label{rem:agent-dual-roles}
In practice, agents often occupy multiple roles simultaneously.
A university professor is both an institution (producing research) and an intermediary (translating it for students) and an audience member (consuming media).
The type classification is analytical, not existential: it describes the \emph{role} an agent plays in a specific doxonomic transaction, not a permanent identity.
\end{remark}

The key structural asymmetry is that funders optimize over $\beliefspace$ (they have preferences about what the population believes and resources to act on those preferences), while audience members are typically modeled as updating agents without strategic control over the information environment.
This asymmetry is the source of doxonomic power.

\begin{example}[Agent Asymmetry in Climate Discourse]
\label{ex:climate-agents}
In the climate-belief system:
\begin{itemize}[nosep]
  \item \textbf{Funders}: fossil fuel companies (ExxonMobil, Koch Industries) invested hundreds of millions in climate-doubt production \autocite{oreskes2010}.
  \item \textbf{Institutions}: think tanks (Heartland Institute, Global Warming Policy Foundation) processed funding into contrarian research, op-eds, and congressional testimony.
  \item \textbf{Intermediaries}: cable news producers, op-ed editors, and social media algorithms amplified doubt narratives.
  \item \textbf{Audience}: the general public, which by 2010 was significantly more skeptical about climate change than the scientific consensus warranted.
\end{itemize}
The funding asymmetry: the fossil fuel industry's doxonomic budget for climate doubt dwarfed the combined science-communication budgets of all climate research institutions.
This did not make the doubt narrative \emph{true}; it made it \emph{ubiquitous}.
\end{example}


\section{Institutions as Belief-Processing Functions}
\label{sec:ch4-institutions}

Institutions are the central machinery of doxonomic production.
They are not mere conduits; they \emph{transform} inputs.
A think tank does not simply relay its donor's wishes.
It translates them into credentialed, publishable, quotable form, adding the epistemic capital that makes the output persuasive.

\begin{definition}[Institutional Processing Function]
\label{def:institution-function}
\index{institution!processing function}
An institution $\iota \in \insteco$ is a function
\[
  \iota: \R_{\geq 0} \times \narspace_{\text{in}} \to \narspace_{\text{out}} \times [0,1]
\]
mapping a funding level and input narrative material to an output narrative and a credibility multiplier.
The credibility multiplier reflects the institution's accumulated epistemic capital $\ecap(\iota)$.
\end{definition}

The processing function captures three things:
\begin{enumerate}[nosep]
  \item \textbf{Content transformation}: the input (a donor's policy preference) becomes output (a research paper, a textbook chapter, an op-ed). The content may be identical, but the form is different, and form determines doxonomic force.
  \item \textbf{Credibility attachment}: the institutional brand is affixed to the output. ``A study by the Heritage Foundation'' and ``a study by Harvard'' carry different credibility weights, even if the methodology is identical.
  \item \textbf{Audience routing}: different institutions have access to different distribution networks. A think tank's output reaches policymakers and journalists; a university's output reaches students and academics; an advertising agency's output reaches consumers.
\end{enumerate}

\begin{example}[Credibility Differential]
\label{ex:think-tank-credibility}
A think tank $\iota$ receives funding $\funding{k} = \$2\text{M/year}$ and raw policy positions from its donors.
It outputs research papers, op-eds, and expert testimony.
Its credibility multiplier $\credibility{\iota} = 0.7$ reflects its reputation among journalists and policymakers.
A university economics department performing similar work might have $\credibility{\iota'} = 0.9$.
The difference in $\credibility{}$ means the same proposition, backed by the same evidence, has different doxonomic force depending on its institutional origin.

This credibility differential is what makes \emph{credibility laundering}\index{credibility laundering} possible: a low-credibility funder channels money through a high-credibility institution, gaining the credibility premium without producing the underlying quality.
The institution's brand is the product being purchased.
\end{example}

\begin{definition}[Institutional Capacity]
\label{def:capacity}
\index{institutional capacity}
The \emph{capacity} of an institution $\iota$ is the maximum narrative throughput it can produce per unit time:
\[
  \Theta^{\max}(\iota) = \max_{\funding{}} \left\{ \reach{\iota}(\funding{}) \cdot \credibility{\iota} \cdot \persistence{\iota} \right\}
\]
Capacity depends on staffing, infrastructure, accumulated expertise, and institutional reputation.
It cannot be created instantly: building a credible institution takes years or decades, which is why doxonomic systems exhibit strong path dependence.
\end{definition}

\begin{remark}[Institutional Autonomy]
\label{rem:autonomy}
Institutions are not perfectly obedient to their funders.
A university department funded by a corporation may produce research that contradicts the funder's preferences.
A media outlet owned by a billionaire may employ journalists who resist editorial direction.
This relative autonomy is an important feature, not a bug, of the model: it means the production function $\iota$ has its own internal logic that may partially resist the input signal.
The degree of institutional autonomy varies: a PR firm has near-zero autonomy (it produces what the client wants); a tenured university department has considerable autonomy (funding affects what is studied, not necessarily what is concluded).
\end{remark}


\section{Channels}
\label{sec:ch4-channels}

Channels are the transmission media through which institutional outputs reach audiences.
Each channel has characteristic properties that determine its doxonomic profile.

\begin{definition}[Doxonomic Channel]
\label{def:channel}
\index{channel}
A \emph{channel} is a medium of narrative transmission characterized by:
\begin{itemize}[nosep]
  \item \emph{reach} $\reach{k} \in [0,1]$: the fraction of the population exposed per unit time,
  \item \emph{credibility} $\credibility{k} \in [0,1]$: the weight audiences assign to messages from this channel,
  \item \emph{persistence} $\persistence{k} \in \R_{>0}$: the half-life of the narrative in audience memory,
  \item \emph{cost function} $c_k: \R_{\geq 0} \to [0,1]$: the mapping from funding to effective reach (typically concave: diminishing returns to spending).
\end{itemize}
\end{definition}

Different channels occupy different regions of the reach-credibility-persistence space:

\begin{center}
\begin{tikzpicture}[
  every node/.style={font=\small},
  ch/.style={draw, thick, rounded corners, minimum width=2.8cm, minimum height=0.7cm, align=center, fill=doxoblue!8},
]
\matrix[row sep=0.4cm, column sep=0.3cm] {
  \node[ch] {Education}; &
  \node[font=\small, text width=8.5cm] {Low reach/time, high credibility ($\sim0.85$), very high persistence ($\sim$decades). Shapes foundational priors. Most expensive per-person but most durable.}; \\
  \node[ch] {Prestige Media}; &
  \node[font=\small, text width=8.5cm] {Medium reach, high credibility ($\sim0.7$), medium persistence ($\sim$months). Sets the daily agenda for elites and opinion leaders.}; \\
  \node[ch] {Broadcast Media}; &
  \node[font=\small, text width=8.5cm] {High reach, medium credibility ($\sim0.5$), low persistence ($\sim$days). Mass exposure with high turnover.}; \\
  \node[ch] {Social Media}; &
  \node[font=\small, text width=8.5cm] {Very high reach, low credibility ($\sim0.2$), very low persistence ($\sim$hours). Amplifies, fragments, and accelerates.}; \\
  \node[ch] {Think Tanks}; &
  \node[font=\small, text width=8.5cm] {Low direct reach, high credibility in policy circles ($\sim0.75$), medium persistence. Produces expert legitimacy for policy narratives.}; \\
  \node[ch] {Entertainment}; &
  \node[font=\small, text width=8.5cm] {High reach, low perceived persuasive intent, medium persistence. Operates at layers 3--5: normalizes through representation and affect.}; \\
  \node[ch] {Peer Networks}; &
  \node[font=\small, text width=8.5cm] {Variable reach (local), high credibility ($\sim0.8$), high persistence. Most resistant to top-down manipulation but most responsive to social proof.}; \\
};
\end{tikzpicture}
\end{center}

\begin{remark}[The Credibility-Reach Tradeoff]
\label{rem:tradeoff}
There is a general tradeoff between credibility and reach.
Channels with high credibility (education, peer networks) tend to have limited reach per unit time.
Channels with high reach (broadcast, social media) tend to have lower credibility.
This tradeoff shapes doxonomic strategy: a funder seeking deep belief change (common-sense capture) must invest in high-credibility, high-persistence channels; a funder seeking temporary opinion shifts (e.g., before an election) can rely on high-reach, low-persistence channels.
\end{remark}

\begin{definition}[Channel Portfolio]
\label{def:portfolio}
\index{channel portfolio}
A \emph{channel portfolio} for belief $i$ is the set of channels through which funding is directed, together with the allocation across them: $\{(k, \funding{k})\}_{k \in K_i}$ where $K_i \subseteq \{1, \ldots, m\}$.
A diversified portfolio spreads investment across multiple channel types to exploit different reach-credibility-persistence profiles simultaneously.
\end{definition}


\section{Budgets as Resource Allocation Vectors}
\label{sec:ch4-budgets}

\begin{definition}[Doxonomic Budget]
\label{def:budget}
\index{budget}
A \emph{doxonomic budget} for belief $i$ is a vector $\fundingvec_i = (\funding{1}, \ldots, \funding{m}) \in \R^m_{\geq 0}$ specifying the resource allocation across $m$ channels directed at promoting belief $i$.
The total doxonomic expenditure on $i$ is $F_i = \|\fundingvec_i\|_1 = \sum_k \funding{k}$.
\end{definition}

The budget allocation problem is: given a total expenditure $F_i$, how should a funder allocate across channels to maximize $\belief{i}$?
This is a constrained optimization problem whose solution depends on the relative cost-effectiveness of each channel.

\begin{proposition}[Optimal Allocation]
\label{prop:optimal-allocation}
Under the first-order model~\eqref{eq:first-order} with concave cost-reach functions $c_k$, the optimal allocation satisfies the equimarginal condition:
\[
  \credibility{k} \cdot \relevance{k}{i} \cdot c_k'(\funding{k}) \cdot \persistence{k} = \lambda \quad \text{for all } k \text{ with } \funding{k} > 0
\]
where $\lambda$ is the Lagrange multiplier on the budget constraint $\sum_k \funding{k} = F_i$.
\end{proposition}

\begin{proof}
The funder maximizes $\belief{i} = \sum_k \funding{k} \cdot \credibility{k} \cdot \relevance{k}{i} \cdot c_k(\funding{k}) / \funding{k} \cdot \persistence{k}$ subject to $\sum_k \funding{k} = F_i$ and $\funding{k} \geq 0$.
Replacing $\reach{k}$ with $c_k(\funding{k})$ (the cost-to-reach function) and forming the Lagrangian $\mathcal{L} = \sum_k \credibility{k} \cdot \relevance{k}{i} \cdot c_k(\funding{k}) \cdot \persistence{k} - \lambda(\sum_k \funding{k} - F_i)$, the first-order condition for interior solutions is $\credibility{k} \cdot \relevance{k}{i} \cdot c_k'(\funding{k}) \cdot \persistence{k} = \lambda$.
Concavity of $c_k$ ensures the second-order conditions are satisfied.
\end{proof}

This is the doxonomic analogue of the equimarginal principle in economics: at the optimum, the marginal belief return per dollar is equalized across all active channels.
Channels where the marginal return exceeds $\lambda$ should receive more funding; channels where it falls below $\lambda$ should be cut.

\begin{example}[Budget Allocation: Think Tank vs.\ Advertising]
\label{ex:allocation}
A funder with $F = \$5\text{M}$ considers two channels:
\begin{itemize}[nosep]
  \item Think tank: $\credibility{1} = 0.75$, $\persistence{1} = 10$ years, $c_1(f) = 0.02\sqrt{f}$ (low reach, high credibility).
  \item Social media advertising: $\credibility{2} = 0.2$, $\persistence{2} = 0.1$ years, $c_2(f) = 0.5\sqrt{f}$ (high reach, low credibility).
\end{itemize}
The marginal return at funding level $f$ is $\credibility{k} \cdot \persistence{k} \cdot c_k'(f)$.
For the think tank: $0.75 \cdot 10 \cdot 0.01 / \sqrt{f} = 0.075/\sqrt{f}$.
For advertising: $0.2 \cdot 0.1 \cdot 0.25/\sqrt{f} = 0.005/\sqrt{f}$.
The think tank's marginal return is 15 times higher.
At the optimum, the funder should invest almost entirely in the think tank, a non-obvious result that reflects the enormous premium of credibility $\times$ persistence over raw reach.
\end{example}

\begin{remark}[Why Doxonomic Budgets Are Hard to Measure]
\label{rem:budget-opacity}
In practice, doxonomic budgets are difficult to estimate because:
\begin{itemize}[nosep]
  \item much funding flows through intermediaries (foundations, donor-advised funds) that obscure the ultimate source,
  \item institutional overhead blurs the distinction between doxonomic spending and general operations,
  \item some channels (education, entertainment) have doxonomic effects that are side effects of their primary function, not their explicit purpose,
  \item non-monetary resources (volunteer labor, social media sharing, cultural prestige) contribute to belief production but are not captured in financial accounting.
\end{itemize}
Chapter~\ref{ch:belief-flow-accounting} develops methods for addressing these measurement challenges.
\end{remark}


\section{Narratives}
\label{sec:ch4-narratives}

Narratives are the \emph{content} of doxonomic systems: what institutions produce, channels transmit, and audiences process.

\begin{definition}[Narrative]
\label{def:narrative}
\index{narrative}
A \emph{narrative} $n \in \narspace$ is a structured proposition or set of propositions equipped with:
\begin{enumerate}[nosep]
  \item a \emph{causal claim} (``$X$ because $Y$''): poverty exists because people don't work hard enough; markets crashed because of government interference; crime rises because of immigration,
  \item a \emph{normative framing} (``this is good/bad/natural/inevitable''): inequality is natural; growth is progress; austerity is responsible,
  \item a \emph{subject position} (who the audience is invited to identify with): the hardworking taxpayer, the innovative entrepreneur, the responsible citizen.
\end{enumerate}
\end{definition}

The same factual content can be embedded in radically different narratives.
``Unemployment is at 8\%'' becomes:
\begin{itemize}[nosep]
  \item $n_1$: ``The economy is recovering: unemployment has fallen from 10\% to 8\%.'' (Causal: recovery. Frame: progress. Subject: optimistic citizen.)
  \item $n_2$: ``Eight million people cannot find work, and the government has failed them.'' (Causal: policy failure. Frame: injustice. Subject: abandoned worker.)
  \item $n_3$: ``Welfare dependency keeps people from seeking employment.'' (Causal: moral failure. Frame: personal responsibility. Subject: taxpayer.)
\end{itemize}
Same statistic, three different doxonomic effects.
The narrative, not the fact, is the doxonomic product.

\begin{definition}[Narrative Distance]
\label{def:narrative-distance}
\index{narrative distance}
Given a vector-space representation of narratives (e.g., via document embeddings), the \emph{narrative distance} between two narratives $n_1, n_2 \in \narspace$ is
\[
  d(n_1, n_2) = 1 - \frac{\langle \bm{v}_{n_1}, \bm{v}_{n_2} \rangle}{\|\bm{v}_{n_1}\| \cdot \|\bm{v}_{n_2}\|}
\]
where $\bm{v}_n$ is the embedding vector of narrative $n$.
Large narrative distance between the dominant frame and an alternative frame indicates that the alternative is not merely a different opinion but a \emph{different way of seeing}, harder to reach by incremental argument, requiring more doxonomic investment to install.
\end{definition}

\begin{remark}[Narratives vs.\ Propositions]
\label{rem:narrative-proposition}
A narrative is more than a proposition.
The proposition ``inequality is increasing'' has a truth value.
The narrative ``inequality is increasing because the lazy are being rewarded'' has a truth value, a moral judgment, a causal attribution, and a subject position.
Doxonomic systems rarely produce naked propositions; they produce narratives, because narratives engage identity, emotion, and action in ways that propositions alone cannot.
This is why doxonomic production operates at multiple layers of the ontology simultaneously (Principle~\ref{prin:layered-ontology}).
\end{remark}


\section{Audiences and Prior Distributions}
\label{sec:ch4-audiences}

Audiences are not homogeneous.
A population consists of subgroups $\{G_1, \ldots, G_L\}$ with different prior belief distributions $\mu_{G_l}$, different media exposure patterns, different resistance profiles, and different identity commitments.
Effective doxonomic strategy requires targeting: knowing which audiences are persuadable and through which channels.

\begin{definition}[Audience Heterogeneity]
\label{def:audience-heterogeneity}
\index{audience}
The population-level belief $\belief{i}$ is a mixture:
\[
  \belief{i} = \sum_{l=1}^{L} w_l \cdot \belief{i}^{(l)}
\]
where $w_l$ is the population share of group $l$ and $\belief{i}^{(l)}$ is the average credence in group $l$.
A doxonomic system may target some groups more than others.
\end{definition}

\begin{definition}[Audience Resistance]
\label{def:resistance}
\index{audience resistance}
The \emph{resistance} of audience group $G_l$ to narrative $n$ is a function $\rho_l(n) \in [0,1]$ reflecting the degree to which the group's prior beliefs, identity commitments, lived experience, and critical resources reduce the narrative's doxonomic effect.
The effective belief change in group $l$ from exposure to $n$ through channel $k$ is
\[
  \Delta\belief{i}^{(l)} = (1 - \rho_l(n)) \cdot \credibility{k} \cdot \relevance{k}{i} \cdot \persistence{k}
\]
High resistance ($\rho_l \to 1$) means the group is effectively immune to the narrative; low resistance ($\rho_l \to 0$) means full susceptibility.
\end{definition}

Resistance is not random.
It is shaped by:
\begin{itemize}[nosep]
  \item \textbf{Lived experience}: a worker who has been laid off despite ``working hard'' is more resistant to meritocratic narratives than a comfortably employed professional.
  \item \textbf{Counter-institutions}: groups with access to alternative information sources (unions, community organizations, alternative media) have higher resistance.
  \item \textbf{Identity}: narratives that threaten core identity commitments provoke backlash rather than persuasion \autocite{kahan2012}.
  \item \textbf{Prior exposure}: groups that have been ``inoculated'' against a narrative (Chapter~\ref{ch:counter-doxonomic}) show elevated resistance to subsequent exposure.
\end{itemize}

\begin{example}[Differential Resistance to Meritocratic Narratives]
\label{ex:resistance}
Survey data consistently shows that belief in meritocracy is negatively correlated with direct experience of structural barriers.
Black Americans score lower on meritocratic belief scales than white Americans; women in male-dominated workplaces are more skeptical of ``effort determines success'' than men in the same workplaces; first-generation college students are less convinced of equal opportunity than their legacy peers.
In doxonomic terms, lived experience of structural disadvantage raises $\rho_l$ for meritocratic narratives, a form of experiential resistance that no amount of doxonomic investment can fully overcome.
\end{example}


\section{Feedback Loops}
\label{sec:ch4-feedback}

The belief production chain is not unidirectional.
Beliefs, once internalized, alter behavior, which alters institutions, which alters funding.
This feedback structure is what makes doxonomic systems self-sustaining, and what makes them so difficult to disrupt.

\begin{definition}[Doxonomic Feedback]
\label{def:feedback}
\index{feedback loop}
A \emph{doxonomic feedback loop} is a dynamical coupling in which the belief state $\beliefvec(t)$ influences the funding vector $\fundingvec(t+1)$ through behavioral and institutional responses:
\[
  \fundingvec(t+1) = g(\beliefvec(t), \fundingvec(t))
\]
for some function $g$.
The system is \emph{self-reinforcing} if $\partial g_k / \partial \belief{i} > 0$ for the belief $i$ that channel $k$ promotes.
The system is \emph{self-undermining} if $\partial g_k / \partial \belief{i} < 0$ (success reduces future investment).
\end{definition}

\begin{example}[Self-Reinforcing: The Tax Cut Cycle]
\label{ex:tax-cycle}
\begin{enumerate}[nosep]
  \item Think tanks promote the belief that tax cuts stimulate growth ($\belief{T}$ increases).
  \item Population supports tax-cutting politicians (behavioral consequence).
  \item Tax cuts are enacted; wealthy beneficiaries gain additional disposable income.
  \item Beneficiaries donate to the same think tanks (funding increases).
  \item Think tanks produce more tax-cut advocacy ($\belief{T}$ increases further).
\end{enumerate}
This is a self-reinforcing loop: $\partial g_k / \partial \belief{T} > 0$.
Each iteration strengthens the cycle.
\end{example}

\begin{example}[Self-Undermining: The Backlash Effect]
\label{ex:backlash}
\begin{enumerate}[nosep]
  \item A corporation promotes the belief that its product is environmentally friendly (``greenwashing'').
  \item The population initially accepts the claim ($\belief{G}$ increases).
  \item Investigative journalists or activists expose the deception.
  \item Public trust in the corporation's claims collapses ($\credibility{k}$ drops).
  \item Further spending on the same message becomes counterproductive.
\end{enumerate}
This is a self-undermining loop: success (increased scrutiny) reduces future effectiveness.
\end{example}

Self-reinforcing feedback is the mechanism by which temporary doxonomic investments become permanent features of a society's common sense.
Self-undermining feedback is the mechanism by which doxonomic systems can fail, a topic developed in Chapter~\ref{ch:how-belief-regimes-fracture}.

\begin{definition}[Doxonomic Stability]
\label{def:stability}
\index{doxonomic stability}
A doxonomic system is \emph{locally stable} at belief state $\beliefvec^*$ if small perturbations (temporary counter-narratives, one-time contrary evidence) are absorbed by the feedback structure and the system returns to $\beliefvec^*$.
Formally, $\beliefvec^*$ is a stable fixed point of the composed dynamical system:
\[
  \beliefvec^* = \Psi(\Phi(g(\beliefvec^*, \fundingvec^*)), \beliefvec^*)
\]
where $\Phi$ is the production function, $\Psi$ is the internalization function, and $g$ is the feedback function.
\end{definition}


\section{Assembling the Units: A Toy Model}
\label{sec:ch4-toy}

To see how the units work together, consider a minimal doxonomic system with one funder, one institution, one channel, and a homogeneous audience.

\begin{model}[Minimal Doxonomic System]
\label{mod:minimal}
\begin{itemize}[nosep]
  \item Funder: budget $F$, objective $\max \belief{}$.
  \item Institution: processing function $\iota(F) = (n, \credibility{})$ where $\credibility{} = 0.8$.
  \item Channel: reach $c(F) = 1 - e^{-F/F_0}$ (saturating), persistence $T = 5$ years.
  \item Audience: $N = 1000$ agents, prior $\pi_j^{(0)} = 0.3$, Bayesian updating.
  \item Feedback: $F(t+1) = F_0 + \beta \belief{}(t)$ (self-reinforcing with rate $\beta$).
\end{itemize}
The dynamics are:
\[
  \belief{}(t+1) = \belief{}(t) + \alpha \cdot \credibility{} \cdot c(F(t)) \cdot T \cdot (1 - \belief{}(t)) - \delta \cdot \belief{}(t)
\]
where $\alpha$ is the persuasion rate and $\delta$ is natural decay.

This system has at most two stable fixed points: a low-belief equilibrium (the narrative fails to take hold) and a high-belief equilibrium (common-sense capture).
The basin of attraction depends on the initial funding $F_0$ and the feedback strength $\beta$.
\end{model}

This toy model is developed into a full dynamical framework in Chapters~\ref{ch:formal-models} and~\ref{ch:belief-production-chain}.


\section{Notes and References}
\label{sec:ch4-notes}

Agent-based models of belief formation draw on \textcite{jackson2008} and \textcite{granovetter1978}.
The institutional processing framework extends \textcite{north1990} on institutions as ``rules of the game.''
On institutional autonomy and the relative independence of cultural production, see \textcite{bourdieu1984}.
Optimal allocation of persuasion resources connects to \textcite{downs1957} and \textcite{tullock1967}.
Audience heterogeneity models are treated in \textcite{manski1993}; on identity-protective cognition, see \textcite{kahan2012}.
Feedback loops in social systems are central to \textcite{schelling1978}.
The climate-doubt example draws on \textcite{oreskes2010}.
On the economics of attention as a scarce resource, see \textcite{kahneman2011}.
Credibility laundering through institutional intermediaries is documented in \textcite{mayer2016}.
Narrative as a structured object with causal, normative, and identity components draws on framing theory; see \textcite{stanley2015} for a philosophical treatment.


\section{Exercises}

\begin{exercise}
Classify the following as agents, institutions, channels, or narratives: (a)~the Koch Foundation, (b)~the \textit{New York Times}, (c)~Twitter/X, (d)~``free markets create prosperity,'' (e)~a university economics department, (f)~a YouTube algorithm, (g)~a public school curriculum, (h)~``hard work is its own reward.''
Note cases where an entity plays multiple roles.
\end{exercise}

\begin{exercise}
\label{ex:ch4-allocation}
Compute the optimal budget allocation for a funder with $F_i = \$10\text{M}$ across two channels with $\credibility{1} = 0.9$, $\credibility{2} = 0.5$, and cost-reach functions $c_1(f) = \sqrt{f/10^6}$, $c_2(f) = 2\sqrt{f/10^6}$.
Assume $\relevance{k}{i} = \persistence{k} = 1$ for both.
How does the allocation change if $\persistence{1} = 20$ and $\persistence{2} = 0.5$?
\end{exercise}

\begin{exercise}
Model a self-reinforcing feedback loop in which belief in meritocracy increases tolerance for inequality, which increases political donations from wealthy actors, which increases funding for meritocratic narratives.
Write the system as a pair of coupled difference equations and identify the conditions for a stable high-meritocracy equilibrium.
\end{exercise}

\begin{exercise}
Consider a population with two groups: $G_1$ (30\% of population, high media exposure, low prior credence in $p$, low resistance $\rho_1 = 0.2$) and $G_2$ (70\%, low media exposure, high prior credence in $p$, high resistance $\rho_2 = 0.8$).
Under what conditions does a doxonomic investment targeting $G_1$ change the population-level $\belief{i}$ more than one targeting $G_2$?
\end{exercise}

\begin{exercise}
For the minimal doxonomic system (Model~\ref{mod:minimal}), find the fixed points of the dynamics as a function of $\alpha$, $\delta$, $\beta$, $F_0$, and $\credibility{}$.
Under what conditions does only one fixed point exist (inevitable capture or inevitable failure)?
Under what conditions do two stable fixed points coexist?
\end{exercise}

\begin{exercise}
Design a doxonomic system from scratch: choose a belief $p$, specify funders, institutions, channels, narratives, target audiences, and feedback structure.
Estimate plausible values for $\credibility{k}$, $\reach{k}$, $\persistence{k}$, and $\rho_l$.
Compute the expected $\belief{i}$ using the first-order model.
Then identify the weakest link in your system, the point where counter-doxonomic intervention would be most effective.
\end{exercise}

\begin{exercise}
The chapter notes that entertainment operates at layers 3--5 of the ontology.
Choose a popular television series or film franchise and analyze it as a doxonomic channel.
What anthropological regime does it normalize?
What subject positions does it offer?
Estimate its $\reach{}$, $\credibility{}$ (noting that entertainment credibility works differently from news credibility), and $\persistence{}$.
\end{exercise}

\chapter{The Belief Production Chain}
\label{ch:belief-production-chain}

\epigraph{It is difficult to get a man to understand something when his salary depends upon his not understanding it.}{Upton Sinclair}

\noindent
The previous chapters introduced the components.
This chapter assembles them into the central object of doxonomic theory: the \emph{belief production chain}\index{belief production chain}, formalized as a composition of operators.

The belief production chain is to doxonomics what the supply chain is to manufacturing, what the food chain is to ecology, and what the epidemiological chain of infection is to public health.
It is the structured sequence through which raw resources become durable public belief.
Understanding its stages, their interconnections, and their failure modes is the core analytical task of the field.


\section{The Chain as Operator Composition}
\label{sec:ch5-composition}

\begin{definition}[Belief Production Chain]
\label{def:bpc-formal}
The belief production chain is a composition of six operators:
\[
  \bpc = \mathcal{R} \circ \mathcal{B} \circ \mathcal{I} \circ \mathcal{C} \circ \mathcal{P} \circ \mathcal{F}
\]
where:
\begin{enumerate}[nosep]
  \item $\mathcal{F}: \R^m_{\geq 0} \to \R^m_{\geq 0}$ is the \emph{funding allocation operator},
  \item $\mathcal{P}: \R^m_{\geq 0} \to \narspace$ is the \emph{institutional processing operator},
  \item $\mathcal{C}: \narspace \to \narspace \times [0,1]^N$ is the \emph{circulation operator} (narrative plus exposure vector),
  \item $\mathcal{I}: \narspace \times [0,1]^N \times \beliefspace \to \beliefspace$ is the \emph{internalization operator},
  \item $\mathcal{B}: \beliefspace \to \mathcal{X}$ is the \emph{behavioral consequence operator} (mapping beliefs to actions),
  \item $\mathcal{R}: \mathcal{X} \to \R^m_{\geq 0}$ is the \emph{reproduction operator} (mapping behaviors back to funding).
\end{enumerate}
\end{definition}

The full chain maps an initial funding allocation to a new funding allocation via beliefs and behavior.
When $\bpc$ has a fixed point, the doxonomic system is in equilibrium: belief reproduces the conditions of its own production.

\begin{remark}
Each operator is itself a complex process with internal structure.
The production function $\mathcal{P}$ is not a single equation but an ensemble of institutions each with its own logic (Chapter~\ref{ch:core-units}).
The internalization operator $\mathcal{I}$ involves Bayesian updating, motivated reasoning, identity protection, and social conformity simultaneously (Chapter~\ref{ch:ontology-belief}).
The composition notation $\bpc = \mathcal{R} \circ \cdots \circ \mathcal{F}$ is a deliberate abstraction: it allows us to reason about the chain's global properties (fixed points, stability, multiplier effects) without specifying the internal structure of each stage.
The internal structure is developed in subsequent chapters.
\end{remark}


\section{Stage 1: Funding Allocation}
\label{sec:ch5-funding}

The funding allocation operator $\mathcal{F}$ distributes total resources $F$ across channels according to a strategy $\sigma$:
\[
  \mathcal{F}(F, \sigma) = (\funding{1}, \ldots, \funding{m}) \quad \text{subject to} \quad \sum_k \funding{k} = F
\]

The strategy $\sigma$ may be the solution to the optimization problem in Proposition~\ref{prop:optimal-allocation}, or it may reflect institutional inertia, path dependence, or donor ideology.
In practice, most doxonomic funding is not optimally allocated; it follows organizational routines.

\begin{example}[Funding Strategies]
\label{ex:funding-strategies}
Consider three funders, each with \$10M/year, pursuing the same belief:
\begin{itemize}[nosep]
  \item \textbf{Optimizer}: allocates using the equimarginal principle, concentrating on high-$\credibility{} \times \persistence{}$ channels. Expected belief impact: high.
  \item \textbf{Inertia funder}: allocates the same proportions as last year, regardless of changed conditions. Common in large foundations with entrenched grantmaking patterns.
  \item \textbf{Ideological funder}: allocates based on personal conviction (``I believe in education'') rather than cost-effectiveness analysis. May over-invest in a favored channel and under-invest in others.
\end{itemize}
The gap between the optimizer and the other two is the \emph{allocative inefficiency} of the doxonomic system.
Real doxonomic systems are typically far from optimal, which, depending on one's normative position, may be a relief.
\end{example}

\begin{remark}[The Dark Budget]
\label{rem:dark-budget}
Not all doxonomic funding is visible.
``Dark money''\index{dark money} (funding channeled through intermediaries that do not disclose their donors) is a significant fraction of total doxonomic investment in some countries.
\textcite{mayer2016} estimates that hundreds of millions of dollars flow annually through donor-advised funds and pass-through foundations in the United States alone, with the ultimate source legally obscured.
The funding allocation operator $\mathcal{F}$ is thus partially unobservable, complicating the doxonomic audit.
Chapter~\ref{ch:belief-flow-accounting} develops methods for inferring hidden funding from observable outputs.
\end{remark}


\section{Stage 2: Institutional Processing}
\label{sec:ch5-processing}

Each institution $\iota_k$ transforms funding and raw material into credentialed narrative output:
\[
  \mathcal{P}_k(\funding{k}) = (n_k, \credibility{k})
\]
where $n_k \in \narspace$ is the narrative output and $\credibility{k} \in [0,1]$ is its credibility weight.

The processing stage is where raw money becomes persuasive content.
A donor's preference for deregulation enters Heritage Foundation as a grant; it exits as a policy brief with citations, data tables, and expert endorsement.
The \emph{content} may not change much, but the \emph{form} is transformed from private preference to public-facing knowledge claim.

\begin{proposition}[Credibility Premium]
\label{prop:credibility-premium}
If two channels produce the same narrative $n$ but with credibility weights $\credibility{1} > \credibility{2}$, then the doxonomic effect satisfies
\[
  \Delta\belief{i}(\text{channel } 1) = \frac{\credibility{1}}{\credibility{2}} \cdot \Delta\belief{i}(\text{channel } 2)
\]
under the first-order model.
High-credibility institutions thus provide a \emph{credibility premium}\index{credibility premium}: the same narrative, processed through a prestigious institution, has proportionally greater doxonomic force.
\end{proposition}

\begin{example}[The Credibility Premium in Practice]
\label{ex:credibility-practice}
In 2010, Carmen Reinhart and Kenneth Rogoff published a paper claiming that government debt exceeding 90\% of GDP leads to sharply lower growth.
The paper was processed through high-credibility channels: Harvard affiliation ($\credibility{} \approx 0.9$), NBER working paper series, prestigious journal publication.
It was cited in congressional testimony, European Commission policy documents, and IMF reports.
It justified austerity programs affecting hundreds of millions of people.

In 2013, a graduate student discovered a spreadsheet error that materially changed the results.
The corrected analysis showed no sharp threshold at 90\%.
But by then, the narrative had been internalized: three years of high-credibility circulation had made ``90\% debt threshold'' into policy common sense.
The \emph{correction} circulated through lower-credibility channels (a working paper from UMass Amherst) and took years to propagate.

This episode illustrates the credibility premium's asymmetry: high-credibility errors propagate faster and stick longer than lower-credibility corrections.
\end{example}

\begin{definition}[Processing Loss]
\label{def:processing-loss}
\index{processing loss}
The \emph{processing loss} of an institution is the difference between the funder's intended message and the institution's actual output.
Formally, if the funder intends narrative $n_{\text{in}}$ and the institution produces $n_{\text{out}}$, the loss is $d(n_{\text{in}}, n_{\text{out}})$ in the narrative distance metric (Definition~\ref{def:narrative-distance}).
Institutions with high autonomy (Remark~\ref{rem:autonomy}) have higher processing loss; institutions with low autonomy (PR firms) have near-zero processing loss.
Processing loss is not necessarily undesirable from a social perspective: it is one mechanism by which institutional norms can resist funder capture.
\end{definition}


\section{Stage 3: Circulation}
\label{sec:ch5-circulation}

The circulation operator models how narratives reach populations.
Let $e_{jk} \in \{0,1\}$ indicate whether individual $j$ is exposed to channel $k$.
The exposure matrix $E = (e_{jk})_{N \times m}$ determines who hears what.

\begin{definition}[Exposure Distribution]
\label{def:exposure}
\index{exposure distribution}
The \emph{exposure distribution} for belief $i$ is the vector $\bm{v}_i \in [0,1]^N$ where
\[
  v_{i,j} = 1 - \prod_{k=1}^{m} (1 - e_{jk} \cdot \relevance{k}{i})
\]
is the probability that individual $j$ encounters at least one narrative relevant to belief $i$.
\end{definition}

The exposure matrix $E$ is the most consequential object in the circulation stage.
It determines the information environment of every individual. In the platform era, it is increasingly determined by algorithms rather than by editorial judgment or audience choice.

\begin{definition}[Exposure Asymmetry]
\label{def:exposure-asymmetry}
\index{exposure asymmetry}
The \emph{exposure asymmetry} for belief $i$ and its negation $\neg i$ is
\[
  A_j(i) = v_{i,j} - v_{\neg i, j}
\]
measuring the degree to which individual $j$'s information environment is skewed toward $i$.
If $A_j(i) > 0$ for most $j$, the population is in an asymmetric information environment favoring $i$, the precondition for Bayesian divergence (Proposition~\ref{prop:unequal-updating}).
\end{definition}

\begin{remark}[Algorithmic Circulation]
\label{rem:algorithmic}
In pre-digital media systems, the exposure matrix $E$ was determined by geography (which newspapers were available), economics (who could afford a subscription), and editorial decisions (what was published).
These factors produced systematic asymmetries, but they were relatively stable and partially observable.

In platform-mediated circulation, $E$ is determined by engagement-maximizing algorithms that personalize exposure at the individual level.
The matrix is no longer stable, observable, or uniform across the population.
Each user receives a different information environment, optimized not for truth or public welfare but for attention capture.
This makes the exposure asymmetry $A_j(i)$ both larger (more personalized skew) and harder to measure (proprietary algorithms).
Chapter~\ref{ch:text-media} discusses methods for partially reconstructing algorithmic exposure.
\end{remark}


\section{Stage 4: Internalization}
\label{sec:ch5-internalization}

Internalization is the process by which exposure becomes credence.
It is the stage where the doxonomic system meets the individual mind, and where the simplifying assumptions of formal models are most strained.

\subsection{The Bayesian Baseline}

As a first approximation, we model internalization as Bayesian updating:
\begin{equation}
\label{eq:bayes-update}
  \pi_j^{(t+1)}(p_i) = \frac{\pi_j^{(t)}(p_i) \cdot L_j(e_j \mid p_i)}{\pi_j^{(t)}(p_i) \cdot L_j(e_j \mid p_i) + (1 - \pi_j^{(t)}(p_i)) \cdot L_j(e_j \mid \neg p_i)}
\end{equation}
where $L_j$ is the likelihood function shaped by $j$'s exposure.
When exposure is heavily skewed (when $j$ sees many more arguments for $p_i$ than against), the posterior converges toward $p_i$ regardless of the prior.

\begin{conjecture}[Asymptotic Capture]
\label{conj:asymptotic-capture}
\index{asymptotic capture}
Let agent $j$ receive a sequence of conditionally independent signals $(s_1, s_2, \ldots)$ where each $s_t$ supports $p_i$ with probability $q > 1/2$, drawn from a well-specified likelihood model.
Assume the agent updates via Bayes' rule, does not discount sources, and does not revise the model itself.
Then $\pi_j^{(t)}(p_i) \to 1$ almost surely as $t \to \infty$, regardless of the prior $\pi_j^{(0)}(p_i) > 0$.
\end{conjecture}

\begin{remark}
The assumptions are strong and worth stating explicitly: conditional independence, well-specified likelihoods, no source discounting, no model revision, no endogenous stopping.
In practice, agents discount unfamiliar sources, update models, experience motivated reasoning, and face correlated rather than independent signals.
The conjecture identifies an idealized mechanism (asymmetric exposure can overwhelm priors), not a universal law.
Repetition alone does not mathematically guarantee capture; it does so only under conditions that doxonomic systems are designed to approximate.
\end{remark}

\subsection{Beyond Bayes: How Real People Internalize}

The Bayesian model is a useful benchmark, but real internalization involves several additional mechanisms:

\begin{enumerate}[nosep]
  \item \textbf{Motivated reasoning}\index{motivated reasoning}: agents update asymmetrically, weighting evidence that confirms identity-consistent beliefs more heavily than evidence that threatens them \autocite{kahan2012}.
  In doxonomic terms, motivated reasoning amplifies the effect of identity-aligned narratives and attenuates the effect of identity-threatening ones.

  \item \textbf{Source credibility weighting}: agents do not treat all sources equally.
  A message from a trusted in-group member has higher effective $\credibility{}$ than the same message from an out-group source, regardless of content.
  This means the credibility weight is not a property of the channel alone but of the \emph{channel $\times$ audience} interaction.

  \item \textbf{Social proof}\index{social proof}: agents use perceived peer behavior as evidence.
  If ``everyone seems to believe $p$,'' this is treated as independent evidence for $p$, even when it is not independent but is itself the product of shared doxonomic exposure \autocite{cialdini2006}.
  Social proof creates a positive feedback loop: doxonomic production increases expressed belief, which increases perceived consensus, which increases genuine belief.

  \item \textbf{Repetition and mere exposure}: repeated exposure to a claim increases its perceived truth, independent of evidential content, the ``illusory truth effect'' \autocite{kahneman2011}.
  This is the most direct mechanism of doxonomic production: sheer volume of exposure substitutes for evidence.

  \item \textbf{Narrative coherence}: agents are more likely to internalize a proposition that fits into an existing narrative structure (Chapter~\ref{ch:ontology-belief}, Definition~\ref{def:ideology}).
  An ideologically bundled proposition benefits from the coherence of the surrounding ideology: accepting one element lowers the threshold for accepting others.

  \item \textbf{Non-propositional internalization}: at layers 3--5 of the ontology (Principle~\ref{prin:layered-ontology}), internalization does not involve conscious belief updating at all.
  Practical schemas are shaped by prolonged immersion, affective orientations by repeated association, identity commitments by social belonging.
  These processes are slower than propositional updating but more durable.
\end{enumerate}

\begin{model}[Augmented Internalization]
\label{mod:augmented-internalization}
\index{internalization!augmented model}
Combining the Bayesian baseline with the mechanisms above, the effective update rule is:
\begin{equation}
\label{eq:augmented-update}
  \pi_j^{(t+1)}(p_i) = \pi_j^{(t)}(p_i) + \underbrace{\alpha_{\text{Bayes}} \cdot \Delta_{\text{Bayes}}}_{\text{evidence}} + \underbrace{\alpha_{\text{social}} \cdot (\hat{\mu}_j^{(t)} - \pi_j^{(t)})}_{\text{social proof}} + \underbrace{\alpha_{\text{rep}} \cdot v_{i,j}}_{\text{repetition}} - \underbrace{\rho_j(n_i) \cdot |\Delta|}_{\text{resistance}}
\end{equation}
where $\Delta_{\text{Bayes}}$ is the Bayesian update, $\hat{\mu}_j^{(t)}$ is the perceived social norm, $v_{i,j}$ is exposure intensity, and $\rho_j$ is resistance.
The relative weights $\alpha$ determine whether internalization is primarily evidence-driven, socially driven, or repetition-driven.
In most real-world settings, $\alpha_{\text{social}}$ and $\alpha_{\text{rep}}$ dominate $\alpha_{\text{Bayes}}$.
\end{model}


\section{Stage 5: Behavioral Consequence}
\label{sec:ch5-behavior}

Internalized beliefs alter behavior.
Let $x_j \in \mathcal{X}$ be the action of agent $j$.
We model behavior as a function of beliefs:
\[
  x_j = \arg\max_{x \in \mathcal{X}} \, u_j(x, \pi_j)
\]

The behavioral consequences of doxonomic production are diverse:
\begin{itemize}[nosep]
  \item \textbf{Voting}: belief in meritocracy reduces support for redistributive policy.
  \item \textbf{Consumption}: belief in consumer identity drives purchasing decisions.
  \item \textbf{Cooperation}: belief in selfish human nature reduces cooperative behavior in public goods settings \autocite{fischbacher2001}.
  \item \textbf{Labor compliance}: belief in the moral value of work increases tolerance for poor working conditions.
  \item \textbf{Tolerance of inequality}: belief that wealth reflects merit reduces demand for redistribution \autocite{piketty2014}.
  \item \textbf{Resistance}: belief that ``there is no alternative'' reduces engagement with counter-narratives.
\end{itemize}

\begin{definition}[Self-Fulfilling Belief]
\label{def:self-fulfilling}
\index{self-fulfilling belief}
A belief $p$ is \emph{self-fulfilling} if the behavioral consequences of holding $p$ generate evidence that appears to confirm $p$.
Formally, $p$ is self-fulfilling if the mapping $\pi \mapsto \text{evidence}(\mathcal{B}(\pi))$ is such that $\text{evidence}(\mathcal{B}(\pi))$ is more consistent with $p$ when $\pi(p)$ is high.
\end{definition}

Self-fulfilling beliefs are the most dangerous doxonomic products because they generate their own evidence.
The belief that ``people are selfish'' leads to institutions designed around selfishness (incentive contracts, competitive structures, monitoring systems), which elicit selfish behavior, which appears to confirm the belief.
Breaking a self-fulfilling cycle requires intervening at the \emph{institutional design} level, not just at the narrative level.


\section{Stage 6: Reproduction}
\label{sec:ch5-reproduction}

The final stage closes the loop.
Behavioral patterns generate institutional demands, electoral outcomes, market structures, and philanthropic priorities that feed back into the funding vector.

\begin{definition}[Doxonomic Equilibrium]
\label{def:doxonomic-equilibrium}
\index{doxonomic equilibrium}
A \emph{doxonomic equilibrium} is a fixed point of the belief production chain:
\[
  \fundingvec^* = \bpc(\fundingvec^*)
\]
At equilibrium, the beliefs produced by the system generate exactly the behaviors and institutions that sustain those beliefs.
\end{definition}

\begin{theorem}[Existence of Equilibrium]
\label{thm:existence}
If each operator in the chain is continuous and the composed map $\bpc: K \to K$ for a compact convex set $K \subset \R^m_{\geq 0}$, then at least one doxonomic equilibrium exists.
\end{theorem}

\begin{proof}
By Brouwer's fixed-point theorem, every continuous map from a compact convex set to itself has a fixed point.
The continuity assumption requires that small changes in funding produce small changes in narrative output, exposure, belief, behavior, and reproduction, reasonable for aggregate quantities, though potentially violated at phase transitions.
The compactness assumption requires bounded funding ($K \subset [0, F_{\max}]^m$), which is physically realistic.
\end{proof}

\begin{conjecture}[Multiplicity of Equilibria]
\label{conj:multiplicity}
For most doxonomic systems with nonlinear internalization (thresholds, social proof, self-fulfilling dynamics), multiple equilibria exist.
A society may be locked into a ``selfish human nature'' equilibrium or a ``cooperative human nature'' equilibrium, each self-sustaining.
Transitions between equilibria require coordinated disruption of the belief production chain at multiple stages simultaneously.
\end{conjecture}

\begin{remark}[Equilibrium Selection]
\label{rem:selection}
When multiple equilibria exist, the question becomes: which one does a society end up in?
This is the \emph{equilibrium selection problem}\index{equilibrium selection}.
History, path dependence, and the timing of doxonomic investments determine selection.
A society that happened to receive large early investments in selfish-anthropology institutions (e.g., through the Mont P\`{e}lerin network in the 1950s--1970s) may be locked into that equilibrium, even if a cooperative equilibrium would be stable too, had the initial investment gone the other way.
This is the doxonomic version of ``history matters.''
\end{remark}


\section{Generational Reproduction as a Fixed-Point Problem}
\label{sec:ch5-generational}

The most powerful form of reproduction is generational.
Each generation's beliefs are partly determined by the institutional environment they grow up in, which was built by the previous generation's beliefs.

\begin{model}[Generational Fixed Point]
\label{mod:generational}
\index{generational reproduction}
Let $\beliefvec_t$ be the belief distribution of generation $t$ and $\insteco_t$ the institutional ecology inherited from the previous generation.
The dynamics are:
\begin{align}
  \insteco_{t+1} &= h(\beliefvec_t, \insteco_t) \quad \text{(institutions reflect current beliefs + inertia)} \\
  \beliefvec_{t+1} &= \Psi(\Phi(\insteco_{t+1}), \beliefvec_0^{(t+1)}) \quad \text{(new generation's beliefs shaped by inherited institutions)}
\end{align}
where $\beliefvec_0^{(t+1)}$ is the prior distribution of the incoming generation (partly shaped by family, partly by innate tendencies).
A \emph{generational equilibrium} satisfies $\beliefvec^* = \Psi(\Phi(h(\beliefvec^*, \insteco^*)), \beliefvec_0)$: the beliefs that institutions transmit to the new generation are the same beliefs that sustain those institutions.
\end{model}

Generational reproduction is why doxonomic equilibria can persist for centuries.
The belief that ``monarchy is divinely ordained'' persisted for a thousand years not because each generation was independently persuaded but because each generation grew up in institutions (churches, courts, feudal structures) built on that assumption.
When the institutional environment changed (print revolution, rising merchant class, Enlightenment universities), the generational fixed point shifted.


\section{Comparison with Input-Output Models}
\label{sec:ch5-leontief}

The belief production chain bears structural resemblance to Leontief input-output models in economics.
In Leontief's framework, each sector's output is an input to other sectors, and equilibrium is a fixed point of the production matrix.

\begin{model}[Doxonomic Input-Output]
\label{mod:input-output}
\index{input-output model}
Let $\bm{A}$ be a matrix where $A_{ij}$ represents the fraction of channel $j$'s output that serves as input to channel $i$.
For example, $A_{21} = 0.3$ means 30\% of think tank output ($j=1$) is picked up by media ($i=2$) as source material.
The total narrative output vector $\bm{n}$ satisfies
\[
  \bm{n} = \bm{A}\bm{n} + \bm{d}
\]
where $\bm{d}$ is exogenous demand (direct funder injection).
The solution, when $(\bm{I} - \bm{A})$ is invertible, is
\[
  \bm{n} = (\bm{I} - \bm{A})^{-1}\bm{d}
\]
The matrix $(\bm{I} - \bm{A})^{-1}$ is the \emph{doxonomic multiplier}\index{doxonomic multiplier}: it shows how a unit injection of narrative through one channel amplifies through the system.
\end{model}

\begin{example}[Multiplier Interpretation]
\label{ex:multiplier}
Consider a three-channel system: think tanks ($k=1$), media ($k=2$), education ($k=3$).
Suppose
\[
  \bm{A} = \begin{pmatrix} 0 & 0.1 & 0 \\ 0.4 & 0 & 0.1 \\ 0.2 & 0.3 & 0 \end{pmatrix}
\]
This means: 40\% of think-tank output is picked up by media; 20\% eventually enters education; 30\% of media output enters education; and so on.
The Leontief inverse $(\bm{I} - \bm{A})^{-1}$ reveals that a \$1 injection into think tanks produces more than \$1 of total narrative output, because the think-tank output amplifies through media and education.
The $(3,1)$ entry of the inverse gives the \emph{total} effect on education of a unit injection into think tanks, including all indirect pathways.
\end{example}


\section{Where Chains Break: Failure Modes}
\label{sec:ch5-failure}

Not every doxonomic investment succeeds.
The chain can break at any stage:

\begin{enumerate}[nosep]
  \item \textbf{Funding failure}: insufficient resources, donor fatigue, or political change cuts funding.
  \item \textbf{Processing failure}: the institution produces low-quality output, is caught fabricating evidence, or faces internal rebellion (researchers refuse to produce the desired conclusion).
  \item \textbf{Circulation failure}: the narrative fails to reach the target audience (gatekeepers block it, algorithms deprioritize it, counter-narratives dominate attention).
  \item \textbf{Internalization failure}: the audience rejects the narrative (identity resistance, lived experience contradicts the claim, counter-inoculation is effective).
  \item \textbf{Behavioral failure}: the belief is internalized but does not produce the expected behavior (people believe meritocracy but vote for redistribution anyway, because material self-interest overrides).
  \item \textbf{Reproduction failure}: behavioral changes do not feed back into funding (the beneficiaries of the belief do not reinvest in its production).
\end{enumerate}

\begin{principle}[Weakest Link]
\label{prin:weakest-link}
\index{weakest link}
The effectiveness of a belief production chain is determined by its weakest stage.
A brilliantly funded and processed narrative that cannot circulate (Stage 3 failure) is as ineffective as a well-circulated narrative that the audience rejects (Stage 4 failure).
Counter-doxonomic strategy (Chapter~\ref{ch:counter-doxonomic}) should target the weakest link of the dominant chain.
\end{principle}


\section{The Full Chain Diagram}

\begin{center}
\begin{tikzpicture}[
  node distance=2cm and 2.5cm,
  block/.style={draw, thick, rounded corners, minimum height=1.2cm, minimum width=2.8cm, align=center, font=\small},
  op/.style={font=\footnotesize\itshape, midway, above},
  arrow/.style={-Stealth, thick}
]
  \node[block] (F) {$\fundingvec \in \R^m_{\geq 0}$\\Funding};
  \node[block, right=of F] (P) {$\insteco$\\Processing};
  \node[block, right=of P] (N) {$\narspace$\\Narratives};
  \node[block, below=of N] (C) {$E$\\Circulation};
  \node[block, left=of C] (I) {$\beliefspace$\\Internalization};
  \node[block, left=of I] (B) {$\mathcal{X}$\\Behavior};

  \draw[arrow] (F) -- (P) node[op] {$\mathcal{F}$};
  \draw[arrow] (P) -- (N) node[op] {$\mathcal{P}$};
  \draw[arrow] (N) -- (C) node[op, right] {$\mathcal{C}$};
  \draw[arrow] (C) -- (I) node[op] {$\mathcal{I}$};
  \draw[arrow] (I) -- (B) node[op] {$\mathcal{B}$};
  \draw[arrow, dashed] (B) -- (F) node[op] {$\mathcal{R}$};
\end{tikzpicture}
\end{center}


\section{Notes and References}
\label{sec:ch5-notes}

The operator-composition framework is original to this text.
The Leontief input-output analogy draws on standard economic input-output analysis; for the original framework, see Leontief (1941).
Self-fulfilling beliefs in cooperation are documented in \textcite{fischbacher2001} and \textcite{henrich2001}.
Threshold models of collective behavior originate with \textcite{granovetter1978} and \textcite{schelling1978}.
Cascade models appear in \textcite{banerjee1992} and \textcite{bikhchandani1992}.
The concept of doxonomic equilibrium as a fixed point connects to \textcite{north1990} on institutional equilibria.
Brouwer's fixed-point theorem is standard; for social science applications, see \textcite{jackson2008}.

On motivated reasoning and identity-protective cognition, see \textcite{kahan2012} and \textcite{mercier2017}.
The illusory truth effect is reviewed in \textcite{kahneman2011}.
Social proof as a persuasion mechanism is analyzed in \textcite{cialdini2006}.
The Reinhart-Rogoff episode is documented in Herndon, Ash, and Pollin (2014).
On dark money and funding opacity, see \textcite{mayer2016}.
Generational reproduction of belief through institutional transmission draws on \textcite{bourdieu1984} and \textcite{bowles1976}.
On performativity (models that reshape the reality they describe), see \textcite{mackenzie2006}.

Equilibrium selection in games with multiple equilibria is a deep problem in economic theory; see \textcite{schelling1978} for focal-point selection and \textcite{north1990} for path-dependent institutional selection.


\section{Exercises}

\begin{exercise}
Compute the doxonomic multiplier $(\bm{I} - \bm{A})^{-1}$ for the three-channel system in Example~\ref{ex:multiplier}.
Interpret the $(2,1)$ and $(3,1)$ entries: how much total media output and education output results from a unit injection into think tanks?
\end{exercise}

\begin{exercise}
\label{ex:ch5-existence}
Verify the conditions of Theorem~\ref{thm:existence} for a specific two-channel system.
Define explicit operators $\mathcal{F}$, $\mathcal{P}$, $\mathcal{C}$, $\mathcal{I}$, $\mathcal{B}$, $\mathcal{R}$ with concrete functional forms, verify continuity and compactness, and conclude that an equilibrium exists.
\end{exercise}

\begin{exercise}
Consider two competing belief production chains (one promoting ``human nature is selfish,'' the other promoting ``human nature is cooperative'') sharing the same population.
Model this as a two-player game over funding allocations and characterize the Nash equilibrium.
Under what conditions does the better-funded chain always win?
Under what conditions can the less-funded chain prevail?
\end{exercise}

\begin{exercise}
For the augmented internalization model~\eqref{eq:augmented-update}, show that if $\alpha_{\text{social}} > \alpha_{\text{Bayes}}$, a doxonomic system can maintain high $\belief{i}$ even when new evidence against $p_i$ is available, provided the perceived social norm $\hat{\mu}_j$ remains high.
Interpret: when does social proof override evidence?
\end{exercise}

\begin{exercise}
Design a disruption strategy that breaks a self-reinforcing doxonomic equilibrium.
At which stage of the chain is intervention most cost-effective?
Argue formally using the weakest-link principle and the relative costs of intervening at each stage.
\end{exercise}

\begin{exercise}
\label{ex:ch5-self-fulfilling}
The text defines self-fulfilling beliefs (Definition~\ref{def:self-fulfilling}).
Give three examples of beliefs that are self-fulfilling and three that are self-undermining (holding them generates evidence against them).
What determines which type a belief is?
\end{exercise}

\begin{exercise}
\label{ex:ch5-generational}
Using Model~\ref{mod:generational}, consider a society where the current generation believes strongly in meritocracy ($\beliefvec_t = 0.8$).
They build institutions (schools, labor markets) around this belief.
Under what conditions does the next generation inherit the same belief level?
Under what conditions does generational reproduction fail (the next generation's beliefs diverge from the previous generation's)?
\end{exercise}

\begin{exercise}
Identify a historical case where a belief production chain broke at each of the six failure modes listed in Section~\ref{sec:ch5-failure}.
For each case, explain what caused the failure and what happened to the belief afterward.
\end{exercise}


\begin{interlude}
You now have the foundational vocabulary: \emph{doxa}, the belief production chain, epistemic capital, narrative infrastructure, and the doxonomic system as a formal object.
You can identify the core units of analysis (funding, institutions, narratives, exposure, belief) and trace them through a production chain.
Part~II develops these concepts into formal models.
\end{interlude}

% --- Part II: Concepts and Models ---
\part{Concepts and Models}
\label{part:concepts}

\chapter{Epistemic Capital}
\label{ch:epistemic-capital}

\epigraph{Every established order tends to produce the naturalization of its own arbitrariness.}{Pierre Bourdieu, \textit{Outline of a Theory of Practice}}

\noindent
Not all dollars are equal in the market for belief.
A dollar spent through a prestigious university carries more doxonomic force than a dollar spent on a billboard.
The difference is \emph{epistemic capital}\index{epistemic capital}: the accumulated stock of credibility, authority, and attention that an institution or agent can deploy in the production of belief.

Epistemic capital is the currency of the doxonomic system.
It determines the conversion rate at which economic resources become belief effects.
A think tank with high epistemic capital can shift policy debates with a single report; one with low epistemic capital can produce a hundred reports that no one reads.
Understanding how epistemic capital is accumulated, converted, deployed, destroyed, and distributed is central to doxonomic analysis.


\section{Definition and Structure}
\label{sec:ch6-definition}

\begin{definition}[Epistemic Capital]
\label{def:epistemic-capital}
\index{epistemic capital}
The \emph{epistemic capital} of an agent or institution $a$ is a vector
\[
  \ecap(a) = \bigl(\ecap_{\text{cred}}(a),\; \ecap_{\text{att}}(a),\; \ecap_{\text{auth}}(a)\bigr) \in \R_{\geq 0}^3
\]
where:
\begin{itemize}[nosep]
  \item $\ecap_{\text{cred}}$ is \emph{credibility capital}: the degree to which $a$'s outputs are treated as reliable by audiences.
  Sources: track record of accuracy, institutional prestige, peer endorsement, professional credentials.
  \item $\ecap_{\text{att}}$ is \emph{attention capital}: the audience share $a$ commands, how many people will encounter $a$'s outputs.
  Sources: subscriber base, media citations, social media following, search ranking.
  \item $\ecap_{\text{auth}}$ is \emph{authority capital}: the degree to which $a$ can set the terms of legitimate discourse, defining what questions are askable, what methods are acceptable, what conclusions are serious.
  Sources: institutional gatekeeping positions, editorial control, curriculum authority, regulatory standing.
\end{itemize}
\end{definition}

This extends Bourdieu's analysis of capital forms \autocite{bourdieu1984, bourdieu1991}.
Where Bourdieu distinguished economic, cultural, social, and symbolic capital qualitatively, we provide a quantitative decomposition oriented specifically toward belief production.

The three components are partially independent.
An institution can have high credibility but low attention (an obscure but respected research lab), high attention but low credibility (a viral social media account), or high authority but low attention (a regulatory body whose rulings shape an industry but whose reports are unread by the public).
The doxonomic force of an institution depends on the \emph{product} of relevant components, not on any single one.

\begin{definition}[Doxonomic Force]
\label{def:doxonomic-force}
\index{doxonomic force}
The \emph{doxonomic force} of an institution $a$ with respect to belief $i$ is
\[
  \Phi_i(a) = \ecap_{\text{cred}}(a) \cdot \ecap_{\text{att}}(a) \cdot \relevance{a}{i}
\]
where $\relevance{a}{i}$ is the relevance of $a$'s domain to belief $i$.
Authority capital enters indirectly: $\ecap_{\text{auth}}$ shapes the boundaries of $\narspace$ (which narratives are ``in play'') rather than the force of any particular narrative.
\end{definition}

\begin{example}[Epistemic Capital Profiles]
\label{ex:profiles}
\index{epistemic capital!profiles}
Consider five institutions and their approximate capital vectors (illustrative scores on a 0--1 scale, not empirically calibrated):

\begin{center}
\begin{tabular}{lccc}
\toprule
\textbf{Institution} & $\ecap_{\text{cred}}$ & $\ecap_{\text{att}}$ & $\ecap_{\text{auth}}$ \\
\midrule
Harvard Economics Dept. & 0.92 & 0.15 & 0.85 \\
\textit{New York Times} & 0.70 & 0.60 & 0.55 \\
Heritage Foundation & 0.55 & 0.35 & 0.50 \\
Joe Rogan podcast & 0.20 & 0.80 & 0.05 \\
Anonymous Twitter account & 0.05 & 0.40 & 0.01 \\
\bottomrule
\end{tabular}
\end{center}

Harvard has high credibility and authority but limited direct public reach.
The \textit{Times} has a balanced profile.
Heritage has moderate credibility but strong policy-world authority.
Joe Rogan has enormous attention capital but low credibility among experts and near-zero authority capital.
An anonymous Twitter account can command surprising attention but has negligible credibility or authority.

The doxonomic strategy of channeling narratives through complementary institutions exploits the comparative advantage of each: Heritage (credibility + authority), the \textit{Times} (attention + credibility), social media (attention).
\end{example}


\section{Conversion Functions}
\label{sec:ch6-conversion}

Epistemic capital is not created from nothing.
It is converted from other forms of capital, primarily economic capital, through specific institutional mechanisms.
The conversion is neither instant nor guaranteed: building credibility takes years, building authority takes decades, and both can be destroyed in days.

\begin{definition}[Capital Conversion]
\label{def:conversion}
\index{capital conversion}
A \emph{capital conversion function} is a map $\phi: \R_{\geq 0} \to \R_{\geq 0}^3$ such that
\[
  \ecap(a, t+1) = \ecap(a, t) + \phi(f_a(t)) - \delta \cdot \ecap(a, t)
\]
where $f_a(t)$ is the economic investment at time $t$ and $\delta \in (0,1)$ is a depreciation rate.
The function $\phi$ is typically concave: the marginal epistemic capital gained per dollar decreases with scale.
\end{definition}

The concavity of $\phi$ is important.
It means that the first million dollars invested in building a new think tank buys more credibility than the hundredth million.
At some point, additional spending on institutional trappings (nicer offices, more conferences, shinier publications) produces negligible additional credibility.
This is the doxonomic analogue of diminishing returns.

\begin{example}[Conversion Pathways]
\label{ex:conversion-pathways}
\index{capital conversion!pathways}
Different investments convert economic capital into different components of $\ecap$:

\begin{center}
\begin{tabular}{lp{6cm}ccc}
\toprule
\textbf{Investment} & \textbf{Mechanism} & $\Delta\ecap_{\text{cred}}$ & $\Delta\ecap_{\text{att}}$ & $\Delta\ecap_{\text{auth}}$ \\
\midrule
Endow a university chair (\$5M) & Institutional affiliation, peer network & High & Low & Medium \\
Fund a think tank (\$5M/yr) & Reports, op-eds, testimony & Medium & Medium & Medium \\
Buy advertising (\$5M) & Audience reach, repetition & Low & High & None \\
Fund a journal (\$2M/yr) & Gatekeeping, peer review & Medium & Low & High \\
Sponsor a conference (\$500K) & Network effects, prestige association & Medium & Low & Medium \\
Hire a PR firm (\$1M) & Media placement, message optimization & Low & Medium & None \\
\bottomrule
\end{tabular}
\end{center}

The most cost-effective credibility conversion is institutional affiliation: attaching a narrative to a university, a professional body, or a credentialed expert.
The most cost-effective attention conversion is advertising or viral content.
Authority capital is the hardest to purchase directly; it accumulates through sustained institutional presence over time.
\end{example}

\begin{proposition}[Steady-State Epistemic Capital]
\label{prop:steady-state}
Under constant investment $f$ and depreciation rate $\delta$, the epistemic capital converges to a steady state:
\[
  \ecap^* = \frac{\phi(f)}{\delta}
\]
Institutions that stop investing ($f = 0$) decay exponentially: $\ecap(t) = \ecap(0) \cdot e^{-\delta t}$ with half-life $t_{1/2} = \ln 2 / \delta$.
For institutions with $\delta \approx 0.02$ (elite universities, legacy media), $t_{1/2} \approx 35$ years.
For institutions with $\delta \approx 0.2$ (social media influencers, digital outlets), $t_{1/2} \approx 3.5$ years.
\end{proposition}

\begin{remark}[The Time Asymmetry of Credibility]
\label{rem:time-asymmetry}
Credibility is slow to build and fast to destroy.
The accumulation timescale is $\tau_{\text{build}} \sim 1/\phi'(f)$, which for credibility capital is typically years to decades.
The destruction timescale can be instantaneous: a single scandal, a single fabrication, a single revealed conflict of interest can reduce $\ecap_{\text{cred}}$ by a large fraction overnight.
This asymmetry has a doxonomic consequence: institutions with high credibility capital have a strong incentive to maintain quality, not out of virtue, but because their most valuable asset is fragile.
When this incentive is undermined (e.g., by pressure from funders who value message over accuracy), the institution's credibility begins to erode.
\end{remark}


\section{Epistemic Rents}
\label{sec:ch6-rents}

\begin{definition}[Epistemic Rent]
\label{def:epistemic-rent}
\index{epistemic rent}
An \emph{epistemic rent} is the excess doxonomic return earned by an agent due to a monopoly or near-monopoly position in the production of legitimate knowledge on a topic.
Formally, if the competitive doxonomic return per dollar is $r^*$, an agent with epistemic rent earns $r > r^*$ due to barriers to entry in the credibility market.
The rent is $\rho = r - r^*$.
\end{definition}

Epistemic rents explain why certain institutions (elite universities, credentialed professions, legacy media) exert disproportionate influence on public belief relative to their expenditure.
They have accumulated credibility that new entrants cannot quickly replicate.

\begin{example}[Sources of Epistemic Rent]
\label{ex:rent-sources}
\begin{itemize}[nosep]
  \item \textbf{Credential barriers}: only holders of an MD can authoritatively pronounce on medical treatment; only holders of a PhD in economics are taken seriously in policy debates.
  The credential system creates a barrier to entry that protects incumbents' epistemic rent.
  \item \textbf{Brand accumulation}: the \textit{New York Times} has accumulated credibility over 170 years.
  A new publication cannot replicate this overnight, regardless of quality.
  The brand acts as a moat around the epistemic rent.
  \item \textbf{Network effects}: a journal's authority depends on who publishes in it, and who publishes in it depends on its authority.
  This circularity creates a self-reinforcing advantage for established journals.
  \item \textbf{Institutional interlocking}: elite institutions cite each other, hire each other's graduates, and appear on each other's boards.
  This network structure distributes epistemic rent across the cluster while making it difficult for outsiders to penetrate.
\end{itemize}
\end{example}

\begin{proposition}[Rent Persistence]
\label{prop:rent-persistence}
If the depreciation rate $\delta$ is low and barriers to credibility accumulation are high, epistemic rents can persist for decades.
Formally, the half-life of an epistemic rent is $t_{1/2} = \ln 2 / \delta$, which for institutions with $\delta \approx 0.02$ gives $t_{1/2} \approx 35$ years.
Rent persistence means that the current credibility landscape is substantially determined by investments made a generation ago.
\end{proposition}


\section{Credibility Destruction}
\label{sec:ch6-destruction}

If epistemic capital is the most valuable asset in doxonomic production, then credibility destruction is the most potent weapon in doxonomic warfare.

\begin{definition}[Credibility Attack]
\label{def:credibility-attack}
\index{credibility attack}
A \emph{credibility attack} on institution $a$ is an action that reduces $\ecap_{\text{cred}}(a)$.
Types include:
\begin{enumerate}[nosep]
  \item \textbf{Exposure}: revealing hidden funding, conflicts of interest, or fabrication.
  \item \textbf{Delegitimation}: framing the institution as biased, politicized, or captured (``liberal academia,'' ``mainstream media'').
  \item \textbf{Flooding}: producing so much contradictory content that the audience cannot distinguish credible from non-credible sources (``epistemic pollution'').
  \item \textbf{Co-optation}: infiltrating the institution to degrade its output quality from within.
\end{enumerate}
\end{definition}

\begin{example}[The War on Expertise]
\label{ex:war-expertise}
The systematic delegitimation of scientific expertise on climate change is a case of organized credibility attack \autocite{oreskes2010}.
The fossil fuel industry did not need to build higher credibility than climate scientists.
It needed only to \emph{reduce} the credibility of climate science enough that the public perceived a ``debate,'' a strategy of epistemic pollution rather than epistemic competition.
The cost of reducing public trust in climate science was far lower than the cost of building a credible alternative scientific consensus.
This asymmetry (destruction is cheaper than construction) is a fundamental property of epistemic capital.
\end{example}

\begin{proposition}[Destruction-Construction Asymmetry]
\label{prop:asymmetry}
\index{credibility!destruction-construction asymmetry}
Let $c_{\text{build}}(\Delta\ecap)$ be the cost of increasing an institution's credibility by $\Delta\ecap$ and $c_{\text{destroy}}(\Delta\ecap)$ the cost of reducing a rival's credibility by the same amount.
Generically, $c_{\text{destroy}} \ll c_{\text{build}}$.
One scandal can destroy what took decades to build.
This asymmetry implies that doxonomic systems with high epistemic capital are structurally vulnerable to well-funded credibility attacks.
\end{proposition}


\section{Interaction Effects Between Capital Types}
\label{sec:ch6-interactions}

The three components of epistemic capital do not operate independently.
They interact in ways that create both synergies and vulnerabilities.

\begin{model}[Capital Interaction Matrix]
\label{mod:interaction}
\index{epistemic capital!interactions}
The evolution of epistemic capital is governed by:
\[
  \frac{d}{dt}\ecap(a) = \phi(f_a) - \delta \cdot \ecap(a) + \bm{M} \cdot \ecap(a)
\]
where $\bm{M}$ is a $3 \times 3$ interaction matrix.
Key interactions:
\begin{itemize}[nosep]
  \item $M_{\text{att},\text{cred}} > 0$: credibility attracts attention (credible sources get cited, shared, invited).
  \item $M_{\text{auth},\text{cred}} > 0$: credibility enables authority (credible institutions are invited to set standards, write curricula, advise governments).
  \item $M_{\text{cred},\text{att}} < 0$ (potentially): excessive attention can reduce credibility (``selling out,'' ``dumbing down,'' ``clickbait'').
  \item $M_{\text{att},\text{auth}} > 0$: authority generates attention (authoritative pronouncements are newsworthy).
\end{itemize}
\end{model}

The negative interaction $M_{\text{cred},\text{att}} < 0$ is doxonomically significant.
It means there is a tension between reaching a mass audience and maintaining credibility among discriminating audiences.
Institutions that ``go viral'' sometimes find their credibility among experts diminished.
This is the \emph{populism-credibility tradeoff}: you can speak to everyone or be believed by the serious people, but doing both is difficult.

\begin{example}[The Academic-Public Dilemma]
\label{ex:academic-public}
An economist who writes a popular book (high $\Delta\ecap_{\text{att}}$) may find their standing among peers reduced (negative $\Delta\ecap_{\text{cred}}$ in academic circles).
An academic department that aggressively markets itself on social media may be perceived by peer departments as unserious.
This creates an institutional incentive against public engagement. From a doxonomic perspective, the highest-credibility institutions systematically underinvest in public communication, leaving the public information environment to lower-credibility sources.
\end{example}


\section{Epistemic Capital Inequality}
\label{sec:ch6-inequality}

Epistemic capital, like economic capital, is unequally distributed.
A small number of institutions control a disproportionate share of credibility, attention, and authority.

\begin{definition}[Epistemic Gini Coefficient]
\label{def:epistemic-gini}
\index{epistemic Gini coefficient}
The \emph{epistemic Gini coefficient} for a population of $n$ institutions with doxonomic forces $\Phi_1, \ldots, \Phi_n$ is
\[
  G_{\ecap} = \frac{\sum_{i=1}^{n}\sum_{j=1}^{n} |\Phi_i - \Phi_j|}{2n\sum_{i=1}^{n}\Phi_i}
\]
$G_{\ecap} = 0$ indicates perfect equality; $G_{\ecap} = 1$ indicates total concentration.
\end{definition}

\begin{remark}
In most media systems, the epistemic Gini is very high, likely exceeding the economic Gini.
A handful of outlets (major newspapers, broadcast networks, dominant platforms) command the majority of public attention, while thousands of smaller outlets share the remainder.
The implication: the effective ``electorate'' of the belief-production system is far more concentrated than the political electorate.
This concentration is the structural basis of doxonomic power (Chapter~\ref{ch:doxonomic-power}).
\end{remark}


\section{The Digital Disruption of Epistemic Capital}
\label{sec:ch6-digital}

The internet disrupted traditional epistemic capital hierarchies in several ways:

\begin{itemize}[nosep]
  \item \textbf{Attention democratization}: anyone can publish, reducing the barriers to $\ecap_{\text{att}}$.
  Blogs, podcasts, YouTube channels, and social media accounts compete with legacy media for audience.
  \item \textbf{Credibility unbundling}: traditional media bundled attention and credibility together (the \textit{Times} had both).
  Digital media unbundled them: a viral tweet has enormous attention but no institutional credibility.
  \item \textbf{Authority erosion}: gatekeeping functions (editorial selection, peer review, professional licensing) are partially circumvented by direct-to-audience publishing.
  \item \textbf{Speed advantage}: digital channels can produce and circulate narrative faster than traditional institutions, meaning low-credibility sources often set the initial frame that high-credibility sources then react to.
\end{itemize}

\begin{model}[Two-Tier Epistemic Capital]
\label{mod:two-tier}
In the digital era, epistemic capital bifurcates:
\begin{itemize}[nosep]
  \item \textbf{Institutional capital} ($\ecap^I$): slow to build, high credibility, high authority, moderate attention.
  Still dominates policy, academia, and elite discourse.
  \item \textbf{Network capital} ($\ecap^N$): fast to build, low credibility, low authority, potentially very high attention.
  Dominates popular discourse, social media, and cultural conversation.
\end{itemize}
The doxonomic competition between institutional and network capital is one of the defining features of the contemporary information environment.
Legacy institutions have $\ecap^I \gg \ecap^N$; digital-native actors have $\ecap^N \gg \ecap^I$.
Neither can easily convert their capital into the other type.
\end{model}

This bifurcation creates a distinctive doxonomic dynamic.
Policy elites and academic experts retain institutional capital and shape law and regulation.
Digital influencers and populist media retain network capital and shape public mood and electoral outcomes.
The two systems operate semi-independently, producing the conditions for epistemic fragmentation (Chapter~\ref{ch:contradiction-coherence}).


\section{The Credibility Market}
\label{sec:ch6-market}

We can model the production of credibility as a market in which institutions compete for audience trust.

\begin{model}[Credibility Market]
\label{mod:credibility-market}
Let $n$ institutions compete by choosing quality $q_k$ (costly) and signaling $s_k$ (also costly).
The audience assigns credibility $\credibility{k} = h(q_k, s_k, \ecap_{\text{cred}}(k))$ where $h$ is increasing in all arguments.
Equilibrium credibility levels satisfy
\[
  \max_{q_k, s_k} \; \credibility{k} \cdot \reach{k} \cdot \pi_k - c_q(q_k) - c_s(s_k)
\]
where $\pi_k$ is the revenue from doxonomic output and $c_q, c_s$ are cost functions for quality and signaling respectively.
\end{model}

This model reveals a market failure: because credibility has a public-good character (audiences benefit from accurate information regardless of whether they pay for it), the market tends to under-produce genuine quality and over-produce cheap signals.

\begin{proposition}[The Lemons Problem in Credibility]
\label{prop:lemons}
\index{credibility!lemons problem}
If audiences cannot perfectly distinguish quality from signaling ($h$ depends on $s_k$ as well as $q_k$), then in equilibrium, some institutions will substitute signaling for quality.
As audiences learn that signals are noisy, they discount all sources, including high-quality ones.
This is the \emph{epistemic lemons problem}\index{epistemic lemons problem}: low-quality producers drive down the perceived value of the entire market, just as bad used cars drive good ones off the lot \autocite{akerlof2015}.
The result is a race to the bottom in credibility, especially in domains where quality is hard to assess (complex policy, long-run predictions, statistical claims).
\end{proposition}


\section{Notes and References}
\label{sec:ch6-notes}

The concept of epistemic capital formalizes ideas in \textcite{bourdieu1984, bourdieu1991}.
The three-component decomposition (credibility, attention, authority) extends Bourdieu's capital typology to the specific needs of doxonomic analysis.
The economics of credibility connects to \textcite{stigler1961} on the economics of information and \textcite{akerlof2015} on manipulation and deception.
Epistemic rents parallel the rent concept in \textcite{stiglitz2012}.
The market for ideas is analyzed in \textcite{coase1974}.
For the role of expertise as a form of capital, see \textcite{eyal2019}.

On the destruction of credibility through organized doubt, see \textcite{oreskes2010}.
The digital disruption of epistemic hierarchies is analyzed in \textcite{benkler2018} and \textcite{zuboff2019}.
On the tension between popular communication and expert credibility, see \textcite{eyal2019}.
The epistemic lemons problem extends Akerlof's (1970) market-for-lemons argument to the information domain.
On epistemic injustice as a structural barrier to credibility accumulation, see \textcite{fricker2007}: members of marginalized groups face systematic credibility deficits regardless of the quality of their testimony, which is a form of epistemic capital inequality not reducible to institutional investment.

For the interaction between different forms of capital, see Bourdieu's general theory of capital conversion in \textcite{bourdieu1984}.
The concept of network capital as distinct from institutional capital draws on \textcite{benkler2006} and the literature on networked public spheres.


\section{Exercises}

\begin{exercise}
Estimate the epistemic capital vector $\ecap$ for three institutions of your choice (e.g., a national newspaper, a social media influencer, a university department).
Justify your estimates.
Compute the doxonomic force $\Phi_i(a)$ for each with respect to a specific belief.
\end{exercise}

\begin{exercise}
Prove that if $\phi$ is strictly concave, there exists an optimal investment level $f^*$ beyond which additional spending on credibility conversion has diminishing returns below the depreciation rate.
Formally: find $f^*$ such that $\phi'(f^*) = \delta$.
What happens if the institution invests $f > f^*$?
\end{exercise}

\begin{exercise}
Model a ``credibility laundering'' operation in which a low-credibility funder channels money through a high-credibility institution.
What is the net credibility gain?
Under what conditions does the institution's credibility itself depreciate from the association?
Use the capital conversion dynamics to model the institution's credibility trajectory when it accepts tainted funding.
\end{exercise}

\begin{exercise}
\label{ex:ch6-gini}
Compute the epistemic Gini coefficient for a hypothetical media system with five outlets having doxonomic forces $\Phi = (100, 50, 20, 5, 1)$.
Is this system more or less concentrated than a typical economic wealth distribution?
What would a policy intervention to reduce $G_{\ecap}$ look like?
\end{exercise}

\begin{exercise}
Analyze the credibility attack strategy used against climate science.
Which of the four attack types (exposure, delegitimation, flooding, co-optation) were employed?
Estimate the cost of the attack relative to the cost of building the scientific credibility it targeted.
\end{exercise}

\begin{exercise}
For the two-tier model (Model~\ref{mod:two-tier}), consider an institution that tries to build both $\ecap^I$ and $\ecap^N$ simultaneously.
Under what conditions is this possible without the populism-credibility tradeoff destroying institutional capital?
Give a real-world example of an institution that has succeeded (or failed) at this.
\end{exercise}

\begin{exercise}
\label{ex:ch6-interaction}
Write out the full $3 \times 3$ interaction matrix $\bm{M}$ for a university, a think tank, and a social media influencer.
For each, which off-diagonal entries are positive, negative, or zero?
How do the interaction effects shape the long-run evolution of each institution's epistemic capital?
\end{exercise}

\chapter{Narrative Infrastructure}
\label{ch:narrative-infrastructure}

\epigraph{Who controls the past controls the future. Who controls the present controls the past.}{George Orwell, \textit{Nineteen Eighty-Four}}

\noindent
Beliefs do not float freely.
They travel through infrastructure: schools, media organizations, platforms, think tanks, entertainment industries, professional training systems, and religious institutions.
Just as an economy requires physical infrastructure (roads, ports, power grids) to move goods, a doxonomic system requires \emph{narrative infrastructure} to move beliefs from production to population.

This chapter models narrative infrastructure as a directed network of transformation functions with measurable capacity, throughput, concentration, and vulnerability.
The infrastructure is not neutral: its topology determines which narratives can travel far and which are stillborn.


\section{Infrastructure as Transformation Network}
\label{sec:ch7-transformation}

\begin{definition}[Narrative Infrastructure]
\label{def:narrative-infrastructure}
\index{narrative infrastructure}
A \emph{narrative infrastructure} $\narr$ is a directed weighted graph $G = (V, E, w)$ where:
\begin{itemize}[nosep]
  \item each vertex $v \in V$ is an institution with processing function $\iota_v: \R_{\geq 0} \times \narspace \to \narspace$,
  \item each directed edge $(u, v) \in E$ carries a flow weight $w_{uv} \in [0,1]$ representing the fraction of $u$'s output that serves as input to $v$,
  \item the network has source nodes (funders injecting resources) and sink nodes (audience segments absorbing narratives).
\end{itemize}
\end{definition}

The key insight is that narrative infrastructure is not a flat landscape.
It has topology: some paths are wide and well-maintained (the pipeline from elite think tanks through prestige media to policy debate), while others are narrow or blocked (the path from grassroots organizations to national attention).
The infrastructure determines \emph{which} narratives are structurally advantaged, independent of their truth or popularity.

\begin{example}[The Conservative Narrative Pipeline]
\label{ex:conservative-pipeline}
\index{narrative pipeline}
The American conservative movement built one of the most effective narrative pipelines in modern history, documented in \textcite{mayer2016} and \textcite{mirowski2009}:
\[
  \text{Donors} \xrightarrow{0.8} \text{Foundations} \xrightarrow{0.9} \text{Think tanks} \xrightarrow{0.5} \text{Op-ed pages} \xrightarrow{0.6} \text{Cable news} \xrightarrow{0.7} \text{Public}
\]
The flow weights indicate the estimated fraction of output that propagates to the next stage.
The end-to-end propagation probability is $\prod_i w_i = 0.8 \times 0.9 \times 0.5 \times 0.6 \times 0.7 \approx 0.15$, seemingly low, but applied to hundreds of reports per year across dozens of institutions, it produces a constant stream of narrative output reaching the public through high-credibility channels.

A parallel pipeline runs through academia:
\[
  \text{Donors} \xrightarrow{0.7} \text{Endowed programs} \xrightarrow{0.6} \text{Textbooks} \xrightarrow{0.8} \text{Curricula} \xrightarrow{0.95} \text{Students}
\]
This pipeline has lower throughput per year but vastly higher persistence ($\persistence{} \sim$ decades).
\end{example}

\begin{remark}[Infrastructure vs.\ Content]
\label{rem:infra-content}
A crucial distinction: narrative infrastructure is about the \emph{channels}, not the \emph{messages}.
The same infrastructure that carries ``markets are efficient'' could, in principle, carry ``markets are prone to failure.''
But infrastructure is not neutral: institutions built to produce one type of narrative develop staff, expertise, donor relationships, and organizational culture that make producing alternative narratives costly.
A think tank funded for 30 years to argue against regulation does not easily pivot to arguing for it.
Infrastructure has inertia.
\end{remark}


\section{Institutional Types and Their Infrastructure Roles}
\label{sec:ch7-types}

Different institutions occupy structurally different positions in the narrative infrastructure.
Their roles can be classified by where they sit in the production chain and what transformation they perform.

\begin{definition}[Infrastructure Role Taxonomy]
\label{def:role-taxonomy}
\index{infrastructure roles}
\begin{description}[nosep]
  \item[Generators] produce original narrative content: research universities, think tanks, investigative journalists, research labs.
  \item[Amplifiers] increase the reach of existing narratives without substantially transforming them: broadcast media, social media platforms, wire services, syndication networks.
  \item[Translators] transform narratives from one register to another: textbook authors (research $\to$ pedagogy), popularizers (academic $\to$ public), lobbyists (policy research $\to$ legislative language).
  \item[Validators] attach credibility to narratives: peer review systems, professional associations, credentialing bodies, editorial boards.
  \item[Gatekeepers] control access to high-value channels: editors, algorithm designers, curriculum committees, conference organizers.
\end{description}
\end{definition}

A single institution may perform multiple roles.
The \textit{New York Times} is simultaneously a generator (original reporting), an amplifier (wire service redistribution), a translator (making complex events accessible), a validator (the prestige of publication implies quality), and a gatekeeper (deciding what is ``fit to print'').
This multi-role character gives major media institutions outsized structural power.

\begin{center}
\begin{tikzpicture}[
  every node/.style={font=\small},
  box/.style={draw, thick, rounded corners, minimum width=3cm, minimum height=0.8cm, align=center},
]
\matrix[row sep=0.5cm, column sep=0.3cm] {
  \node[box, fill=doxoblue!8] {Education}; &
  \node[font=\small, text width=7.5cm] {\textbf{Generator + Translator + Validator.} High persistence ($\sim$decades), high credibility, low throughput per unit time. Shapes foundational priors. Most consequential long-run infrastructure.}; \\
  \node[box, fill=doxoblue!8] {Prestige Media}; &
  \node[font=\small, text width=7.5cm] {\textbf{Generator + Amplifier + Gatekeeper.} Medium persistence ($\sim$months), high credibility, medium throughput. Sets the daily agenda and frames for elites and opinion leaders.}; \\
  \node[box, fill=doxoblue!8] {Social Media}; &
  \node[font=\small, text width=7.5cm] {\textbf{Amplifier + (sometimes) Generator.} Very low persistence ($\sim$hours), low credibility, very high throughput. Amplifies, fragments, and accelerates. Algorithms act as invisible gatekeepers.}; \\
  \node[box, fill=doxoblue!8] {Think Tanks}; &
  \node[font=\small, text width=7.5cm] {\textbf{Generator + Validator.} High credibility in policy circles, low direct public reach. Produces the ``studies show'' citations that legitimate policy narratives.}; \\
  \node[box, fill=doxoblue!8] {Entertainment}; &
  \node[font=\small, text width=7.5cm] {\textbf{Translator + Amplifier.} High reach, low perceived persuasive intent. Normalizes through narrative representation, not argument. Operates at layers 3--5.}; \\
  \node[box, fill=doxoblue!8] {Religious Institutions}; &
  \node[font=\small, text width=7.5cm] {\textbf{Generator + Validator + Community.} Very high persistence, high credibility within congregations, medium reach. Unique access to moral framing and identity formation.}; \\
  \node[box, fill=doxoblue!8] {Professional Training}; &
  \node[font=\small, text width=7.5cm] {\textbf{Translator + Validator.} Converts academic knowledge into operational assumptions. MBA programs, law schools, medical training embed anthropological regimes into professional practice.}; \\
};
\end{tikzpicture}
\end{center}


\section{Throughput and Capacity}
\label{sec:ch7-throughput}

\begin{definition}[Narrative Throughput]
\label{def:throughput}
\index{narrative throughput}
The \emph{throughput} of an institution $v$ is the quantity
\[
  \Theta(v) = \funding{v} \cdot \reach{v} \cdot \persistence{v}
\]
measuring the total belief-exposure-hours produced per unit time.
The \emph{capacity} $\Theta^{\max}(v)$ is the maximum throughput achievable given institutional constraints (staffing, infrastructure, brand).
\end{definition}

Throughput combines three dimensions: how much is spent, how many people are reached, and how long the effect lasts.
This composite measure allows comparison across radically different channel types.

\begin{example}[Throughput Comparison]
\label{ex:throughput-comparison}
Consider three channels, each receiving \$1M/year:
\begin{center}
\begin{tabular}{lcccc}
\toprule
\textbf{Channel} & $\reach{}$ & $\persistence{}$ (years) & $\Theta$ (M\$·person·yr) \\
\midrule
University economics course & 0.001 & 30 & 0.03 \\
Cable news segment & 0.05 & 0.02 & 0.001 \\
Textbook chapter & 0.005 & 20 & 0.1 \\
\bottomrule
\end{tabular}
\end{center}
The textbook has 100 times the throughput of cable news, despite reaching far fewer people per day, because its persistence is three orders of magnitude higher.
This is why doxonomic strategists who understand the mathematics invest in education and publishing, not just media.
\end{example}

\begin{remark}[Why Throughput Is Not Impact]
\label{rem:throughput-impact}
Throughput measures exposure, not belief change.
A high-throughput channel with low credibility (social media) may produce less actual belief change than a low-throughput channel with high credibility (a mentor's advice).
The full impact depends on the internalization operator $\mathcal{I}$ (Chapter~\ref{ch:belief-production-chain}), which weights exposure by credibility, resistance, and audience characteristics.
Throughput is a necessary but not sufficient condition for doxonomic effect.
\end{remark}


\section{Network Topology and Structural Advantage}
\label{sec:ch7-topology}

The topology of narrative infrastructure creates structural advantages and disadvantages for different narratives.

\begin{definition}[Narrative Pathway]
\label{def:pathway}
\index{narrative pathway}
A \emph{narrative pathway} from source $s$ to audience segment $t$ is a directed path $s = v_0, v_1, \ldots, v_l = t$ in the infrastructure graph $G$.
The \emph{pathway efficiency} is
\[
  \eta(s \to t) = \prod_{i=0}^{l-1} w_{v_i, v_{i+1}} \cdot \credibility{v_i}
\]
measuring the fraction of the original narrative that reaches the audience, weighted by the credibility accumulated along the way.
\end{definition}

\begin{proposition}[Structural Advantage]
\label{prop:structural-advantage}
\index{structural advantage}
A narrative $n_1$ has a \emph{structural advantage} over narrative $n_2$ if the maximum pathway efficiency from any source to the target audience is higher for $n_1$:
\[
  \max_{s} \eta_{n_1}(s \to t) > \max_{s} \eta_{n_2}(s \to t)
\]
Structural advantage arises from the infrastructure, not from the content.
A narrative with better-connected institutional sponsors, more pathways to high-credibility amplifiers, and shorter routes to the target audience has a structural advantage regardless of its truth value.
\end{proposition}

\begin{example}[Structural Advantage: Deregulation vs.\ Public Ownership]
\label{ex:structural-advantage}
The narrative ``deregulation promotes efficiency'' has extensive infrastructure: well-funded think tanks (Heritage, Cato, AEI) connected to prestigious op-ed pages (\textit{WSJ}, \textit{FT}), cable news panels, congressional testimony channels, and economics curricula.
The maximum pathway efficiency to U.S.\ policymakers is high.

The narrative ``public ownership can be efficient'' has minimal infrastructure: a few academic advocates, limited think-tank support, almost no connections to major op-ed pages or cable news.
The maximum pathway efficiency to the same policymakers is low.

This structural difference exists \emph{independently} of the empirical merits of either claim.
It is a property of the infrastructure, not of the evidence.
\end{example}


\section{Concentration Metrics}
\label{sec:ch7-concentration}

\begin{definition}[Epistemic Herfindahl Index]
\label{def:ehhi-infra}
\index{epistemic Herfindahl index}
The \emph{epistemic Herfindahl index} for narrative infrastructure is
\[
  \ehhi = \sum_{v \in V} \left(\frac{\Theta(v)}{\sum_{u \in V} \Theta(u)}\right)^2
\]
Values near $1$ indicate monopoly in narrative production; values near $1/|V|$ indicate deconcentration.
\end{definition}

\begin{conjecture}[Concentration and Capture]
\label{conj:concentration-capture}
If $\ehhi > \bar{h}$ for a critical threshold $\bar{h}$, then a single institution or small coalition has the \emph{capacity} to achieve common-sense capture (Definition~\ref{def:common-sense}) on propositions within its narrative domain, provided sufficient funding and absent strong counter-institutions, resistant audiences, or contradicting lived experience.
The threshold $\bar{h}$ is not a universal constant; it depends on audience receptivity, competing credibility structures, and proposition-content fit.
\end{conjecture}

\begin{example}[Concentration in U.S.\ Media]
\label{ex:us-concentration}
As of 2020, five companies (Comcast/NBCUniversal, Disney, ViacomCBS, WarnerMedia, News Corp/Fox) controlled approximately 90\% of U.S.\ media consumption \autocite{bagdikian2004}.
If we weight by throughput (reach $\times$ persistence), the $\ehhi$ for U.S.\ national media is approximately 0.20--0.25, near the threshold typically associated with high concentration in market economics.

But this understates doxonomic concentration, because the five companies share substantial overlap in:
\begin{itemize}[nosep]
  \item advertising dependence (same advertisers, same incentive structure),
  \item sourcing patterns (same wire services, same government press conferences),
  \item professional culture (journalism schools produce a shared worldview),
  \item ownership class interests (all publicly traded, all dependent on deregulatory environment).
\end{itemize}
Effective doxonomic concentration is higher than $\ehhi$ alone suggests, because the ``competing'' outlets share structural filters (Chapter~\ref{ch:prehistory}, Section~\ref{sec:ch3-herman-chomsky}).
\end{example}

\begin{definition}[Effective Concentration]
\label{def:effective-concentration}
\index{effective concentration}
The \emph{effective concentration} adjusts $\ehhi$ for narrative correlation:
\[
  \ehhi_{\text{eff}} = \ehhi + (1 - \ehhi) \cdot \bar{\rho}
\]
where $\bar{\rho} \in [0,1]$ is the average pairwise narrative correlation between outlets (measured by text similarity, framing overlap, or source sharing).
When all outlets produce identical narratives ($\bar{\rho} = 1$), $\ehhi_{\text{eff}} = 1$ regardless of market structure: a perfectly correlated oligopoly is doxonomically equivalent to a monopoly.
\end{definition}


\section{Infrastructure Resilience and Vulnerability}
\label{sec:ch7-resilience}

\begin{definition}[Infrastructure Vulnerability]
\label{def:vulnerability}
\index{infrastructure vulnerability}
A narrative infrastructure is \emph{vulnerable} at node $v$ if removing $v$ significantly reduces the maximum pathway efficiency for a class of narratives.
The \emph{vulnerability index} of $v$ is
\[
  \nu(v) = 1 - \frac{\max_s \eta_{-v}(s \to t)}{\max_s \eta(s \to t)}
\]
where $\eta_{-v}$ is the pathway efficiency with $v$ removed.
High $\nu(v)$ means $v$ is a critical bottleneck.
\end{definition}

\begin{example}[Bottleneck Institutions]
\label{ex:bottleneck}
In many countries, a small number of editorial gatekeepers control access to the prestige media pipeline.
If three editors at major newspapers decide that a topic is not newsworthy, the pathway from think-tank research to public attention is effectively blocked.
These editors have high vulnerability indices: removing them from the network (or changing their editorial stance) would significantly alter which narratives reach the public.

Similarly, a curriculum committee that approves textbooks for a state with millions of students is a high-vulnerability node.
Texas's State Board of Education, which approves textbooks used by 5.5 million students, has historically been a bottleneck node with outsized doxonomic influence on national textbook content (publishers produce one version for Texas and sell it nationally).
\end{example}

\begin{proposition}[Infrastructure Fragility Under Concentration]
\label{prop:fragility-concentration}
Highly concentrated infrastructure ($\ehhi$ high) is more efficient for narrative transmission but more fragile: disruption of a single high-throughput node has a larger system-wide effect.
Formally, if $\Theta(v) / \sum_u \Theta(u) > 1/2$ for some $v$, then $\nu(v) > 1/2$ and the infrastructure is critically dependent on $v$.
\end{proposition}

This creates a design tension: concentrated infrastructure is powerful but fragile; distributed infrastructure is resilient but less efficient.
Chapter~\ref{ch:epistemic-democracy} considers what institutional design best balances efficiency and resilience.


\section{Infrastructure Construction and Path Dependence}
\label{sec:ch7-construction}

Narrative infrastructure does not appear overnight.
Building it requires sustained investment over years or decades.

\begin{example}[Building Conservative Infrastructure]
\label{ex:building-infrastructure}
The conservative narrative infrastructure in the United States was consciously constructed over 40 years \autocite{mayer2016, phillipsfein2009}:
\begin{itemize}[nosep]
  \item \textbf{1947}: Mont P\`{e}lerin Society founded (intellectual seed).
  \item \textbf{1971}: Lewis Powell memo urges American business to invest in ideological infrastructure.
  \item \textbf{1973}: Heritage Foundation founded (\$250K initial; \$80M+ budget by 2020).
  \item \textbf{1977}: Cato Institute founded.
  \item \textbf{1980s}: Olin, Scaife, and Bradley foundations dramatically increase grants to conservative law-and-economics programs at universities.
  \item \textbf{1990s}: Fox News launched (1996), completing the broadcast amplification layer.
  \item \textbf{2000s}: Digital expansion (Daily Wire, Turning Point USA, PragerU).
\end{itemize}
Total estimated investment: several billion dollars over four decades.
The result: a vertically integrated pipeline from academic seed-planting through policy production, media amplification, and political mobilization.

No comparable left-of-center infrastructure was built during the same period \autocite{mudge2018}.
This infrastructure asymmetry, not a difference in the quality of ideas, explains much of the rightward shift in American policy common sense between 1980 and 2020.
\end{example}

\begin{principle}[Infrastructure Path Dependence]
\label{prin:path-dependence}
\index{path dependence!infrastructure}
Once built, narrative infrastructure is self-reinforcing: institutions produce alumni who staff other institutions, generate citations that boost credibility, create demand for their own output, and develop donor relationships that ensure continued funding.
Replicating an established infrastructure from scratch requires investment of comparable scale, which is why infrastructure asymmetries persist across decades.
\end{principle}


\section{Platform Infrastructure: The Algorithmic Layer}
\label{sec:ch7-platforms}

Digital platforms add a new layer to narrative infrastructure: the algorithmic layer, which sits between institutional production and audience reception.

\begin{model}[Algorithmic Infrastructure]
\label{mod:algorithmic-infra}
\index{algorithmic infrastructure}
A platform algorithm is a function $A: \narspace \times \mathcal{U} \to [0,1]$ mapping each narrative-user pair $(n, j)$ to a display probability $A(n, j)$.
The algorithm determines the effective exposure matrix:
\[
  e_{jk}^{\text{eff}} = e_{jk} \cdot A(n_k, j)
\]
The platform's objective function (typically engagement maximization) determines $A$ endogenously:
\[
  A^* = \arg\max_A \sum_{j,k} A(n_k, j) \cdot \text{engagement}(n_k, j)
\]
Because high-engagement content tends toward moral outrage, identity threat, and simplistic narratives, engagement maximization systematically distorts the effective narrative infrastructure \autocite{zuboff2019}.
\end{model}

\begin{remark}[Platforms as Invisible Gatekeepers]
\label{rem:invisible-gatekeepers}
Traditional gatekeepers (editors, curriculum committees) are visible and accountable.
Algorithmic gatekeepers are invisible and optimized for a private objective (engagement/revenue).
This means the most consequential infrastructure decisions (which narratives reach which audiences) are made by systems that are neither transparent nor democratically accountable.
From a doxonomic perspective, platform algorithms may be the single most powerful gatekeeper institution in the contemporary information environment, yet they are the least understood and least regulated.
\end{remark}


\section{Notes and References}
\label{sec:ch7-notes}

Media concentration is studied in \textcite{bagdikian2004} and \textcite{mcchesney1999}.
Comparative media systems analysis is developed in \textcite{hallin2004}.
The role of think tanks in narrative production is documented in \textcite{mayer2016}, \textcite{mirowski2009}, and \textcite{phillipsfein2009}.
Platform infrastructure is analyzed in \textcite{zuboff2019} and \textcite{benkler2018}.
The network model of media draws on \textcite{jackson2008}.
The Herfindahl index originates in \textcite{herfindahl1950}.
For public-interest alternatives, see \textcite{benkler2006}.
On the propaganda model's structural filters as infrastructure properties, see \textcite{herman1988}.
The concept of narrative correlation adjusting effective concentration is new to this text but draws on the literature on media parallelism in \textcite{hallin2004}.
On the construction of conservative intellectual infrastructure, see \textcite{mayer2016} and \textcite{phillipsfein2009}; on the absence of a comparable left infrastructure, see \textcite{mudge2018}.


\section{Exercises}

\begin{exercise}
Compute $\ehhi$ for U.S.\ national television news using publicly available audience share data.
Then estimate the average narrative correlation $\bar{\rho}$ (using a qualitative assessment of framing similarity) and compute $\ehhi_{\text{eff}}$.
How much does the correction change the concentration assessment?
\end{exercise}

\begin{exercise}
Model the narrative pipeline from the Atlas Network to public opinion on regulation in a specific country.
Draw the infrastructure graph, estimate the flow weights $w_{uv}$ at each stage, and compute the end-to-end pathway efficiency.
Where is the pipeline most likely to fail?
\end{exercise}

\begin{exercise}
Prove that $\ehhi \geq 1/|V|$ with equality if and only if all institutions have equal throughput.
\end{exercise}

\begin{exercise}
A new social media platform enters the narrative infrastructure with $\reach{} = 0.4$, $\credibility{} = 0.15$, $\persistence{} = 0.05$.
How does this change the overall $\ehhi$?
Does it increase or decrease the conditions for common-sense capture?
(Consider both the direct effect on concentration and the indirect effect on narrative correlation.)
\end{exercise}

\begin{exercise}
\label{ex:ch7-bottleneck}
Identify the three highest-vulnerability nodes in the narrative infrastructure of your country (using the vulnerability index $\nu(v)$).
For each, describe what would happen if that node were captured by a single funder or if it ceased to function.
\end{exercise}

\begin{exercise}
Compare the conservative narrative pipeline (Example~\ref{ex:building-infrastructure}) with the infrastructure available to labor movements or environmental movements.
Draw both pipeline graphs.
Where are the structural asymmetries?
What would it cost to close the infrastructure gap?
\end{exercise}

\begin{exercise}
\label{ex:ch7-algorithm}
Model a platform algorithm that maximizes engagement.
Show that if outrage-inducing content generates higher engagement, the algorithm will systematically amplify narratives with high emotional valence and low nuance.
What does this imply for the narrative entropy $H$ (Definition~\ref{def:narrative-entropy}) of the platform's output?
\end{exercise}

\begin{exercise}
The throughput comparison (Example~\ref{ex:throughput-comparison}) suggests that textbooks are 100 times more cost-effective than cable news.
Why, then, do doxonomic funders invest heavily in cable news?
Give at least three reasons why the throughput calculation does not capture the full strategic picture.
\end{exercise}

\chapter{Common-Sense Capture}
\label{ch:common-sense-capture}

\epigraph{The smart way to keep people passive and obedient is to strictly limit the spectrum of acceptable opinion, but allow very lively debate within that spectrum.}{Noam Chomsky}

\noindent
The most powerful doxonomic achievement is not persuading people to hold a belief.
It is making a belief so pervasive that it ceases to be recognized as a belief at all. A funded narrative becomes invisible infrastructure, indistinguishable from the texture of reality.

This chapter formalizes \emph{common-sense capture}\index{common-sense capture}: the phase transition by which a contested proposition becomes unquestioned background, the mechanisms that produce it, the conditions that sustain it, and the signatures by which it can be detected.


\section{Capture as Phase Transition}
\label{sec:ch8-phase-transition}

\begin{definition}[Common-Sense Capture]
\label{def:capture}
\index{common-sense capture!formal definition}
A belief $p$ is \emph{captured} in a population if:
\begin{enumerate}[nosep]
  \item $\mu(p) > \theta$ for a high threshold $\theta$ (widespread adoption),
  \item $\mathrm{Var}(\pi_j(p)) < \varepsilon$ for small $\varepsilon$ (tight concentration),
  \item the proposition $p$ is not salient as a \emph{claim}. It is experienced as a feature of reality.
\end{enumerate}
Conditions (1) and (2) are measurable.
Condition (3) is the distinctively doxonomic condition: $p$ has exited the space of contestation.
\end{definition}

The transition from ``widely held opinion'' to ``common sense'' is not gradual.
It exhibits threshold behavior analogous to phase transitions in physics.
Below the critical threshold, a belief is debated, contested, and recognized as a position.
Above it, the belief disappears into the background: it is no longer argued for because it no longer needs to be.

\begin{example}[Stages of Capture]
\label{ex:stages-capture}
Consider the proposition ``economic growth is the primary objective of government policy.''
\begin{enumerate}[nosep]
  \item \textbf{Contested (pre-1950s)}: alternative objectives were actively debated, including full employment, social solidarity, moral improvement, imperial expansion, and class harmony.
  \item \textbf{Dominant (1950s--1970s)}: GDP growth became the primary metric, but alternatives (quality of life, distributional equity) were still voiced in mainstream discourse.
  \item \textbf{Common sense (1980s--2000s)}: growth became so dominant that questioning it sounded eccentric. Policy debates assumed growth as the goal and argued only about means.
  \item \textbf{Captured/doxic (2000s--present)}: growth is not argued for; it is the unstated premise of virtually all policy discussion. A politician who says ``we should not pursue growth'' is not engaging in a legitimate policy debate; they are committing a category error.
\end{enumerate}
The transition from stage 2 to stage 4 is what doxonomics calls capture.
\end{example}


\section{Threshold Models}
\label{sec:ch8-thresholds}

Following \textcite{granovetter1978} and \textcite{schelling1978}, we model adoption as a threshold process.
The core insight is that individual adoption decisions depend on how many others have already adopted: people are more willing to hold a belief when they perceive it as socially supported.

\begin{model}[Granovetter-Schelling Belief Adoption]
\label{mod:granovetter}
\index{threshold model}
Each agent $j$ has a threshold $\tau_j \in [0,1]$ drawn from a distribution $G(\tau)$.
Agent $j$ adopts belief $p$ if and only if the fraction of the population already holding $p$ exceeds $\tau_j$:
\[
  \pi_j(p) = 1 \iff \mu(p) \geq \tau_j
\]
The equilibrium population fraction is a fixed point of
\[
  \mu^* = G(\mu^*)
\]
\end{model}

The threshold $\tau_j$ can be interpreted as agent $j$'s ``social proof requirement'': how many people must hold a belief before $j$ is willing to adopt it.
Agents with low thresholds are \emph{early adopters} (they adopt without much social support); agents with high thresholds are \emph{holdouts} (they adopt only when almost everyone else already has).

\begin{theorem}[Tipping Point]
\label{thm:tipping}
\index{tipping point}
If $G$ is continuous and $G(0) > 0$ (some agents adopt unconditionally), then there exists at least one stable equilibrium $\mu^*$.
A doxonomic investment that shifts $G$ downward (lowering thresholds by increasing social pressure, institutional endorsement, or reducing the perceived cost of adoption) can cause a discontinuous jump in $\mu^*$, a \emph{tipping point}.
\end{theorem}

\begin{proof}[Proof sketch]
The function $h(\mu) = G(\mu) - \mu$ is continuous with $h(0) = G(0) > 0$ and $h(1) = G(1) - 1 \leq 0$.
By the intermediate value theorem, $h$ has a zero.
Stability requires $G'(\mu^*) < 1$.
A downward shift in $G$ (uniform reduction in thresholds) raises $G(\mu)$ for all $\mu$, which can push $\mu^*$ past a saddle-node bifurcation, causing the low-adoption equilibrium to disappear and the system to jump to the high-adoption equilibrium.
\end{proof}

\begin{example}[Tipping in Belief Adoption]
\label{ex:tipping}
\index{tipping point!example}
Consider a population with threshold distribution $G(\tau) = \tau^{1/2}$ (many agents have low thresholds; few are extreme holdouts).
The fixed points satisfy $\mu = \mu^{1/2}$, giving $\mu^* = 0$ and $\mu^* = 1$.
The interior dynamics are: for any $\mu > 0$, we have $G(\mu) = \mu^{1/2} > \mu$, so adoption grows.
This system has only one stable equilibrium: $\mu^* = 1$ (complete adoption).
Any nonzero seed of adoption cascades to unanimity.

Now consider $G(\tau) = \tau^2$ (most agents have high thresholds).
The fixed points satisfy $\mu = \mu^2$, giving $\mu^* = 0$ and $\mu^* = 1$.
For $\mu \in (0,1)$, $G(\mu) = \mu^2 < \mu$, so adoption shrinks.
This system has only one stable equilibrium: $\mu^* = 0$ (no adoption).
Even substantial initial adoption erodes to zero.

The doxonomic implication: the \emph{shape} of the threshold distribution (determined by institutional environment, prior narrative exposure, and social structure) determines whether a belief can cascade.
Doxonomic investment works by reshaping $G$: lowering thresholds through institutional endorsement, repetition, and social proof until the system tips.
\end{example}

\begin{remark}[What Determines Thresholds?]
\label{rem:thresholds}
Individual thresholds $\tau_j$ are not fixed parameters.
They depend on:
\begin{itemize}[nosep]
  \item \textbf{Institutional endorsement}: when authoritative institutions adopt $p$, individual thresholds drop (``if Harvard says so, it must be reasonable'').
  \item \textbf{Perceived cost of dissent}: in environments where questioning $p$ is socially costly (mockery, professional risk, ostracism), thresholds drop (people adopt to avoid punishment, not because they are convinced).
  \item \textbf{Narrative coherence}: when $p$ fits into an existing ideological bundle (Definition~\ref{def:ideology}), thresholds drop because adoption is cognitively cheap.
  \item \textbf{Prior exposure}: repeated exposure to $p$ through multiple channels lowers thresholds via the illusory truth effect.
  \item \textbf{Identity alignment}: when $p$ is framed as aligned with a valued identity, thresholds drop; when framed as threatening, thresholds rise.
\end{itemize}
Doxonomic investment works on thresholds through all five mechanisms simultaneously.
\end{remark}


\section{The Doxonomic Investment Needed for Capture}
\label{sec:ch8-investment}

How much does capture cost?
This depends on the initial threshold distribution, the available channels, and the target belief.

\begin{model}[Capture Cost]
\label{mod:capture-cost}
\index{capture cost}
Let $G_0(\tau)$ be the initial threshold distribution and $G_F(\tau)$ the distribution after doxonomic investment $F$.
Assume investment uniformly lowers thresholds: $G_F(\tau) = G_0(\tau - \alpha F)$ for some effectiveness parameter $\alpha$.
The minimum investment $F^*$ needed for capture is the smallest $F$ such that the high-adoption fixed point of $\mu = G_F(\mu)$ exists and is globally stable.

At the tipping point, $G_F$ must be tangent to the 45-degree line from below:
\[
  G_F(\mu^*) = \mu^* \quad \text{and} \quad G_F'(\mu^*) = 1
\]
These two equations determine $(\mu^*, F^*)$.
\end{model}

\begin{remark}
The capture cost $F^*$ depends critically on the initial threshold distribution.
In populations with many low-threshold agents (already predisposed by prior narrative exposure or identity alignment), $F^*$ is small: a nudge suffices.
In populations with many high-threshold agents (resistant due to counter-institutions, lived experience, or strong alternative narratives), $F^*$ can be enormous.
This explains why the same doxonomic operation can succeed in one context and fail in another: the target population's threshold distribution matters as much as the budget.
\end{remark}


\section{Mechanisms of Capture}
\label{sec:ch8-mechanisms}

Capture does not happen through a single mechanism.
It typically involves the coordinated operation of several:

\subsection{Institutional Saturation}

When enough institutions treat $p$ as an operational premise, the cumulative effect is overwhelming.
A student who encounters $p$ in economics class, in management training, in the workplace, in advertising, in entertainment, and in political discourse receives $p$ from so many independent-seeming sources that it appears to be a natural fact rather than a specific claim.

\begin{definition}[Institutional Saturation]
\label{def:saturation}
\index{institutional saturation}
A belief $p$ is \emph{institutionally saturated} if the fraction of institutional contexts in which $p$ functions as an operational premise exceeds a threshold $\sigma$:
\[
  \frac{|\{U_i \in \mathcal{U} : p \in \sigma_i\}|}{|\mathcal{U}|} > \sigma
\]
where $\mathcal{U}$ is the cover of social life by institutional contexts and $\sigma_i$ is the local belief system of context $U_i$ (Chapter~\ref{ch:contradiction-coherence}).
\end{definition}

Institutional saturation is related to but distinct from population-level adoption.
A belief can be institutionally saturated ($p$ is embedded in workplaces, schools, and media) without being unanimously held (many individuals privately doubt $p$).
This is the pluralistic-ignorance configuration (Definition~\ref{def:pluralistic-ignorance}): the institutional environment enforces $p$ as operational reality even where private belief is weak.

\subsection{Crowding Out Alternatives}

Capture involves not just amplifying $p$ but \emph{suppressing alternatives}.
When the narrative infrastructure is saturated with $p$, there is no institutional space for $\neg p$: no funding, no platforms, no expert endorsement, no curriculum presence.
The alternatives do not need to be explicitly censored; they are simply not produced at sufficient volume to be encountered.

\begin{definition}[Narrative Crowding Out]
\label{def:crowding-out}
\index{crowding out}
A belief $p$ \emph{crowds out} alternative $q$ if the doxonomic investment in $p$ reduces the effective throughput for $q$:
\[
  \frac{\partial \Theta(q)}{\partial F(p)} < 0
\]
This occurs when $p$ and $q$ compete for the same infrastructure (media attention, curriculum space, expert endorsement, public attention) and $p$'s dominance leaves insufficient residual capacity for $q$.
\end{definition}

\subsection{Normalization Through Repetition}

Sheer repetition makes a claim feel true (the illusory truth effect).
When $p$ is encountered daily across multiple channels, it acquires the phenomenological quality of familiarity, which the cognitive system interprets as evidence for truth.
This is not a rational process; it is a cognitive heuristic that doxonomic systems exploit.

\subsection{Social Proof Cascades}

Once a critical mass of people (or, more importantly, a critical mass of \emph{visible} people) endorses $p$, social proof dynamics take over.
Each new adopter makes $p$ appear more normal, which lowers the threshold for the next adopter, which makes $p$ appear even more normal, a positive feedback cascade that accelerates toward capture.

\begin{model}[Social Proof Cascade]
\label{mod:cascade}
\index{social proof cascade}
Let $\mu(t)$ be the fraction holding $p$ at time $t$, with dynamics:
\[
  \frac{d\mu}{dt} = \alpha\phi(t)(1-\mu) + \beta\mu(1-\mu) - \delta\mu
\]
where $\alpha\phi(t)$ is institutional production, $\beta$ is the social proof coefficient, and $\delta$ is natural decay.
Setting $d\mu/dt = 0$, the non-trivial equilibria satisfy:
\[
  \mu^* = \frac{(\alpha\phi + \beta - \delta) \pm \sqrt{(\alpha\phi + \beta - \delta)^2 + 4\beta\alpha\phi}}{2\beta}
\]
When $\beta$ is large (strong social proof), the system is bistable: there is a low-$\mu$ equilibrium and a high-$\mu$ equilibrium, separated by an unstable threshold.
Doxonomic investment (increasing $\phi$) pushes the system across the threshold and into the high-$\mu$ basin.
\end{model}


\section{Path Dependence and Lock-In}
\label{sec:ch8-lock-in}

Once a belief crosses the tipping point, it becomes self-reinforcing through the feedback mechanisms described in Chapter~\ref{ch:belief-production-chain}: the belief generates behaviors that generate institutions that generate funding that generates more of the same belief.

\begin{definition}[Doxonomic Lock-In]
\label{def:lock-in}
\index{lock-in}
A belief $p$ is \emph{locked in} if the doxonomic equilibrium $\mu^*(p)$ is locally asymptotically stable and the basin of attraction is large, meaning that small perturbations (counter-narratives, contrary evidence, temporary funding cuts) do not dislodge it.
\end{definition}

\begin{proposition}[Lock-In Condition]
\label{prop:lock-in}
Lock-in occurs when the self-reinforcement gain exceeds the natural decay rate:
\[
  \frac{\partial \mu^{(t+1)}}{\partial \mu^{(t)}} \bigg|_{\mu = \mu^*} > 1 - \delta
\]
where $\delta$ is the rate at which belief decays absent reinforcement.
\end{proposition}

Lock-in has a temporal dimension.
The longer a belief has been locked in, the harder it is to dislodge, because:
\begin{itemize}[nosep]
  \item institutional infrastructure has been built around it (path dependence),
  \item professional careers depend on it (vested interests),
  \item curricula have transmitted it to new generations (generational reproduction),
  \item alternative beliefs have atrophied for lack of institutional support (crowding out),
  \item the belief has shaped behavior in ways that generate confirming evidence (self-fulfillment).
\end{itemize}

\begin{example}[Lock-In: Homo Economicus]
\label{ex:lock-in-homo-economicus}
The selfish-rational-agent model (homo economicus) has been locked in as the operational anthropology of economic policy since approximately the 1980s.
Dislodging it would require simultaneous intervention at multiple points:
\begin{itemize}[nosep]
  \item \textbf{Curricula}: rewriting economics textbooks to foreground conditional cooperation, behavioral complexity, and institutional context.
  \item \textbf{Professional norms}: changing what counts as ``rigorous'' economic analysis in journals, tenure review, and policy consulting.
  \item \textbf{Policy design}: replacing incentive-based regulation with trust-based institutional design.
  \item \textbf{Management theory}: replacing principal-agent models with cooperative organizational theories.
  \item \textbf{Popular culture}: producing narratives in which cooperative behavior is normal, not exceptional.
\end{itemize}
No single intervention suffices because the lock-in operates at all five levels simultaneously.
This is why lock-in is so durable: disrupting one level is compensated by the other four.
\end{example}


\section{Saturation and Diminishing Returns}
\label{sec:ch8-saturation-returns}

\begin{proposition}[Diminishing Returns Post-Capture]
\label{prop:diminishing}
Once common-sense capture is achieved ($\mu(p) > \theta$), the marginal belief change per additional dollar of funding approaches zero:
\[
  \lim_{\mu \to 1} \frac{\partial \mu}{\partial F} = 0
\]
Post-capture, doxonomic spending shifts from persuasion to \emph{maintenance}: preventing erosion rather than producing new converts.
\end{proposition}

This has several practical implications:

\begin{enumerate}[nosep]
  \item \textbf{Maintenance is cheaper than capture.}
  The cost of maintaining a captured belief is typically an order of magnitude lower than the cost of capturing it.
  Once the infrastructure is built, the institutions staffed, and the curricula written, the system runs on inertia plus modest annual investment.
  This asymmetry explains why dominant narratives persist long after the original funding rationale has changed.

  \item \textbf{Over-investment is common.}
  Funders often continue spending at capture-level rates even after capture has been achieved, because they cannot easily measure the current belief level and err on the side of caution.
  This produces \emph{doxonomic waste}: resources spent reinforcing beliefs that are already unquestioned.

  \item \textbf{The marginal dollar is better spent on new frontiers.}
  Once $p$ is captured, a rational doxonomic investor should redirect resources to adjacent beliefs ($p_2, p_3, \ldots$) that extend the ideological bundle, rather than pouring more money into $p$ itself.
  This is the doxonomic analogue of market saturation in product economics.
\end{enumerate}


\section{Detecting Capture: Empirical Signatures}
\label{sec:ch8-detection}

How does a researcher identify that a belief has been captured rather than merely being popular?
The detection methods from Section~\ref{sec:ch2-detecting-doxa} apply, but we can add quantitative signatures:

\begin{definition}[Capture Signature]
\label{def:capture-signature}
\index{capture!detection}
A belief $p$ shows a \emph{capture signature} if:
\begin{enumerate}[nosep]
  \item \textbf{Asymmetric argumentation}: $p$ appears overwhelmingly as a premise rather than a conclusion in public discourse (measurable via corpus analysis).
  \item \textbf{Low narrative entropy}: the narrative space around $p$'s domain has collapsed to a small number of frames ($H(t)$ below baseline; see Definition~\ref{def:narrative-entropy}).
  \item \textbf{High institutional saturation}: $p$ functions as an operational assumption in a large fraction of institutional contexts (Definition~\ref{def:saturation}).
  \item \textbf{Anomaly response}: challenges to $p$ are met with confusion or ridicule rather than counter-argument.
  \item \textbf{Funding-belief decoupling}: $p$'s prevalence persists even during periods of reduced doxonomic investment (the belief runs on institutional inertia, not active funding).
  \item \textbf{Cross-domain consistency}: $p$ appears in economics, management, law, media, and everyday speech without any domain citing the others as its source.
\end{enumerate}
\end{definition}

\begin{remark}
Not all captured beliefs are false, and not all uncaptured beliefs are true.
The capture signature identifies a \emph{structural} property of the belief's social position, not an epistemic property of its content.
``The earth orbits the sun'' is also captured (near-universal, low variance, invisible as a claim), but it is captured \emph{because it is true and well-evidenced}, not because of doxonomic investment.
The researcher must distinguish evidence-driven capture from investment-driven capture, which requires the separation of audit and evaluation (Principle~\ref{prin:separation}).
\end{remark}


\section{Historical Case: The Capture of ``Markets Know Best''}
\label{sec:ch8-case}

The belief that markets are generally superior to government in allocating resources provides a detailed case study of the capture process.

\textbf{Phase 1: Contestation (1930s--1960s).}
In the aftermath of the Great Depression, market-skeptical views dominated: Keynesian macroeconomics, strong regulation, public ownership of utilities, and active industrial policy were mainstream.
The belief ``markets know best'' was a minority position, associated with a small network of Austrian and Chicago economists.

\textbf{Phase 2: Infrastructure Building (1947--1980).}
The Mont P\`{e}lerin Society (1947), the Foundation for Economic Education, and later the Heritage Foundation (1973), Cato Institute (1977), and dozens of smaller think tanks built the narrative infrastructure.
Funding from Olin, Scaife, Koch, and Bradley foundations supported law-and-economics programs at universities, endowed chairs, and produced a generation of economists, lawyers, and policy intellectuals committed to market-first thinking \autocite{mirowski2009, slobodian2018}.

\textbf{Phase 3: Political Capture (1979--1990).}
Thatcher and Reagan brought the narrative to political power.
Deregulation, privatization, and tax cuts were enacted.
The behavioral consequences (rising stock markets, new consumer products, visible wealth creation for some) appeared to confirm the narrative.

\textbf{Phase 4: Common-Sense Capture (1990s--2008).}
By the 1990s, even center-left parties (Clinton's Democrats, Blair's New Labour) adopted market-first assumptions.
The ``Washington Consensus'' made market liberalization the default prescription for developing countries.
Alternatives (industrial policy, public ownership, cooperative models) were no longer debated; they were simply absent from serious policy discussion.

\textbf{Phase 5: Partial Fracture (2008--present).}
The financial crisis disrupted the capture by generating a massive contradiction (markets spectacularly failed to self-correct).
But the fracture was partial: despite the crisis, the policy response (austerity, quantitative easing for banks, limited re-regulation) largely preserved market-first assumptions.
The belief was shaken but not dislodged: evidence of deep lock-in.

\begin{remark}
This case illustrates the full capture trajectory: from minority position to infrastructure building to political implementation to common-sense invisibility to partial fracture under contradiction.
The total doxonomic investment over 60 years was plausibly in the billions of dollars.
The capture was not permanent (the 2008 crisis opened cracks), but the lock-in was strong enough to survive even a dramatic empirical refutation.
\end{remark}


\section{Escape from Capture}
\label{sec:ch8-escape}

If common-sense capture is so durable, how does it ever end?
Chapter~\ref{ch:how-belief-regimes-fracture} treats this in detail, but a preview is useful here.

\begin{conjecture}[Conditions for De-Capture]
\label{conj:de-capture}
\index{de-capture}
A captured belief $p$ can be dislodged when:
\begin{enumerate}[nosep]
  \item \textbf{Contradiction accumulates}: the gap between $p$ and lived experience becomes too large for institutional maintenance to bridge (e.g., ``markets self-correct'' during a prolonged recession).
  \item \textbf{Counter-infrastructure exists}: alternative institutions have been built that can supply a replacement narrative at sufficient volume (e.g., the environmental movement built enough infrastructure to challenge growth-as-progress).
  \item \textbf{Generational turnover}: new cohorts arrive without the formative exposure that locked in $p$ for the previous generation.
  \item \textbf{Elite defection}: high-credibility insiders publicly break with $p$, reducing the social cost of dissent (e.g., prominent economists questioning homo economicus after 2008).
  \item \textbf{Crisis opens the Overton window}: an exogenous shock temporarily suspends the normal inertia of common sense (Chapter~\ref{ch:crisis-doxonomics}).
\end{enumerate}
Typically, multiple conditions must coincide.
Single-condition de-capture is rare; this is why lock-in is so durable.
\end{conjecture}


\section{Notes and References}
\label{sec:ch8-notes}

Threshold models originate with \textcite{granovetter1978} and \textcite{schelling1978}.
Cascade and herding models appear in \textcite{banerjee1992} and \textcite{bikhchandani1992}.
Path dependence in institutional contexts is treated in \textcite{north1990}.
The concept of lock-in from technology economics \autocite{arthur1989} transfers to the belief domain.
On the invisibility of dominant frames, see \textcite{lukes2005}.

The capture of market-first ideology is documented in \textcite{mirowski2009}, \textcite{harvey2005}, \textcite{slobodian2018}, and \textcite{phillipsfein2009}.
On the illusory truth effect and repetition-based persuasion, see \textcite{kahneman2011}.
Social proof dynamics are analyzed in \textcite{cialdini2006}.
On the partial fracture of market fundamentalism after 2008, see \textcite{streeck2016}.
For the concept of crowding out in media and attention markets, see \textcite{mcchesney1999}.
On narrative entropy as a measure of discursive openness, see the information-theoretic framework in Chapter~\ref{ch:text-media}.
The concept of institutional saturation extends the Althusserian notion of ideological state apparatuses \autocite{althusser1970} into a measurable condition.


\section{Exercises}

\begin{exercise}
Given a threshold distribution $G(\tau) = \tau^2$ on $[0,1]$, find all fixed points of $\mu^* = G(\mu^*)$.
Determine their stability.
What is the minimum uniform threshold reduction $\Delta$ needed to create a stable high-adoption equilibrium?
Plot $G(\tau)$ and the shifted $G(\tau - \Delta)$ to illustrate.
\end{exercise}

\begin{exercise}
For the social proof cascade model (Model~\ref{mod:cascade}) with $\alpha = 0.3$, $\beta = 0.5$, $\delta = 0.1$, find the minimum institutional investment $\phi^*$ needed for a stable high-adoption equilibrium to exist.
Plot the phase portrait for $\phi$ just below and just above $\phi^*$.
\end{exercise}

\begin{exercise}
Estimate the maintenance cost of the belief ``economic growth is always desirable'' by identifying three institutions that continuously reinforce it and estimating their annual budgets.
Compare the maintenance cost to a rough estimate of the original capture cost (infrastructure built from the 1940s onward).
\end{exercise}

\begin{exercise}
Design a formal test to distinguish a belief that is common sense (Definition~\ref{def:common-sense}) from one that is captured (Definition~\ref{def:capture}).
What empirical evidence would differentiate the two?
Apply your test to ``democracy is the best form of government.''
\end{exercise}

\begin{exercise}
\label{ex:ch8-crowding}
Model the crowding-out effect formally.
Let two beliefs $p$ and $q$ compete for a shared narrative infrastructure with total capacity $\Theta_{\max}$.
Show that if $\Theta(p) + \Theta(q) \leq \Theta_{\max}$, increased investment in $p$ necessarily reduces throughput for $q$.
Under what conditions can $q$ survive despite being outspent?
\end{exercise}

\begin{exercise}
Apply the capture signature (Definition~\ref{def:capture-signature}) to the belief ``a university degree is necessary for a good career.''
Assess each of the six criteria.
Is this belief captured?
If so, who benefits from its capture?
\end{exercise}

\begin{exercise}
\label{ex:ch8-de-capture}
Choose a belief that has been partially de-captured in the past 20 years (e.g., ``smoking is glamorous,'' ``homosexuality is a disorder,'' ``colonialism was civilizing'').
Trace the de-capture process through the five conditions in Conjecture~\ref{conj:de-capture}.
Which conditions were most important?
Was the de-capture complete or partial?
\end{exercise}

\begin{exercise}
The chapter claims that ``maintenance is cheaper than capture.''
Formalize this claim using the social proof cascade model.
Show that the per-period funding needed to maintain $\mu^*$ near 1 is lower than the total funding needed to reach $\mu^*$ from $\mu = 0$.
Under what conditions is this \emph{not} true?
\end{exercise}

\chapter{Doxonomic Power}
\label{ch:doxonomic-power}

\epigraph{The most effective way to destroy people is to deny and obliterate their own understanding of their history.}{George Orwell}

\noindent
Power in doxonomics is not the capacity to coerce but the capacity to set the conditions under which people form their beliefs.
This chapter formalizes doxonomic power, extends Lukes' three dimensions, and shows how distributed systems can produce coordinated ideological effects without central planning.

\section{Three Dimensions, Extended}
\label{sec:ch9-lukes}

\textcite{lukes2005} distinguished three faces of power.
We extend each to the doxonomic domain.

\begin{definition}[Dimensions of Doxonomic Power]
\label{def:doxonomic-power}
\index{doxonomic power}
\begin{enumerate}[nosep]
  \item \textbf{First dimension}: the capacity to win arguments in public debate (persuasion within existing frames).
  \item \textbf{Second dimension}: the capacity to control which arguments reach public debate (agenda-setting, narrative infrastructure control).
  \item \textbf{Third dimension}: the capacity to shape what people want, expect, and consider possible (prior formation, anthropological regime installation).
  \item \textbf{Fourth dimension} (doxonomic extension): the capacity to determine the \emph{epistemic rules}---what counts as evidence, who counts as an expert, which methods are legitimate.
\end{enumerate}
\end{definition}

The fourth dimension is the distinctively doxonomic contribution.
It is power over the \emph{meta-level}: not just shaping beliefs, but shaping the processes by which beliefs are evaluated.

\begin{proposition}[Power Ordering]
\label{prop:power-ordering}
The dimensions form a hierarchy of cost-effectiveness.
Fourth-dimensional power is the most cost-effective because it shapes all subsequent belief formation:
\[
  \iroi_4 > \iroi_3 > \iroi_2 > \iroi_1
\]
where $\iroi_d$ is the ideological return on investment at dimension $d$.
\end{proposition}

\section{Distributed Coordination Without Conspiracy}
\label{sec:ch9-distributed}

A common objection to doxonomic analysis is that it requires conspiracy: a secret group deliberately coordinating belief production.
This is a misunderstanding.

\begin{principle}[Structural Alignment]
\label{prin:structural-alignment}
\index{structural alignment}
A doxonomic system can produce coordinated ideological output without centralized planning if:
\begin{enumerate}[nosep]
  \item funders share material interests (e.g., reduced regulation),
  \item institutions face similar incentive structures (e.g., dependence on the same funding sources),
  \item professional norms align outputs (e.g., economics training produces a shared worldview),
  \item feedback loops reward convergence (e.g., media outlets that deviate lose audience or advertisers).
\end{enumerate}
\end{principle}

\textcite{herman1988} demonstrated this with the propaganda model: five structural filters---ownership, advertising, sourcing, flak, and ideology---produce systematic media bias without requiring any explicit coordination.

\begin{model}[Emergent Alignment]
\label{mod:emergent}
Let $n$ institutions each independently choose a narrative position $n_k \in \narspace$ to maximize revenue $r_k(n_k, \bm{n}_{-k})$ given the positions of others.
If revenue functions are correlated (because funders, advertisers, and audiences overlap), the Nash equilibrium $\bm{n}^*$ will cluster in a narrow region of $\narspace$---producing apparent coordination from independent profit maximization.
\end{model}

\section{Power as Prior-Setting}
\label{sec:ch9-priors}

The deepest form of doxonomic power is the ability to set the \emph{prior distribution} from which a population evaluates new information.

\begin{definition}[Prior-Setting Power]
\label{def:prior-setting}
\index{prior-setting power}
An agent $a$ has \emph{prior-setting power} over a population if $a$ can influence the initial distribution $\pi_j^{(0)}$ for agents $j$ in the population, typically through control of educational institutions, childhood media, or cultural narratives.
\end{definition}

Prior-setting power is the most durable form of doxonomic power because Bayesian agents with different priors will interpret the same evidence differently, and prior formation occurs early in life when resistance is lowest.

\section{Epistemic Injustice as Doxonomic Power}
\label{sec:ch9-epistemic-injustice}

\textcite{fricker2007} identified a form of power closely related to the fourth dimension: \emph{epistemic injustice}\index{epistemic injustice}, in which certain speakers are systematically denied credibility due to their social identity.
In doxonomic terms, epistemic injustice is a mechanism by which the credibility weights $\credibility{k}$ are allocated not by evidence quality but by the speaker's position in hierarchies of race, gender, class, or nationality.

This means that some agents face a structural credibility deficit: their testimony about lived experience is discounted, their counter-narratives are dismissed as ``anecdotal,'' and their epistemic labor is uncompensated.
Epistemic injustice is a doxonomic power that operates through the fourth dimension: it does not merely suppress specific beliefs but shapes the rules of whose evidence counts.

\section{Notes and References}
\label{sec:ch9-notes}

\textcite{lukes2005} provides the three-dimensional framework.
The propaganda model is in \textcite{herman1988}.
Structural power without conspiracy is analyzed in \textcite{gramsci1971} and formalized in network terms by \textcite{jackson2008}.
On the role of economics in shaping priors, see \textcite{bowles2011} and \textcite{mirowski2009}.
On epistemic injustice, see \textcite{fricker2007} and \textcite{medina2013}.
On colonial knowledge production as an exercise of doxonomic power, see \textcite{said1978}.

\section{Exercises}

\begin{exercise}
Classify the following exercises of power by dimension: (a)~a pharmaceutical company funding a medical journal, (b)~a government classifying documents, (c)~an economics curriculum teaching rational self-interest as human nature, (d)~a platform algorithm determining what counts as ``trending.''
\end{exercise}

\begin{exercise}
Using Model~\ref{mod:emergent}, show that if all institutions depend on the same advertising market, the equilibrium narrative range is narrower than if institutions have diverse funding sources.
\end{exercise}

\begin{exercise}
Formalize the claim that prior-setting power compounds over generations.
If each generation's priors are a function of the previous generation's posteriors and the institutional environment, write the recurrence relation and identify conditions for convergence.
\end{exercise}

\chapter{The Economics of Belief Production}
\label{ch:economics-belief}

\epigraph{There is no such thing as a free market.}{Ha-Joon Chang}

\noindent
Belief production has costs, returns, market structures, and welfare effects.
This chapter applies economic analysis to the doxonomic system itself, asking: what does belief cost, who earns the returns, and what are the welfare consequences of concentrated reality production?

\section{Cost Functions for Narrative Repetition}
\label{sec:ch10-costs}

\begin{model}[Stylized Doxonomic Cost Function]
\label{mod:cost-function}
\index{cost function}
A simple candidate cost function for achieving belief level $\mu$ for proposition $p$ through channel $k$ is
\[
  c_k(\mu) = \frac{-\ln(1-\mu)}{\credibility{k} \cdot \reach{k}}
\]
This is a heuristic, not an empirically derived function.
Its justification is qualitative: the last few percentage points of adoption are exponentially more expensive than the first, because holdouts are more resistant, more isolated, or more committed to alternatives.
The functional form is chosen for tractability.
Empirical estimation of actual cost curves is an open problem for the field.
\end{model}

The total cost of achieving population-level belief $\mu(p) = \bar{\mu}$ using $m$ channels is the solution to
\[
  \min_{\mu_1, \ldots, \mu_m} \sum_{k=1}^m c_k(\mu_k) \quad \text{subject to} \quad \sum_k w_k \mu_k = \bar{\mu}
\]
where $w_k$ is the population weight of group $k$.

\section{Ideological Return on Investment}
\label{sec:ch10-iroi}

\begin{definition}[Ideological ROI]
\label{def:iroi}
\index{ideological ROI}
The \emph{ideological return on investment} for belief $p$ through channel $k$ is
\[
  \iroi_k(p) = \frac{\Delta \belief{p}}{\Delta \funding{k}} = \frac{\credibility{k} \cdot \relevance{k}{p} \cdot \reach{k} \cdot \persistence{k}}{c_k'(\mu)}
\]
evaluated at the current belief level.
The aggregate $\iroi$ is the weighted average across all active channels.
\end{definition}

\begin{proposition}[Decreasing IROI]
\label{prop:decreasing-iroi}
$\iroi_k(p)$ is decreasing in $\mu(p)$: beliefs that are already widely held offer lower marginal returns to additional investment.
This creates an incentive to diversify across beliefs or to shift from persuasion to maintenance once capture is achieved.
\end{proposition}

\section{Belief as Public Good or Public Bad}
\label{sec:ch10-public-good}

Beliefs are non-excludable (you cannot prevent someone from adopting a freely circulating narrative) and partially non-rival (one person's belief does not deplete the narrative).
This gives beliefs the structure of a \emph{public good}---or, when harmful, a \emph{public bad}.

\begin{definition}[Doxonomic Externality]
\label{def:doxonomic-externality}
\index{doxonomic externality}
A \emph{doxonomic externality} arises when a subsidized belief $p$ imposes costs on agents who did not choose to produce or receive it.
Examples include: reduced social trust from selfishness narratives, increased tolerance for inequality from meritocratic myths, or reduced collective action capacity from individualist anthropologies.
\end{definition}

\begin{theorem}[Welfare Loss from Doxonomic Externalities]
\label{thm:welfare-loss}
If the social cost of belief $p$ is $D(\mu(p))$ (increasing and convex in adoption level) and private funders do not internalize this cost, the equilibrium adoption level $\mu^*$ exceeds the social optimum $\mu^{**}$:
\[
  \mu^* > \mu^{**} \quad \text{where} \quad \mu^{**} = \arg\min_{\mu} \left[ c(\mu) + D(\mu) \right]
\]
The welfare loss is $\int_{\mu^{**}}^{\mu^*} [D'(\mu) - c'(\mu)] \, d\mu > 0$.
\end{theorem}

\section{Epistemic Herfindahl Index}
\label{sec:ch10-ehhi}

\begin{definition}[EHI]
\label{def:ehhi}
\index{epistemic Herfindahl index!formal}
The \emph{epistemic Herfindahl index} is
\[
  \ehhi = \sum_{k=1}^{m} s_k^2 \qquad \text{where} \qquad s_k = \frac{\Theta_k}{\sum_{j} \Theta_j}
\]
and $\Theta_k = \funding{k} \cdot \credibility{k} \cdot \reach{k}$ is the effective doxonomic throughput of channel $k$.
\end{definition}

High $\ehhi$ indicates a concentrated market for reality production, in which a small number of institutions control most of the population's narrative exposure.
In standard antitrust economics, $\mathrm{HHI} > 0.25$ signals high concentration.
Whether the same threshold applies to belief production is an empirical question---narrative markets differ from commodity markets in credibility dynamics, network effects, and audience resistance.
We adopt $\ehhi > 0.25$ as a provisional warning indicator, not as a settled regulatory standard.

\section{Subsidy Structures}
\label{sec:ch10-subsidies}

Much doxonomic spending is invisible because it is structured as tax-exempt philanthropy, corporate sponsorship, or institutional cross-subsidy.

\begin{model}[Hidden Doxonomic Transfer]
\label{mod:hidden-transfer}
Let a corporation spend $S$ on think-tank funding.
The effective doxonomic transfer is $S \cdot (1 - t)$ where $t$ is the tax rate on the deduction.
Society effectively co-finances the doxonomic investment through foregone tax revenue.
The net public cost is $S \cdot t$ for a narrative chosen by the funder.
\end{model}

\section{Notes and References}
\label{sec:ch10-notes}

The economics of public goods and externalities draws on \textcite{stiglitz2012} and \textcite{arrow1963}.
The market for ideas framework is from \textcite{coase1974}.
Media market concentration is analyzed in \textcite{bagdikian2004}.
The Herfindahl index originates in \textcite{herfindahl1950}.
On manipulation as an economic phenomenon, see \textcite{akerlof2015}.
The political economy of think-tank funding is documented in \textcite{mayer2016}.

\section{Exercises}

\begin{exercise}
Compute $\ehhi$ for the following hypothetical market: three channels with throughputs $\Theta_1 = 50$, $\Theta_2 = 30$, $\Theta_3 = 20$.
Is this market concentrated by doxonomic standards?
\end{exercise}

\begin{exercise}
A funder can invest \$10M in advertising ($\credibility{} = 0.2$, $\reach{} = 0.6$) or in university research ($\credibility{} = 0.85$, $\reach{} = 0.05$).
Compute $\iroi$ for each.
Under what conditions is the university investment more effective despite lower reach?
\end{exercise}

\begin{exercise}
Model the welfare loss from the belief ``poverty reflects personal failure'' using a specific damage function $D(\mu) = \alpha \mu^2$.
Find the socially optimal belief level and the deadweight loss.
\end{exercise}

\begin{exercise}
Estimate the hidden public subsidy to ideological production through the U.S.\ 501(c)(3) tax exemption.
What total doxonomic spending does this structure enable?
\end{exercise}

\chapter{Formal Models of Doxonomic Systems}
\label{ch:formal-models}

\epigraph{A theory has only the alternative of being right or wrong. A model has a third possibility: it may be right, but irrelevant.}{Manfred Eigen}

\noindent
This chapter develops the mathematical core of doxonomics: stock-and-flow models, epidemiological diffusion, Bayesian updating under unequal exposure, cascade models, and differential equation frameworks for belief dynamics.

\section{Stock-and-Flow Models}
\label{sec:ch11-stock-flow}

\begin{definition}[Belief Stock and Narrative Flow]
\label{def:stock-flow}
\index{stock-and-flow model}
The \emph{belief stock} $\belief{i}(t) \in [0,1]$ is the fraction of the population holding belief $i$ at time $t$.
The \emph{narrative flow} $\phi_i(t) \in \R_{\geq 0}$ is the rate at which narrative content supporting $i$ is produced and circulated.
The dynamics are:
\begin{equation}
\label{eq:stock-flow}
  \frac{d\belief{i}}{dt} = \alpha \cdot \phi_i(t) \cdot (1 - \belief{i}(t)) - \delta \cdot \belief{i}(t) + \beta \cdot \belief{i}(t)(1 - \belief{i}(t))
\end{equation}
where $\alpha$ is the persuasion rate, $\delta$ is the natural decay rate, and $\beta$ is the social reinforcement rate (peer effects).
\end{definition}

\begin{proposition}[Steady States]
\label{prop:steady-states}
Equation~\eqref{eq:stock-flow} has at most three steady states.
Setting $d\belief{i}/dt = 0$, the trivial steady state $\belief{i}^* = 0$ exists when $\phi_i = 0$.
For $\phi_i > 0$ and $\beta > 0$, there exist two non-trivial steady states:
\[
  \belief{i}^* = \frac{(\alpha\phi_i + \beta - \delta) \pm \sqrt{(\alpha\phi_i + \beta - \delta)^2 + 4\beta\alpha\phi_i}}{2\beta}
\]
The system exhibits bistability when $\beta$ is sufficiently large relative to $\delta$, corresponding to the lock-in phenomenon of Chapter~\ref{ch:common-sense-capture}.
\end{proposition}

\section{Epidemiological Diffusion Models}
\label{sec:ch11-sir}

Belief spread shares structural features with disease transmission.

\begin{model}[SIR Model for Belief]
\label{mod:sir-belief}
\index{SIR model}
Partition the population into three compartments:
\begin{itemize}[nosep]
  \item $S(t)$: susceptible (have not encountered the narrative),
  \item $I(t)$: infected (actively hold and transmit the belief),
  \item $R(t)$: recovered (have encountered and rejected the belief, or lost interest).
\end{itemize}
The dynamics are:
\begin{align}
  \frac{dS}{dt} &= -\beta S I - \alpha \phi(t) S \label{eq:sir-s}\\
  \frac{dI}{dt} &= \beta S I + \alpha \phi(t) S - \gamma I \label{eq:sir-i}\\
  \frac{dR}{dt} &= \gamma I \label{eq:sir-r}
\end{align}
where $\beta$ is the peer transmission rate, $\alpha\phi(t)$ is the institutional transmission rate (from narrative infrastructure), and $\gamma$ is the recovery rate (skepticism, counter-narrative exposure).
\end{model}

\begin{definition}[Basic Reproduction Number]
\label{def:R0}
\index{basic reproduction number}
The \emph{doxonomic reproduction number} is
\[
  R_0 = \frac{\beta + \alpha\phi}{\gamma}
\]
If $R_0 > 1$, the belief can spread epidemically.
If $R_0 < 1$, it dies out without sustained institutional support ($\phi > 0$).
\end{definition}

A crucial difference from epidemiology: institutions can maintain $\phi > 0$ indefinitely, sustaining beliefs that would otherwise die out.
This is the formal basis of the claim that some beliefs persist not because they are persuasive but because they are subsidized.

\section{Bayesian Updating Under Unequal Exposure}
\label{sec:ch11-bayesian}

\begin{model}[Asymmetric Information Environment]
\label{mod:asymmetric-bayes}
\index{Bayesian updating!asymmetric}
Agent $j$ observes a sequence of signals $\{s_t^{(j)}\}_{t=1}^{T}$ where each $s_t^{(j)}$ is drawn from
\[
  s_t^{(j)} \sim \begin{cases} \nu^+_j & \text{with probability } q_j \\ \nu^-_j & \text{with probability } 1 - q_j \end{cases}
\]
and $q_j$ depends on $j$'s media environment.
Agents in channel $k$ have $q_j = q_k$, determined by the channel's editorial stance.
After $T$ signals, the posterior is
\[
  \pi_j^{(T)}(p) = \frac{\pi_j^{(0)}(p) \cdot \prod_{t=1}^T L(s_t^{(j)} \mid p)}{\pi_j^{(0)}(p) \cdot \prod_{t=1}^T L(s_t^{(j)} \mid p) + (1-\pi_j^{(0)}(p)) \cdot \prod_{t=1}^T L(s_t^{(j)} \mid \neg p)}
\]
\end{model}

\begin{theorem}[Polarization Under Asymmetric Exposure]
\label{thm:polarization}
\index{polarization}
If two groups $A$ and $B$ have $q_A > 1/2$ and $q_B < 1/2$ (their channels systematically favor opposite conclusions), then as $T \to \infty$:
\[
  \pi_A^{(T)}(p) \to 1 \quad \text{and} \quad \pi_B^{(T)}(p) \to 0
\]
almost surely, even if both groups are perfectly rational Bayesian updaters.
\end{theorem}

This theorem is the formal foundation for the claim that doxonomic polarization does not require irrationality---only asymmetric evidence supply.

\section{Cascade Models}
\label{sec:ch11-cascades}

\begin{model}[Information Cascade]
\label{mod:cascade}
\index{information cascade}
Following \textcite{banerjee1992} and \textcite{bikhchandani1992}, agents act sequentially, observing predecessors' actions but not their private signals.
An \emph{information cascade} occurs when agent $j$'s optimal action is independent of $j$'s private signal---the public evidence overwhelms private information.
In doxonomic terms, a cascade occurs when the accumulated weight of institutional narrative is so strong that individual experience becomes irrelevant to expressed belief.
\end{model}

\section{Differential Equation Framework}
\label{sec:ch11-ode}

We can model competing beliefs as a dynamical system.
Let $\belief{1}(t)$ and $\belief{2}(t) = 1 - \belief{1}(t)$ be two competing beliefs.

\begin{model}[Lotka-Volterra Belief Competition]
\label{mod:lotka-volterra}
\index{Lotka-Volterra model}
\begin{align}
  \frac{d\belief{1}}{dt} &= \belief{1}(\alpha_1 \phi_1 - \gamma_1 - \beta_{12}\belief{2}) \\
  \frac{d\belief{2}}{dt} &= \belief{2}(\alpha_2 \phi_2 - \gamma_2 - \beta_{21}\belief{1})
\end{align}
where $\phi_k$ is the narrative flow for belief $k$, $\gamma_k$ is the intrinsic decay, and $\beta_{ij}$ is the competitive displacement rate.
The funding asymmetry $\phi_1 \gg \phi_2$ can ensure dominance of belief 1 regardless of intrinsic merit.
\end{model}

\section{Notes and References}
\label{sec:ch11-notes}

Stock-and-flow models adapt systems dynamics \autocite{jackson2008}.
The SIR framework for information is developed in \textcite{watts2002}.
Bayesian polarization under asymmetric information is treated in the economics of media bias \autocite{gentzkow2010}.
Cascade models originate with \textcite{banerjee1992} and \textcite{bikhchandani1992}.
Lotka-Volterra models for cultural competition appear in cultural evolution literature \autocite{nowak2006}.

\section{Exercises}

\begin{exercise}
For the stock-and-flow model~\eqref{eq:stock-flow} with $\alpha = 0.3$, $\delta = 0.1$, $\beta = 0.5$, find the minimum narrative flow $\phi^*$ needed for a stable high-adoption equilibrium to exist.
\end{exercise}

\begin{exercise}
Compute $R_0$ for a belief spread through social media ($\beta = 0.05$, $\gamma = 0.2$) with institutional support $\alpha\phi = 0.1$.
Can the belief sustain itself without institutional support?
\end{exercise}

\begin{exercise}
Simulate the Bayesian polarization theorem: generate $T = 100$ signals for two groups with $q_A = 0.7$ and $q_B = 0.3$, starting from the same prior $\pi^{(0)} = 0.5$.
Plot the posterior trajectories.
\end{exercise}

\begin{exercise}
In the Lotka-Volterra model, find the conditions under which belief 2 survives despite $\phi_1 > \phi_2$.
Interpret economically: what makes a counter-narrative resilient?
\end{exercise}

\begin{exercise}
Extend the SIR model to include a ``vaccinated'' compartment $V(t)$ representing people who have received counter-narrative inoculation.
Write the modified equations and find the new $R_0$.
\end{exercise}

\chapter{Contradiction, Fragmentation, and Sheaf-Like Coherence}
\label{ch:contradiction-coherence}

\epigraph{The test of a first-rate intelligence is the ability to hold two opposed ideas in the mind at the same time, and still retain the ability to function.}{F.\ Scott Fitzgerald}

\noindent
Real societies do not hold a single coherent ideology.
They hold fragments: the workplace teaches competition, the family teaches care, the church teaches sacrifice, the market teaches acquisition.
These local belief systems often contradict each other.
Yet the social order functions as if they were compatible.
This chapter uses sheaf-theoretic language to formalize how local coherence coexists with global contradiction, and how ideology manages the seams.

\section{Local Belief Systems}
\label{sec:ch12-local}

\begin{definition}[Local Belief System]
\label{def:local-belief}
\index{local belief system}
Let $\mathcal{U} = \{U_1, \ldots, U_n\}$ be a cover of social life by institutional contexts (workplace, family, school, market, polity, etc.).
A \emph{local belief system} on $U_i$ is a consistent set of propositions $\sigma_i \in \Gamma(U_i, \sheaf)$---a section of a sheaf $\sheaf$ over $U_i$---specifying the operational norms, anthropological assumptions, and behavioral expectations within that context.
\end{definition}

\begin{example}
\label{ex:local-contradiction}
Consider two contexts:
\begin{itemize}[nosep]
  \item $U_1$ (workplace): $\sigma_1 = $ ``maximize personal advantage, compete, individuals are self-interested''
  \item $U_2$ (family): $\sigma_2 = $ ``sacrifice for others, cooperate, persons are constitutively relational''
\end{itemize}
Each is internally consistent.
But on the overlap $U_1 \cap U_2$ (e.g., work-life balance decisions), they prescribe contradictory behavior.
\end{example}

\section{Sheaf-Theoretic Framework}
\label{sec:ch12-sheaf}

\begin{definition}[Belief Sheaf]
\label{def:belief-sheaf}
\index{sheaf!belief sheaf}
A \emph{belief sheaf} is a sheaf $\sheaf$ on the topological space $(\mathcal{U}, \tau)$ of institutional contexts, where:
\begin{itemize}[nosep]
  \item to each open set $U_i$, $\sheaf$ assigns a set $\sheaf(U_i)$ of consistent local belief systems,
  \item for $U_j \subset U_i$, there is a restriction map $\rho_{ij}: \sheaf(U_i) \to \sheaf(U_j)$,
  \item the sheaf condition holds: if local sections $\sigma_i \in \sheaf(U_i)$ agree on overlaps ($\rho_{ij}(\sigma_i) = \rho_{ij}(\sigma_j)$ on $U_i \cap U_j$), they glue to a global section $\sigma \in \sheaf(\bigcup U_i)$.
\end{itemize}
\end{definition}

\begin{definition}[Gluing Obstruction]
\label{def:obstruction}
\index{gluing obstruction}
A \emph{gluing obstruction} exists when local sections are individually consistent but cannot be glued into a global section.
Formally, the obstruction lies in the first cohomology $H^1(\mathcal{U}, \sheaf)$.
When $H^1 \neq 0$, the society maintains local coherence without global consistency.
\end{definition}

\begin{theorem}[Ideology as Obstruction Management]
\label{thm:ideology-obstruction}
\index{ideology!as obstruction management}
A dominant ideology is a mechanism for managing non-trivial $H^1(\mathcal{U}, \sheaf)$.
It does not eliminate contradictions.
It provides narrative bridges---\emph{transition maps}---that allow agents to move between contradictory local belief systems without experiencing the contradiction as a crisis.
\end{theorem}

\begin{proof}[Proof sketch]
Define transition maps $g_{ij}: U_i \cap U_j \to \mathrm{Aut}(\sheaf)$ that ``translate'' between local systems on overlaps.
These maps form a \v{C}ech 1-cocycle.
The ideology is effective if the cocycle is perceived as a coboundary---i.e., agents experience the transition as natural rather than contradictory.
The obstruction class $[g] \in H^1(\mathcal{U}, \sheaf)$ measures the irreducible contradiction that no ideology can fully dissolve.
\end{proof}

\section{Managed Contradiction}
\label{sec:ch12-managed}

\begin{principle}[Productive Incoherence]
\label{prin:productive-incoherence}
A doxonomic system benefits from moderate incoherence.
Complete coherence would expose the system's commitments to scrutiny.
Complete incoherence would collapse social coordination.
The optimal regime maintains enough local coherence for institutional functioning and enough global incoherence to prevent systematic critique.
\end{principle}

\begin{model}[Optimal Incoherence]
\label{mod:optimal-incoherence}
Let $h = \dim H^1(\mathcal{U}, \sheaf)$ measure the level of global incoherence.
The stability of the doxonomic regime is
\[
  S(h) = \begin{cases}
    \text{low} & \text{if } h = 0 \text{ (fully coherent: visible and critiquable)} \\
    \text{high} & \text{if } 0 < h < h^* \text{ (managed contradiction)} \\
    \text{low} & \text{if } h > h^* \text{ (fragmentation: social coordination fails)}
  \end{cases}
\]
The regime optimizes at intermediate $h$, maintaining contradictions that are locally invisible.
\end{model}

\section{Fragmented Publics}
\label{sec:ch12-fragmentation}

When the transition maps break down---when agents can no longer smoothly move between contradictory local systems---the sheaf fails to cohere even locally.
This is the condition of \emph{epistemic fragmentation}: multiple publics inhabiting incommensurable belief worlds.

\begin{definition}[Epistemic Fragmentation]
\label{def:fragmentation}
\index{epistemic fragmentation}
A population is \emph{epistemically fragmented} if the cover $\mathcal{U}$ can be partitioned into disconnected components $\mathcal{U}_1, \ldots, \mathcal{U}_r$ such that no transition maps exist between components.
In sheaf terms, $\sheaf$ decomposes as a direct sum of sheaves on disconnected spaces.
\end{definition}

Fragmentation is distinct from pluralism.
Pluralism means different beliefs coexist within a shared epistemic framework.
Fragmentation means the frameworks themselves are incommensurable.

\section{Notes and References}
\label{sec:ch12-notes}

Sheaf theory is introduced in any standard algebraic topology text; for applications to data and social systems, see \textcite{ghrist2014} and \textcite{curry2014}.
The concept of managed contradiction draws on \textcite{gramsci1971} and \textcite{bourdieu1984}.
\v{C}ech cohomology and obstruction theory are standard; we use them here metaphorically but with formal precision.
On epistemic fragmentation, see \textcite{sunstein2017} and the literature on filter bubbles.
The idea that ideology manages rather than eliminates contradiction is central to \textcite{adorno1944}.

\section{Exercises}

\begin{exercise}
Construct a concrete belief sheaf for a society with three institutional contexts: market, family, and polity.
Assign local sections to each and identify the overlaps where contradictions arise.
Compute whether a global section exists.
\end{exercise}

\begin{exercise}
For the two-context example (workplace vs.\ family), write explicit transition maps $g_{12}$ that a common ideology might provide.
Examples: ``providing for my family requires competing at work'' or ``the market rewards everyone fairly, so competition is care.''
Verify that these transition maps reduce the perceived obstruction.
\end{exercise}

\begin{exercise}
Model epistemic fragmentation in a polarized society.
Define a cover with two disconnected components and show that no global section exists.
What would ``reunification'' require in sheaf-theoretic terms?
\end{exercise}

\begin{exercise}
Argue that the $\dim H^1$ measure of incoherence is related to the number of independent ``contradictions'' a society must manage.
Give three examples of irreducible contradictions in contemporary capitalism.
\end{exercise}


\begin{interlude}
You now have the conceptual toolkit: epistemic capital, narrative infrastructure, common-sense capture, doxonomic power, belief markets, dynamical models, and the sheaf-theoretic framework for coherence and contradiction.
These tools let you \emph{model} doxonomic systems formally.
Part~III asks: how would you \emph{measure} them empirically?
The methods chapters are the scientific backbone of the book: where the framework meets data.
\end{interlude}

% --- Part III: Methods ---
\part{Methods}
\label{part:methods}

\chapter{Doxometric Measurement}
\label{ch:doxometric-measurement}

\epigraph{Not everything that counts can be counted, and not everything that can be counted counts.}{William Bruce Cameron}

\noindent
Doxonomics requires measurement.
Without it, the field remains assertion---a collection of plausible claims about belief production that cannot be tested, compared, or refuted.
This chapter develops the theory of \emph{doxometrics}\index{doxometrics}: the quantitative measurement of public belief, its intensity, distribution, and institutional embeddedness.

The central difficulty is that beliefs are \emph{latent variables}.
We do not observe them directly.
We observe indicators---survey responses, behaviors, choices, texts, institutional practices---and must infer the underlying belief structure.
This inference is never perfect, and the measurement challenges are at least as severe as those in any social science.
Doxometrics must therefore be psychometrically serious: attentive to reliability, validity, bias, and the gap between what we measure and what we claim to measure.


\section{What Counts as a Measurable Belief?}
\label{sec:ch13-operationalization}

\begin{definition}[Operationalized Belief]
\label{def:operationalized}
\index{operationalized belief}
A belief $p$ is \emph{doxometrically operationalized} when there exists a measurement function $m: \text{Pop} \to [0,1]$ such that $m(j)$ approximates $\pi_j(p)$ with known reliability and validity.
The population-level measure is $\hat{\mu}(p) = \frac{1}{N}\sum_{j=1}^N m(j)$.
\end{definition}

Not all beliefs are equally measurable.
The layered ontology (Principle~\ref{prin:layered-ontology}) implies that different layers present different measurement challenges:

\begin{center}
\begin{tabular}{lp{4.5cm}p{5cm}}
\toprule
\textbf{Layer} & \textbf{Measurement method} & \textbf{Difficulty} \\
\midrule
1. Propositional belief & Surveys, polls & Moderate (social desirability bias, framing effects) \\
2. Norm perception & Norm-elicitation surveys, vignettes & Moderate (perception $\neq$ reality) \\
3. Practical schema & Behavioral observation, reaction time & High (must be inferred from action, not report) \\
4. Identity commitment & Identity scales, in-group/out-group tasks & Moderate-high (self-report may be performative) \\
5. Affective orientation & Implicit association tests, physiological measures & High (requires specialized equipment; validity debated) \\
\bottomrule
\end{tabular}
\end{center}

Most doxometric work focuses on layer~1 because it is the most tractable.
But the most consequential doxonomic products---deep anthropological regimes, embodied dispositions, affective prejudices---operate at layers 3--5, where measurement is hardest.
This gap between what is most measurable and what is most important is the field's central methodological challenge.


\section{Indicators and Proxies}
\label{sec:ch13-indicators}

Direct measurement of $\pi_j(p)$ is rarely possible.
We distinguish four types of indicators, each with characteristic strengths and biases:

\begin{description}[nosep]
  \item[Stated belief:] direct survey response (``Do you agree that human nature is selfish?'').
  \textit{Strength}: easy to collect at scale.
  \textit{Bias}: social desirability (respondents give ``acceptable'' answers), acquiescence (tendency to agree with any statement), framing effects (wording changes responses dramatically).

  \item[Behavioral marker:] observed action consistent with $p$ (e.g., defection rates in public goods games as a proxy for selfishness belief; willingness to cross a picket line as a proxy for individualist values).
  \textit{Strength}: resistant to social desirability bias (actions are costly, so they reveal more than words).
  \textit{Bias}: behavior is multiply determined (defection in a game may reflect strategic calculation, not belief about human nature).

  \item[Revealed preference:] choices that imply $p$ (e.g., voting for policies consistent with meritocratic belief; purchasing choices consistent with consumerist identity).
  \textit{Strength}: based on high-stakes real-world decisions.
  \textit{Bias}: preferences are not beliefs (I may vote for redistribution because it benefits me, not because I reject meritocracy).

  \item[Textual indicator:] frequency and framing of $p$ in media, social media, institutional documents, or personal communication.
  \textit{Strength}: available at massive scale through computational methods.
  \textit{Bias}: discourse frequency is not belief prevalence (a proposition may be frequently discussed precisely because it is contested, not because it is believed).
\end{description}

\begin{proposition}[Divergence of Indicators]
\label{prop:divergence}
Stated belief and behavioral markers may diverge systematically.
A population may report cooperative values while behaving competitively (or vice versa).
The gap $|m_{\text{stated}}(p) - m_{\text{behavioral}}(p)|$ is itself a doxonomic observable, measuring the \emph{depth of internalization}.
\end{proposition}

\begin{remark}[The Divergence as Data]
\label{rem:divergence-as-data}
When stated belief exceeds behavioral compliance, the belief has not been internalized to layer~3 (practical schema): people profess it but do not enact it.
When behavioral compliance exceeds stated belief, the belief may operate doxically (layer~5): people enact it without explicitly endorsing it.
The gap between layers is not noise to be eliminated; it is signal about the depth and mode of doxonomic penetration.
\end{remark}

\begin{example}[Measuring Belief in Meritocracy]
\label{ex:measuring-meritocracy}
Four indicators for the belief ``success reflects effort and talent'':
\begin{enumerate}[nosep]
  \item \textbf{Stated}: Likert-scale agreement with ``people who work hard almost always get ahead'' (World Values Survey item).
  \item \textbf{Behavioral}: willingness to assign unequal rewards in an experimental dictator game after observing performance differences.
  \item \textbf{Revealed preference}: voting for politicians who oppose redistributive taxation.
  \item \textbf{Textual}: frequency of meritocratic framing (``earned,'' ``deserved,'' ``self-made'') in media coverage of wealth.
\end{enumerate}
Each captures a different facet.
A complete doxometric assessment uses all four and reports both the individual estimates and the divergences.
\end{example}


\section{Measurement Bias and Correction}
\label{sec:ch13-bias}

Several systematic biases plague doxometric measurement.

\begin{definition}[Social Desirability Bias]
\label{def:social-desirability}
\index{social desirability bias}
\emph{Social desirability bias} (SDB) is the tendency of respondents to answer survey questions in a way that will be viewed favorably by others, rather than truthfully.
In doxometric contexts, SDB causes under-reporting of socially disapproved beliefs (racism, selfish motivation) and over-reporting of socially approved beliefs (cooperation, fairness).
\end{definition}

Correction strategies:
\begin{itemize}[nosep]
  \item \textbf{List experiments}\index{list experiment}: respondents report the number of statements they agree with from a list, rather than identifying which ones.
  The treatment group's list includes the sensitive item; the control group's does not.
  The difference in average counts estimates the prevalence of the sensitive belief without any individual revealing their answer.
  \item \textbf{Endorsement experiments}: respondents evaluate a policy endorsed by a figure associated with the belief.
  The endorsement effect measures implicit alignment without direct self-report.
  \item \textbf{Implicit measures}: reaction-time tasks (Implicit Association Test) measure automatic associations, bypassing deliberate self-presentation.
  \textit{Caveat}: the reliability and validity of IATs for measuring beliefs (as opposed to associations) remains debated.
  \item \textbf{Behavioral validation}: compare survey responses with behavior in incentivized tasks.
  Respondents who claim to value cooperation but defect in public goods games are flagged as exhibiting SDB.
\end{itemize}

\begin{definition}[Framing Bias]
\label{def:framing-bias}
\index{framing bias}
\emph{Framing bias} occurs when the wording of a survey question influences the response, independent of the underlying belief.
``Do you believe in free enterprise?'' and ``Do you think corporations should be unregulated?'' may target the same proposition but elicit different responses.
Doxometric instruments should include multiple framings of the same target belief and report the framing sensitivity as a measure of measurement precision.
\end{definition}

\begin{definition}[Measurement Invariance]
\label{def:invariance}
\index{measurement invariance}
A doxometric instrument satisfies \emph{measurement invariance} across groups $G_1, \ldots, G_L$ if the relationship between latent belief and observed indicators is the same in each group.
Formally: the factor loadings $\lambda_k$ and intercepts are equal across groups (metric and scalar invariance).
Without invariance, cross-group comparisons of belief levels are not meaningful---a higher score in group $G_1$ might reflect a stronger belief or a different relationship between the belief and the indicator.
\end{definition}

\begin{remark}[Why Invariance Matters for Doxonomics]
\label{rem:invariance-matters}
Doxonomic analysis frequently compares belief levels across countries, classes, racial groups, or media environments.
If the measurement instrument is not invariant across these groups---if ``hard work leads to success'' means something different to a white professional in Connecticut than to a Black factory worker in Detroit---then the comparison is spurious.
Every cross-group doxometric comparison should be accompanied by an invariance test.
\end{remark}


\section{Latent Variable Models}
\label{sec:ch13-latent}

\begin{model}[Factor Model for Belief]
\label{mod:factor}
\index{latent variable model}
Let $\bm{y}_j = (y_{j1}, \ldots, y_{jK})$ be $K$ observable indicators for agent $j$.
The latent belief $\pi_j$ generates observables via
\[
  y_{jk} = \lambda_k \pi_j + \varepsilon_{jk}
\]
where $\lambda_k$ is the factor loading for indicator $k$ (how strongly the indicator reflects the latent belief) and $\varepsilon_{jk} \sim \mathcal{N}(0, \sigma_k^2)$ is measurement error.
The latent belief is recovered as the weighted average:
\[
  \hat{\pi}_j = \frac{\sum_k \lambda_k y_{jk} / \sigma_k^2}{\sum_k \lambda_k^2 / \sigma_k^2}
\]
This is the minimum-variance unbiased estimator under the model.
\end{model}

\begin{proposition}[Reliability]
\label{prop:reliability}
\index{reliability}
The \emph{reliability} of the latent belief estimate is
\[
  \text{Rel}(\hat{\pi}) = 1 - \frac{\mathrm{Var}(\hat{\pi} - \pi)}{\mathrm{Var}(\hat{\pi})} = \frac{\sum_k \lambda_k^2 / \sigma_k^2}{\sum_k \lambda_k^2 / \sigma_k^2 + 1}
\]
Reliability increases with the number of indicators $K$ (more items → more signal) and with the factor loadings $\lambda_k$ (stronger items → more signal), and decreases with measurement error variance $\sigma_k^2$.
The standard threshold for acceptable reliability in psychometrics is $\text{Rel} > 0.70$.
\end{proposition}

\begin{remark}[Beyond Single-Factor Models]
\label{rem:multi-factor}
The single-factor model assumes that all indicators load on a single latent belief.
Many doxonomic beliefs are multi-dimensional.
``Meritocratic belief'' may decompose into: belief in effort → success, belief in equal opportunity, belief in desert-based distribution, and attribution of poverty to individual failure.
These may be correlated but not identical.
Exploratory factor analysis (EFA) and confirmatory factor analysis (CFA) should be used to assess dimensionality before constructing composite indices.
Treating a multi-dimensional construct as unidimensional produces unreliable and uninterpretable results.
\end{remark}

\subsection{Item Response Theory}

For more refined measurement, item response theory (IRT) models the probability of endorsing each item as a function of the latent belief and item characteristics:

\begin{model}[Two-Parameter IRT]
\label{mod:irt}
\index{item response theory}
The probability that agent $j$ endorses item $k$ is
\[
  P(y_{jk} = 1 \mid \pi_j) = \frac{1}{1 + \exp(-a_k(\pi_j - b_k))}
\]
where $a_k$ is the \emph{discrimination parameter} (how sharply the item distinguishes between high and low $\pi_j$) and $b_k$ is the \emph{difficulty parameter} (the belief level at which endorsement probability is 50\%).
\end{model}

IRT improves on the factor model by allowing each item to have a different measurement precision at different points on the belief continuum.
A well-designed doxometric instrument should include items with different difficulty parameters to measure accurately across the full range of belief levels---from strong disbelief through ambivalence to strong endorsement.


\section{Composite Indices of Belief Dominance}
\label{sec:ch13-composite}

\begin{definition}[Doxonomic Dominance Index]
\label{def:dominance-index}
\index{dominance index}
The \emph{doxonomic dominance index} (DDI) for belief $p$ is
\[
  \mathrm{DDI}(p) = \alpha_1 \hat{\mu}(p) + \alpha_2 \cdot (1 - \mathrm{Var}(\hat{\pi})) + \alpha_3 \cdot \mathrm{InstEmbed}(p) + \alpha_4 \cdot \mathrm{MediaFreq}(p)
\]
where:
\begin{itemize}[nosep]
  \item $\hat{\mu}(p)$ is the population-mean belief (how widely held),
  \item $1 - \mathrm{Var}(\hat{\pi})$ is the concentration of belief (how uniformly held),
  \item $\mathrm{InstEmbed}(p)$ measures institutional embeddedness (how many institutions treat $p$ as operational truth),
  \item $\mathrm{MediaFreq}(p)$ is the normalized frequency of $p$ in public discourse.
\end{itemize}
The weights $\alpha_i$ satisfy $\sum \alpha_i = 1$.
\end{definition}

\begin{remark}[DDI Limitations]
\label{rem:ddi-limitations}
The DDI is a \emph{dashboard sketch}, not a validated index.
The weights $\alpha_i$ are currently arbitrary; different weights produce different rankings.
The components mix unlike quantities (survey proportions, variance measures, frequency counts) without a clear scaling doctrine.
Before the DDI can serve as a serious research tool, it requires:
\begin{itemize}[nosep]
  \item empirical validation (does DDI predict policy outcomes, behavioral patterns, or other independently measurable consequences?),
  \item sensitivity analysis (how much do results change with different weight specifications?),
  \item normalization (each component should be on a comparable scale).
\end{itemize}
We include it here as a conceptual prototype, not as a finished instrument.
\end{remark}


\section{Measuring Institutional Embeddedness}
\label{sec:ch13-embeddedness}

The DDI requires a measure of institutional embeddedness---how deeply a belief is woven into the operational assumptions of institutions.

\begin{definition}[Institutional Embeddedness]
\label{def:inst-embed}
\index{institutional embeddedness}
The \emph{institutional embeddedness} of belief $p$ is
\[
  \mathrm{InstEmbed}(p) = \frac{1}{|\mathcal{U}|} \sum_{U_i \in \mathcal{U}} \mathbf{1}[p \in \sigma_i]
\]
where $\mathcal{U}$ is the set of institutional contexts and $\mathbf{1}[p \in \sigma_i] = 1$ if $p$ functions as an operational assumption in context $U_i$.
\end{definition}

Operationalizing this requires coding institutional practices, curricula, regulations, and organizational structures for the presence of $p$ as an operational assumption---a qualitative judgment that can be systematized through:
\begin{itemize}[nosep]
  \item textbook analysis (does the economics curriculum assume rational self-interest?),
  \item organizational document analysis (do HR policies assume shirking? do compensation structures assume competitive motivation?),
  \item legal analysis (does regulatory framework assume market efficiency?),
  \item expert survey (do practitioners in a field report that $p$ is a taken-for-granted assumption in their work?).
\end{itemize}

\begin{example}[Embeddedness of ``Rational Self-Interest'']
\label{ex:embeddedness}
Coding for the belief ``agents are rational self-interest maximizers'' across U.S.\ institutional contexts:
\begin{center}
\begin{tabular}{lcc}
\toprule
\textbf{Institution} & \textbf{Embedded?} & \textbf{Evidence} \\
\midrule
Economics curricula & Yes & Textbook analysis \\
Corporate HR / compensation & Yes & Incentive-based pay design \\
Antitrust regulation & Yes & Consumer welfare standard \\
Criminal justice & Partial & Deterrence theory; also retributive \\
Public education & No & Assumes cooperation, development \\
Religious institutions & No & Assumes altruism, duty \\
Military & Partial & Hierarchy + unit solidarity \\
Family law & No & Assumes care, obligation \\
\bottomrule
\end{tabular}
\end{center}
$\mathrm{InstEmbed} \approx 3.5/8 = 0.44$.
The belief is institutionally embedded in roughly half of major contexts---far from total, but substantial enough to function as a default assumption in public life.
\end{example}


\section{Measuring Doxonomic Change Over Time}
\label{sec:ch13-temporal}

Doxonomic analysis requires not just snapshot measurement but the tracking of belief trajectories over time.

\begin{definition}[Belief Trajectory]
\label{def:trajectory}
\index{belief trajectory}
The \emph{belief trajectory} for proposition $p$ is the time series $\{\hat{\mu}(p, t)\}_{t=0}^{T}$.
A doxonomic event (a major campaign, a crisis, a curriculum reform) produces a detectable shift in the trajectory.
\end{definition}

Sources for longitudinal belief data:
\begin{itemize}[nosep]
  \item \textbf{Repeated cross-sections}: World Values Survey (waves every 5 years since 1981), General Social Survey (annual/biennial since 1972), Eurobarometer (biannual since 1973).
  \item \textbf{Panel surveys}: tracking the same individuals over time (more informative but more expensive; subject to attrition).
  \item \textbf{Text time series}: media corpus analysis tracking narrative frequency and framing over decades (Google Ngrams, GDELT, Media Cloud).
  \item \textbf{Behavioral time series}: cooperation rates in experimental games, charitable giving rates, labor market indicators.
\end{itemize}

\begin{remark}[Identifying Doxonomic Events in Time Series]
\label{rem:event-identification}
A sudden shift in $\hat{\mu}(p, t)$ coinciding with a known doxonomic investment (a new think tank, a major ad campaign, a curriculum reform, a media ownership change) is \emph{suggestive} but not \emph{causal} evidence of doxonomic effect.
Establishing causality requires the methods of Chapter~\ref{ch:causal-inference}: natural experiments, difference-in-differences, or synthetic controls that rule out confounding events.
Time-series correlation between funding and belief is a starting point, not a conclusion.
\end{remark}


\section{Notes and References}
\label{sec:ch13-notes}

Measurement theory for social science draws on \textcite{king1994} and \textcite{gelman2020}.
Latent variable models and factor analysis are standard in psychometrics; see Bollen (1989) for the structural equation modeling framework.
Item response theory is developed in Embretson and Reise (2000).
Behavioral measures of beliefs about cooperation are developed in \textcite{fischbacher2001} and \textcite{henrich2001}.
Survey methodology for the digital age connects to \textcite{salganik2018}.
On the divergence between stated and revealed preferences, see \textcite{kahneman2011}.

List experiments and other techniques for sensitive questions are reviewed in Blair, Imai, and Zhou (2015).
The Implicit Association Test is from Greenwald, McGhee, and Schwartz (1998); for a critical assessment, see Oswald et al.\ (2013).
Measurement invariance testing is covered in Meredith (1993) and Vandenberg and Lance (2000).
For topic models as tools for measuring latent constructs in text, see \textcite{blei2003} and \textcite{grimmer2013}.
On the distinction between discourse frequency and belief prevalence, see \textcite{zaller1992}: media coverage shapes the \emph{accessibility} of considerations, not directly the \emph{beliefs} of the audience.


\section{Exercises}

\begin{exercise}
Design a five-item survey instrument to measure belief in selfish human nature.
Include at least two reverse-coded items.
Specify the response format, argue for content validity, and describe how you would assess reliability and dimensionality.
\end{exercise}

\begin{exercise}
Using the factor model (Model~\ref{mod:factor}), show that the reliability of the latent belief estimate $\hat{\pi}_j$ increases with the number of indicators $K$ and decreases with measurement error variance $\sigma_k^2$.
Compute the minimum $K$ needed for $\text{Rel} > 0.80$ when $\lambda_k = 0.7$ and $\sigma_k = 0.5$ for all $k$.
\end{exercise}

\begin{exercise}
Compute the DDI for the belief ``markets are efficient'' using plausible estimates for each component.
Then repeat with different weight specifications ($\alpha_1 = 0.5$ vs.\ $\alpha_1 = 0.1$).
How sensitive is the DDI to the choice of weights?
What does this imply about its usefulness?
\end{exercise}

\begin{exercise}
\label{ex:ch13-list}
Design a list experiment to measure the prevalence of the belief ``most poor people are poor because they are lazy'' in a population where this belief is socially disapproved.
Specify the control and treatment lists, the sample size calculation, and the estimator.
\end{exercise}

\begin{exercise}
\label{ex:ch13-invariance}
A doxometric instrument measuring meritocratic belief is administered in the U.S.\ and Sweden.
Describe how you would test for measurement invariance across the two countries.
What would you conclude if configural invariance holds but scalar invariance fails?
Can you compare mean belief levels across the two countries?
\end{exercise}

\begin{exercise}
\label{ex:ch13-layers}
Choose a specific belief (e.g., ``inequality is natural'').
For each of the five layers in the ontology (Principle~\ref{prin:layered-ontology}), propose a concrete measurement method.
Which layers are most important for doxonomic analysis?
Which are most measurable?
What is the gap between importance and measurability, and what does it imply for the field?
\end{exercise}

\begin{exercise}
Using publicly available World Values Survey data, plot the trajectory of belief in ``hard work leads to success'' in three countries over the past 30 years.
Do you observe any doxonomic events (sudden shifts)?
If so, what was happening in each country at the time of the shift?
\end{exercise}

\begin{exercise}
The chapter notes that discourse frequency is not the same as belief prevalence.
Give an example where $\mathrm{MediaFreq}(p)$ is high but $\hat{\mu}(p)$ is low, and vice versa.
What does each configuration imply doxonomically?
\end{exercise}

\chapter{Belief-Flow Accounting}
\label{ch:belief-flow-accounting}

\epigraph{Follow the money.}{Deep Throat (attr.)}

\noindent
If doxonomics is to be a science rather than a commentary, it needs accounting.
The claim that ``beliefs are funded'' is empty without a ledger: who spent how much, through which institutions, producing which narratives, reaching which audiences, with what measurable effects.
This chapter develops the framework for tracing financial flows into narrative outputs: the doxonomic ledger.

Belief-flow accounting is to doxonomics what national income accounting is to macroeconomics: a systematic method for measuring aggregate flows.
Just as GDP accounting tracks the flow of goods and services through an economy, doxonomic accounting tracks the flow of resources through the belief production chain.
The analogy is not perfect (narrative production is harder to measure than automobile production), but the aspiration is the same: replace impressionistic claims with documented flows.


\section{The Accounting Framework}
\label{sec:ch14-framework}

\begin{definition}[Doxonomic Account]
\label{def:doxonomic-account}
\index{belief-flow accounting}
A \emph{doxonomic account} for belief $p$ in period $[0, T]$ is a table
\[
  \mathcal{L}_p = \{(s_k, \iota_k, c_k, \funding{k}, \Theta_k)\}_{k=1}^{m}
\]
where $s_k$ is the funding source, $\iota_k$ is the processing institution, $c_k$ is the channel, $\funding{k}$ is the expenditure, and $\Theta_k$ is the estimated doxonomic throughput.
The \emph{total doxonomic investment} in $p$ is $F_p = \sum_k \funding{k}$.
\end{definition}

A complete doxonomic account answers five questions:
\begin{enumerate}[nosep]
  \item \textbf{Who paid?} (source identification)
  \item \textbf{Who processed?} (institutional intermediation)
  \item \textbf{What was produced?} (narrative output)
  \item \textbf{Who was reached?} (audience and exposure)
  \item \textbf{What was the effect?} (belief change, behavioral consequence)
\end{enumerate}

In practice, complete accounts are rarely achievable.
Funding opacity, institutional complexity, and the difficulty of measuring belief effects mean that most doxonomic accounts are partial. Partial accounting is vastly better than no accounting.

\begin{example}[Partial Account: Climate Doubt]
\label{ex:climate-account}
\index{belief-flow accounting!climate example}
A partial doxonomic account for the belief ``climate science is uncertain'' in the U.S., 1990--2010 \autocite{oreskes2010}:

\begin{center}
\begin{tabular}{p{2.3cm}p{2.5cm}p{2.5cm}rp{1.5cm}}
\toprule
\textbf{Source} & \textbf{Institution} & \textbf{Channel} & \textbf{\$ (est.)} & \textbf{Reach} \\
\midrule
ExxonMobil & Heartland Inst. & Reports, op-eds & \$8M & Policy \\
Koch Industries & Americans for Prosperity & TV ads, rallies & \$50M+ & Mass \\
Coal industry & Global Climate Coalition & Lobbying, PR & \$13M & Congress \\
Multiple & George C. Marshall Inst. & Congressional testimony & \$2M & Elite \\
Multiple & Competitive Enterprise Inst. & Media appearances & \$5M & Broadcast \\
\bottomrule
\end{tabular}
\end{center}

This is a fragment of the total.
\textcite{oreskes2010} and subsequent investigations (Frumhoff et al., 2015; Supran and Oreskes, 2017) have documented hundreds of millions in total spending across dozens of institutions.
The account is partial because:
\begin{itemize}[nosep]
  \item not all funding is traceable (dark money, donor-advised funds),
  \item institutional overhead is not separated from narrative production,
  \item throughput estimates are rough (how many people were actually reached and influenced?),
  \item the causal link from spending to belief change is not established by the account alone (that requires Chapter~\ref{ch:causal-inference}).
\end{itemize}
\end{example}


\section{Source Identification}
\label{sec:ch14-sources}

Tracing funding to belief requires navigating layers of intermediation.
A typical chain:
\[
  \text{Corporation} \to \text{Foundation} \to \text{Donor-advised fund} \to \text{Think tank} \to \text{Publication}
\]

Each intermediary adds opacity.
The corporation may be legally required to disclose its contributions to the foundation, but the foundation may route grants through a donor-advised fund (which does not disclose the ultimate recipient), and the think tank may not disclose its funders.

\begin{definition}[Funding Chain]
\label{def:funding-chain}
\index{funding chain}
A \emph{funding chain} for belief $p$ is a directed path in the financial network:
\[
  s_0 \xrightarrow{f_1} s_1 \xrightarrow{f_2} \cdots \xrightarrow{f_l} \iota \xrightarrow{\text{produces}} n_p
\]
where $s_0$ is the ultimate source, $s_1, \ldots, s_{l-1}$ are intermediaries, $\iota$ is the narrative-producing institution, and $f_i$ is the flow at each stage.
The \emph{chain length} $l$ measures the number of intermediation layers.
Longer chains are more opaque.
\end{definition}

\subsection{Data Sources for Source Identification}

\begin{itemize}[nosep]
  \item \textbf{IRS Form 990 filings} (U.S.): tax-exempt organizations must file annual returns listing revenues, expenditures, and grants made.
  Available through ProPublica Nonprofit Explorer and Guidestar.
  \textit{Limitation}: does not require disclosure of donors to 501(c)(4) organizations; donor-advised funds disclose the fund sponsor but not the donor.

  \item \textbf{Lobbying disclosure records}: the U.S.\ Lobbying Disclosure Act requires registration of lobbyists and quarterly reports of expenditures.
  Available through OpenSecrets.
  \textit{Limitation}: covers only direct lobbying, not grassroots campaigns, astroturfing, or think-tank-mediated advocacy.

  \item \textbf{Corporate annual reports and SEC filings}: publicly traded companies disclose political spending and trade-association memberships.
  \textit{Limitation}: disclosure requirements vary; many expenditures are classified as ``education,'' ``research,'' or ``community engagement'' without specifying ideological content.

  \item \textbf{Investigative journalism databases}: ProPublica, InfluenceWatch, LittleSis, SourceWatch.
  \textit{Limitation}: coverage is selective (focused on controversial or politically salient cases); data quality varies.

  \item \textbf{Academic forensic accounting}: scholars have reconstructed funding networks through painstaking archival work.
  \textcite{mayer2016} traced Koch network funding; \textcite{oreskes2010} traced fossil fuel funding of doubt.
  \textit{Limitation}: labor-intensive, not systematic, not reproducible at scale.
\end{itemize}

\begin{definition}[Funding Opacity]
\label{def:opacity}
\index{funding opacity}
The \emph{opacity} of a funding chain is a measure of how difficult it is to trace the chain from narrative output back to ultimate source.
We define:
\[
  \text{Opacity}(\text{chain}) = 1 - \prod_{i=1}^{l} d_i
\]
where $d_i \in [0,1]$ is the disclosure probability at intermediary $i$ (the probability that the flow $f_i$ is publicly documented).
A fully transparent chain (all $d_i = 1$) has opacity 0.
A chain with one perfectly opaque intermediary ($d_i = 0$ for some $i$) has opacity 1.
\end{definition}

\begin{remark}[The Opacity Arms Race]
\label{rem:opacity-arms-race}
As doxonomic accounting improves, funders develop more sophisticated opacity mechanisms: pass-through entities, foreign intermediaries, anonymized cryptocurrency transfers, and corporate structures designed to frustrate tracing.
Doxonomic accounting is therefore an adversarial exercise: the accounting methodology must evolve in response to evasion strategies, much as tax enforcement evolves in response to tax avoidance.
\end{remark}


\section{Ideological Input-Output Tables}
\label{sec:ch14-leontief}

The belief production chain involves institutions that are both producers and consumers of narrative.
A think tank produces reports that are consumed by media; media produces coverage that is consumed by politicians; politicians produce policies that are consumed by think tanks as evidence for their next round of reports.
This circular flow is captured by the input-output framework.

\begin{model}[Doxonomic Input-Output Table]
\label{mod:io-table}
\index{input-output table}
Let $n$ institutional sectors produce narrative output.
The input-output matrix $\bm{A}$ has entries $a_{ij} \in [0,1]$ representing the fraction of institution $j$'s narrative output consumed as input by institution $i$.
The total output vector $\bm{n}$ satisfies
\[
  \bm{n} = \bm{A}\bm{n} + \bm{d}
\]
where $\bm{d}$ is the exogenous injection (direct funder input).
The solution, when $(\bm{I} - \bm{A})$ is invertible and the spectral radius of $\bm{A}$ is less than 1, is
\[
  \bm{n} = (\bm{I} - \bm{A})^{-1}\bm{d}
\]
The matrix $(\bm{I} - \bm{A})^{-1}$ is the \emph{doxonomic multiplier}\index{doxonomic multiplier}: it shows how a unit injection of narrative through one sector amplifies through the system.
\end{model}

\begin{example}[Constructing an Input-Output Table]
\label{ex:io-construction}
Consider four sectors: think tanks ($T$), prestige media ($M$), broadcast news ($B$), and education ($E$).
Estimating the inter-sector flows from citation patterns, sourcing data, and content analysis:

\begin{center}
\begin{tabular}{l|cccc|c}
\toprule
\textbf{To $\downarrow$ / From $\to$} & $T$ & $M$ & $B$ & $E$ & Exog.\ $\bm{d}$ \\
\midrule
Think tanks $T$ & 0 & 0.05 & 0.02 & 0.10 & \$50M \\
Prestige media $M$ & 0.30 & 0 & 0.10 & 0.05 & \$20M \\
Broadcast news $B$ & 0.15 & 0.40 & 0 & 0.02 & \$10M \\
Education $E$ & 0.20 & 0.10 & 0.05 & 0 & \$30M \\
\bottomrule
\end{tabular}
\end{center}

Reading the table: 30\% of think-tank output is picked up by prestige media ($a_{MT} = 0.30$); 40\% of prestige media output is recycled into broadcast news ($a_{BM} = 0.40$); 20\% of think-tank output eventually enters educational materials ($a_{ET} = 0.20$).

The Leontief inverse $(\bm{I} - \bm{A})^{-1}$ reveals:
\begin{itemize}[nosep]
  \item a \$1M injection into think tanks produces approximately \$1.5M in total narrative output (50\% amplification through the system),
  \item the largest multiplier path is $T \to M \to B$: think-tank research is picked up by prestige media, which is then amplified by broadcast news,
  \item the education sector is a ``slow amplifier'': it absorbs think-tank and media output and releases it over decades through curricula.
\end{itemize}
\end{example}

\begin{remark}[Estimating the $\bm{A}$ Matrix]
\label{rem:estimating-A}
The inter-institutional flow coefficients $a_{ij}$ are not directly observable.
They must be estimated from:
\begin{itemize}[nosep]
  \item \textbf{Citation analysis}: how often does institution $i$ cite or reference institution $j$'s outputs?
  \item \textbf{Sourcing studies}: content analyses of media output identifying the institutional origins of claims, data, and expert quotes.
  \item \textbf{Network analysis}: co-membership, co-citation, and personnel-flow networks that trace idea movement between institutions (Chapter~\ref{ch:network-analysis}).
  \item \textbf{Survey of practitioners}: journalists report which sources they rely on; policymakers report which think-tank outputs influence their thinking.
\end{itemize}
All methods are imprecise.
The input-output table should be accompanied by uncertainty bounds on each coefficient.
\end{remark}


\section{Accounting Identities}
\label{sec:ch14-identities}

\begin{principle}[Doxonomic Accounting Identity]
\label{prin:accounting-identity}
\index{accounting identity}
For each institution $\iota$ in steady state:
\[
  \text{Total funding in} = \text{Narrative output} \times \text{Unit cost} + \text{Institutional overhead}
\]
In symbols: $\funding{\iota} = \Theta_{\iota} \cdot c_{\text{unit}} + O_{\iota}$.
The \emph{doxonomic efficiency} is $\eta_{\iota} = \Theta_{\iota} / \funding{\iota} = 1 - O_{\iota}/\funding{\iota}$.
\end{principle}

Doxonomic efficiency varies widely across institution types:
\begin{itemize}[nosep]
  \item \textbf{High efficiency} ($\eta > 0.6$): small advocacy organizations, online media outlets, social media campaigns.
  Low overhead, direct narrative output.
  \item \textbf{Medium efficiency} ($\eta \approx 0.3$--$0.5$): think tanks, PR firms.
  Moderate overhead (buildings, staff, credibility maintenance), substantial narrative output.
  \item \textbf{Low efficiency} ($\eta < 0.2$): universities, government agencies.
  High overhead (research infrastructure, administration, teaching, non-doxonomic functions), but the doxonomic output has very high credibility and persistence.
\end{itemize}

\begin{remark}[Efficiency vs.\ Effectiveness]
\label{rem:efficiency-effectiveness}
Low doxonomic efficiency does not mean low doxonomic effectiveness.
A university with $\eta = 0.1$ spends 90\% of its budget on non-doxonomic functions (research, teaching, administration), but the 10\% that produces belief effects has $\credibility{} \approx 0.9$ and $\persistence{} \approx 30$ years.
A social media campaign with $\eta = 0.8$ converts most of its budget to reach, but with $\credibility{} = 0.15$ and $\persistence{} = 0.01$ years.
The long-run doxonomic impact (throughput $\times$ persistence $\times$ credibility) may be higher for the university despite its lower efficiency.
This is why the equimarginal principle (Proposition~\ref{prop:optimal-allocation}) and the IROI calculation (Definition~\ref{def:iroi}) account for credibility and persistence, not just efficiency.
\end{remark}


\section{Uncertainty in Doxonomic Accounting}
\label{sec:ch14-uncertainty}

Doxonomic accounts are never precise.
Every entry involves estimation, and honest accounting must report the uncertainty.

\begin{definition}[Account Uncertainty]
\label{def:account-uncertainty}
\index{accounting uncertainty}
For each line in a doxonomic account, the uncertainty is characterized by:
\begin{itemize}[nosep]
  \item \textbf{Source uncertainty}: the ultimate funder may be unknown or partially traced (opacity $> 0$).
  Report as a confidence interval: $s_k \in \{$known sources$\}$ with probability $p_s$, or ``unknown'' with probability $1-p_s$.
  \item \textbf{Amount uncertainty}: the expenditure $\funding{k}$ is often estimated from partial data (e.g., grant totals that include both doxonomic and non-doxonomic spending).
  Report as $\funding{k} \in [F_L, F_U]$ with 90\% confidence.
  \item \textbf{Throughput uncertainty}: the estimated doxonomic throughput $\Theta_k$ depends on reach, credibility, and persistence estimates, each with its own error.
  Report the propagated uncertainty.
  \item \textbf{Attribution uncertainty}: the fraction of an institution's output attributable to belief $p$ (vs.\ other beliefs or non-doxonomic functions) is a judgment call.
  Report the sensitivity of $F_p$ to this attribution fraction.
\end{itemize}
\end{definition}

\begin{principle}[Honest Accounting]
\label{prin:honest-accounting}
A doxonomic account without uncertainty bounds is not a doxonomic account.
It is advocacy with numbers.
The field's credibility (its own epistemic capital) depends on honest reporting of what is known, what is estimated, and what is guessed.
\end{principle}


\section{A Worked Example: Accounting for Austerity Belief}
\label{sec:ch14-austerity}

To illustrate the full accounting framework, consider a partial account for the belief ``government debt is dangerous and austerity is necessary'' in the UK, 2009--2013.

\textbf{Step 1: Source identification.}
The austerity narrative was promoted by:
\begin{itemize}[nosep]
  \item Centre for Policy Studies (CPS): founded by Thatcher and Keith Joseph, annual budget $\sim$\pounds2M.
  \item Institute of Economic Affairs (IEA): annual budget $\sim$\pounds3M.
  \item Taxpayers' Alliance: annual budget $\sim$\pounds1M.
  \item The Treasury itself: the largest single source of austerity framing, but with public rather than private funding.
  \item Financial press (\textit{Financial Times}, \textit{Economist}): editorial stances aligned with austerity, funded by subscriptions and advertising from the financial sector.
\end{itemize}

\textbf{Step 2: Processing.}
Think tanks produced policy briefs, opinion pieces, and parliamentary testimony.
The Treasury produced budget documents, forecasts, and press conferences.
Financial media produced daily coverage framing debt as crisis.

\textbf{Step 3: Throughput estimation.}
\begin{center}
\begin{tabular}{lrrcc}
\toprule
\textbf{Source} & \textbf{\pounds/yr} & \textbf{Reach} & \textbf{Credibility} & \textbf{$\Theta$} \\
\midrule
CPS & \pounds2M & 0.01 & 0.6 & 12K \\
IEA & \pounds3M & 0.01 & 0.7 & 21K \\
Taxpayers' Alliance & \pounds1M & 0.05 & 0.3 & 15K \\
Treasury (austerity share) & \pounds10M est. & 0.3 & 0.8 & 2.4M \\
Financial press & \pounds50M est. & 0.2 & 0.7 & 7M \\
\bottomrule
\end{tabular}
\end{center}

\textbf{Step 4: Total.}
The private think-tank investment was modest ($\sim$\pounds6M/year), but the amplification through the Treasury and financial press produced throughput orders of magnitude larger.
The multiplier effect (Model~\ref{mod:io-table}) was enormous: a small think-tank injection was amplified by the most credible institutions in the system.

\textbf{Step 5: Uncertainty.}
The Treasury's ``austerity share'' is an estimate (what fraction of Treasury communication was devoted to the austerity narrative vs.\ other policy); the financial press estimate is rough (advertising revenue is not doxonomic spending, but editorial alignment is).
The total investment is somewhere in the range \pounds50M--\pounds100M/year, with enormous amplification from institutional credibility.

\begin{remark}
This example illustrates a general pattern: private doxonomic investment is often small relative to the total narrative output, because public institutions and commercial media provide massive amplification.
The think tanks do not need to reach the public directly; they need to reach the \emph{institutions} that reach the public.
The multiplier, not the direct investment, is the story.
\end{remark}


\section{Accounting for Non-Monetary Flows}
\label{sec:ch14-nonmonetary}

Not all doxonomic resources are monetary.
Several important flows resist financial accounting:

\begin{itemize}[nosep]
  \item \textbf{Volunteer labor}: social movements, religious organizations, and community groups produce narrative through unpaid work.
  The monetary equivalent can be estimated (hours $\times$ wage rate), but the institutional dynamics differ from funded production.
  \item \textbf{Attention as currency}: on social media, sharing, liking, and commenting are forms of doxonomic production with zero monetary cost to the producer but significant narrative effect.
  \item \textbf{Institutional cross-subsidy}: a corporation's marketing budget (spent on product advertising) also reinforces the consumer-identity anthropological regime; a military's recruitment budget also reinforces nationalist narratives.
  These doxonomic effects are by-products, not intentional investments, but they are real.
  \item \textbf{Architecture and space}: the built environment communicates beliefs without any ongoing narrative production (monumental state buildings communicate power; open-plan offices communicate transparency/surveillance; shopping malls communicate consumption-as-activity).
\end{itemize}

\begin{definition}[Expanded Doxonomic Account]
\label{def:expanded-account}
\index{belief-flow accounting!expanded}
An \emph{expanded doxonomic account} adds non-monetary flows to the standard account:
\[
  \mathcal{L}_p^+ = \mathcal{L}_p \cup \{(\text{source}, \text{type}, \text{estimated equivalent}, \Theta)\}_{\text{non-monetary}}
\]
The ``estimated equivalent'' converts non-monetary contributions to dollar values where possible and flags them as imputed rather than observed.
\end{definition}


\section{Notes and References}
\label{sec:ch14-notes}

The accounting approach draws on national income accounting methods \autocite{stiglitz2012}.
The Leontief input-output framework originates in Leontief (1941); for a modern introduction, see Miller and Blair (2009).
Dark money flows are documented in \textcite{mayer2016}.
The application of input-output accounting to narrative flows is original to this text.
For lobbying data, see OpenSecrets; for think-tank funding, see InfluenceWatch and ProPublica Nonprofit Explorer.
On media ownership structures, see \textcite{bagdikian2004}.

On the specific case of climate-doubt funding, see \textcite{oreskes2010} and Supran and Oreskes (2017).
For the UK austerity case, see \textcite{streeck2016} and Blyth (2013), \textit{Austerity: The History of a Dangerous Idea}.
On non-monetary doxonomic production (architecture, space, material culture), see Latour (2005) and \textcite{scott1998}.
The opacity concept connects to the literature on dark money and political finance regulation \autocite{mayer2016}.


\section{Exercises}

\begin{exercise}
Using publicly available IRS 990 data (via ProPublica Nonprofit Explorer), construct a partial doxonomic account for the belief ``regulation harms economic growth'' by tracing funding from three major foundations to think tanks to publications.
Report the funding chain length and opacity for each path.
\end{exercise}

\begin{exercise}
Compute the Leontief inverse for the four-sector system in Example~\ref{ex:io-construction}.
What is the total system-wide narrative output from a \$1M injection into think tanks?
Which indirect pathway produces the largest multiplier effect?
\end{exercise}

\begin{exercise}
Estimate the doxonomic efficiency $\eta$ of: (a)~a university economics department, (b)~a social media campaign, (c)~a think tank.
For each, estimate the total budget, the fraction devoted to doxonomic output, and the resulting $\eta$.
Which has the highest efficiency?
Which has the highest long-run effectiveness?
Why might these differ?
\end{exercise}

\begin{exercise}
\label{ex:ch14-uncertainty}
Take the austerity worked example (Section~\ref{sec:ch14-austerity}) and construct 90\% confidence intervals for each throughput estimate.
How does the uncertainty in the ``austerity share'' of Treasury communication affect the total investment estimate?
Perform a sensitivity analysis: report $F_p$ under optimistic, central, and pessimistic attribution assumptions.
\end{exercise}

\begin{exercise}
\label{ex:ch14-nonmonetary}
Identify three non-monetary doxonomic flows in your daily life (e.g., architectural, spatial, routine, performative).
Estimate their doxonomic throughput in dollar-equivalent terms.
How do they compare in magnitude to the monetary flows you encounter through media and advertising?
\end{exercise}

\begin{exercise}
Construct a doxonomic account for a belief of your choice.
Report: (a)~at least three funding sources with amounts, (b)~at least two processing institutions, (c)~at least two channels with reach and credibility estimates, (d)~total $F_p$ with uncertainty bounds, (e)~a candid assessment of what you could not trace and why.
\end{exercise}

\begin{exercise}
\label{ex:ch14-opacity}
The chapter defines funding opacity as $1 - \prod_i d_i$.
For a chain of length $l = 4$ where each intermediary has disclosure probability $d_i = 0.7$, compute the opacity.
What opacity would the chain have if a donor-advised fund ($d = 0.1$) were inserted as one intermediary?
What does this imply about the doxonomic accounting challenge posed by anonymous philanthropy?
\end{exercise}

\chapter{Network Analysis of Belief Systems}
\label{ch:network-analysis}

\epigraph{Tell me who your friends are, and I will tell you who you are.}{Proverb}

\noindent
Doxonomic systems are networks: funders connect to institutions, institutions connect to media, media connect to audiences.
This chapter applies graph theory and network science to map and measure these structures.

\section{Doxonomic Networks}
\label{sec:ch15-networks}

\begin{definition}[Doxonomic Network]
\label{def:doxonomic-network}
\index{doxonomic network}
A \emph{doxonomic network} is a weighted directed graph $G = (V, E, w)$ where vertices $V$ are doxonomic agents (funders, institutions, media, audiences), edges $E$ represent influence or resource flows, and weights $w: E \to \R_{\geq 0}$ measure the magnitude of the flow (dollars, citations, audience share).
\end{definition}

\section{Centrality Measures}
\label{sec:ch15-centrality}

\begin{definition}[Doxonomic Centrality]
\label{def:centrality}
\index{centrality}
The \emph{doxonomic centrality} of node $v$ is its eigenvector centrality in the weighted network:
\[
  x_v = \frac{1}{\lambda} \sum_{u: (u,v) \in E} w_{uv} \cdot x_u
\]
where $\lambda$ is the largest eigenvalue of the weighted adjacency matrix.
Nodes with high doxonomic centrality are those that receive large weighted inputs from other central nodes.
\end{definition}

High-centrality nodes are the bottlenecks of belief production.
Disrupting or capturing a high-centrality institution has disproportionate effect on the entire narrative ecosystem.

\section{Community Detection}
\label{sec:ch15-community}

\begin{definition}[Narrative Cluster]
\label{def:narrative-cluster}
\index{narrative cluster}
A \emph{narrative cluster} is a community in the doxonomic network: a subset $S \subset V$ with denser internal connections than external ones.
Formally, $S$ maximizes the modularity:
\[
  Q = \frac{1}{2W} \sum_{i,j \in S} \left(w_{ij} - \frac{s_i s_j}{2W}\right)
\]
where $s_i = \sum_j w_{ij}$ and $W = \frac{1}{2}\sum_{ij} w_{ij}$.
\end{definition}

Narrative clusters reveal which institutions form coherent ideological coalitions, even without explicit coordination.

\section{Brokerage}
\label{sec:ch15-brokerage}

\begin{definition}[Epistemic Broker]
\label{def:broker}
\index{epistemic broker}
An \emph{epistemic broker} is a node $v$ with high betweenness centrality:
\[
  g(v) = \sum_{s \neq v \neq t} \frac{\sigma_{st}(v)}{\sigma_{st}}
\]
where $\sigma_{st}$ is the number of shortest paths from $s$ to $t$ and $\sigma_{st}(v)$ is the number passing through $v$.
Brokers bridge otherwise disconnected narrative clusters and can translate ideas between communities.
\end{definition}

\section{Example: A Donor-Institution Network}

\begin{center}
\begin{tikzpicture}[
  funder/.style={circle, draw, thick, fill=doxoblue!20, minimum size=1cm, font=\footnotesize},
  inst/.style={circle, draw, thick, fill=doxogreen!20, minimum size=1cm, font=\footnotesize},
  media/.style={circle, draw, thick, fill=doxored!20, minimum size=1cm, font=\footnotesize},
  edge/.style={-Stealth, thick},
]
  \node[funder] (f1) at (0,2) {$F_1$};
  \node[funder] (f2) at (0,0) {$F_2$};
  \node[inst] (i1) at (3,3) {$I_1$};
  \node[inst] (i2) at (3,1) {$I_2$};
  \node[inst] (i3) at (3,-1) {$I_3$};
  \node[media] (m1) at (6,2) {$M_1$};
  \node[media] (m2) at (6,0) {$M_2$};

  \draw[edge] (f1) -- node[above, font=\scriptsize] {\$3M} (i1);
  \draw[edge] (f1) -- node[above, font=\scriptsize] {\$1M} (i2);
  \draw[edge] (f2) -- node[below, font=\scriptsize] {\$2M} (i2);
  \draw[edge] (f2) -- node[below, font=\scriptsize] {\$4M} (i3);
  \draw[edge] (i1) -- (m1);
  \draw[edge] (i2) -- (m1);
  \draw[edge] (i2) -- (m2);
  \draw[edge] (i3) -- (m2);
\end{tikzpicture}
\end{center}

In this network, institution $I_2$ is the epistemic broker: it bridges both funders and both media outlets.

\section{Notes and References}
\label{sec:ch15-notes}

Network analysis draws on \textcite{jackson2008}.
Community detection via modularity is due to \textcite{newman2006}.
Betweenness centrality is from \textcite{freeman1977}.
For applications to media and political networks, see \textcite{watts2002}.
Donor networks in U.S.\ politics are mapped in \textcite{mayer2016}.

\section{Exercises}

\begin{exercise}
Compute the eigenvector centrality of each node in the example network above.
Which node has the highest doxonomic centrality?
\end{exercise}

\begin{exercise}
Using citation data from a policy domain of your choice, construct a donor-institution-media network and identify narrative clusters.
\end{exercise}

\begin{exercise}
Prove that removing a node with high betweenness centrality from a connected doxonomic network increases the number of connected components.
Under what conditions is this guaranteed?
\end{exercise}

\chapter{Causal Inference in Doxonomics}
\label{ch:causal-inference}

\epigraph{Correlation is not causation, but it sure is a hint.}{Edward Tufte}

\noindent
Doxonomics makes causal claims: funding \emph{causes} belief shift, institutions \emph{produce} narratives, exposure \emph{changes} behavior.
These claims require rigorous causal methods, yet the doxonomic setting presents distinctive challenges that standard econometric toolkits do not fully address.
Beliefs are partially observable, treatments are continuous and multi-dimensional, interference between units is the norm rather than the exception, and the actors who generate the ``treatment'' have strategic reasons to obscure the causal chain.
This chapter adapts the modern causal inference toolkit (potential outcomes, difference-in-differences, instrumental variables, causal graphs, synthetic controls, and regression discontinuity) to these distinctive challenges.

We proceed from a methodological conviction: doxonomics cannot content itself with establishing that correlations exist between funding flows and belief distributions.
The field must identify \emph{mechanisms} and estimate \emph{magnitudes}.
Without credible causal identification, doxonomic claims reduce to conspiracy theory dressed in Greek notation.

\section{The Potential Outcomes Framework}
\label{sec:ch16-potential}

\subsection{Basic setup}

\begin{definition}[Doxonomic Treatment Effect]
\label{def:treatment-effect}
\index{treatment effect}
Let $Y_j(1)$ be agent $j$'s belief state under exposure to narrative treatment and $Y_j(0)$ under control.
The \emph{individual doxonomic treatment effect} (IDTE) is
\[
  \tau_j = Y_j(1) - Y_j(0).
\]
The \emph{average doxonomic treatment effect} (ADTE) is
\[
  \mathrm{ADTE} = \E[Y_j(1) - Y_j(0)].
\]
The fundamental problem of causal inference: we observe either $Y_j(1)$ or $Y_j(0)$, never both \autocite{rubin2005}.
\end{definition}

The notation is deceptively simple.
In practice, three features of doxonomic settings complicate the standard Rubin framework.

\begin{enumerate}
  \item \textbf{Multi-valued, continuous treatment.}
  Narrative exposure is not binary but varies in intensity, frequency, source credibility, and framing.
  Let $D_j \in [0, \infty)$ represent a dose of exposure.
  We need dose-response functions:
  \[
    Y_j(d) = g(d, X_j, U_j)
  \]
  where $X_j$ are observed covariates and $U_j$ unobserved heterogeneity.
  The average dose-response function is $\mu(d) = \E[Y_j(d)]$, and the marginal effect at dose $d$ is $\mu'(d)$.

  \item \textbf{Multi-dimensional belief outcomes.}
  A single narrative intervention may shift beliefs along multiple dimensions simultaneously.
  If $\mathbf{Y}_j(d) \in \R^k$ is a vector of $k$ belief dimensions, the treatment effect becomes a vector $\boldsymbol{\tau}_j = \mathbf{Y}_j(1) - \mathbf{Y}_j(0)$, raising multiple-testing issues and requiring that researchers specify which dimension is primary ex ante.

  \item \textbf{Interference.}
  The Stable Unit Treatment Value Assumption (SUTVA) requires that $j$'s outcome depends only on $j$'s own treatment.
  In doxonomic settings, agent $j$'s belief depends on the beliefs (and hence treatments) of $j$'s social contacts.
  Formally, $Y_j = Y_j(D_j, \mathbf{D}_{-j})$, where $\mathbf{D}_{-j}$ includes others' exposures.
  We address interference below in Section~\ref{sec:ch16-interference}.
\end{enumerate}

\subsection{Conditional treatment effects}

The ADTE masks important heterogeneity.
In doxonomic settings, treatment effects are likely to differ systematically across groups.

\begin{definition}[Conditional Average Doxonomic Treatment Effect]
\label{def:cate}
\index{conditional treatment effect}
For a subgroup defined by covariates $X = x$:
\[
  \mathrm{CADTE}(x) = \E[Y_j(1) - Y_j(0) \mid X_j = x].
\]
Of particular interest in doxonomics:
\begin{itemize}[nosep]
  \item \emph{Prior-belief heterogeneity}: do those with strong prior beliefs respond differently than the uncommitted?
  \item \emph{Educational heterogeneity}: does the effect differ by formal education level?
  \item \emph{Source-trust heterogeneity}: does the treatment effect depend on whether the subject trusts the source institution?
\end{itemize}
\end{definition}

A key doxonomic hypothesis is that treatment effects are \emph{non-monotone} in prior belief: agents with moderate prior beliefs are most susceptible, while those at the extremes resist new information (cf.\ the Bayesian updating results in Chapter~\ref{ch:formal-models}).
This generates a U-shaped resistance function that causal methods must be able to detect.

\subsection{Worked example: estimating a think-tank effect}

\begin{example}
\label{ex:think-tank-ate}
A think tank begins publishing policy briefs targeting state legislators.
We model the treatment as binary: $D_j = 1$ if legislator $j$ receives the briefs, $D_j = 0$ otherwise.
The outcome $Y_j$ is the legislator's roll-call voting score on the policy domain the briefs address (e.g., deregulation), measured on a 0--100 scale.

If assignment were random, identification would be trivial:
\[
  \widehat{\mathrm{ADTE}} = \bar{Y}_{\text{treated}} - \bar{Y}_{\text{control}} = 62.3 - 54.7 = 7.6 \text{ points}.
\]
But assignment is not random: the think tank targets swing-district legislators.
Selection bias means $\E[Y_j(0) \mid D_j = 1] \neq \E[Y_j(0) \mid D_j = 0]$: swing legislators may differ from safe-seat legislators in baseline ideology.
The naive estimator confounds the treatment effect with selection.

We need one of the identification strategies developed below.
\end{example}

\section{Natural Experiments in Doxonomic Settings}
\label{sec:ch16-natural}

Natural experiments arise when narrative exposure varies for reasons unrelated to individual characteristics, what \textcite{angrist2009} call ``instruments of fate.''
The doxonomic setting offers several distinctive sources of quasi-random variation.

\subsection{Media market boundaries}

Following \textcite{gentzkow2010}, adjacent counties on opposite sides of a media market boundary receive different television content despite sharing demographic and economic characteristics.
If county $A$ receives Fox News at channel position 2 and county $B$ receives it at channel position 73, channel position creates quasi-random variation in exposure.

\begin{model}[Media Boundary Design]
\label{mod:media-boundary}
\index{media boundary design}
For counties $c$ near boundary $b$:
\[
  Y_c = \alpha + \tau D_c + f(\text{dist}_c) + \gamma X_c + \varepsilon_c
\]
where $D_c$ is exposure, $\text{dist}_c$ is distance to the boundary, $f(\cdot)$ is a smooth function (typically linear or polynomial), and $X_c$ are county-level covariates.
The identifying assumption is that unobserved determinants of $Y$ vary smoothly across the boundary, a form of regression discontinuity in geographic space.
\end{model}

\subsection{Outlet entry and exit}

The introduction or removal of a media outlet provides a before/after comparison.
Historical examples include:
\begin{itemize}[nosep]
  \item the rollout of television across US counties in the 1940s--50s,
  \item the entry of Fox News into cable markets in the late 1990s \autocite{dellavigna2007},
  \item the shutdown of local newspapers (``news deserts'') and subsequent shifts in political knowledge and participation,
  \item the introduction of broadband internet, which shifted the information mix from broadcast to on-demand.
\end{itemize}

\subsection{Platform algorithm changes}

When a social media platform changes its content recommendation algorithm, the resulting shift in exposure is exogenous to any individual user's characteristics.
These events create a useful before/after contrast, though identification requires that no other changes occurred simultaneously, a condition that platform companies rarely guarantee.

\subsection{Policy shocks}

Regulatory changes that affect narrative production or distribution provide natural experiments:
\begin{itemize}[nosep]
  \item the repeal of the Fairness Doctrine (1987) and subsequent growth of partisan talk radio,
  \item foreign-ownership restrictions on media,
  \item campaign finance rulings (e.g., \emph{Citizens United}) that alter the volume of political advertising,
  \item educational curriculum mandates that change what students encounter in schools.
\end{itemize}

Each provides variation in the ``supply side'' of the doxonomic system that is plausibly independent of individual-level belief determinants.

\section{Difference-in-Differences}
\label{sec:ch16-did}

\subsection{Basic design}

\begin{model}[Doxonomic DiD]
\label{mod:did}
\index{difference-in-differences}
Let a new think tank begin operating in region $A$ but not in region $B$ at time $t_0$.
The doxonomic effect is estimated as:
\[
  \hat{\tau} = (\bar{Y}_{A,\text{post}} - \bar{Y}_{A,\text{pre}}) - (\bar{Y}_{B,\text{post}} - \bar{Y}_{B,\text{pre}})
\]
under the parallel trends assumption: absent treatment, beliefs in $A$ and $B$ would have evolved identically.
\end{model}

In regression form:
\[
  Y_{it} = \alpha + \beta \cdot \text{Treat}_i + \gamma \cdot \text{Post}_t + \tau(\text{Treat}_i \times \text{Post}_t) + \varepsilon_{it}
\]
where $\text{Treat}_i = 1$ for units in the treated group, $\text{Post}_t = 1$ for post-treatment periods, and $\tau$ is the DiD estimate of the ADTE.

\subsection{Parallel trends and pre-treatment tests}

The parallel trends assumption is untestable: it concerns the counterfactual path of the treated group.
However, we can probe its plausibility by examining whether treated and control groups followed parallel trajectories \emph{before} treatment.

\begin{proposition}[Pre-Trend Test]
\label{prop:pre-trend}
\index{parallel trends}
If pre-treatment outcomes satisfy
\[
  \E[Y_{it} - Y_{it-1} \mid \text{Treat}_i = 1] = \E[Y_{it} - Y_{it-1} \mid \text{Treat}_i = 0] \quad \text{for all } t < t_0,
\]
this is consistent with (but does not prove) parallel trends.
Rejection provides evidence against the identifying assumption.
\end{proposition}

\begin{center}
\begin{tikzpicture}[scale=0.9]
  \begin{axis}[
    width=10cm, height=6cm,
    xlabel={Time},
    ylabel={Belief measure $\bar{Y}$},
    xmin=-4.5, xmax=4.5,
    ymin=30, ymax=75,
    xtick={-4,-3,-2,-1,0,1,2,3,4},
    xticklabels={$t_{-4}$,$t_{-3}$,$t_{-2}$,$t_{-1}$,$t_0$,$t_1$,$t_2$,$t_3$,$t_4$},
    legend pos=north west,
    legend style={font=\small},
    grid=major,
    grid style={dotted, gray!50},
  ]
    % Control group
    \addplot[blue, thick, mark=square*] coordinates {
      (-4,40) (-3,42) (-2,44) (-1,46) (0,48) (1,50) (2,52) (3,54) (4,56)
    };
    \addlegendentry{Control}

    % Treated group - actual
    \addplot[red, thick, mark=*] coordinates {
      (-4,45) (-3,47) (-2,49) (-1,51) (0,53) (1,59) (2,63) (3,66) (4,68)
    };
    \addlegendentry{Treated (observed)}

    % Treated group - counterfactual
    \addplot[red, dashed, thick] coordinates {
      (0,53) (1,55) (2,57) (3,59) (4,61)
    };
    \addlegendentry{Treated (counterfactual)}

    % Treatment effect arrow
    \draw[-Stealth, thick, doxogreen] (axis cs:4,61) -- (axis cs:4,68);
    \node[right, doxogreen, font=\small] at (axis cs:4.1,64.5) {$\hat{\tau}$};

    % Vertical line at treatment
    \draw[dashed, gray] (axis cs:0,30) -- (axis cs:0,75);
    \node[above, font=\small] at (axis cs:0,75) {Treatment};
  \end{axis}
\end{tikzpicture}
\end{center}

\subsection{Staggered adoption and modern DiD}

In many doxonomic settings, treatment rolls out at different times across units: a think-tank network opens regional offices sequentially, or a media platform enters different markets over several years.
Classical two-way fixed effects (TWFE) regression can produce biased estimates under staggered treatment timing because early-treated units serve as ``controls'' for later-treated units.

Recent methodological advances address this problem.
Following \textcite{callaway2021}, define the \emph{group-time average treatment effect}:
\[
  \mathrm{ATT}(g, t) = \E[Y_t(g) - Y_t(\infty) \mid G_i = g]
\]
where $G_i = g$ denotes units first treated at time $g$, $Y_t(g)$ is the potential outcome under treatment beginning at $g$, and $Y_t(\infty)$ is the never-treated potential outcome.
These group-time effects can be aggregated into an overall summary:
\[
  \hat{\tau}^{\text{agg}} = \sum_{g} \sum_{t \geq g} w_{g,t} \cdot \widehat{\mathrm{ATT}}(g, t)
\]
with weights $w_{g,t}$ proportional to group size.

\subsection{Worked example: think-tank regional entry}

\begin{example}
\label{ex:did-think-tank}
The Prosperity Institute opens offices in three states at staggered intervals: Ohio in 2015, Michigan in 2017, and Pennsylvania in 2019.
Using survey data on beliefs about economic regulation from 2012--2022, we estimate group-time ATTs:

\begin{center}
\begin{tabular}{lccc}
\toprule
& Ohio ($g = 2015$) & Michigan ($g = 2017$) & Pennsylvania ($g = 2019$) \\
\midrule
1 year post  & $+3.2$ (1.4) & $+2.8$ (1.6) & $+4.1$ (1.8) \\
2 years post & $+5.1$ (1.7) & $+4.4$ (1.9) & $+5.8$ (2.3) \\
3 years post & $+6.3$ (2.0) & $+5.9$ (2.2) & -- \\
\bottomrule
\end{tabular}
\end{center}
Standard errors in parentheses, clustered by state.
The aggregated ATT across all groups and post-periods is $\hat{\tau}^{\text{agg}} = 4.6$ (SE $= 1.1$), suggesting the think tank shifted beliefs about economic regulation by roughly 4.6 points on a 100-point scale.
Effects appear to grow over time, consistent with cumulative exposure models from Chapter~\ref{ch:formal-models}.

Note: This example uses simulated data for illustration.
Real applications require actual survey instruments, pre-registration, and sensitivity analyses.
\end{example}

\section{Instrumental Variables}
\label{sec:ch16-iv}

\subsection{The doxonomic endogeneity problem}

In most doxonomic settings, exposure to narratives is endogenous: people select into media, educational, and social environments based on pre-existing beliefs.
This \emph{selective exposure} means that naive regression of belief on exposure confounds the causal effect with self-selection.

\begin{definition}[Doxonomic Instrument]
\label{def:instrument}
\index{instrumental variable}
An \emph{instrument} $Z$ for doxonomic exposure $D$ is a variable that satisfies:
\begin{enumerate}[nosep]
  \item \textbf{Relevance}: $\operatorname{Cov}(Z, D) \neq 0$, meaning the instrument predicts exposure.
  \item \textbf{Exclusion}: $Z$ affects belief $Y$ only through exposure $D$, with no direct effect.
  \item \textbf{Independence}: $Z \perp\!\!\!\perp U$, meaning the instrument is uncorrelated with unobserved confounders.
\end{enumerate}
Under these conditions, the IV estimator identifies the Local Average Treatment Effect (LATE):
\[
  \tau^{\text{LATE}} = \frac{\operatorname{Cov}(Y, Z)}{\operatorname{Cov}(D, Z)} = \frac{\text{reduced form}}{\text{first stage}}
\]
which is the average treatment effect among \emph{compliers}, agents whose exposure is changed by the instrument.
\end{definition}

\subsection{Candidate instruments}

Finding valid instruments is the central challenge of IV estimation.
In doxonomic research, several categories of instruments have been proposed or used:

\begin{enumerate}
  \item \textbf{Geographic instruments.}
  Distance to a university (affects exposure to academic narratives), distance to a think-tank office, proximity to media production centers.
  The concern: geography correlates with urbanization, education, and wealth, creating potential violations of independence.

  \item \textbf{Technological instruments.}
  Cable channel position (lower channel numbers get more viewers), internet connection speed (affects online news consumption), television signal strength.
  These are often more credible because the assignment mechanism is more clearly exogenous.

  \item \textbf{Regulatory instruments.}
  Changes in media ownership rules, campaign finance limits, or educational mandates that shift exposure.
  These affect the supply side of the doxonomic system in ways plausibly independent of individual demand.

  \item \textbf{Weather and timing instruments.}
  Rain on election day (affects turnout and hence exposure to campaign messaging at polling places), natural disasters that preempt regular media programming, scheduling of major sporting events that shift audience attention.
\end{enumerate}

\subsection{Two-stage least squares}

The practical workhorse for IV estimation is two-stage least squares (2SLS):

\medskip
\noindent\textbf{First stage:}
\[
  D_j = \pi_0 + \pi_1 Z_j + \pi_2 X_j + v_j
\]

\noindent\textbf{Second stage:}
\[
  Y_j = \beta_0 + \tau \hat{D}_j + \beta_1 X_j + \varepsilon_j
\]
where $\hat{D}_j$ is the fitted value from the first stage.
The coefficient $\tau$ is the 2SLS estimate of the causal effect.

A critical diagnostic is first-stage strength.
The rule of thumb is that the first-stage $F$-statistic should exceed 10 \autocite{staiger1997}.
In doxonomic applications, weak instruments are a persistent concern because most plausible instruments have only modest effects on narrative exposure.

\subsection{Worked example: channel position instrument}

\begin{example}
\label{ex:iv-channel}
We instrument Fox News viewership with channel position, following \textcite{dellavigna2007}.
In cable systems, channels assigned to lower positions receive more viewers due to habitual channel-surfing patterns.
Channel position is determined by negotiations between cable providers and networks, plausibly independent of local political preferences.

\medskip
\noindent\textbf{First stage}: regress Fox News viewership share on channel position (controlling for demographics):
\[
  \text{FoxShare}_c = 0.23 - 0.0018 \cdot \text{Position}_c + \gamma X_c + v_c
\]
Each 10-position increase reduces Fox viewership by 1.8 percentage points.
First-stage $F = 23.4 > 10$.

\medskip
\noindent\textbf{Second stage}: regress Republican vote share on instrumented Fox viewership:
\[
  \text{RepVote}_c = \alpha + \hat{\tau} \cdot \widehat{\text{FoxShare}}_c + \beta X_c + \varepsilon_c
\]
The 2SLS estimate $\hat{\tau} = 0.28$ (SE $= 0.12$): each percentage point of Fox viewership causes a 0.28 percentage point increase in Republican vote share.
This is larger than the OLS estimate of 0.15, suggesting that OLS is attenuated by measurement error in viewership (or that compliers are more responsive than the average viewer).
\end{example}

\section{Causal Directed Acyclic Graphs}
\label{sec:ch16-dags}

\subsection{The core doxonomic DAG}

Causal DAGs make identification assumptions transparent by encoding them as graphical structures \autocite{pearl2009}.
The core doxonomic DAG links funding, exposure, belief, and behavior:

\begin{center}
\begin{tikzpicture}[
  node distance=2cm,
  every node/.style={draw, circle, thick, minimum size=1.2cm, font=\small},
  arrow/.style={-Stealth, thick}
]
  \node (F) {$F$};
  \node[right=of F] (E) {$E$};
  \node[right=of E] (B) {$B$};
  \node[above=1cm of E] (U) {$U$};
  \node[right=of B] (X) {$X$};

  \draw[arrow] (F) -- (E) node[midway, below, draw=none, font=\footnotesize] {funding};
  \draw[arrow] (E) -- (B) node[midway, below, draw=none, font=\footnotesize] {exposure};
  \draw[arrow] (B) -- (X) node[midway, below, draw=none, font=\footnotesize] {action};
  \draw[arrow, dashed] (U) -- (E);
  \draw[arrow, dashed] (U) -- (B);
\end{tikzpicture}
\end{center}

Here $F$ is funding, $E$ is exposure, $B$ is belief, $X$ is behavior, and $U$ represents unobserved confounders (e.g., pre-existing ideology that affects both media choice and beliefs).
The confounding path $E \leftarrow U \to B$ means naive regression of $B$ on $E$ is biased.
An instrument that affects $E$ but not $U$ identifies the causal effect \autocite{pearl2009, angrist2009}.

\subsection{An expanded DAG}

The core DAG is too simple for most real doxonomic settings.
An expanded version includes institutional structure, social influence, and feedback loops:

\begin{center}
\begin{tikzpicture}[
  node distance=1.8cm,
  every node/.style={draw, rounded corners, thick, minimum size=1cm, font=\small},
  arrow/.style={-Stealth, thick},
  dashed arrow/.style={-Stealth, thick, dashed}
]
  \node (W) {Wealth $W$};
  \node[right=of W] (F) {Funding $F$};
  \node[right=of F] (I) {Institution $I$};
  \node[right=of I] (N) {Narrative $N$};
  \node[below=1.5cm of N] (E) {Exposure $E$};
  \node[left=of E] (S) {Social net $S$};
  \node[below=1.5cm of E] (B) {Belief $B$};
  \node[right=1.5cm of B] (X) {Behavior $X$};
  \node[above=1cm of S] (U) {$U$ (unobs.)};

  \draw[arrow] (W) -- (F);
  \draw[arrow] (F) -- (I);
  \draw[arrow] (I) -- (N);
  \draw[arrow] (N) -- (E);
  \draw[arrow] (S) -- (E);
  \draw[arrow] (E) -- (B);
  \draw[arrow] (B) -- (X);
  \draw[dashed arrow] (U) -- (E);
  \draw[dashed arrow] (U) -- (B);
  \draw[arrow, bend left=20] (B) to (S);
  \draw[arrow, bend right=30] (X) to (W);
\end{tikzpicture}
\end{center}

Two feedback loops appear: $B \to S$ (beliefs shape social networks through homophily) and $X \to W$ (behavior affects wealth, which funds future narrative production).
These feedback loops mean the system is not a DAG in the strict sense. It is a \emph{dynamic causal model} that unfolds over time.
We recover a DAG by indexing variables by time period: $B_t$ can cause $S_{t+1}$ but not $S_t$.

\subsection{Identification via d-separation}

Pearl's d-separation criterion determines which causal effects are identifiable from observational data.

\begin{proposition}[D-Separation for Doxonomic Effects]
\label{prop:d-sep}
\index{d-separation}
In the expanded doxonomic DAG, the causal effect of $F$ on $B$ is identifiable by conditioning on $\{U, S\}$.
Since $U$ is unobserved, we need either:
\begin{enumerate}[nosep]
  \item an instrument for $F$ (affecting $F$ but not $U$ or $S$ directly), or
  \item a proxy for $U$ that satisfies the conditions of proxy variable estimation, or
  \item a design that eliminates $U$ (e.g., panel data with individual fixed effects, if $U$ is time-invariant).
\end{enumerate}
\end{proposition}

The DAG framework clarifies why doxonomic causal inference is hard: the chain $W \to F \to I \to N \to E \to B$ has many links, and confounders can enter at each one.
But it also clarifies the \emph{modularity} of the problem: if we can credibly identify $E \to B$ (the exposure-to-belief link), we can combine it with descriptive evidence on $F \to I \to N \to E$ (the production chain) to construct a complete causal narrative, even if the full chain is not identified in a single study.

\section{Synthetic Control Methods}
\label{sec:ch16-synth}

\subsection{The setup}

When only one or a few units are treated, difference-in-differences lacks statistical power.
Synthetic control methods \autocite{abadie2010} construct a weighted combination of untreated units that closely matches the treated unit's pre-treatment trajectory.

\begin{definition}[Synthetic Control Estimator]
\label{def:synth-control}
\index{synthetic control}
Let unit 1 be treated at time $T_0$.
The synthetic control is a weighted average of $J$ donor units:
\[
  \hat{Y}_{1t}^{(0)} = \sum_{j=2}^{J+1} w_j^* Y_{jt}, \quad t > T_0
\]
where the weights $\mathbf{w}^* = (w_2^*, \ldots, w_{J+1}^*)$ minimize the pre-treatment discrepancy:
\[
  \mathbf{w}^* = \arg\min_{\mathbf{w}} \sum_{t=1}^{T_0} \left( Y_{1t} - \sum_{j=2}^{J+1} w_j Y_{jt} \right)^2 \quad \text{s.t.\ } w_j \geq 0, \; \sum_j w_j = 1.
\]
The treatment effect estimate is $\hat{\tau}_{1t} = Y_{1t} - \hat{Y}_{1t}^{(0)}$ for $t > T_0$.
\end{definition}

\subsection{Doxonomic applications}

Synthetic control is well-suited to studying country-level or state-level doxonomic interventions where the number of treated units is small:
\begin{itemize}[nosep]
  \item the doxonomic effect of a country banning political advertising (treated: Norway; donors: other Nordic/European countries),
  \item the effect of a media mogul acquiring a national newspaper (treated: the newspaper's audience region; donors: comparable regions),
  \item the effect of a new educational curriculum on long-run beliefs (treated: the adopting state; donors: non-adopting states).
\end{itemize}

Inference relies on \emph{placebo tests}: applying the synthetic control procedure to each donor unit in turn and checking whether the treated unit's gap exceeds the placebo distribution.
If the treated unit's post-treatment gap is larger than all or nearly all placebo gaps, the effect is deemed significant.

\subsection{Worked example: advertising ban}

\begin{example}
\label{ex:synth-adban}
Country $A$ bans political television advertising in 2010.
We construct a synthetic version of $A$ from a donor pool of 15 comparable democracies, matching on pre-2010 values of: trust in government, media consumption patterns, GDP per capita, internet penetration, and education levels.

The synthetic $A$ is: $0.31 \times \text{Country B} + 0.24 \times \text{Country C} + 0.22 \times \text{Country D} + 0.23 \times \text{Country E}$.
Pre-treatment fit: root mean squared prediction error $= 1.3$ on a 100-point trust scale.

Post-ban, actual trust in government in $A$ exceeds synthetic $A$ by 4.2 points on average over 2011--2018.
Placebo tests show that $A$'s gap ranks 1st out of 16 (all countries), yielding a placebo $p$-value of $1/16 = 0.0625$, suggestive but not conventionally significant.
This illustrates a common limitation: with few donor units, synthetic control has limited statistical power.
\end{example}

\section{Regression Discontinuity}
\label{sec:ch16-rd}

\subsection{Sharp and fuzzy designs}

Regression discontinuity (RD) exploits a threshold rule that creates a jump in treatment probability at a cutoff value of a running variable.

\begin{definition}[Doxonomic RD Design]
\label{def:rd}
\index{regression discontinuity}
Let $R_j$ be a running variable with cutoff $c$.
In a \emph{sharp} design, treatment is deterministic:
\[
  D_j = \mathbf{1}(R_j \geq c).
\]
In a \emph{fuzzy} design, the probability of treatment jumps at $c$:
\[
  \lim_{r \downarrow c} P(D = 1 \mid R = r) - \lim_{r \uparrow c} P(D = 1 \mid R = r) > 0.
\]
The RD estimator is
\[
  \hat{\tau}_{\text{RD}} = \lim_{r \downarrow c} \E[Y \mid R = r] - \lim_{r \uparrow c} \E[Y \mid R = r]
\]
for the sharp case, and $\hat{\tau}_{\text{fuzzy}} = \hat{\tau}_{\text{RD}} / \hat{\Delta}_D$ for the fuzzy case, where $\hat{\Delta}_D$ is the first-stage jump in treatment.
\end{definition}

\subsection{Doxonomic RD settings}

Several doxonomic settings feature natural thresholds:
\begin{enumerate}[nosep]
  \item \textbf{Electoral thresholds}: parties that barely win elections gain control of government communications infrastructure; parties that barely lose do not.
  Comparing beliefs in jurisdictions where a party barely won vs.\ barely lost identifies the effect of controlling institutional narrative channels.

  \item \textbf{Funding thresholds}: grant programs with score cutoffs create sharp discontinuities in institutional funding.
  Organizations just above the threshold receive funding; those just below do not.

  \item \textbf{Platform thresholds}: content that crosses a virality threshold (e.g., appearing on a ``trending'' list) receives vastly more exposure than content just below.
  This creates a discontinuity in narrative reach.

  \item \textbf{Age thresholds}: legal ages for voting, drinking, or military service create discontinuities in the types of political messaging individuals receive.
\end{enumerate}

\section{Interference and Spillovers}
\label{sec:ch16-interference}

\subsection{Why SUTVA fails in doxonomics}

The Stable Unit Treatment Value Assumption (SUTVA) requires that each unit's outcome depends only on its own treatment assignment.
In doxonomic settings, this assumption is routinely violated: beliefs are socially contagious, and treating agent $j$ may shift the beliefs of $j$'s contacts, their contacts' contacts, and so on.

\begin{definition}[Interference Structure]
\label{def:interference}
\index{interference}
Let $\mathcal{G} = (V, E)$ be the social network.
Agent $j$'s potential outcome depends on the full treatment vector:
\[
  Y_j = Y_j(D_j, \mathbf{D}_{\mathcal{N}(j)})
\]
where $\mathcal{N}(j) = \{k : (j,k) \in E\}$ is $j$'s neighborhood.
Under the \emph{partial interference} assumption, $j$'s outcome depends on $D_j$ and a summary statistic of neighbors' treatments:
\[
  Y_j = Y_j(D_j, g(\mathbf{D}_{\mathcal{N}(j)}))
\]
where $g(\cdot)$ might be the fraction of treated neighbors.
\end{definition}

\subsection{Direct and indirect effects}

With interference, we distinguish:
\begin{align}
  \text{Direct effect:} &\quad \tau^{\text{dir}} = \E[Y_j(1, g) - Y_j(0, g)] \\
  \text{Indirect (spillover) effect:} &\quad \tau^{\text{ind}} = \E[Y_j(d, g') - Y_j(d, g)] \\
  \text{Total effect:} &\quad \tau^{\text{tot}} = \E[Y_j(1, g'(1)) - Y_j(0, g'(0))]
\end{align}
where $g'(d)$ is the neighborhood treatment summary that would obtain if $j$ were assigned treatment $d$.

In doxonomic settings, indirect effects may exceed direct effects: a media campaign's largest impact may come not from those who directly see it, but from the social diffusion of the beliefs it generates.
This has profound implications for cost-effectiveness analysis of doxonomic interventions (cf.\ Chapter~\ref{ch:belief-flow-accounting}).

\subsection{Cluster randomized designs}

One practical solution to interference is to randomize at the cluster level, treating entire communities, media markets, or network components:
\[
  Y_{jc} = \alpha + \tau D_c + \gamma X_{jc} + \varepsilon_{jc}
\]
where $D_c$ is the cluster-level treatment.
Within-cluster spillovers are absorbed into the treatment effect; between-cluster spillovers remain a threat to validity.
The effective sample size is the number of clusters, not the number of individuals, which reduces statistical power substantially.

\section{Sensitivity Analysis and Robustness}
\label{sec:ch16-sensitivity}

\subsection{Omitted variable bias bounds}

No observational study can rule out all confounders.
Following \textcite{oster2019}, we can bound the bias from unobservable confounders:

\begin{proposition}[Oster Bound]
\label{prop:oster}
\index{sensitivity analysis}
Let $\tilde{\tau}$ be the coefficient from a regression of $Y$ on $D$ without controls, $\hat{\tau}$ the coefficient after adding observed controls, and $\tilde{R}^2, \hat{R}^2$ the corresponding $R$-squared values.
Under proportional selection (unobservables relate to treatment in the same proportion as observables), the bias-adjusted estimate is:
\[
  \tau^* = \hat{\tau} - \delta \cdot (\tilde{\tau} - \hat{\tau}) \cdot \frac{R_{\max} - \hat{R}^2}{\hat{R}^2 - \tilde{R}^2}
\]
where $\delta$ measures the relative importance of unobservables and $R_{\max}$ is the hypothetical $R^2$ if all confounders were observed.
If $\tau^*$ remains meaningfully different from zero for reasonable values of $\delta$ and $R_{\max}$, the finding is robust to omitted variable bias.
\end{proposition}

\subsection{Rosenbaum bounds}

For matched designs, Rosenbaum bounds assess how strong an unmeasured confounder would need to be to explain away the observed effect.
The sensitivity parameter $\Gamma$ represents the ratio of treatment odds for two units with the same observed covariates but different values of the unobserved confounder.
If the effect survives at $\Gamma = 2$, a confounder would need to more than double the odds of treatment to nullify the finding.

\subsection{Reporting standards}

We propose that doxonomic causal studies report:
\begin{enumerate}[nosep]
  \item the primary causal estimate with confidence intervals,
  \item at least one sensitivity analysis (Oster bounds, Rosenbaum bounds, or placebo tests),
  \item a discussion of the most plausible threat to identification,
  \item the LATE-vs-ATE distinction: for whom does the estimate apply?
  \item a pre-analysis plan (ideally pre-registered) specifying outcomes, subgroups, and model specifications.
\end{enumerate}

\section{Doxonomic-Specific Identification Challenges}
\label{sec:ch16-challenges}

Beyond the standard econometric challenges, doxonomic causal inference faces several distinctive difficulties.

\subsection{Strategic obscuration}

Unlike natural phenomena, doxonomic ``treatments'' are often designed by strategic actors who have incentives to obscure the causal chain.
Dark money in political spending, astroturf organizations that mimic grassroots movements, and undisclosed native advertising all make it difficult to measure treatment accurately.
Measurement error in the treatment variable biases IV toward zero and DiD toward the null, making it harder to detect real effects.

\subsection{Belief measurement}

Beliefs are not directly observable.
Survey measures are subject to social desirability bias, framing effects, and the distinction between expressed and held beliefs (cf.\ Chapter~\ref{ch:doxometric-measurement}).
If measurement error is \emph{classical} (random noise), it attenuates treatment effect estimates.
If it is \emph{non-classical} (e.g., people misreport beliefs more when politically sensitive), it can bias estimates in either direction.

\subsection{Long causal chains}

The doxonomic causal chain $W \to F \to I \to N \to E \to B \to X$ spans many links.
End-to-end identification (from wealth to behavior) is extremely difficult in a single study.
The modular approach described in Section~\ref{sec:ch16-dags} provides a practical alternative: identify individual links credibly and combine them.
But this requires an additional ``no interaction'' assumption, that the causal effect of each link does not depend on how other links are configured.

\subsection{Feedback and general equilibrium}

Doxonomic systems feature feedback: beliefs affect behavior, behavior affects economic outcomes, economic outcomes affect the resources available for narrative production.
Standard partial-equilibrium causal estimates may miss these amplification (or dampening) effects.
Structural estimation methods (which embed the causal model within a general equilibrium framework) can address this, but at the cost of stronger parametric assumptions.

\section{Notes and References}
\label{sec:ch16-notes}

The potential outcomes framework originates with \textcite{rubin2005}; see \textcite{imbens2015} for a comprehensive treatment.
Difference-in-differences is covered in \textcite{angrist2009}; for the staggered-treatment extensions, see \textcite{callaway2021} and \textcite{dechaisemartin2020}.
Instrumental variables are developed in \textcite{angrist2009}; the weak-instruments problem is analyzed in \textcite{staiger1997}.

DAGs and structural causal models are in \textcite{pearl2009}.
Synthetic control methods are introduced in \textcite{abadie2010}.
Media market instruments for causal identification are used in \textcite{gentzkow2010} and \textcite{dellavigna2007}.

Sensitivity analysis frameworks include \textcite{oster2019} for coefficient stability and \textcite{rosenbaum2002} for matched designs.
The interference literature is surveyed in \textcite{hudgens2008}.
For general overviews of modern causal methods, see \textcite{athey2017} and \textcite{morgan2015}.

On the specific challenges of causal inference in media effects research, see \textcite{dellavigna2010}.
The problem of strategic obscuration in political spending is documented in \textcite{mayer2016}.

\section{Exercises}

\begin{exercise}
Design a difference-in-differences study to estimate the effect of a new cable news channel on beliefs about immigration.
State the parallel trends assumption explicitly.
Discuss at least two plausible threats to parallel trends and how you would test for them using pre-treatment data.
\end{exercise}

\begin{exercise}
Draw a causal DAG for the relationship between economics education, belief in selfish human nature, and cooperative behavior.
Include at least two unobserved confounders.
Identify all confounding paths using d-separation.
Propose two distinct identification strategies and discuss their relative strengths.
\end{exercise}

\begin{exercise}
Propose an instrumental variable for the effect of social media exposure on political beliefs.
Verify the three IV conditions in detail, discussing the most likely violation.
Calculate the first-stage $F$-statistic you would need for the instrument to be credible.
\end{exercise}

\begin{exercise}
Using synthetic control, design a study to estimate the doxonomic effect of a country banning political advertising.
Specify the donor pool, matching variables, and pre-treatment period.
Describe how you would conduct placebo tests and interpret the results.
\end{exercise}

\begin{exercise}
\label{ex:interference}
Consider a randomized experiment where half of a social network is exposed to a media literacy intervention.
\begin{enumerate}[nosep]
  \item Explain why SUTVA is violated.
  \item Define the direct and indirect effects formally.
  \item Propose a cluster-randomized design that addresses interference.
  \item Discuss what the cluster-level estimate captures vs.\ what it misses.
\end{enumerate}
\end{exercise}

\begin{exercise}
A researcher finds that exposure to economic think-tank publications is correlated with pro-deregulation beliefs ($r = 0.35$, $p < 0.001$, $N = 2{,}000$).
\begin{enumerate}[nosep]
  \item List three confounders that could produce this correlation without any causal effect.
  \item For each confounder, describe a research design that would control for it.
  \item Apply the Oster bound: if adding demographics reduces $\hat{\tau}$ from 0.35 to 0.28 and raises $R^2$ from 0.04 to 0.12, how robust is the estimate to further unobservable confounders?
  \item Assume $R_{\max} = 0.5$ and $\delta = 1$. Calculate $\tau^*$.
\end{enumerate}
\end{exercise}

\begin{exercise}[Computational]
Simulate a doxonomic DiD study with the following data-generating process:
\begin{itemize}[nosep]
  \item 200 units, 10 time periods, treatment at $t = 6$ for half the units,
  \item $Y_{it} = 2 + 0.5t + 3 \cdot \text{Treat}_i + \tau \cdot \text{Treat}_i \cdot \text{Post}_t + \varepsilon_{it}$, with $\tau = 5$ and $\varepsilon_{it} \sim N(0, 4)$.
\end{itemize}
\begin{enumerate}[nosep]
  \item Estimate $\hat{\tau}$ using the DiD regression. How close is it to the true value?
  \item Introduce a violation of parallel trends: add $0.3t \cdot \text{Treat}_i$ to the DGP. Re-estimate $\hat{\tau}$. How biased is it?
  \item Plot the pre-treatment trends for treated and control groups in both scenarios. Can you visually detect the violation?
\end{enumerate}
\end{exercise}

\chapter{Text, Media, and Discourse Analysis}
\label{ch:text-media}

\epigraph{Language is not merely a tool for communication but a medium in which we live.}{Hans-Georg Gadamer}

\noindent
Narratives are made of text, images, and speech.
Computational methods allow us to measure their frequency, framing, drift, and concentration at scale.
This chapter introduces the text-as-data methods most relevant to doxonomics.

\section{Corpus Construction}
\label{sec:ch17-corpus}

\begin{definition}[Doxonomic Corpus]
\label{def:corpus}
\index{corpus}
A \emph{doxonomic corpus} for belief $p$ is a collection of documents $\mathcal{D} = \{d_1, \ldots, d_N\}$ drawn from the narrative infrastructure, annotated with metadata: source institution, date, channel, estimated audience, and funding source (where traceable).
\end{definition}

Corpus construction is not neutral.
What you include determines what you can measure.
A corpus of newspaper editorials captures elite narrative production; a corpus of social media posts captures peer-level circulation.
Both are needed for a complete doxonomic analysis.

\section{Framing Analysis}
\label{sec:ch17-framing}

\begin{definition}[Frame]
\label{def:frame}
\index{frame}
A \emph{frame} for proposition $p$ is a function $f: p \to (C, N, S)$ mapping the proposition to a causal story $C$, a normative valence $N \in \{-1, 0, +1\}$, and a subject position $S$.
Different frames for the same proposition can produce opposite doxonomic effects.
\end{definition}

\begin{example}
The proposition ``unemployment is at 8\%'' can be framed as:
\begin{itemize}[nosep]
  \item $f_1$: ``the economy is recovering'' ($C$: trend is improving, $N = +1$),
  \item $f_2$: ``millions remain without work'' ($C$: structural failure, $N = -1$).
\end{itemize}
Same fact, different doxonomic force.
\end{example}

\section{Topic Models}
\label{sec:ch17-topic-models}

\begin{model}[LDA for Doxonomic Analysis]
\label{mod:lda}
\index{topic model}
Latent Dirichlet Allocation \autocite{blei2003} models each document $d$ as a mixture of topics $\theta_d \sim \mathrm{Dir}(\alpha)$, and each topic $k$ as a distribution over words $\phi_k \sim \mathrm{Dir}(\beta)$.
In doxonomic application, topics correspond to narrative themes.
Changes in topic prevalence $\theta_k(t)$ over time reveal which narratives are gaining or losing institutional support.
\end{model}

\section{Semantic Drift}
\label{sec:ch17-drift}

\begin{definition}[Narrative Drift]
\label{def:drift}
\index{semantic drift}
The \emph{semantic drift} of a term or concept $w$ over period $[t_1, t_2]$ is
\[
  D(w) = 1 - \cos(\bm{v}_w(t_1), \bm{v}_w(t_2))
\]
where $\bm{v}_w(t)$ is the word embedding vector trained on the corpus from period $t$.
High drift indicates that the meaning or connotation of $w$ has shifted---a potential marker of doxonomic intervention.
\end{definition}

\section{Information-Theoretic Measures}
\label{sec:ch17-info-theory}

\begin{definition}[Narrative Entropy]
\label{def:narrative-entropy}
\index{narrative entropy}
The \emph{narrative entropy} of a discourse at time $t$ is
\[
  H(t) = -\sum_{k=1}^{K} p_k(t) \log p_k(t)
\]
where $p_k(t)$ is the prevalence of narrative frame $k$.
Low entropy indicates narrative concentration (a few frames dominate); high entropy indicates pluralism \autocite{shannon1948}.
\end{definition}

\begin{proposition}[Entropy and Capture]
\label{prop:entropy-capture}
Common-sense capture (Definition~\ref{def:capture}) implies $H(t) < H^*$ for a critical threshold $H^*$: the narrative space has collapsed to a small number of dominant frames.
Monitoring $H(t)$ over time provides an early warning of doxonomic capture.
\end{proposition}

\section{Notes and References}
\label{sec:ch17-notes}

Text-as-data methods are surveyed in \textcite{grimmer2013} and \textcite{gentzkow2019}.
LDA is introduced in \textcite{blei2003}.
Media slant measurement is developed in \textcite{gentzkow2010}.
Information-theoretic approaches to discourse draw on \textcite{shannon1948}.
For corpus-based social science, see \textcite{salganik2018}.

\section{Exercises}

\begin{exercise}
Construct a small corpus of 20 editorials on a topic of your choice from two newspapers with different editorial stances.
Apply framing analysis to identify the dominant frames in each.
\end{exercise}

\begin{exercise}
Compute the narrative entropy $H(t)$ for U.S.\ media coverage of ``human nature'' using LDA topic proportions from a suitable corpus.
How has entropy changed over the past 30 years?
\end{exercise}

\begin{exercise}
Using word embeddings, measure the semantic drift of the term ``freedom'' between 1950 and 2020 in American political discourse.
Interpret the drift doxonomically.
\end{exercise}

\chapter{Historical and Archival Methods}
\label{ch:historical-archival}

\epigraph{History is the memory of states.}{Henry Kissinger}

\noindent
Doxonomic systems operate over decades and centuries.
Computational methods capture the present; historical methods reveal how we arrived here.
This chapter develops the tools for doxonomic history: genealogy of concepts, institutional archives, long-duration narrative change, and the integration of historical and computational evidence.

\section{Genealogy of Concepts}
\label{sec:ch18-genealogy}

\begin{definition}[Doxonomic Genealogy]
\label{def:genealogy}
\index{genealogy}
A \emph{doxonomic genealogy} of belief $p$ is a reconstruction of:
\begin{enumerate}[nosep]
  \item when $p$ was first articulated as an explicit claim,
  \item which agents promoted it and with what resources,
  \item through which institutional channels it gained credibility,
  \item what alternative beliefs it displaced,
  \item at what point it achieved common-sense status (if ever).
\end{enumerate}
\end{definition}

This extends \textcite{foucault1971} from discourse analysis to funded discourse analysis: not just \emph{who said what when}, but \emph{who paid for it to be said, through what institutions, with what effect}.

\section{Institutional Archives}
\label{sec:ch18-archives}

Key archival sources for doxonomic history:
\begin{itemize}[nosep]
  \item foundation annual reports and grant records,
  \item corporate archives (advertising, internal strategy),
  \item government propaganda agency records,
  \item educational curriculum archives,
  \item media ownership transfer records,
  \item think-tank publication archives.
\end{itemize}

\section{Long-Duration Narrative Change}
\label{sec:ch18-longue-duree}

\begin{model}[Periodization of Belief Regimes]
\label{mod:periodization}
\index{periodization}
A doxonomic period is defined by the dominance of an anthropological regime $\alpha$.
Transitions between periods are marked by:
\begin{itemize}[nosep]
  \item shifts in funding patterns (new donors, new institutional vehicles),
  \item institutional realignment (new think tanks, curriculum reforms),
  \item narrative crises (events that destabilize the dominant frame),
  \item measurable changes in population belief distributions.
\end{itemize}
\end{model}

\begin{example}
The shift from Keynesian common sense (``government manages the economy'') to neoliberal common sense (``markets are self-correcting'') can be periodized as: intellectual groundwork (1947--1970, Mont P\`{e}lerin Society), institutional build-out (1970--1980, Heritage Foundation, Cato Institute, AEI), political capture (1979--1983, Thatcher and Reagan), and normalization (1990--2008, Third Way, Washington Consensus) \autocite{mirowski2009, harvey2005}.
\end{example}

\section{Notes and References}
\label{sec:ch18-notes}

Genealogical method draws on \textcite{foucault1971}.
The institutional history of neoliberalism is documented in \textcite{mirowski2009} and \textcite{mayer2016}.
On long-duration institutional change, see \textcite{north1990} and \textcite{acemoglu2012}.
Archival methods for social science are covered in \textcite{king1994}.

\section{Exercises}

\begin{exercise}
Construct a doxonomic genealogy of the belief ``taxation is theft'' from its philosophical origins to its current institutional support.
\end{exercise}

\begin{exercise}
Using the periodization framework, identify the major doxonomic periods for beliefs about gender roles in the 20th century.
What institutional shifts mark the transitions?
\end{exercise}

\begin{exercise}
Design an archival research project to estimate the total spending on neoliberal narrative production from 1970 to 2000.
What sources would you use?
What are the main obstacles?
\end{exercise}

\chapter{Experimental Doxonomics}
\label{ch:experimental}

\epigraph{The great tragedy of science---the slaying of a beautiful hypothesis by an ugly fact.}{Thomas Henry Huxley}

\noindent
Observational methods identify correlations; experiments identify causes.
This chapter develops experimental designs for testing doxonomic hypotheses in the laboratory and the field.

\section{Cooperation Games as Belief Measurement}
\label{sec:ch19-games}

\begin{definition}[Doxonomic Game]
\label{def:doxonomic-game}
\index{doxonomic game}
A \emph{doxonomic game} is an experimental economic game (public goods, dictator, trust, ultimatum) administered after exposure to a doxonomic treatment (narrative about human nature, institutional framing, etc.).
The outcome variable is the degree of cooperative or self-interested behavior, interpreted as a \emph{behavioral proxy} for internalized belief.
\end{definition}

\textcite{fischbacher2001} found that most people are conditional cooperators---not the selfish maximizers predicted by standard economic models.
\textcite{henrich2001} showed this holds cross-culturally.
The doxonomic question is: does exposure to selfishness narratives \emph{reduce} cooperation?

\section{Priming and Framing Experiments}
\label{sec:ch19-priming}

\begin{model}[Narrative Priming]
\label{mod:priming}
\index{priming experiment}
Randomly assign subjects to:
\begin{itemize}[nosep]
  \item \textbf{Treatment}: read a text asserting that human nature is fundamentally self-interested,
  \item \textbf{Control}: read a neutral text or one asserting cooperative human nature.
\end{itemize}
Then measure behavior in a public goods game.
The treatment effect $\hat{\tau} = \bar{Y}_{\text{treat}} - \bar{Y}_{\text{control}}$ estimates the causal effect of a single narrative exposure on cooperative behavior.
\end{model}

\section{Norm Perception Studies}
\label{sec:ch19-norms}

\begin{definition}[Perceived Social Norm]
\label{def:perceived-norm}
\index{perceived norm}
The \emph{perceived social norm} for behavior $x$ is agent $j$'s belief about the population's behavior: $\hat{\mu}_j(x)$.
Doxonomic systems can shift behavior by shifting norm perceptions---making people think others are more selfish than they actually are.
\end{definition}

\begin{proposition}[Norm Misperception Effect]
\label{prop:norm-misperception}
If agents are conditional cooperators (cooperate when they believe others cooperate), then a doxonomic intervention that increases $\hat{\mu}_j(\text{selfish})$ without changing actual behavior will reduce cooperation in subsequent interactions.
The belief becomes self-fulfilling.
\end{proposition}

\section{Ethical Limits}
\label{sec:ch19-ethics}

Experimental doxonomics faces distinctive ethical constraints:
\begin{itemize}[nosep]
  \item exposing subjects to false narratives about human nature may have lasting effects,
  \item deception about the social environment may erode trust,
  \item experiments that increase selfishness ``work'' by making people worse off.
\end{itemize}

These constraints limit the scope of laboratory doxonomics but do not eliminate it.
Natural experiments and field studies with minimal deception remain viable.

\section{Notes and References}
\label{sec:ch19-notes}

Public goods experiments are reviewed in \textcite{fischbacher2001}.
Cross-cultural behavioral games are in \textcite{henrich2001}.
Priming effects on economic behavior are studied in the behavioral economics literature; see \textcite{kahneman2011} for an overview.
On norm misperception, see \textcite{cialdini2006}.
Experimental ethics in social science is discussed in \textcite{salganik2018}.

\section{Exercises}

\begin{exercise}
Design a priming experiment to test whether exposure to ``greed is good'' narratives increases defection in a public goods game.
Specify sample size, treatment arms, and outcome measures.
\end{exercise}

\begin{exercise}
A norm perception study finds that 80\% of people believe ``most people are selfish,'' but behavioral games show only 30\% defect.
Compute the norm misperception gap and predict its doxonomic consequences.
\end{exercise}

\begin{exercise}
Discuss the ethics of an experiment that exposes children to competing anthropological narratives (selfish vs.\ cooperative) and measures long-term behavioral effects.
Under what conditions (if any) would this be justifiable?
\end{exercise}

\chapter{Model Validation and Epistemic Humility}
\label{ch:model-validation}

\epigraph{The only useful function of a statistician is to make predictions, and thus to provide a basis for action.}{W.\ Edwards Deming}

\noindent
Doxonomics makes strong claims about the causal structure of belief: that material resources shape public opinion through identifiable institutional mechanisms.
These claims must be testable, falsifiable, and honest about their limitations.
A field that studies how power corrupts knowledge production must be especially vigilant about how its own knowledge production might be corrupted---by ideological commitment, confirmation bias, unfalsifiable theorizing, or the temptation to see doxonomic causation in every correlation between money and belief.

This chapter develops the tools for validating doxonomic models---uncertainty quantification, out-of-sample prediction, sensitivity analysis, triangulation, and Bayesian model comparison---and the intellectual discipline to admit when those models fail.
It is the methodological conscience of the book.

\section{Uncertainty Quantification}
\label{sec:ch20-uncertainty}

\subsection{The uncertainty imperative}

\begin{principle}[Doxonomic Uncertainty]
\label{prin:uncertainty}
\index{uncertainty quantification}
Every doxonomic estimate---of funding flows, belief prevalence, causal effects, model parameters---must be accompanied by a quantified uncertainty.
Point estimates without uncertainty intervals are not doxonomic science; they are doxonomic assertion.
\end{principle}

The principle sounds banal---every statistics course teaches confidence intervals---but in practice, doxonomic research faces distinctive challenges in uncertainty quantification:
\begin{itemize}[nosep]
  \item funding data is incomplete (dark money, international flows, in-kind contributions), so the inputs to doxonomic models carry large measurement error,
  \item belief data comes from surveys with known biases (social desirability, framing effects), so the outputs carry systematic error,
  \item causal identification is rarely bulletproof, so structural uncertainty (``is this the right model?'') dominates statistical uncertainty (``given this model, how precise is the estimate?'').
\end{itemize}

\subsection{Frequentist and Bayesian frameworks}

\begin{definition}[Doxonomic Uncertainty Interval]
\label{def:uncertainty-interval}
\index{uncertainty interval}
A $100(1-\alpha)\%$ uncertainty interval for a doxonomic parameter $\theta$ can be constructed in two frameworks:
\begin{itemize}[nosep]
  \item \textbf{Frequentist (confidence interval)}: the procedure $[\hat{\theta}_L, \hat{\theta}_U]$ is calibrated so that, over repeated sampling, the interval contains the true $\theta$ at least $100(1-\alpha)\%$ of the time.
  The probability statement is about the procedure, not about $\theta$ lying in any particular realized interval.
  \item \textbf{Bayesian (credible interval)}: given a prior $\pi(\theta)$ and data $y$, the interval $[\theta_L, \theta_U]$ satisfies $P(\theta_L \leq \theta \leq \theta_U \mid y) \geq 1 - \alpha$ under the posterior.
  The probability statement is about $\theta$ conditional on the observed data.
\end{itemize}
Doxonomic research should specify which framework is in use.
\end{definition}

Given the typically small samples, strong prior information, and complex hierarchical structures in doxonomic studies, Bayesian methods are often more natural.
A Bayesian framework allows the analyst to incorporate prior knowledge about funding magnitudes, belief distributions, and institutional behavior while updating transparently as data arrives.

\subsection{Propagating uncertainty through the doxonomic chain}

The doxonomic causal chain $F \to I \to N \to E \to B$ involves multiple estimated quantities at each link.
Uncertainty at each link compounds:

\begin{proposition}[Uncertainty Propagation]
\label{prop:uncertainty-propagation}
\index{uncertainty propagation}
If the doxonomic effect of funding $F$ on belief $B$ is the product of link-specific effects:
\[
  \frac{\partial B}{\partial F} = \frac{\partial B}{\partial E} \cdot \frac{\partial E}{\partial N} \cdot \frac{\partial N}{\partial I} \cdot \frac{\partial I}{\partial F}
\]
and each link is estimated with independent uncertainty, then the total variance of the end-to-end effect is (by the delta method):
\[
  \operatorname{Var}\!\left(\frac{\partial B}{\partial F}\right) \approx \left(\frac{\partial B}{\partial F}\right)^2 \sum_{k} \left(\frac{\sigma_k}{\hat{\beta}_k}\right)^2
\]
where $\hat{\beta}_k$ and $\sigma_k$ are the estimate and standard error at link $k$.
The relative uncertainty of the total effect is approximately the square root of the sum of squared relative uncertainties at each link.
\end{proposition}

This means that even if each individual link is estimated with moderate precision (say, a coefficient of variation of 30\%), the end-to-end estimate may have enormous uncertainty.
For a four-link chain with 30\% relative uncertainty at each link, the total relative uncertainty is $\sqrt{4 \times 0.09} = 0.60$---a 60\% coefficient of variation, yielding a very wide confidence interval.
This is a fundamental constraint on doxonomic measurement, not a failure of technique.

\subsection{Reporting standards}

We propose minimum reporting standards for doxonomic uncertainty:

\begin{enumerate}[nosep]
  \item \textbf{Primary estimates}: point estimate $\pm$ 95\% interval, with the framework (frequentist or Bayesian) stated.
  \item \textbf{Data uncertainty}: document missing data rates for funding, belief, and exposure variables.
  \item \textbf{Model uncertainty}: report results under at least two plausible model specifications.
  \item \textbf{Identification uncertainty}: state the identifying assumptions and assess their plausibility.
  \item \textbf{Propagated uncertainty}: when combining estimates across links, report the compounded interval.
\end{enumerate}

\section{Overfitting and the Ideological Hazard}
\label{sec:ch20-overfitting}

\subsection{The problem}

\begin{definition}[Doxonomic Overfitting]
\label{def:overfitting}
\index{overfitting}
A doxonomic model \emph{overfits} if it explains observed patterns in a specific dataset (e.g., media coverage correlates with funding) but fails to predict out-of-sample---because the apparent pattern reflects noise, selection, or post hoc rationalization rather than genuine causal structure.
\end{definition}

Doxonomics is especially vulnerable to overfitting for three reasons:

\begin{enumerate}
  \item \textbf{The dependent variable is over-determined.}
  Public belief is influenced by funding, institutions, media, education, personal experience, peer effects, economic conditions, historical memory, and random events.
  Any model that includes only doxonomic variables will have low $R^2$, tempting analysts to mine the specification until something ``works.''

  \item \textbf{Funding data is incomplete and noisy.}
  Missing data on dark money, in-kind contributions, and indirect funding creates measurement error that can generate spurious patterns when combined with flexible model specifications.

  \item \textbf{The analyst has priors.}
  A researcher drawn to doxonomics is likely to believe, prior to analysis, that power shapes belief.
  This prior creates a confirmation bias that makes it easier to accept evidence consistent with doxonomic explanations and harder to accept null results.
\end{enumerate}

\subsection{The ideological hazard}

\begin{principle}[Ideological Hazard]
\label{prin:ideological-hazard}
\index{ideological hazard}
A discipline that studies how ideology distorts knowledge production must acknowledge the possibility that it is itself subject to ideological distortion.
Doxonomics is not immune to the forces it studies.
The ideological hazard is the temptation to:
\begin{itemize}[nosep]
  \item see doxonomic causation in every correlation between money and belief,
  \item attribute all belief to manipulation while ignoring genuine persuasion, rational updating, and independent thought,
  \item treat one's own beliefs as exempt from doxonomic influence,
  \item present speculative hypotheses with the confidence appropriate to well-identified causal claims.
\end{itemize}
\end{principle}

The ideological hazard does not invalidate doxonomics---the same hazard exists in any social science that studies power---but it requires institutional safeguards: pre-registration, adversarial collaboration, publication of null results, and the kind of epistemic humility developed in Section~\ref{sec:ch20-humility}.

\section{Out-of-Sample Prediction}
\label{sec:ch20-prediction}

\subsection{The prediction test}

The most powerful safeguard against overfitting is out-of-sample prediction.
A model that can predict doxonomic outcomes in data it has not seen has demonstrated that it has captured genuine structure, not noise.

\begin{definition}[Predictive Validation]
\label{def:predictive-validation}
\index{predictive validation}
Split the available data into training set $\mathcal{D}_{\text{train}}$ and test set $\mathcal{D}_{\text{test}}$ (or use $k$-fold cross-validation).
Fit the doxonomic model on $\mathcal{D}_{\text{train}}$ and evaluate on $\mathcal{D}_{\text{test}}$ using:
\begin{itemize}[nosep]
  \item \textbf{Root mean squared prediction error} (RMSPE) for continuous outcomes:
  \[
    \text{RMSPE} = \sqrt{\frac{1}{|\mathcal{D}_{\text{test}}|} \sum_{i \in \mathcal{D}_{\text{test}}} (Y_i - \hat{Y}_i)^2}
  \]
  \item \textbf{Classification accuracy} (or AUC) for binary belief outcomes.
  \item \textbf{Log predictive density} for Bayesian models: $\sum_{i \in \mathcal{D}_{\text{test}}} \log p(Y_i \mid \mathcal{D}_{\text{train}})$.
\end{itemize}
Compare the doxonomic model's prediction against a \emph{null model} that uses only non-doxonomic predictors (demographics, economic conditions, prior beliefs).
The doxonomic model must outperform the null to justify its additional complexity.
\end{definition}

\subsection{Temporal out-of-sample tests}

In doxonomic settings, temporal splits are more informative than random splits.
Train the model on data up to year $t$ and predict outcomes for years $t+1, t+2, \ldots$
This tests whether the model's structural relationships are stable over time---a stronger test than cross-sectional prediction.

\begin{example}
\label{ex:temporal-prediction}
A model relates think-tank funding to public opinion on climate change, estimated on 1990--2010 data.
The model predicts that the decline in fossil-fuel think-tank funding after 2015 should lead to a modest increase in climate concern.
Checking this prediction against 2015--2022 survey data:

\begin{center}
\begin{tabular}{lcc}
\toprule
Year & Predicted climate concern & Actual climate concern \\
\midrule
2016 & 62.3\% & 64.1\% \\
2018 & 64.7\% & 67.2\% \\
2020 & 66.1\% & 70.8\% \\
2022 & 67.4\% & 72.3\% \\
\bottomrule
\end{tabular}
\end{center}
The model correctly predicts the direction and approximate magnitude of the change but systematically underpredicts, suggesting that factors beyond think-tank funding (e.g., direct experience of extreme weather, youth activism) contributed to the belief shift.
The prediction error is informative: it reveals the limits of a funding-only doxonomic model.
(Illustrative data.)
\end{example}

\subsection{Cross-national out-of-sample tests}

An even stronger test: estimate the model on one country and predict outcomes in another.
If the doxonomic model captures generalizable mechanisms, a model estimated on US data should have some predictive power for UK or Australian outcomes (after adjusting for institutional differences).
Failure to generalize cross-nationally suggests that the model has overfit to country-specific features.

\section{Sensitivity Analysis}
\label{sec:ch20-sensitivity}

\subsection{Rosenbaum bounds revisited}

\begin{model}[Sensitivity Analysis for Doxonomic Claims]
\label{mod:sensitivity}
\index{sensitivity analysis}
Following \textcite{rosenbaum2002}, a sensitivity analysis asks: how large would an unobserved confounder $U$ have to be to explain away the estimated doxonomic effect?
The sensitivity parameter $\Gamma$ represents the maximum ratio of treatment odds for two units with the same observed covariates:
\[
  \frac{1}{\Gamma} \leq \frac{P(D=1 \mid X, U=1) / P(D=0 \mid X, U=1)}{P(D=1 \mid X, U=0) / P(D=0 \mid X, U=0)} \leq \Gamma.
\]
At $\Gamma = 1$, there is no unobserved confounding.
As $\Gamma$ increases, the bounds on the treatment effect widen.
The \emph{breaking point} $\Gamma^*$ is the smallest $\Gamma$ at which the doxonomic effect is no longer statistically significant.
A claim with $\Gamma^* = 1.2$ is fragile; one with $\Gamma^* = 3.0$ is robust.
\end{model}

\subsection{The Oster test}

\textcite{oster2019} provides an alternative sensitivity framework based on coefficient stability:

\begin{proposition}[Coefficient Stability Test]
\label{prop:oster-test}
\index{coefficient stability}
Let $\tilde{\tau}$ be the doxonomic effect estimate from a short regression (without controls), $\hat{\tau}$ the estimate from a long regression (with observable controls), and $\tilde{R}^2, \hat{R}^2$ the corresponding $R$-squared values.
Under the assumption that observable and unobservable confounders are equally important ($\delta = 1$) and that the maximum attainable $R^2$ is $R_{\max}$, the bias-adjusted estimate is:
\[
  \tau^* = \hat{\tau} - (\tilde{\tau} - \hat{\tau}) \cdot \frac{R_{\max} - \hat{R}^2}{\hat{R}^2 - \tilde{R}^2}.
\]
If $\tau^* > 0$ (or retains the hypothesized sign) for $R_{\max} = 1.3 \hat{R}^2$ and $\delta = 1$, the estimate is considered robust.
\end{proposition}

\begin{example}
\label{ex:oster-application}
A study finds that think-tank proximity predicts pro-deregulation beliefs:
\begin{itemize}[nosep]
  \item Short regression (no controls): $\tilde{\tau} = 0.42$, $\tilde{R}^2 = 0.05$.
  \item Long regression (with demographics, income, education, urban/rural): $\hat{\tau} = 0.28$, $\hat{R}^2 = 0.18$.
  \item With $R_{\max} = 1.3 \times 0.18 = 0.234$ and $\delta = 1$:
  \[
    \tau^* = 0.28 - (0.42 - 0.28) \cdot \frac{0.234 - 0.18}{0.18 - 0.05} = 0.28 - 0.14 \cdot \frac{0.054}{0.13} = 0.28 - 0.058 = 0.222.
  \]
\end{itemize}
The bias-adjusted estimate $\tau^* = 0.222$ remains positive and substantively meaningful.
Even if unobserved confounders are as important as observed ones, the doxonomic effect survives.
\end{example}

\subsection{Placebo tests}

\begin{definition}[Doxonomic Placebo Test]
\label{def:placebo}
\index{placebo test}
A \emph{placebo test} applies the doxonomic estimation procedure to an outcome that should \emph{not} be affected by the hypothesized mechanism.
If a significant effect is found on the placebo outcome, the original finding is likely driven by confounding rather than the doxonomic mechanism.
\end{definition}

\begin{example}
\label{ex:placebo}
A study claims that fossil-fuel funding shifted beliefs about climate change.
Placebo tests:
\begin{enumerate}[nosep]
  \item Does fossil-fuel funding predict beliefs about \emph{evolution}?
  If yes, the ``effect'' may reflect general conservative ideology, not a specific doxonomic mechanism.
  \item Does fossil-fuel funding predict beliefs about climate change \emph{before} the funding began?
  If yes, the relationship is driven by pre-existing differences (selection), not the funding.
  \item Does fossil-fuel funding predict climate beliefs in countries where the funded institutions have no presence?
  If yes, the correlation reflects global trends, not institutional influence.
\end{enumerate}
Passing all three placebo tests substantially strengthens the doxonomic claim.
\end{example}

\section{Bayesian Model Comparison}
\label{sec:ch20-bayesian}

\subsection{The model comparison problem}

Doxonomic claims typically involve choosing among competing explanations for the same observed pattern:
\begin{itemize}[nosep]
  \item $M_1$: funding causes belief change (the doxonomic explanation),
  \item $M_2$: ideology drives both funding and belief (confounding),
  \item $M_3$: belief change drives funding (reverse causation---funders invest where they see receptive audiences),
  \item $M_4$: no systematic relationship (coincidence).
\end{itemize}

Bayesian model comparison provides a principled framework for adjudicating among these.

\begin{definition}[Bayes Factor]
\label{def:bayes-factor}
\index{Bayes factor}
The Bayes factor comparing model $M_1$ to model $M_2$ is:
\[
  \text{BF}_{12} = \frac{P(\text{data} \mid M_1)}{P(\text{data} \mid M_2)} = \frac{\int P(\text{data} \mid \theta_1, M_1) \pi(\theta_1 \mid M_1) \, d\theta_1}{\int P(\text{data} \mid \theta_2, M_2) \pi(\theta_2 \mid M_2) \, d\theta_2}.
\]
A Bayes factor of $\text{BF}_{12} > 10$ provides ``strong'' evidence for $M_1$; $\text{BF}_{12} > 100$ provides ``decisive'' evidence (Jeffreys' scale).
The Bayes factor automatically penalizes model complexity through the prior-weighted integral, implementing a natural Occam's razor.
\end{definition}

\subsection{Comparing doxonomic and non-doxonomic models}

The critical comparison is $M_1$ (doxonomic model: funding $\to$ belief) vs.\ $M_2$ (confounding model: ideology $\to$ funding and ideology $\to$ belief):

\begin{example}
\label{ex:bayes-factor}
Using panel data on think-tank funding and regional beliefs about deregulation:

\begin{center}
\begin{tabular}{lcc}
\toprule
Model & Log marginal likelihood & Bayes factor vs.\ $M_4$ \\
\midrule
$M_1$: funding $\to$ belief & $-3{,}412.7$ & $245.3$ \\
$M_2$: ideology confounds & $-3{,}415.1$ & $22.8$ \\
$M_3$: reverse causation & $-3{,}418.9$ & $0.5$ \\
$M_4$: no relationship & $-3{,}419.6$ & $1.0$ (reference) \\
\bottomrule
\end{tabular}
\end{center}

Both $M_1$ and $M_2$ are strongly supported over no-relationship.
The Bayes factor of $M_1$ over $M_2$ is $245.3 / 22.8 = 10.8$, providing strong (but not decisive) evidence for the doxonomic over the confounding explanation.
The reverse-causation model is not supported.
(Illustrative data.)
\end{example}

\subsection{Model averaging}

When multiple models are plausible, Bayesian model averaging (BMA) provides predictions that account for model uncertainty:
\[
  P(\theta \mid \text{data}) = \sum_{m=1}^{M} P(\theta \mid \text{data}, M_m) \cdot P(M_m \mid \text{data})
\]
where $P(M_m \mid \text{data})$ is the posterior model probability.
This is especially appropriate for doxonomics because the correct model is genuinely uncertain---confounding, reverse causation, and doxonomic causation may all operate simultaneously.

\section{Triangulation}
\label{sec:ch20-triangulation}

\subsection{The convergence criterion}

\begin{principle}[Doxonomic Triangulation]
\label{prin:triangulation}
\index{triangulation}
A doxonomic claim is strengthened when it is supported by multiple independent lines of evidence.
The key categories:
\begin{enumerate}[nosep]
  \item \textbf{Accounting evidence}: money was spent (foundation records, tax filings).
  \item \textbf{Production evidence}: narratives were produced (publications, media content).
  \item \textbf{Network evidence}: institutional connections exist between funder and producer.
  \item \textbf{Causal evidence}: exposure changes beliefs (experiments, natural experiments).
  \item \textbf{Textual evidence}: narrative content shifted in the predicted direction (text analysis).
  \item \textbf{Historical evidence}: the timing of funding precedes the timing of belief change.
  \item \textbf{Comparative evidence}: similar funding produces similar effects in different contexts.
\end{enumerate}
No single line is sufficient.
The claim is credible when multiple independent lines converge on the same conclusion.
\end{principle}

\subsection{Formal triangulation assessment}

\begin{model}[Triangulation Scorecard]
\label{mod:triangulation-score}
\index{triangulation scorecard}
For a specific doxonomic claim, rate the evidence on each line as: strong ($++$), moderate ($+$), weak ($\circ$), absent ($-$), or contradictory ($--$).
The overall assessment:
\begin{itemize}[nosep]
  \item \textbf{Established}: 4+ lines at strong or moderate, none contradictory.
  \item \textbf{Supported}: 2--3 lines at strong or moderate, none contradictory.
  \item \textbf{Suggestive}: 1--2 lines at moderate, others absent.
  \item \textbf{Speculative}: only weak evidence or single line.
  \item \textbf{Challenged}: at least one contradictory line.
\end{itemize}
\end{model}

\begin{example}
\label{ex:triangulation}
Scorecard for ``fossil-fuel funding produced climate doubt'':

\begin{center}
\begin{tabular}{lcc}
\toprule
Evidence line & Rating & Key source \\
\midrule
Accounting & $++$ & Tax filings, litigation discovery \\
Production & $++$ & Think-tank publications, op-eds \\
Network & $++$ & Board overlaps, funding flows \\
Causal & $+$ & Natural experiments, some RCTs \\
Textual & $++$ & ``Uncertainty'' frame six times more prevalent \\
Historical & $++$ & Internal memos date strategy to 1989 \\
Comparative & $+$ & Weaker in countries without think-tank network \\
\bottomrule
\end{tabular}
\end{center}
Assessment: \textbf{Established}.
Six lines at strong or moderate, none contradictory.
This is one of the best-documented doxonomic cases because litigation discovery provided archival evidence that is normally inaccessible.
\end{example}

\section{Falsifiability and Degenerating Research Programs}
\label{sec:ch20-falsifiability}

\subsection{What would falsify doxonomics?}

A field that cannot specify conditions under which its core claims would be falsified is not science.
Doxonomics should state explicitly what observations would count as evidence against its hypotheses:

\begin{enumerate}
  \item \textbf{Funding-belief independence}: if careful causal studies consistently find no relationship between funding and belief (after correcting for confounding), the core doxonomic hypothesis is falsified.

  \item \textbf{Symmetry failure}: if well-funded narratives are no more prevalent than poorly funded narratives of equal quality, the power-shapes-belief mechanism is not operative.

  \item \textbf{Resilience of reason}: if populations consistently resist funded narratives and converge on true beliefs regardless of the institutional environment, the doxonomic model of belief formation is wrong.

  \item \textbf{Reverse dominance}: if belief change consistently precedes (rather than follows) funding and institutional activity, the causal arrow runs the other direction.
\end{enumerate}

None of these falsifications would be established by a single study---they would require a pattern of findings across multiple contexts and methods.

\subsection{Lakatos and degenerating programs}

Following \textcite{lakatos1978}, a research program \emph{progresses} if it successfully predicts novel facts and \emph{degenerates} if it can only accommodate observations post hoc.
Doxonomics should be evaluated by this standard:

\begin{itemize}[nosep]
  \item A \emph{progressive} sign: a doxonomic model estimated on 1990--2010 data correctly predicts the direction and approximate magnitude of belief change in 2010--2020.
  \item A \emph{degenerating} sign: every counter-example to a doxonomic claim is explained away by ad hoc auxiliary hypotheses (``the funding was dark money we can't trace,'' ``the effect operates at a level we can't measure,'' ``the alternative explanation is itself a doxonomic product'').
\end{itemize}

The danger of degeneration is especially acute for doxonomics because its subject matter---the production of belief by power---provides a built-in excuse for any failure: ``the reason you don't see the doxonomic effect is that the doxonomic system has prevented you from seeing it.''
This is the epistemic equivalent of an unfalsifiable conspiracy theory, and guarding against it requires the formal validation tools of this chapter.

\section{Epistemic Humility in Practice}
\label{sec:ch20-humility}

\subsection{The humility checklist}

\begin{principle}[Epistemic Humility]
\label{prin:humility}
\index{epistemic humility}
Every doxonomic publication should explicitly state:
\begin{enumerate}[nosep]
  \item \textbf{What it has shown}: the specific empirical finding, described in modest terms (correlation, association, estimated causal effect under stated assumptions).
  \item \textbf{What it assumes}: the identifying assumptions, measurement assumptions, and theoretical priors that underpin the analysis.
  \item \textbf{What it has not shown}: the causal claims that the evidence does not support, even if the analyst suspects they are true.
  \item \textbf{How it could be wrong}: the most plausible alternative explanation for the observed pattern.
  \item \textbf{What would change the conclusion}: the specific evidence that would cause the author to abandon or substantially revise the claim.
\end{enumerate}
Over-claiming undermines the field's credibility---its own epistemic capital.
A field that studies the corruption of epistemic authority must guard its own epistemic authority with unusual care.
\end{principle}

\subsection{Graduated confidence language}

We propose a graduated vocabulary for doxonomic claims, tied to the triangulation scorecard:

\begin{center}
\begin{tabular}{ll}
\toprule
Triangulation level & Appropriate language \\
\midrule
Established & ``The evidence strongly suggests that\ldots'' \\
Supported & ``There is substantial evidence that\ldots'' \\
Suggestive & ``Preliminary evidence is consistent with\ldots'' \\
Speculative & ``It is plausible that\ldots'' or ``We hypothesize that\ldots'' \\
Challenged & ``The evidence is mixed; some findings suggest\ldots'' \\
\bottomrule
\end{tabular}
\end{center}

This vocabulary is not merely stylistic.
It encodes real epistemic information and allows readers to calibrate their confidence appropriately.
A field-wide norm of graduated language would do more to protect doxonomics from overclaiming than any number of statistical tests.

\subsection{The symmetry obligation}

\begin{principle}[Analytical Symmetry]
\label{prin:symmetry}
\index{analytical symmetry}
Doxonomic analysis should be applied symmetrically to all belief-producing institutions, not selectively to those the researcher opposes.
If the tools developed in this book can show that fossil-fuel funding shaped climate beliefs, the same tools should be applied to environmental foundation funding, government science communication, and academic knowledge production.
The asymmetric application of doxonomic analysis---scrutinizing only one side's narrative infrastructure---is itself a form of doxonomic bias.
\end{principle}

This is not a ``both sides'' equivalence.
Different institutions may have very different levels of doxonomic power, and the analysis may well find that one side's infrastructure is vastly larger or more effective.
But the \emph{method} of analysis must be applied symmetrically; only the \emph{results} should be asymmetric.

\section{Common Pitfalls and How to Avoid Them}
\label{sec:ch20-pitfalls}

We conclude with a catalog of the most common errors in doxonomic reasoning, distilled from the methodological principles developed throughout this part of the book.

\begin{enumerate}
  \item \textbf{Post hoc ergo propter hoc.}
  ``Think tank was founded in 1985; beliefs shifted in 1990; therefore the think tank caused the shift.''
  Temporal precedence is necessary but not sufficient for causation.
  Remedy: use the causal methods of Chapter~\ref{ch:causal-inference}.

  \item \textbf{Guilt by association.}
  ``Institution $X$ receives funding from corporation $Y$; therefore $X$'s conclusions are wrong.''
  Funding creates a \emph{suspicion} of bias, not proof of it.
  The doxonomic question is whether funding predicts framing \emph{conditional on the evidence}---not whether funded conclusions are false.
  Remedy: distinguish between the genetic fallacy (dismissing claims because of their origin) and legitimate doxonomic analysis (studying how origin correlates with content).

  \item \textbf{Monocausal explanation.}
  ``Belief $p$ exists because of doxonomic campaign $X$.''
  Real beliefs have multiple causes: personal experience, peer influence, institutional exposure, economic conditions, and yes, funded narrative production.
  The doxonomic contribution is one factor among many.
  Remedy: estimate the \emph{partial} effect of doxonomic variables, controlling for other factors.

  \item \textbf{Unfalsifiable theorizing.}
  ``If you can't see the doxonomic influence, that's because the influence is so total that it has become invisible.''
  This renders the theory immune to evidence.
  Remedy: specify observable implications and test them (Section~\ref{sec:ch20-falsifiability}).

  \item \textbf{The totalizing error.}
  ``All belief is socially produced; therefore no belief is genuinely held.''
  Showing that a belief has social origins does not show that it is false or that the believer is a puppet.
  People are agents who process, filter, resist, and transform the narratives they encounter.
  Remedy: model agent heterogeneity in belief response (Chapter~\ref{ch:formal-models}).

  \item \textbf{Asymmetric scrutiny.}
  Applying doxonomic analysis only to beliefs one disagrees with.
  Remedy: the symmetry obligation (Principle~\ref{prin:symmetry}).
\end{enumerate}

\section{Notes and References}
\label{sec:ch20-notes}

Model selection and validation are covered in \textcite{freedman2009} and \textcite{gelman2020}.
Sensitivity analysis via Rosenbaum bounds is developed in \textcite{rosenbaum2002}; the Oster test is in \textcite{oster2019}.
On the limits of causal inference in social science, see \textcite{manski1993} and \textcite{imbens2015}.

Bayesian model comparison and Bayes factors are treated in \textcite{kass1995} and Gelman et al.\ (2013).
Model averaging is developed in Hoeting et al.\ (1999).
The philosophy of falsifiability and research programs draws on \textcite{lakatos1978} and Popper (1963).

For triangulation methodology, see \textcite{king1994} and \textcite{morgan2015}.
On the replication crisis and its implications for social science methodology, see Open Science Collaboration (2015).
The ethics of social science research are discussed in \textcite{salganik2018}.

On the problem of reflexivity in social science---the fact that studying social phenomena can change them---see \textcite{hacking1999} and Bourdieu (2004).

\section{Exercises}

\begin{exercise}
Take a published claim of the form ``think-tank funding increased public support for deregulation.''
\begin{enumerate}[nosep]
  \item Design a Rosenbaum sensitivity analysis. For what values of $\Gamma$ does the claim survive?
  \item Apply the Oster coefficient stability test. Compute $\tau^*$ for $R_{\max} = 1.3\hat{R}^2$ and $\delta = 1$.
  \item Design three placebo tests that would strengthen or weaken the claim.
\end{enumerate}
\end{exercise}

\begin{exercise}
Describe a doxonomic hypothesis that is unfalsifiable in practice.
\begin{enumerate}[nosep]
  \item Explain why it is unfalsifiable.
  \item Reformulate it as a testable hypothesis with specific observable implications.
  \item Identify the data and methods needed to test the reformulated hypothesis.
  \item Describe what results would constitute falsification.
\end{enumerate}
\end{exercise}

\begin{exercise}
Construct a triangulation scorecard (Model~\ref{mod:triangulation-score}) for the claim that economics education increases belief in selfish human nature.
Rate each of the seven evidence lines and provide an overall assessment.
Where is the evidence weakest?
What study would most improve the overall assessment?
\end{exercise}

\begin{exercise}
\label{ex:bayes-factor-exercise}
Consider two models for the correlation between media ownership concentration and political polarization:
\begin{itemize}[nosep]
  \item $M_1$: concentration causes polarization (doxonomic model).
  \item $M_2$: polarization causes audiences to self-sort, leading to concentrated viewership that \emph{looks like} media power but is actually demand-driven.
\end{itemize}
\begin{enumerate}[nosep]
  \item What data would discriminate between these models?
  \item If you had panel data on both ownership changes and polarization measures, how would you compute the Bayes factor for $M_1$ vs.\ $M_2$?
  \item What prior over model parameters would you use, and how sensitive is the Bayes factor to this choice?
\end{enumerate}
\end{exercise}

\begin{exercise}
Apply the common pitfalls catalog (Section~\ref{sec:ch20-pitfalls}) to a real-world doxonomic claim from public discourse (e.g., ``Big Pharma controls what doctors believe,'' ``the liberal media brainwashes viewers,'' ``tech companies radicalize users'').
\begin{enumerate}[nosep]
  \item Which pitfalls does the claim commit?
  \item How would you reformulate the claim to avoid those pitfalls?
  \item What evidence would you need to evaluate the reformulated claim?
\end{enumerate}
\end{exercise}

\begin{exercise}
Design a pre-registration protocol for a doxonomic study that is maximally protected against the ideological hazard (Principle~\ref{prin:ideological-hazard}).
Include:
\begin{enumerate}[nosep]
  \item specific hypotheses with predicted directions and magnitudes,
  \item primary and secondary outcomes,
  \item the analysis plan (including what constitutes a null result),
  \item an adversarial critique section: write the strongest objection a skeptic could raise, and describe how the study design addresses it.
\end{enumerate}
\end{exercise}

\begin{exercise}[Computational]
Simulate the uncertainty propagation problem (Proposition~\ref{prop:uncertainty-propagation}):
\begin{enumerate}[nosep]
  \item Create a four-link doxonomic chain where each link has a true effect of $\beta = 0.5$ and is estimated with noise $\varepsilon \sim N(0, 0.15^2)$.
  \item Compute the end-to-end effect $\hat{\beta}_{\text{total}} = \prod_{k=1}^{4} \hat{\beta}_k$ for 10,000 Monte Carlo replications.
  \item Plot the sampling distribution of $\hat{\beta}_{\text{total}}$.
  \item What fraction of replications yield an estimate within 50\% of the true total effect ($0.5^4 = 0.0625$)?
  \item How does this fraction change if you increase the per-link noise to $\sigma = 0.25$?
\end{enumerate}
\end{exercise}


\begin{interlude}
You now have both the models and the methods.
You can define a doxonomic system, measure its components, trace causal effects, analyze texts, reconstruct histories, design experiments, and assess the robustness of your conclusions.
Part~IV applies this toolkit to the belief regimes that matter most: selfishness, meritocracy, consumerism, inequality, race, austerity, work, nation, technology, and crisis.
These chapters are case studies, not proofs.
The framework illuminates each case; each case tests the framework.
Where the fit is strong, it increases our confidence.
Where it is weak, it marks the boundaries of what doxonomics can explain.
\end{interlude}

% --- Part IV: Domains of Application ---
\part{Domains of Application}
\label{part:applications}

\chapter{Human Nature as a Managed Narrative}
\label{ch:human-nature}

\epigraph{It is not that people are either good or bad by nature, but that they can be made more selfish or more cooperative by the stories they are told about themselves.}{Samuel Bowles}

\noindent
This is the signature case study of the book.
The belief that ``human nature is selfish'' is among the most consequential doxonomic products of the modern era: an anthropological regime reinforced through centuries of institutional investment, whose self-fulfilling properties make it appear to confirm itself.
Whether it is \emph{the} most consequential---more than racial hierarchy, national myth, or religious cosmology---is a comparative question this chapter does not attempt to settle.
What it does attempt is a detailed doxonomic analysis of one regime whose mechanisms are well-documented and whose effects are measurable.

The analysis proceeds in stages: first the empirical evidence on what human nature actually looks like (Section~\ref{sec:ch21-evidence}); then the institutional genealogy of the selfishness regime (Sections~\ref{sec:ch21-genealogy}--\ref{sec:ch21-media}); the self-fulfilling mechanism by which the belief generates its own confirmation (Section~\ref{sec:ch21-self-fulfilling}); the policy feedback loop (Section~\ref{sec:ch21-policy}); the formal competition model (Section~\ref{sec:ch21-competition}); counter-narratives and their institutional base (Section~\ref{sec:ch21-counter}); and cross-cultural comparison (Section~\ref{sec:ch21-cross-cultural}).

\section{The Claim and the Evidence}
\label{sec:ch21-evidence}

\subsection{The dominant anthropological regime}

The dominant anthropological regime in Western market societies holds that:
\begin{enumerate}[nosep]
  \item humans are fundamentally self-interested,
  \item cooperation requires external enforcement (contracts, laws, incentives),
  \item altruism is either disguised self-interest or evolutionarily marginal,
  \item competition is the natural state; solidarity must be artificially maintained.
\end{enumerate}

This regime is not presented as ideology; it is presented as science---as the sober conclusion of economics, evolutionary biology, and rational choice theory.
Its power lies precisely in this scientific self-presentation: to question it is to be naive, sentimental, or ignorant of ``how the world really works.''

\subsection{What the experimental evidence actually shows}

The experimental evidence does not support these claims as a complete anthropology, though it does not deny that self-interest, status competition, and opportunism are real human tendencies.

\paragraph{Public goods games.}
\textcite{fischbacher2001} showed that the majority of participants in public goods games are \emph{conditional cooperators}: they cooperate when they expect others to cooperate.
Only about 30\% behave as consistent free riders.
The remaining 20\% display mixed or erratic patterns.
This means the modal human is not selfish---the modal human is \emph{conditionally pro-social}, adjusting behavior to perceived social norms.

\paragraph{Cross-cultural evidence.}
\textcite{henrich2001} demonstrated generosity in ultimatum and dictator games across 15 small-scale societies on four continents.
People in diverse societies routinely reject pure self-interest, though the degree varies with market integration and institutional context.
No society tested exhibited the pure self-interest predicted by the homo economicus model.

\paragraph{Commons management.}
\textcite{ostrom1990} documented hundreds of communities successfully managing common-pool resources---fisheries, irrigation systems, forests, grazing lands---without privatization or state enforcement, many for centuries.
The ``tragedy of the commons'' \autocite{hardin1968} turns out to be not a law of nature but a prediction about behavior under specific institutional conditions (open access with no communication), conditions that real communities routinely avoid through institutional design.

\paragraph{Developmental evidence.}
Children as young as 14 months display helping behavior (picking up dropped objects for an adult) without training or reward.
By age 3, children enforce fairness norms on others and protest unfair distributions.
By age 5, they engage in costly punishment of cheaters.
These behaviors appear before cultural training has had much time to operate, suggesting an evolved pro-social foundation.

\paragraph{Neuroscience.}
Cooperation activates reward centers in the brain (striatum, ventromedial prefrontal cortex) even when the cooperator receives no material benefit.
Rejecting unfair offers in ultimatum games activates the anterior insula, a region associated with disgust.
The neural architecture supports cooperation as intrinsically rewarding, not merely instrumentally advantageous.

\paragraph{The synthetic picture.}
\textcite{bowles2011} synthesized this evidence: cooperation is a fundamental, evolved human capacity, coexisting with competition and conflict rather than replacing them.
Humans are \emph{conditional cooperators with a taste for punishing cheaters}---a combination that supports cooperation in groups but also generates inter-group conflict.
The best available anthropology is not ``humans are selfish,'' not ``humans are cooperative,'' but \emph{humans are flexibly social}, with the balance between cooperation and competition depending heavily on institutional design, narrative environment, and perceived norms \autocite{gintis2005}.

The doxonomic point is not that the selfishness regime is ``wrong'' in every particular---self-interest is real, free-riding occurs, enforcement helps.
The point is that the regime \emph{selectively amplifies} one real tendency while suppressing evidence of others, producing a distorted anthropology that serves specific institutional interests.

\section{The Institutional Genealogy of Homo Economicus}
\label{sec:ch21-genealogy}

\subsection{Intellectual precursors (1600--1900)}

\begin{definition}[Homo Economicus as Anthropological Regime]
\label{def:homo-economicus}
\index{homo economicus}
\emph{Homo economicus} is the anthropological regime $\alpha_{\text{HE}}$ defined by:
motivational theory = rational self-interest;
social ontology = isolated individuals;
capacity claim = collective action requires external incentives.
\end{definition}

The intellectual lineage:
\begin{itemize}[nosep]
  \item \textbf{Hobbes} (1651): ``the life of man, solitary, poor, nasty, brutish, and short''---without a sovereign, self-interest produces war.
  \item \textbf{Mandeville} (1714): \emph{The Fable of the Bees}---``private vices, public benefits.'' Self-interest, properly channeled, produces social good.
  \item \textbf{Smith} (1776): ``It is not from the benevolence of the butcher, the brewer, or the baker that we expect our dinner, but from their regard to their own interest.'' Note: Smith also wrote \emph{The Theory of Moral Sentiments} (1759), emphasizing sympathy and social connection, but this text is far less cited in economics curricula.
  \item \textbf{Bentham} (1789): utility maximization as the foundation of rational choice.
  \item \textbf{Spencer} (1864): ``survival of the fittest'' applied to human societies---Social Darwinism as a naturalization of competition.
\end{itemize}

The doxonomic observation: none of these thinkers simply ``discovered'' that humans are selfish.
Each was writing in a specific institutional context---the rise of market society, the dissolution of feudal obligations, colonial expansion---and their theories served specific functions in legitimating the emerging order.

\subsection{Formalization (1900--1960)}

The selfishness assumption was transformed from a philosophical claim into a formal axiom:
\begin{itemize}[nosep]
  \item \textbf{Pareto} (1906): ordinal utility theory---agents have well-ordered preferences, interpreted as self-regarding.
  \item \textbf{Robbins} (1932): economics as the science of allocating scarce means among competing ends---no room for altruism, solidarity, or non-instrumental motivation.
  \item \textbf{Von Neumann and Morgenstern} (1944): expected utility theory---the rational agent maximizes expected payoff.
  \item \textbf{Arrow and Debreu} (1954): general equilibrium---the economy as a system of self-interested agents reaching optimal allocation through markets.
  \item \textbf{Friedman} (1953): the ``as if'' defense---even if people are not self-interested, models that assume they are predict well, so the assumption is ``vindicated.''
\end{itemize}

The formalization was crucial doxonomically: it turned a contestable philosophical claim (``people are selfish'') into a mathematical axiom (``preferences are complete and transitive'') that appeared value-neutral.
Students who might have questioned the philosophy were instead taught the mathematics, absorbing the anthropology as an unstated background assumption.

\subsection{Institutional build-out (1960--2000)}

The second half of the 20th century saw the selfishness regime embedded in an institutional infrastructure of unprecedented scale:

\begin{enumerate}
  \item \textbf{Economics departments.}
  From roughly 200 PhD-granting economics departments in the US alone, graduating $\sim$1,200 PhDs per year, who then teach the self-interest axiom to $\sim$1.5 million undergraduate students annually.

  \item \textbf{Business schools.}
  The MBA program---which reached $\sim$200,000 graduates per year in the US by 2000---teaches principal-agent theory, incentive design, and shareholder value maximization, all premised on self-interest.
  Jensen and Meckling's (1976) ``Theory of the Firm'' paper, which models managers as self-interested agents who must be disciplined by contracts, became the most cited paper in economics.

  \item \textbf{Think tanks.}
  Free-market think tanks (Heritage, Cato, AEI, IEA, Adam Smith Institute, Atlas Network of 500+ affiliated organizations) produce a steady stream of publications, op-eds, and policy recommendations premised on the rational self-interest model.

  \item \textbf{Textbooks.}
  Mankiw's \emph{Principles of Economics}---the dominant introductory textbook from the mid-1990s onward---presents self-interest as Principle \#4 (``People respond to incentives'') and frames market outcomes as the natural result of rational choice.
  Over 4 million copies sold.
\end{enumerate}

\begin{example}
\label{ex:funding-estimate}
A rough doxonomic accounting of the selfishness-regime's annual institutional investment in the United States (circa 2010):

\begin{center}
\begin{tabular}{lrr}
\toprule
Channel & Annual expenditure (est.) & Annual reach (est.) \\
\midrule
Economics departments (teaching) & \$2--4 billion & 1.5 million students \\
Business schools (MBA + exec ed) & \$5--8 billion & 250,000 graduates \\
Free-market think tanks (top 50) & \$0.5--1 billion & Millions (media, policy) \\
Management consulting & \$2--3 billion & 100,000+ corporate clients \\
Business media (CNBC, Bloomberg, etc.) & \$1--2 billion & 50+ million viewers \\
\midrule
\textbf{Total (conservative)} & \textbf{\$10--18 billion} & \\
\bottomrule
\end{tabular}
\end{center}

These figures are rough order-of-magnitude estimates, not precise accounting.
They include only the costs of institutions that teach or assume self-interest as a behavioral foundation; they exclude indirect channels (entertainment media, social media, interpersonal transmission).
Even at the lower bound, the annual investment in reproducing the selfishness regime exceeds the entire budget of the National Science Foundation (\$8.3 billion in 2020).
\end{example}

\section{The Education Channel}
\label{sec:ch21-education}

\subsection{Economics as doxonomic delivery mechanism}

The economics classroom is the single most important delivery mechanism for the selfishness regime.
It reaches more people, with more institutional authority, over a longer period of exposure, than any think tank or media outlet.

\begin{model}[Educational Doxonomic Pipeline]
\label{mod:education-pipeline}
\index{education pipeline}
The pipeline has three stages:
\begin{enumerate}[nosep]
  \item \textbf{Undergraduate economics} ($\sim$1.5M students/year in the US): Intro courses teach self-interest as the motivational axiom. Most students take only the intro course, absorbing the axiom without encountering the behavioral qualifications taught in advanced courses.
  \item \textbf{MBA programs} ($\sim$200K graduates/year): Principal-agent theory, incentive design, and shareholder primacy deepen the self-interest assumption and apply it to organizational design.
  \item \textbf{Corporate practice}: MBA graduates design compensation systems, performance metrics, and organizational structures premised on worker self-interest, creating institutional environments that elicit the behavior the model predicts.
\end{enumerate}
Each stage amplifies the previous one: the undergraduate absorbs the axiom, the MBA applies it, the executive institutionalizes it.
\end{model}

\subsection{The Frank-Gilovich-Regan finding}

\textcite{frank1993} documented that economics students behave more self-interestedly than students in other disciplines---contributing less in public goods games, defecting more in prisoner's dilemma, and expecting others to defect more.
The finding has been replicated across multiple countries and experimental protocols.

The causal question is contested (selection vs.\ indoctrination), but longitudinal evidence suggests both mechanisms operate.
Students who enter economics are somewhat more self-interested than average (selection), but their self-interested behavior increases during the course (indoctrination).
The indoctrination effect is concentrated in the first course, suggesting that the introductory axiom has a larger impact than advanced training.

\begin{proposition}[Economics Education as Doxonomic Treatment]
\label{prop:econ-education}
\index{economics education}
If the average introductory economics course shifts cooperative behavior by $\delta = -3$ to $-5$ percentage points (the range suggested by existing studies), and $\sim$1.5 million US students take the course annually, then the annual doxonomic output is:
\[
  \Delta_{\text{annual}} = 1.5 \times 10^6 \times \delta \approx 4.5\text{--}7.5 \times 10^6 \text{ cooperation-points lost per year.}
\]
Over a 40-year working life, the cumulative effect on the population stock of cooperative disposition is substantial---even before accounting for the multiplier effects of changed behavior influencing peers.
\end{proposition}

\subsection{What textbooks teach vs.\ what the evidence says}

\begin{center}
\begin{tabular}{p{5.5cm}p{5.5cm}}
\toprule
Standard textbook claim & Current evidence \\
\midrule
People are rational self-maximizers & Most people are conditional cooperators; pure maximizers are a minority \\[4pt]
Altruism is anomalous & Pro-social behavior is developmentally early, neurologically grounded, cross-culturally robust \\[4pt]
Commons are tragic without privatization & Hundreds of self-governing commons have persisted for centuries \\[4pt]
Incentives always help & Monetary incentives can crowd out intrinsic motivation and reduce cooperation \\[4pt]
Markets produce optimal outcomes & Markets produce efficient outcomes only under restrictive conditions rarely met in practice \\
\bottomrule
\end{tabular}
\end{center}

The gap between textbook claims and current evidence is not a gap that exists because the evidence is new.
Ostrom published in 1990.
Fischbacher and colleagues published in 2001.
Henrich and colleagues published in 2001.
These findings are a generation old.
Their absence from standard introductory textbooks is itself a doxonomic phenomenon---the institutional inertia of a curriculum controlled by departments whose senior faculty were trained in the self-interest paradigm.

\section{Media and Popular Culture}
\label{sec:ch21-media}

The selfishness regime extends far beyond the classroom.
Popular culture provides continuous reinforcement through narrative channels that reach audiences who never take an economics course.

\paragraph{Financial media.}
CNBC, Bloomberg, the Financial Times, and the Wall Street Journal operate within the self-interest framework as default.
Market movements are explained by ``rational'' responses to incentives; cooperation is a ``market failure''; regulation is ``interference'' with natural market processes.
The aggregate audience of financial media in the US exceeds 50 million.

\paragraph{Entertainment.}
The ``greed is good'' speech in \emph{Wall Street} (1987) became cultural shorthand for a worldview already institutionalized.
Reality television (e.g., \emph{Survivor}, \emph{The Apprentice}) is structured around competition, betrayal, and individual victory---reinforcing the narrative that life is a zero-sum game.
Self-help literature frequently adopts evolutionary-psychology framing: ``your brain is wired for competition.''

\paragraph{Evolutionary psychology popularization.}
Books like Dawkins' \emph{The Selfish Gene} (1976)---a brilliant work of science communication whose title was widely misread as a claim about organisms rather than genes---provided scientific-sounding support for the selfishness regime.
The ``gene's-eye view'' was translated, in popular understanding, into ``selfishness is natural, cooperation is an illusion.''
Dawkins himself has protested this misreading, but the doxonomic damage was done: the phrase ``selfish gene'' became a cultural reference point for the naturalization of self-interest.

\paragraph{The quantitative asymmetry.}
Narratives about human selfishness and competition are vastly more prevalent in mainstream media than narratives about human cooperation and solidarity.
A content analysis of US newspaper opinion pages would reveal that ``incentive'' appears orders of magnitude more frequently than ``solidarity,'' and ``market'' more frequently than ``commons.''
This asymmetry is not a natural reflection of reality; it reflects the institutional infrastructure that produces media content.

\section{The Self-Fulfilling Prophecy}
\label{sec:ch21-self-fulfilling}

\subsection{The mechanism}

Consider a first-year university student taking an introductory economics course.
The professor explains that rational agents maximize their own utility; the textbook models assume self-interest; the problem sets reward strategic thinking.
The student leaves the course believing---not because she has been argued into it, but because it is the unquestioned premise of every exercise she completed---that most people are selfish.
She enters a group project expecting free-riders.
She monitors her partners for shirking.
Her vigilance makes them uncomfortable; trust erodes; cooperation falters.
``See?'' she concludes.
``People really are selfish.''
She has just experienced the most dangerous property of the selfishness regime: it generates the evidence for its own truth.

\begin{proposition}[Doxonomic Self-Fulfillment]\footnote{This result follows from the model assumptions; it is a heuristic prediction, not an empirically validated law. The qualitative mechanism (narrative shifts norm perceptions, which shifts behavior, which confirms the narrative) is well-supported; the specific dynamics are stylized.}
\label{thm:self-fulfilling}
\index{self-fulfilling prophecy}
Let agents be conditional cooperators who cooperate when $\hat{\mu}_j(\text{cooperate}) > \tau_j$, where $\hat{\mu}_j$ is $j$'s perception of the cooperation norm and $\tau_j$ is $j$'s cooperation threshold.
If a doxonomic intervention reduces $\hat{\mu}_j(\text{cooperate})$ by shifting perceived norms (by amount $\delta$), then:
\begin{enumerate}[nosep]
  \item fewer agents exceed their threshold: cooperation drops,
  \item the observed drop confirms the narrative (``see, people are selfish''),
  \item which further reduces $\hat{\mu}_j$ for those who observe behavior,
  \item producing a downward cascade until the system reaches a new (lower) equilibrium.
\end{enumerate}
Formally, the dynamics satisfy:
\[
  \mu^{(t+1)}_{\text{coop}} = G\!\left(\hat{\mu}^{(t)}_{\text{coop}} - \delta\right)
\]
where $G$ is the fraction of the population whose threshold $\tau_j$ falls below the perceived norm.
If $G'(x) > 1$ at the initial equilibrium, the doxonomic shock $\delta$ triggers a cascade to the low-cooperation equilibrium.
\end{proposition}

\subsection{Multiple equilibria}

The self-fulfilling mechanism implies that the same population, with the same underlying preferences, can sustain \emph{different equilibrium cooperation levels} depending on the narrative environment:

\begin{itemize}[nosep]
  \item \textbf{High-cooperation equilibrium}: most people believe others cooperate $\to$ most cooperate $\to$ belief confirmed.
  \item \textbf{Low-cooperation equilibrium}: most people believe others defect $\to$ most defect $\to$ belief confirmed.
\end{itemize}

Both equilibria are self-consistent.
The narrative environment---specifically, which stories about human nature are institutionally amplified---determines which equilibrium the population settles into.
The selfishness regime pushes toward the low-cooperation equilibrium; a cooperative regime would push toward the high-cooperation equilibrium.

\begin{example}
\label{ex:self-fulfilling-numerical}
Let $\tau_j \sim \text{Beta}(3, 2)$ (most people have moderate cooperation thresholds, with a mean of 0.6).
At the high-cooperation equilibrium, $\hat{\mu} = 0.75$: 86\% of agents have $\tau_j < 0.75$ and cooperate.
At the low-cooperation equilibrium, $\hat{\mu} = 0.30$: only 17\% cooperate.

A doxonomic shock of $\delta = 0.15$ (shifting perceived norms from 0.75 to 0.60) triggers a cascade: at $\hat{\mu} = 0.60$, 66\% cooperate.
But the observed 66\% further reduces $\hat{\mu}$, which reduces cooperation further, until the system settles at the low-cooperation equilibrium ($\hat{\mu} \approx 0.30$, 17\% cooperate).

The \emph{same population} has gone from 86\% to 17\% cooperation---not because preferences changed, but because the narrative environment tipped the norm perception past the critical threshold.
\end{example}

\subsection{The doxonomic trap}

\begin{definition}[Doxonomic Trap]
\label{def:doxonomic-trap}
\index{doxonomic trap}
A \emph{doxonomic trap} is a locally stable equilibrium that is Pareto-inferior to an alternative equilibrium but is sustained by the self-fulfilling properties of the dominant narrative.
Escaping the trap requires a coordinated shift in both narrative and norm perception---which is difficult precisely because the current narrative makes the alternative seem naive or impossible.
\end{definition}

The selfishness regime constitutes a doxonomic trap: the low-cooperation equilibrium is Pareto-inferior (everyone would be better off with more cooperation), but the narrative that ``people are selfish'' is self-confirming, and anyone who proposes the cooperative alternative is dismissed as ignoring ``human nature.''

\section{Affect, Precision, and the Moral Basin}
\label{sec:ch21-affect-precision}

The self-fulfilling mechanism of Section~\ref{sec:ch21-self-fulfilling} operates through perceived norms: the regime shifts what agents expect others to do, and behavior follows expectation.
A deeper layer operates through moral feeling itself.
The selfishness regime does not only tell agents that others will defect; it also teaches them which emotional responses count as \emph{moral warrant}.
Contempt for a free-rider reads as moral perception; unease at a cooperator who ``left money on the table'' reads as prudence; sympathy for a diffuse victim reads as sentimentality.
Over time the regime shapes not only propositional belief but the precision assigned to affect---which gut responses are treated as evidence, and which are treated as noise to be overridden.

This section sketches an individual-level mechanism that feeds the population dynamics already developed.
Doxonomics is not psychology (Chapter~\ref{ch:ontology-belief}), but population belief aggregates individual states, and some of the most durable regimes operate by recruiting the machinery of moral affect.
The account is stylized and active-inference-compatible; it is a proposed synthesis, not a canonical result.

\subsection{Feelings as dashboards, not verdicts}

The folk theory of moral cognition treats feeling as perception: disgust reveals wrongness, anger reveals injustice, guilt reveals culpability.
The empirical picture is more constrained.
Moral evaluation can, in some tasks, be reported faster than self-reported felt emotion, and incidental affect (disgust primed by a bad smell, anger carried over from traffic) leaks into unrelated moral judgments, though meta-analysis suggests average incidental-affect effects are small.
The useful compression is that moral feelings are \emph{dashboards}---fast compressions of harm, threat, norm violation, and action urgency---not \emph{verdicts}.

Under active inference---a family of frameworks in which the brain is modeled as continually predicting its sensory inputs and updating to minimize prediction error---emotion is an \emph{interoceptive inference}: the organism's best model of its own bodily state (heart rate, arousal, visceral tone), used to regulate action.
What governs how much that inference shapes downstream judgment is \emph{precision}---the confidence the system assigns to the bodily signal relative to other evidence channels.
When precision on interoception is high and evidence on harm, intent, or victim structure is diffuse, the system falls into a moral basin driven by feeling.
When precision is more evenly distributed, the same feeling becomes one input among several.

\subsection{A minimal state space}

Let the individual moral state be
\[
  x_t = \bigl[\,I,\ \Pi,\ H,\ T,\ N,\ V,\ C,\ R,\ S,\ U\,\bigr]
\]
with $I$ interoceptive arousal, $\Pi$ the precision placed on it, $H$ perceived harm, $T$ inferred intent, $N$ taboo/purity salience, $V$ victim salience, $C$ self-causation, $R$ repair efficacy, $S$ social confirmation, and $U$ residual ambiguity.
The state moves on a potential $\mathcal{V}(x; c)$ shaped by context $c$: threat cues, narrative framing, group signals, institutional environment.
The picture is a landscape: $\mathcal{V}$ has local minima---\emph{attractors}---around which the system settles, each surrounded by a \emph{basin} of nearby states that drain into it.
Attractors here are not physical locations but stable patterns of appraisal and disposition.
Candidate attractors correspond to action dispositions: \emph{avoid} (disgust-like exclusion), \emph{punish} (anger-like sanction), \emph{repair} (guilt-like restitution), \emph{help} (empathic concern), and \emph{suspend} (withholding judgment).
Context $c$ reshapes the landscape: a regime that deepens one basin and flattens others changes which attractor the same facts tend to produce.
The same person, facing the same facts, can land in different basins depending on which variables the environment amplifies.

\begin{definition}[Moral Distortion Index]
\label{def:moral-distortion}
\index{moral distortion index}
\index{distortion index, moral}
For an individual in moral state $x$, the \emph{moral distortion index} is
\[
  D(x) \;=\; \frac{\Pi\,I + S}{H + T + V + C + R + \varepsilon}
\]
for small $\varepsilon > 0$.
$D \ll 1$: judgment is anchored to appraisal evidence (harm, intent, victim structure, causation, repair).
$D \gg 1$: bodily affect weighted by precision, plus social reinforcement, outruns appraisal.
\end{definition}

$D$ is a bookkeeping device, not a measurement instrument.
It names the regime in which ``this feels wrong'' most readily hardens into ``this \emph{is} wrong,'' and---symmetrically---the regime in which moral feeling serves as useful input rather than mistaken revelation.

\subsection{Doxonomic amplification of \texorpdfstring{$D$}{D}}

A regime does not control the full state $x$; it shapes context $c$, and context shapes which variables the environment makes salient.
The selfishness regime raises $D$ asymmetrically through two complementary moves:
\begin{enumerate}[nosep]
  \item \textbf{Amplification.} It elevates $\Pi\,I$ and $S$ on regime-supporting stimuli---betrayal narratives about free-riders, affective framings of welfare recipients, spectacle-driven coverage of individual wrongdoing---so that anger and contempt feel like moral perception.
  \item \textbf{Diffusion.} It suppresses $H$, $V$, $C$, and $R$ on regime-threatening stimuli---institutional harms spread across populations, victims with no camera, causal chains too long to narrate, reparative policies framed as fiscally irresponsible.
\end{enumerate}
The aggregate effect is to grant moral feeling high precision where it supports the regime and low precision where it would challenge it.
This composes with Section~\ref{sec:ch21-self-fulfilling}: narrative shifts norm perception, and it also shifts which feelings are granted moral authority.

\begin{proposition}[Affective Recruitment]\footnote{This proposition follows from the model assumptions; it is a structural claim about how regimes exploit affect, not a calibrated prediction.}
\label{prop:affective-recruitment}
\index{affective recruitment}
A regime that selectively amplifies $\Pi\,I$ and $S$ on regime-supporting stimuli while diffusing $H$, $V$, $C$, and $R$ on regime-threatening stimuli raises population $D$ asymmetrically.
The resulting distribution of moral feeling tracks institutional interest rather than evidential warrant, even when no agent experiences the shift as manipulation.
\end{proposition}

The mechanism matters precisely because it is not experienced as indoctrination.
Agents feel their moral responses as their own.
The regime has shaped not what they think but what feels morally self-evident---which is a deeper lock-in than propositional belief, and correspondingly harder to dislodge.

\subsection{Hysteresis and the cost of exit}

Attractor dynamics imply hysteresis: once a moral basin deepens through repetition, social reinforcement, and enacted judgment, the evidence required to exit it exceeds the evidence that would have been required to prevent entry.
This is consistent with the lock-in properties of doxonomic systems developed in Chapter~\ref{ch:formal-models}.
The policy consequence is that counter-doxonomic work cannot rely on argument alone.
Arguments operate on appraisal variables ($H$, $T$, $V$); lock-in operates on the affect-precision channel $\Pi\,I$ and the social-coupling channel $S$.
Shifting the moral basin requires sustained experiential counter-evidence (cooperative institutions that visibly work), sustained reframing of affect (narratives that re-locate anger toward structural sources of harm and concern toward structurally produced victims), or both.

Whether populations actually exhibit the asymmetric $D$ predicted by Proposition~\ref{prop:affective-recruitment} is an open empirical problem.
Candidate approaches are discussed in Chapter~\ref{ch:doxometric-measurement} (individual-level affect measurement) and Chapter~\ref{ch:text-media} (regime-level amplification patterns in discourse).

\section{The Policy Feedback Loop}
\label{sec:ch21-policy}

The selfishness regime does not merely describe behavior; it shapes the institutional environment in which behavior occurs, creating a feedback loop:

\begin{model}[Policy Feedback]
\label{mod:policy-feedback}
\index{policy feedback}
\begin{enumerate}[nosep]
  \item The selfishness assumption enters policy design: regulators assume that agents will cheat unless monitored; managers assume workers will shirk unless incentivized; policymakers assume citizens will free-ride unless forced to contribute.
  \item Policy instruments are designed accordingly: surveillance systems, performance-based pay, privatization, means-testing.
  \item These instruments signal to citizens that they are not trusted, which reduces intrinsic motivation and cooperative behavior \autocite{frey1997}.
  \item The resulting behavior confirms the selfishness assumption, justifying further monitoring and incentivization.
\end{enumerate}
\end{model}

\textcite{frey1997} documented the ``crowding out'' effect: monetary incentives can reduce intrinsic motivation.
Offering to pay people for blood donations \emph{reduces} donations compared to asking them to donate voluntarily.
Imposing fines for late pickup at daycare centers \emph{increases} lateness (parents treat the fine as a price, replacing a social norm with a market transaction).

The policy feedback loop means that the selfishness regime is not just a belief about behavior---it is a \emph{generator} of behavior.
The regime creates the institutional conditions under which self-interest is the rational response, then points to the resulting behavior as evidence that self-interest is natural.

\section{Two Competing Regimes: A Formal Model}
\label{sec:ch21-competition}

\begin{model}[Anthropological Regime Competition]
\label{mod:regime-competition}
\index{anthropological regime!competition}
Let $\belief{S}(t)$ be the prevalence of the selfish regime and $\belief{C}(t) = 1 - \belief{S}(t)$ the cooperative alternative.
The dynamics are:
\begin{align}
  \frac{d\belief{S}}{dt} &= \belief{S}\bigl(\alpha_S \phi_S - \gamma_S - \beta_{SC}\belief{C}\bigr) \\
  \frac{d\belief{C}}{dt} &= \belief{C}\bigl(\alpha_C \phi_C - \gamma_C - \beta_{CS}\belief{S}\bigr)
\end{align}
where:
\begin{itemize}[nosep]
  \item $\phi_S, \phi_C$ are institutional funding levels ($\phi_S \gg \phi_C$),
  \item $\alpha_S, \alpha_C$ are conversion efficiencies (funding $\to$ belief change),
  \item $\gamma_S, \gamma_C$ are natural decay rates,
  \item $\beta_{SC}, \beta_{CS}$ are competitive interaction coefficients.
\end{itemize}
\end{model}

\begin{proposition}[Funding Threshold for Regime Change]
\label{prop:funding-threshold}
\index{regime change!funding threshold}
The cooperative regime displaces the selfish regime if and only if
\[
  \alpha_C \phi_C > \gamma_C + \beta_{CS} \cdot \belief{S}^*
\]
where $\belief{S}^*$ is the current selfish-regime prevalence.
Given the current dominance of the selfish regime ($\belief{S}^* \approx 0.7$), this implies:
\[
  \phi_C > \frac{\gamma_C + 0.7 \beta_{CS}}{\alpha_C}.
\]
With plausible parameter values ($\gamma_C = 0.05$, $\beta_{CS} = 0.3$, $\alpha_C = 0.001$ per million dollars), the required annual investment is:
\[
  \phi_C > \frac{0.05 + 0.21}{0.001} = \$260 \text{ million/year.}
\]
This is a fraction of the selfishness regime's estimated \$10--18 billion annual investment, suggesting that the cooperative regime is not disadvantaged by intrinsic weakness but by massive resource asymmetry.
\end{proposition}

Note: these parameter values are illustrative, not empirically calibrated.
Rigorous estimation would require the methods of Chapters~\ref{ch:doxometric-measurement}--\ref{ch:causal-inference}.
The point is not the specific number but the \emph{structure} of the analysis: regime change depends on the ratio of institutional investment to natural decay and competitive pressure, and the required investment can in principle be estimated.

\section{Counter-Narratives and Their Institutional Base}
\label{sec:ch21-counter}

The selfishness regime has not gone unchallenged.
Several counter-narratives have emerged, each with its own (much smaller) institutional base:

\paragraph{Behavioral economics.}
Beginning with Kahneman and Tversky in the 1970s and gaining institutional power through Nobel Prizes (2002, 2017), the Journal of Behavioral Economics, and ``nudge units'' in government, behavioral economics challenges the rationality assumption---though it generally preserves self-interest, replacing ``rational selfishness'' with ``biased selfishness'' \autocite{kahneman2011}.

\paragraph{Commons research.}
Ostrom's work, recognized with the 2009 Nobel Prize, demonstrated that communities can manage shared resources without privatization.
The institutional base: the Ostrom Workshop (Indiana University), the International Association for the Study of the Commons, and a network of field researchers.
Funding level: tiny compared to the economics mainstream.

\paragraph{Cooperative economics.}
The cooperative movement (Mondragon, credit unions, worker cooperatives, platform cooperatives) provides empirical counter-evidence: organizations that assume cooperative motivation can be economically successful.
The institutional base is fragmented and underfunded compared to the corporate sector.

\paragraph{The funding asymmetry.}
A rough estimate of the counter-narrative institutional investment:

\begin{center}
\begin{tabular}{lr}
\toprule
Counter-narrative institution & Annual budget (est.) \\
\midrule
Behavioral economics research centers (top 10) & \$20--40 million \\
Commons research network (global) & \$5--10 million \\
Cooperative development organizations (US) & \$10--20 million \\
Progressive economic think tanks (EPI, Roosevelt, etc.) & \$30--50 million \\
\midrule
\textbf{Total counter-narrative investment} & \textbf{\$65--120 million} \\
\bottomrule
\end{tabular}
\end{center}

The funding ratio is approximately 100:1 in favor of the selfishness regime.
This asymmetry, not the relative quality of the evidence, is the primary explanation for the selfishness regime's dominance.

\section{Cross-Cultural Comparison}
\label{sec:ch21-cross-cultural}

The selfishness regime is not universal.
Different societies operate under different anthropological regimes, with different institutional infrastructures supporting different beliefs about human nature.

\paragraph{Japan.}
Japanese management theory emphasizes group harmony (\emph{wa}), lifetime employment, and organizational loyalty---an anthropological regime closer to ``humans are group-oriented'' than ``humans are selfish.''
The institutional base: the Ministry of Economy, Trade and Industry (METI), corporate culture, educational system emphasizing group work.
Note: this regime has been under pressure from neoliberal reform since the 1990s, illustrating the global reach of the selfishness regime's institutional infrastructure.

\paragraph{Nordic countries.}
Scandinavian societies maintain high levels of social trust and cooperation, supported by: universal welfare states, strong labor unions, cooperative traditions, and educational systems that emphasize solidarity.
The doxonomic infrastructure supporting cooperation is embedded in state institutions rather than private foundations.

\paragraph{Indigenous commons regimes.}
Many indigenous societies maintain anthropological regimes centered on reciprocity, kinship obligation, and commons stewardship.
These regimes have been systematically disrupted by colonial and neoliberal institutional intervention (cf.\ Chapter~\ref{ch:race-coloniality-gender}).

The cross-cultural comparison supports the doxonomic thesis: where institutional infrastructure supports cooperation, people cooperate more; where it supports self-interest, people behave more self-interestedly.
Human nature is not fixed at ``selfish'' or ``cooperative''---it is \emph{institutionally mediated}, and the mediating institutions are themselves products of material investment.

\section{Notes and References}
\label{sec:ch21-notes}

This chapter draws heavily on \textcite{bowles2011}, \textcite{fischbacher2001}, \textcite{henrich2001}, and \textcite{ostrom1990}.
On the ``tragedy of the commons'' and its limitations, see \textcite{hardin1968} and Ostrom's (1990) devastating response.
The developmental evidence for pro-social behavior is reviewed in Tomasello (2009).
Neuroscience of cooperation is surveyed in Rilling et al.\ (2002).

On the indoctrination effect of economics education, see \textcite{frank1993} and \textcite{bauman1996}.
The institutional history of homo economicus is traced in \textcite{mirowski2009} and \textcite{fourcade2009}.
Self-fulfilling beliefs are modeled in \textcite{schelling1978}.
On the crowding-out effect of incentives, see \textcite{frey1997}.

On the alternative cooperative anthropology, see \textcite{gintis2005} and \textcite{nowak2006}.
The behavioral economics challenge is surveyed in \textcite{kahneman2011}.
For a popular account of the ``selfish gene'' misreading, see Midgley (1979).

On cross-cultural variation in economic behavior, see the comprehensive survey in \textcite{henrich2001} and Henrich et al.\ (2010).
The neoliberal transformation of Japanese economic policy is discussed in Vogel (2006).
Nordic social democracy is analyzed in Esping-Andersen (1990).

\section{Exercises}

\begin{exercise}
Using Proposition~\ref{prop:funding-threshold}, estimate the annual counter-doxonomic investment needed to shift the U.S.\ anthropological regime from selfish to cooperative.
State your parameter assumptions explicitly.
Conduct a sensitivity analysis: how does the required investment change if $\alpha_C$ is twice as large (the cooperative message is more persuasive per dollar)?
\end{exercise}

\begin{exercise}
Design an experiment to test the self-fulfilling prophecy of selfishness narratives.
\begin{enumerate}[nosep]
  \item Randomly expose groups to different anthropological narratives.
  \item Measure cooperation in subsequent public goods games.
  \item Measure perceived norms (what do subjects think others will do?).
  \item Test whether the effect operates through the norm perception channel (mediation analysis).
\end{enumerate}
Specify sample size, treatment arms, and pre-registration plan.
\end{exercise}

\begin{exercise}
Construct a doxonomic account (Chapter~\ref{ch:belief-flow-accounting}) for the belief ``human nature is selfish,'' identifying at least five funding channels and estimating their annual expenditure.
Use the accounting framework to identify which channels have the highest return on doxonomic investment.
\end{exercise}

\begin{exercise}
Argue that the self-fulfilling property of selfishness beliefs creates a doxonomic trap (Definition~\ref{def:doxonomic-trap}).
\begin{enumerate}[nosep]
  \item Show formally that both the high-cooperation and low-cooperation equilibria are stable.
  \item Compute the basin of attraction for each equilibrium using the Beta(3,2) threshold distribution.
  \item Calculate the minimum doxonomic shock $\delta^*$ needed to tip the system from the low-cooperation to the high-cooperation equilibrium.
  \item Discuss what real-world intervention could produce a shock of this magnitude.
\end{enumerate}
\end{exercise}

\begin{exercise}
Compare the doxonomic investment in ``humans are selfish'' with the investment in ``humans are cooperative.''
\begin{enumerate}[nosep]
  \item Estimate the funding ratio using the institutional accounting approach.
  \item Given this ratio, what does the epistemic status of the dominant regime's dominance tell us---that it reflects the evidence, or that it reflects the funding?
  \item Design a study that could distinguish between these two explanations.
\end{enumerate}
\end{exercise}

\begin{exercise}
Conduct a content analysis of the first three chapters of two introductory economics textbooks.
\begin{enumerate}[nosep]
  \item Count explicit and implicit invocations of self-interest as a motivational assumption.
  \item Count mentions of cooperation, altruism, or pro-social behavior.
  \item Compare with the actual experimental evidence summarized in Section~\ref{sec:ch21-evidence}.
  \item What is the ``textbook-evidence gap,'' and how would you explain it doxonomically?
\end{enumerate}
\end{exercise}

\begin{exercise}
Select a non-Western society with a different anthropological regime (e.g., Japan, Bhutan, an indigenous commons community).
\begin{enumerate}[nosep]
  \item Describe the dominant anthropological assumptions.
  \item Map the institutional infrastructure that supports these assumptions (education, media, policy, religion).
  \item Assess whether neoliberal doxonomic intervention has shifted the regime, and if so, through which channels.
  \item What would the doxonomic model predict about the future trajectory?
\end{enumerate}
\end{exercise}

\begin{exercise}[Computational]
Simulate the self-fulfilling prophecy model from Theorem~\ref{thm:self-fulfilling}:
\begin{enumerate}[nosep]
  \item Generate 1,000 agents with thresholds $\tau_j \sim \text{Beta}(3, 2)$.
  \item Start at the high-cooperation equilibrium ($\hat{\mu} = 0.75$).
  \item Apply a doxonomic shock $\delta = 0.15$ and iterate the norm-updating dynamics for 50 periods.
  \item Plot the cooperation rate over time. Does the system tip to the low-cooperation equilibrium?
  \item Repeat for $\delta = 0.05, 0.10, 0.15, 0.20, 0.25$. Find the critical $\delta^*$ at which tipping occurs.
  \item Now simulate a ``counter-doxonomic'' intervention: starting from the low-cooperation equilibrium, apply a positive shock. What magnitude is required to restore high cooperation?
\end{enumerate}
\end{exercise}

\chapter{Meritocracy}
\label{ch:meritocracy}

\epigraph{The more a society proclaims its meritocracy, the more its winners can justify their position and the more its losers blame themselves.}{Michael Sandel, \textit{The Tyranny of Merit}}

\noindent
Meritocracy is the belief that outcomes reflect effort and talent.
It is also one of the most heavily subsidized narratives in modern capitalism, functioning as the legitimation mechanism for inequality.

\section{The Meritocratic Belief System}
\label{sec:ch22-system}

\begin{definition}[Meritocratic Regime]
\label{def:meritocratic-regime}
\index{meritocracy}
The \emph{meritocratic regime} is an ideology $\mathcal{P}_M$ containing:
\begin{enumerate}[nosep]
  \item success reflects ability and effort ($p_1$),
  \item opportunity is broadly equal ($p_2$),
  \item failure reflects personal deficiency ($p_3$),
  \item inequality is therefore justified ($p_4$).
\end{enumerate}
These propositions are mutually reinforcing: high credence in $p_1$ raises credence in $p_3$ and $p_4$.
\end{definition}

\section{Institutional Repetition}
\label{sec:ch22-repetition}

The meritocratic regime is reinforced through:
\begin{itemize}[nosep]
  \item educational systems (grading, ranking, credentialing),
  \item corporate hiring and promotion (``talent management''),
  \item media (rags-to-riches narratives, billionaire hagiography),
  \item policy (means-testing, workfare, individual responsibility discourse).
\end{itemize}

Each institution treats meritocratic assumptions as operational truth, creating the conditions under which the assumptions appear self-evident.

\section{Credential Systems as Epistemic Capital Laundering}
\label{sec:ch22-credentials}

\begin{proposition}[Credential Laundering]
\label{prop:credential}
\index{credential laundering}
Credential systems (degrees, certifications, rankings) convert inherited advantage into apparent merit.
If access to credentials correlates with parental wealth ($r > 0$), then the credential $c_j$ for agent $j$ satisfies
\[
  c_j = f(\text{ability}_j, \text{effort}_j) + g(\text{wealth}_j) + \varepsilon_j
\]
The meritocratic narrative claims $g \equiv 0$.
The doxonomic analysis estimates $g$ and reveals the laundering function.
\end{proposition}

\section{A Model of Meritocratic Belief}
\label{sec:ch22-model}

\begin{model}[Meritocratic Belief Production]
\label{mod:meritocracy}
The strength of meritocratic belief $\belief{M}$ is a function of inequality visibility $V$ and narrative investment $\phi_M$:
\[
  \belief{M} = \frac{\phi_M}{\phi_M + \kappa \cdot V}
\]
where $\kappa > 0$ measures the sensitivity of belief to visible inequality.
High investment $\phi_M$ can sustain meritocratic belief even at high inequality, but the required investment grows with $V$.
\end{model}

\section{The Doxonomic Genealogy of Meritocracy}
\label{sec:ch22-genealogy}

The term ``meritocracy'' was coined satirically by Michael Young in \textit{The Rise of the Meritocracy} (1958) as a \emph{warning}, not a recommendation.
Its transformation into a term of praise is itself a doxonomic achievement.

Key institutional channels:
\begin{itemize}[nosep]
  \item \textbf{Standardized testing} (SAT, GRE, IQ tests): created the appearance of objective merit measurement while correlating heavily with parental income and educational investment.
  \item \textbf{Business schools}: propagated the idea that managerial success reflects personal talent, justified executive compensation, and trained a professional class that internalized meritocratic assumptions.
  \item \textbf{Self-help and motivational industry}: annual revenues exceeding \$10 billion in the U.S., reinforcing the belief that success is a matter of mindset and effort.
  \item \textbf{Rags-to-riches media narratives}: disproportionate coverage of exceptional upward mobility, creating availability bias about its frequency.
  \item \textbf{Policy design}: means-testing, workfare, and conditional welfare systems that operationalize the assumption that poverty reflects insufficient effort.
\end{itemize}

\section{Negative Cases: When Meritocratic Belief Fails}
\label{sec:ch22-negative}

Not all investments in meritocratic narrative succeed.
Doxonomic analysis requires attention to failure:
\begin{itemize}[nosep]
  \item In highly visible economic crises (2008, COVID-19), meritocratic belief temporarily declines as structural causes become salient.
  \item In societies with strong labor movements or social-democratic traditions, meritocratic belief is weaker despite similar media exposure.
  \item When inequality becomes extreme enough that the gap between narrative and lived experience is too large, the legitimation function (Definition~\ref{def:legitimation}) breaks down.
\end{itemize}

These negative cases constrain the model: doxonomic investment is necessary but not sufficient for common-sense capture.
Lived experience, counter-institutions, and competing narratives limit the reach of any regime.

\section{Racial and Gender Dimensions}
\label{sec:ch22-racial}

Meritocratic belief is not distributed uniformly.
It intersects with racial and gender regimes (Chapter~\ref{ch:race-coloniality-gender}) in specific ways:
\begin{itemize}[nosep]
  \item White respondents in the U.S.\ consistently score higher on meritocratic belief scales than Black respondents---the people for whom the narrative is most evidently contradicted by structural barriers are least persuaded.
  \item Meritocratic narratives often implicitly assume a subject who is male, able-bodied, and free of caregiving responsibilities.
  \item The ``model minority'' narrative applies meritocratic logic selectively to Asian Americans, using their achievements to delegitimize structural claims by other racialized groups.
\end{itemize}

These patterns confirm that doxonomic capture is partial and contested, not total.

\section{Notes and References}
\label{sec:ch22-notes}

On meritocracy as ideology, see \textcite{sandel2012} and \textcite{bourdieu1984}.
Inequality and its legitimation are analyzed in \textcite{piketty2014} and \textcite{stiglitz2012}.
The role of education in reproducing class structure is central to \textcite{bourdieu1984} and \textcite{bowles1976}.
On credentialism, see Collins (1979).
On the racial dimensions of meritocratic belief, see \textcite{hochschild2016}.
The self-help industry as a doxonomic channel connects to \textcite{ehrenreich2009}.

\section{Exercises}

\begin{exercise}
Using Model~\ref{mod:meritocracy}, compute the narrative investment $\phi_M$ required to maintain $\belief{M} = 0.7$ when inequality doubles.
\end{exercise}

\begin{exercise}
Design a doxometric survey to measure meritocratic belief across socioeconomic groups.
Predict which groups will score highest and explain why doxonomically.
\end{exercise}

\begin{exercise}
Analyze the genre of ``billionaire biography'' as a doxonomic product.
Identify the belief production chain from publisher to reader.
\end{exercise}

\chapter{Consumerism and Manufactured Desire}
\label{ch:consumerism}

\epigraph{People recognize themselves in their commodities; they find their soul in their automobile, hi-fi set, split-level home, kitchen equipment.}{Herbert Marcuse, \textit{One-Dimensional Man}}

\noindent
Advertising is the largest doxonomic industry.
Global advertising expenditure surpassed one trillion dollars in 2024, making it the single greatest investment in belief production in human history.
This chapter analyzes consumerism as a doxonomic system.

\section{Advertising as Doxonomic Investment}
\label{sec:ch23-advertising}

\begin{definition}[Advertising as Doxonomic Production]
\label{def:advertising-doxonomic}
\index{advertising}
Advertising is a doxonomic investment whose primary output is not product awareness but \emph{desire infrastructure}: the narrative environment in which consumption appears as the natural expression of identity, freedom, and self-worth.
\end{definition}

The individual advertisement promotes a product.
The advertising \emph{system} promotes a worldview: you are what you buy, happiness is purchasable, identity is consumable \autocite{debord1967, adorno1944}.

\section{The Desire Production Function}
\label{sec:ch23-desire}

\begin{model}[Manufactured Desire]
\label{mod:desire}
\index{manufactured desire}
Let $D_j$ be the desire intensity of agent $j$ for goods beyond subsistence.
Desire is produced by:
\[
  D_j(t) = D_0 + \alpha \sum_k \funding{k}^{\text{ad}} \cdot e_{jk} - \beta \cdot S_j(t)
\]
where $D_0$ is baseline desire, $\funding{k}^{\text{ad}}$ is advertising through channel $k$, $e_{jk}$ is exposure, and $S_j(t)$ is satisfaction with current possessions (which advertising systematically reduces).
The term $-\beta S_j$ captures the mechanism of \emph{induced dissatisfaction} \autocite{schor1998, galbraith1958}.
\end{model}

\section{Consumption as Moral Language}
\label{sec:ch23-moral}

Consumerism is not merely an economic pattern but an anthropological regime: the belief that the good life consists in acquisition, that social worth is legible through possessions, and that consumer choice is the primary form of agency \autocite{fromm1976}.

\section{Notes and References}
\label{sec:ch23-notes}

The critical theory of consumer society draws on \textcite{adorno1944}, \textcite{debord1967}, and \textcite{fromm1976}.
On advertising and identity, see \textcite{veblen1899} and \textcite{galbraith1958}.
Induced dissatisfaction is analyzed in \textcite{schor1998}.
For the economics of manipulation through advertising, see \textcite{akerlof2015}.

\section{Exercises}

\begin{exercise}
Estimate the advertising expenditure per capita in your country and compare it to spending on public education.
Interpret the ratio doxonomically.
\end{exercise}

\begin{exercise}
Using the desire production model, show that if $\beta > 0$, advertising simultaneously increases desire \emph{and} decreases satisfaction, creating a treadmill effect.
\end{exercise}

\begin{exercise}
Analyze a specific advertising campaign as a belief production chain.
Identify funder, institution, channel, narrative, and behavioral consequence.
\end{exercise}

\chapter{Inequality and Legitimacy}
\label{ch:inequality-legitimacy}

\epigraph{The comfort of the rich depends upon an abundant supply of the poor.}{Voltaire}

\noindent
Every unequal society faces a legitimation problem: why do the many accept the privilege of the few?
The question is ancient (Hume noted that the rulers are always outnumbered by the ruled, so government must rest ultimately on opinion), but doxonomics gives it new precision.
The previous chapters analyzed specific regimes (selfishness, meritocracy, consumerism) that contribute to legitimation.
This chapter integrates them into a unified framework: the doxonomic theory of inequality legitimation.

The central argument: inequality persists not (primarily) because of coercion but because of narrative.
The narrative infrastructure that legitimates inequality (the network of institutions, media outlets, educational practices, and cultural forms that make the existing distribution of rewards seem natural, inevitable, or deserved) is itself a product of material investment.
The wealthy fund the narratives that legitimate their wealth: a feedback loop that doxonomics is uniquely equipped to analyze.

\section{The Legitimation Problem}
\label{sec:ch24-problem}

\subsection{Why legitimation is necessary}

In any society where the top 1\% holds 30--40\% of the wealth (as in the contemporary US and UK), the distribution could in principle be challenged by majoritarian political action.
That it is not (that democratic majorities consistently fail to enact redistributive policies commensurate with the level of inequality) requires explanation.

Three explanations are commonly offered:
\begin{enumerate}
  \item \textbf{Coercion}: the wealthy use force (police, military, legal system) to suppress redistributive demands.
  This is occasionally true but insufficient as a general explanation in democracies with contested elections and free speech.

  \item \textbf{Institutional capture}: the political system is structured to block redistribution regardless of majority preferences (gerrymandering, lobbying, campaign finance, filibuster).
  This is substantially true and well-documented, but it is only part of the story.

  \item \textbf{Narrative legitimation}: the majority \emph{accepts} the existing distribution as fair, natural, or inevitable, or at least does not find it sufficiently objectionable to prioritize redistribution.
  This is the doxonomic explanation.
\end{enumerate}

All three mechanisms operate simultaneously.
The doxonomic contribution is not to deny coercion or institutional capture but to analyze the third mechanism, the one that makes the first two unnecessary for much of the time.

\subsection{The legitimation function}

\begin{definition}[Legitimation Function]
\label{def:legitimation}
\index{legitimation}
The \emph{legitimation function} $L: \R_{\geq 0} \times \R_{\geq 0} \to [0,1]$ maps the level of inequality $I$ and doxonomic investment $\phi_L$ to the population's acceptance level:
\[
  L(I, \phi_L) = \frac{\phi_L}{\phi_L + \kappa(I - I_0)^+}
\]
where $I_0$ is a baseline inequality level below which legitimation is ``free'' (inequality is accepted without narrative investment) and $\kappa > 0$ is the sensitivity parameter measuring how effectively visible inequality erodes acceptance.
\end{definition}

Properties of the legitimation function:
\begin{itemize}[nosep]
  \item $L = 1$ when $I \leq I_0$: below the baseline, inequality requires no justification.
  \item $L$ decreases as $I$ increases for fixed $\phi_L$: rising inequality erodes legitimacy.
  \item $L$ increases as $\phi_L$ increases for fixed $I$: more narrative investment increases acceptance.
  \item $\lim_{\phi_L \to \infty} L = 1$: with unlimited narrative investment, any inequality can be legitimated.
  \item $\lim_{I \to \infty} L = 0$ for fixed $\phi_L$: at extreme inequality, finite investment cannot maintain legitimacy.
\end{itemize}

\subsection{The escalating cost of legitimation}

\begin{proposition}[Escalating Cost of Legitimation]
\label{prop:escalation}
\index{legitimation cost}
For a fixed acceptance level $L^*$, the required narrative investment is:
\[
  \phi_L^* = \frac{L^*}{1 - L^*} \cdot \kappa(I - I_0)
\]
which grows linearly in inequality.
Highly unequal societies must spend proportionally more on narrative maintenance.
\end{proposition}

\begin{example}
\label{ex:legitimation-cost-us}
Applying the model to the United States, 1970--2020:

\begin{center}
\begin{tabular}{lcccc}
\toprule
& 1970 & 1990 & 2010 & 2020 \\
\midrule
Top 1\% wealth share ($I$) & 22\% & 30\% & 35\% & 38\% \\
Think-tank spending (est., \$B) & 0.1 & 0.5 & 0.8 & 1.0 \\
Meritocratic belief (\%) & 70 & 75 & 72 & 68 \\
\bottomrule
\end{tabular}
\end{center}

Inequality rose by 73\% (from 22\% to 38\% of wealth held by the top 1\%) while think-tank spending rose tenfold.
Meritocratic belief initially rose despite growing inequality (the model predicts this: investment can outpace erosion), but began declining after 2010 as inequality outstripped the narrative infrastructure's capacity to legitimate it.
The 2008 crisis and post-2016 populism may mark the onset of legitimation failure, a point we return to below.
(Illustrative data; actual estimation requires the methods of Chapters~\ref{ch:doxometric-measurement}--\ref{ch:causal-inference}.)
\end{example}

\section{The Legitimation Repertoire}
\label{sec:ch24-repertoire}

The narrative infrastructure of inequality legitimation draws on several interlocking belief systems.
Each provides a distinct justification for the existing distribution.

\subsection{The wealth-virtue association}

The most efficient legitimating narrative is the association of wealth with virtue: the belief that the rich \emph{deserve} their wealth because they are more talented, hardworking, or morally superior.
This is the intersection of the meritocratic regime (Chapter~\ref{ch:meritocracy}) and the selfishness regime (Chapter~\ref{ch:human-nature}): competition is natural, markets are fair, winners earned their position, losers deserve their fate.

\begin{model}[Wealth-Virtue Feedback]
\label{mod:wealth-virtue}
\index{wealth-virtue association}
The wealth-virtue association creates a feedback loop:
\begin{enumerate}[nosep]
  \item Wealthy individuals and corporations fund institutions that promote the wealth-virtue association (think tanks, business schools, media).
  \item These institutions produce narratives that frame wealth as deserved.
  \item The public internalizes the association, reducing political demand for redistribution.
  \item Reduced redistribution preserves the wealth that funds the narrative infrastructure.
\end{enumerate}
The loop is self-sustaining as long as the narrative investment is sufficient to maintain legitimation above the threshold $L^*$.
\end{model}

\subsection{Trickle-down narratives}

The ``trickle-down'' narrative holds that inequality benefits everyone because wealth at the top generates investment, innovation, and employment that eventually reach the bottom.
The narrative has been a staple of supply-side economics since the 1980s and is actively promoted by business organizations and free-market think tanks.

The empirical evidence is weak: \textcite{piketty2014} documents that periods of high inequality are not associated with faster growth; several studies find the opposite relationship.
But the narrative persists because it serves a specific doxonomic function: it converts a moral problem (is inequality just?) into a technical claim (does inequality promote growth?) that can be debated endlessly without ever reaching the normative question.

\subsection{Aspirational narratives}

Rather than justifying the existing distribution, aspirational narratives redirect attention to future possibility: ``You too could be rich.''
Lottery advertising, rags-to-riches media, and ``American Dream'' rhetoric serve this function.
The doxonomic operation: convert a statistical reality (low intergenerational mobility) into a subjective experience (possibility of upward mobility), making the current distribution feel temporary rather than structural.

\begin{proposition}[Prospect of Upward Mobility]
\label{prop:poum}
\index{prospect of upward mobility}
The ``prospect of upward mobility'' (POUM) hypothesis holds that tolerance of inequality reflects the belief that one's own position may improve.
Formally, agent $j$ tolerates inequality if:
\[
  \E_j[U(Y_j^{\text{future}})] > U(Y_j^{\text{current}} + r_j)
\]
where $r_j$ is the redistribution $j$ would receive.
If $j$ overestimates their probability of upward mobility (because aspirational narratives are more prevalent than accurate mobility data), they rationally oppose redistribution even when it would benefit them in expectation.
This is the ``temporarily embarrassed millionaire'' effect.
\end{proposition}

\subsection{Blame narratives}

Blame narratives locate the cause of poverty in the poor themselves: lack of effort, bad decisions, moral failings, cultural deficiency.
These narratives operationalize the meritocratic regime's third proposition ($p_3$: failure reflects personal deficiency) and direct attention away from structural causes (deindustrialization, discrimination, healthcare costs, housing markets).

\begin{definition}[Narrative Shield]
\label{def:narrative-shield}
\index{narrative shield}
A \emph{narrative shield} is a set of beliefs that deflects attention from structural causes of inequality to individual explanations.
Formally, a narrative shield is a mapping:
\[
  \sigma: \{\text{structural explanations for outcome } O\} \to \{\text{individual explanations for outcome } O\}
\]
that replaces ``$O$ is caused by institutional structures $S$'' with ``$O$ is caused by individual characteristics $C$.''
The shield reduces the perceived need for collective action by framing outcomes as matters of personal responsibility.
\end{definition}

Common narrative shields:
\begin{itemize}[nosep]
  \item ``The poor are lazy'' (shields from: labor market structure, wage stagnation, precarious employment).
  \item ``If you wanted health insurance, you should have gotten a better job'' (shields from: employer-based insurance system, pre-existing condition exclusions).
  \item ``College is available to anyone who works hard enough'' (shields from: tuition costs, school quality segregation, information asymmetries).
  \item ``If immigrants can succeed, why can't you?'' (shields from: selection effects in immigration, survivorship bias in success stories).
\end{itemize}

\section{The Political Economy of Legitimation}
\label{sec:ch24-political-economy}

\subsection{Who funds legitimation?}

The legitimation infrastructure is not funded by ``the economy'' in the abstract; it is funded by identifiable actors with identifiable interests.

\begin{example}
\label{ex:legitimation-funders}
The institutional infrastructure of inequality legitimation in the US includes:
\begin{enumerate}[nosep]
  \item \textbf{Corporate-funded think tanks}: Heritage Foundation (\$85M/year), Cato Institute (\$40M), AEI (\$60M), Manhattan Institute (\$25M), and hundreds of smaller organizations in the Atlas Network.
  \item \textbf{Family foundations of the ultra-wealthy}: Koch network (\$400M+/election cycle), Scaife foundations, Bradley Foundation, Olin Foundation (now defunct but highly influential in building the intellectual infrastructure).
  \item \textbf{Industry associations}: US Chamber of Commerce (\$200M+/year in lobbying and communications), National Association of Manufacturers, Business Roundtable.
  \item \textbf{Corporate PR and communications departments}: Fortune 500 companies collectively spend billions on public affairs, issue advertising, and narrative management.
\end{enumerate}
Total estimated annual spending on inequality-legitimating narrative infrastructure: \$1--3 billion in direct think-tank and advocacy spending, plus an unknown (but much larger) amount in corporate communications and media.
\end{example}

\subsection{Demand-side factors}

Legitimation succeeds not only because it is well-funded but because audiences have psychological incentives to accept it.
Just-world belief, system justification theory \autocite{jost2004}, and cognitive dissonance reduction all predict that people will be motivated to accept narratives that justify the existing order, especially if they perceive themselves as benefiting from it (or hope to benefit in the future).

\begin{proposition}[System Justification Premium]
\label{prop:system-justification}
\index{system justification}
System justification theory predicts that, holding narrative investment constant, legitimation is higher among:
\begin{enumerate}[nosep]
  \item those who perceive themselves as benefiting from the current system,
  \item those with higher need for cognitive closure (preference for certainty),
  \item those with lower socioeconomic status (paradoxically, because accepting the system reduces the psychological cost of their position: it is less painful to believe one's poverty is deserved than to believe one is a victim of injustice).
\end{enumerate}
The third prediction (that the disadvantaged sometimes embrace legitimating narratives more than the advantaged) is empirically documented and represents a distinctive contribution of system justification theory to doxonomic analysis \autocite{jost2004}.
\end{proposition}

\section{Legitimation Failure and Crisis}
\label{sec:ch24-failure}

\subsection{When legitimation breaks down}

The model predicts that legitimation fails when inequality grows faster than narrative investment can compensate.
Observable signatures of legitimation failure:

\begin{enumerate}[nosep]
  \item \textbf{Declining trust in institutions}: if the institutions that produce legitimating narratives lose credibility, their output loses effectiveness.
  \item \textbf{Rise of populism}: populist movements (left and right) are characterized by explicit rejection of elite legitimacy: ``the system is rigged.''
  \item \textbf{Increased protest activity}: strikes, demonstrations, and social movements increase when the population no longer accepts the distribution as legitimate.
  \item \textbf{Narrative fragmentation}: the unified legitimating narrative fragments into competing frames, reducing the coherence that sustains common-sense capture.
  \item \textbf{Elite anxiety}: increased spending on private security, gated communities, and ``prepper'' culture among the wealthy signals awareness that legitimation is failing.
\end{enumerate}

\subsection{The 2008--present legitimation crisis}

\begin{example}
\label{ex:legitimation-crisis}
The 2008 financial crisis produced the most visible legitimation failure in Western democracies since the 1930s.

\paragraph{The triggering event.}
Banks that had been celebrated as innovative were revealed to be reckless.
Executives who had been compensated as geniuses were exposed as incompetent or fraudulent.
The meritocratic narrative, that financial success reflects talent, was visibly falsified.

\paragraph{The legitimation response.}
The policy response (bailouts without prosecution, austerity for the public, continued bonuses for executives) further eroded legitimacy by making the double standard visible.
Think tanks and business groups invested heavily in counter-narratives: ``the crisis was caused by government regulation'' (deflecting blame from the financial sector), ``we need growth, not redistribution'' (the trickle-down narrative), ``too much regulation will prevent recovery'' (the deregulation narrative).

\paragraph{Partial legitimation recovery.}
By 2015, meritocratic belief had partially recovered in survey data: the narrative infrastructure proved resilient.
But new cracks had appeared: the Occupy movement (2011), the Sanders and Trump campaigns (2016), and the ``rigged system'' discourse all reflected persistent legitimation failure among significant segments of the population.

\paragraph{The COVID-19 shock.}
The pandemic produced a second legitimation shock: ``essential workers'' earning minimum wage while billionaire wealth increased by trillions made the wealth-virtue association difficult to sustain.
Whether this produces lasting legitimation failure or another partial recovery depends on the relative investment in legitimating vs.\ delegitimating narratives in the coming years.
\end{example}

\subsection{A model of legitimation dynamics}

\begin{model}[Legitimation Dynamics]
\label{mod:legitimation-dynamics}
\index{legitimation dynamics}
\begin{align}
  \frac{dL}{dt} &= \alpha \phi_L (1 - L) - \beta (I - I_0)^+ L - \delta_{\text{crisis}} \cdot \mathbf{1}[\text{crisis at } t] \label{eq:legit-dynamics} \\
  \frac{dI}{dt} &= \gamma_I \cdot r(L) - \mu \label{eq:ineq-dynamics}
\end{align}
where:
\begin{itemize}[nosep]
  \item $\alpha \phi_L (1-L)$: narrative investment increases legitimation (with diminishing returns),
  \item $\beta(I - I_0)^+ L$: excess inequality erodes legitimation,
  \item $\delta_{\text{crisis}}$: a crisis event produces a discrete legitimation shock,
  \item $r(L)$: the rate of return to capital increases with legitimation (high $L$ blocks redistribution and regulation),
  \item $\mu$: a natural compression force (growth, diffusion, demographic change).
\end{itemize}
The system can exhibit:
\begin{enumerate}[nosep]
  \item \emph{Stable legitimation}: $L$ and $I$ reach a self-sustaining equilibrium.
  \item \emph{Vicious spiral}: rising $I$ erodes $L$, but reduced $L$ fails to produce redistribution because institutional capture blocks the political channel, so $I$ continues to rise, further eroding $L$.
  \item \emph{Crisis-triggered regime change}: a legitimation shock ($\delta_{\text{crisis}}$ large enough) pushes $L$ below a threshold, triggering redistributive political action that reduces $I$.
\end{enumerate}
\end{model}

\section{Inequality, Race, and Compound Legitimation}
\label{sec:ch24-race}

Economic inequality intersects with racial and gender hierarchies (Chapter~\ref{ch:race-coloniality-gender}), creating \emph{compound legitimation}: each hierarchy provides narrative resources for legitimating the others.

\begin{proposition}[Compound Legitimation]
\label{prop:compound}
\index{compound legitimation}
When economic inequality correlates with racial hierarchy, each provides narrative cover for the other:
\begin{itemize}[nosep]
  \item Economic inequality is legitimated by racial narratives: ``Those people are poor because of their culture/behavior.''
  \item Racial hierarchy is legitimated by economic narratives: ``If racism were real, how do you explain successful minorities?''
\end{itemize}
The compound effect means that neither inequality nor racial hierarchy can be fully delegitimated without also challenging the other.
Intersectional doxonomic analysis is therefore not a supplement to economic analysis but a requirement for it.
\end{proposition}

Historically, racial division has been a powerful tool for blocking class-based redistribution.
\textcite{du-bois2007} identified the ``psychological wage of whiteness'': the psychic benefit that poor whites derive from racial superiority, which compensates for economic deprivation and reduces demand for redistribution.
This is a doxonomic mechanism: the racial narrative substitutes for material benefit, reducing the political threat that economic inequality would otherwise face.

\section{Comparative Legitimation Regimes}
\label{sec:ch24-comparative}

Different societies legitimate their inequality structures through different narrative repertoires:

\begin{center}
\begin{tabular}{p{2.5cm}p{3cm}p{3.5cm}p{3cm}}
\toprule
Society & Inequality level & Primary legitimation narrative & Counter-narrative strength \\
\midrule
United States & Very high (Gini $\approx 0.39$) & Meritocracy, aspiration & Moderate (populist, progressive) \\[4pt]
United Kingdom & High (Gini $\approx 0.35$) & Tradition, meritocracy & Moderate (labor movement) \\[4pt]
Germany & Moderate (Gini $\approx 0.31$) & Social market economy & Strong (unions, SPD tradition) \\[4pt]
Sweden & Low (Gini $\approx 0.27$) & Social solidarity & Very strong (universal welfare) \\[4pt]
South Africa & Very high (Gini $\approx 0.63$) & Post-apartheid aspiration & Strong but contested \\
\bottomrule
\end{tabular}
\end{center}

The comparison supports the model: higher inequality requires stronger legitimating narratives, and the specific narrative used depends on the available institutional infrastructure.
The US and Sweden are polar cases: the US has high inequality and a strong meritocratic infrastructure; Sweden has low inequality and a strong solidarity infrastructure.
The causal direction runs both ways: institutions produce narratives that sustain inequality levels, and inequality levels demand institutions that produce legitimating narratives.

\section{Notes and References}
\label{sec:ch24-notes}

On inequality and its justification, see \textcite{piketty2014} and \textcite{stiglitz2012}.
The legitimation function draws on \textcite{lukes2005} on the third dimension of power and Gramsci's theory of hegemony.
System justification theory is developed in \textcite{jost2004}.
The wealth-virtue association is analyzed in \textcite{sandel2012}.

On the prospect of upward mobility, see Benabou and Ok (2001).
On the racial dimensions of inequality legitimation, see \textcite{du-bois2007} on the wages of whiteness and \textcite{hochschild2016} on the American Dream's racial contours.
For the political economy of inequality, see \textcite{acemoglu2012} and Hacker and Pierson (2010).

The 2008 legitimation crisis is analyzed in Streeck (2016).
On populism as a symptom of legitimation failure, see M\"{u}ller (2016).
For the Koch network's institutional infrastructure, see \textcite{mayer2016}.

\section{Exercises}

\begin{exercise}
Using the legitimation function (Definition~\ref{def:legitimation}), estimate $\kappa$ for the United States by comparing inequality growth and think-tank spending over 1980--2020.
\begin{enumerate}[nosep]
  \item Gather data on the top 1\% wealth share (as a proxy for $I$) and total think-tank spending (as a proxy for $\phi_L$).
  \item Estimate $I_0$ (the inequality level at which legitimation becomes necessary) from the historical data.
  \item Calibrate $\kappa$ so that the model reproduces observed meritocratic belief levels.
  \item Does the model predict the post-2008 decline in meritocratic belief?
\end{enumerate}
\end{exercise}

\begin{exercise}
Identify three narrative shields (Definition~\ref{def:narrative-shield}) currently active in public discourse about wealth inequality.
For each:
\begin{enumerate}[nosep]
  \item State the structural explanation being shielded.
  \item State the individual explanation being substituted.
  \item Trace the belief production chain: who produces this shield, through what institutions, with what funding?
  \item How would you empirically test whether the shield is effective (i.e., whether exposure to the shield narrative reduces support for structural explanations)?
\end{enumerate}
\end{exercise}

\begin{exercise}
Model a society where legitimation fails ($L < L^*$).
\begin{enumerate}[nosep]
  \item What observable consequences follow from legitimation failure?
  \item Using Model~\ref{mod:legitimation-dynamics}, identify the parameter conditions under which a crisis event triggers lasting legitimation failure vs.\ temporary disruption followed by recovery.
  \item Connect to Chapter~\ref{ch:how-belief-regimes-fracture}: how does legitimation failure relate to belief regime fracture?
\end{enumerate}
\end{exercise}

\begin{exercise}
Apply the compound legitimation framework (Proposition~\ref{prop:compound}) to a specific historical case.
\begin{enumerate}[nosep]
  \item Choose a society where economic and racial inequality are correlated.
  \item Show how economic narratives legitimate racial hierarchy and vice versa.
  \item Identify a historical moment when the compound was challenged (e.g., a cross-racial labor movement) and explain why it succeeded or failed.
\end{enumerate}
\end{exercise}

\begin{exercise}
Construct the comparative legitimation table (Section~\ref{sec:ch24-comparative}) using actual data.
\begin{enumerate}[nosep]
  \item Find Gini coefficients and top-income shares for at least four countries.
  \item Find survey data on meritocratic or just-world belief (World Values Survey, ISSP).
  \item Map the legitimation infrastructure in each country (think tanks, media, policy design).
  \item Does the model's prediction (higher inequality $\to$ more legitimation investment) hold cross-nationally?
\end{enumerate}
\end{exercise}

\begin{exercise}
Analyze the POUM hypothesis (Proposition~\ref{prop:poum}) using mobility data.
\begin{enumerate}[nosep]
  \item Find actual intergenerational mobility rates for the US (e.g., Chetty et al.).
  \item Find survey data on \emph{perceived} mobility (what people think their chances of upward mobility are).
  \item Compute the ``mobility perception gap'': the difference between perceived and actual mobility.
  \item Estimate the doxonomic investment needed to maintain this gap (i.e., to sustain the belief in mobility despite contrary evidence).
\end{enumerate}
\end{exercise}

\begin{exercise}[Computational]
Simulate the legitimation dynamics model (Model~\ref{mod:legitimation-dynamics}):
\begin{enumerate}[nosep]
  \item Set $I_0 = 20$, $I(0) = 25$, $L(0) = 0.7$, $\alpha = 0.1$, $\phi_L = 3$, $\beta = 0.02$, $\gamma_I = 0.5$, $\mu = 0.3$.
  \item Simulate 200 time steps with no crisis ($\delta_{\text{crisis}} = 0$). Plot $L(t)$ and $I(t)$.
  \item At $t = 100$, introduce a crisis ($\delta_{\text{crisis}} = 0.2$). How does the system respond?
  \item Experiment with different $\phi_L$ values. What is the minimum narrative investment needed to prevent legitimation collapse after the crisis?
  \item Compare scenarios where the crisis triggers redistributive policy ($I$ decreases) vs.\ where institutional capture blocks redistribution ($I$ continues rising despite low $L$).
\end{enumerate}
\end{exercise}

\chapter{Race, Coloniality, and Gender as Doxonomic Systems}
\label{ch:race-coloniality-gender}

\epigraph{The most potent weapon of the oppressor is the mind of the oppressed.}{Steve Biko}

\noindent
Race, coloniality, and gender are among the most consequential doxonomic systems in human history.
Each involves the large-scale production and institutional stabilization of beliefs about what kinds of people exist, what they are capable of, and what they deserve.
A field that claims to study how power shapes belief cannot treat these as optional examples.
They are central cases, arguably the cases where doxonomic production has been most systematic, most durable, and most materially consequential.

This chapter extends the doxonomic framework to analyze these systems, revealing both the framework's strengths (it provides precise tools for tracing funded narrative production) and its limitations (some of the most durable belief systems are reproduced through unpaid labor, embodied practice, and institutional routine rather than through the explicit funding flows that doxonomic accounting tracks most naturally).

\section{Race as Manufactured Social Ontology}
\label{sec:ch24b-race}

\subsection{Definition and mechanism}

\begin{definition}[Racial Regime]
\label{def:racial-regime}
\index{racial regime}
A \emph{racial regime} is a doxonomic system that produces and maintains a categorical division of humanity into hierarchically ordered groups based on perceived phenotypic or ancestral characteristics, where:
\begin{enumerate}[nosep]
  \item the categories are presented as natural rather than constructed,
  \item differential treatment is legitimated by attributed group capacities or deficiencies,
  \item the institutional infrastructure (law, education, media, science) continuously reproduces the classification \autocite{gould1981},
  \item the material consequences of classification (segregation, differential wealth, differential health) generate ``evidence'' that appears to confirm the classification.
\end{enumerate}
\end{definition}

The fourth criterion is distinctively doxonomic: the racial regime produces the conditions that make it appear true.
Segregation produces differential educational outcomes; differential education produces differential economic outcomes; differential economics produces differential health outcomes; and all of these differentials are then cited as evidence that racial categories correspond to genuine differences in ability, motivation, or culture.

\subsection{The institutional genealogy of scientific racism}

The history of race is a history of doxonomic production.
Scientific racism in the 18th--20th centuries was not a marginal pseudoscience; it was a heavily funded, institutionally embedded research program that produced the ``common sense'' of racial hierarchy \autocite{dubois1903}.

\paragraph{The production chain.}
The institutional infrastructure of scientific racism included:
\begin{itemize}[nosep]
  \item \textbf{Universities and research institutes}: craniometry, phrenology, eugenics institutes at leading universities (Cold Spring Harbor, Kaiser Wilhelm Institute, UCL's Galton Laboratory). The Eugenics Record Office (US, 1910--1939) received substantial Rockefeller and Carnegie funding.
  \item \textbf{Census bureaus}: the racial categories of national censuses (which varied dramatically between countries) reified administrative classifications as natural kinds.
  \item \textbf{Legal systems}: anti-miscegenation laws, Jim Crow, apartheid, colonial legal codes, all requiring precise racial classification and producing an institutional apparatus for maintaining it.
  \item \textbf{Educational institutions}: racially segregated schools, textbooks presenting racial hierarchy as scientific fact, curricula omitting the achievements of non-European civilizations.
  \item \textbf{Popular science and media}: National Geographic's century of racial representation, IQ testing popularization, the ``bell curve'' discourse \autocite{herrnstein1994}.
\end{itemize}

\paragraph{The funding dimension.}
Unlike many doxonomic systems, early racial regime production was funded directly by the state (colonial administrations, census bureaus, segregated school systems) and by major philanthropies (Carnegie, Rockefeller, and Harriman foundations funded eugenics research).
The total institutional investment in producing and maintaining racial categories over the period 1800--2000 is difficult to estimate but clearly enormous, encompassing the entire apparatus of colonial governance, segregated institutions, and racialized law enforcement.

\subsection{The racial belief production model}

\begin{model}[Racial Belief Production]
\label{mod:racial-production}
\index{racial regime!production model}
The strength of racial-hierarchy belief $\belief{R}$ is sustained by:
\[
  \frac{d\belief{R}}{dt} = \alpha_{\text{inst}} \phi_{\text{inst}} + \alpha_{\text{seg}} S(t) + \beta \belief{R}(1-\belief{R}) - \delta \belief{R} - \gamma E(t)
\]
where:
\begin{itemize}[nosep]
  \item $\phi_{\text{inst}}$ is institutional investment in racial classification (curricula, media, law),
  \item $S(t)$ is the degree of residential and economic segregation (which generates ``evidence'' for racial difference through differential outcomes),
  \item $\beta$ captures peer reinforcement and social transmission,
  \item $\delta$ is natural decay (generational replacement, individual updating),
  \item $E(t)$ represents counter-evidence (interracial cooperation, anti-racist education, disconfirming experience).
\end{itemize}
\end{model}

The segregation term $S(t)$ is crucial: spatial and economic separation produces the differential outcomes that appear to confirm the racial narrative, creating a self-fulfilling loop analogous to the selfishness trap of Chapter~\ref{ch:human-nature}.

\begin{proposition}[Segregation as Doxonomic Feedback]
\label{prop:segregation-feedback}
\index{segregation!doxonomic feedback}
If segregation $S(t)$ is itself a function of racial belief ($S(t) = h(\belief{R}(t))$ with $h' > 0$, so stronger racial belief leads to more segregation through housing discrimination, white flight, school district formation), then the system has a positive feedback loop:
\[
  \belief{R} \uparrow \;\Rightarrow\; S \uparrow \;\Rightarrow\; \text{differential outcomes} \;\Rightarrow\; \text{``evidence'' for racial hierarchy} \;\Rightarrow\; \belief{R} \uparrow
\]
This loop can sustain racial hierarchy belief \emph{even after the explicit institutional investment} $\phi_{\text{inst}}$ \emph{is withdrawn} (e.g., after the formal end of Jim Crow or apartheid), because segregation continues to generate confirming ``evidence.''
Breaking the loop requires intervening on segregation directly, not merely on belief.
\end{proposition}

\subsection{The post-civil-rights transformation}

The explicit institutional investment in racial hierarchy ($\phi_{\text{inst}}$) declined substantially after the civil rights movement (US), anti-apartheid movement (South Africa), and decolonization (Global South).
But racial belief did not decline proportionally, because:
\begin{enumerate}[nosep]
  \item the segregation feedback loop ($S(t)$ term) continued operating,
  \item new narrative forms emerged: ``colorblind racism'' (denying the relevance of race while maintaining racial structures), ``cultural deficit'' explanations (attributing racial inequality to cultural rather than biological differences, but with the same doxonomic function), and ``reverse racism'' narratives (reframing anti-racist policy as the ``real'' racism) \autocite{bonilla-silva2018},
  \item the institutional infrastructure of racial classification (census categories, residential patterns, educational tracking, criminal justice disparities) remained largely intact.
\end{enumerate}

The doxonomic insight: the racial regime \emph{transformed} rather than disappeared.
The explicit belief ``race X is biologically inferior'' was replaced by the implicit belief ``racial inequality reflects cultural or behavioral differences,'' a narrative shield (Definition~\ref{def:narrative-shield}) that achieves the same legitimation function through different discursive means.

\section{Colonial Civilizational Narratives}
\label{sec:ch24b-coloniality}

\subsection{Colonialism as doxonomic operation}

Colonialism was, among other things, a massive doxonomic operation: the production of beliefs about civilizational hierarchy that justified extraction, governance, and violence.
\textcite{said1978} showed how ``Orientalism'' (a systematic body of knowledge about the East produced by Western institutions) functioned as epistemic infrastructure for colonial rule.

The belief production chain for colonial narratives ran through:
\begin{itemize}[nosep]
  \item \textbf{Missionary education}: replacing indigenous knowledge systems with European curricula, languages, and religious frameworks.
  \item \textbf{Colonial administration}: legal and census categories imposing European classifications on diverse populations; the creation of ``tribes'' and ``castes'' as administrative units that then became social identities.
  \item \textbf{Academic disciplines}: anthropology (classifying ``primitive'' societies), geography (mapping ``empty'' lands for settlement), linguistics (ranking languages by ``complexity''), history (narrating progress as a European monopoly).
  \item \textbf{Museums and exhibitions}: world's fairs, natural history museums, and ethnographic displays that presented non-European peoples as specimens for study rather than agents of history.
  \item \textbf{Economic institutions}: structuring labor markets, land tenure, and trade around racial and civilizational categories.
\end{itemize}

\subsection{Coloniality after colonialism}

The persistence of these narratives after formal decolonization (what decolonial theorists call \emph{coloniality}) is a doxonomic phenomenon: the institutional infrastructure outlasts the political regime that built it.

\begin{definition}[Doxonomic Coloniality]
\label{def:coloniality}
\index{coloniality}
\emph{Doxonomic coloniality} is the persistence of colonial belief structures after formal decolonization, sustained by:
\begin{enumerate}[nosep]
  \item educational institutions that still teach from colonial-era curricula or frameworks (European-centered history, English/French as languages of prestige),
  \item academic disciplines whose categories were shaped by colonial needs (``development economics'' as the study of why non-Western countries are ``behind''),
  \item international institutions (IMF, World Bank) whose policy prescriptions reproduce colonial-era assumptions about governance and economic organization,
  \item media representations that continue to frame the Global South as a site of crisis, backwardness, or exoticism.
\end{enumerate}
\end{definition}

\begin{example}
\label{ex:development-economics}
``Development economics'' is a doxonomically interesting field: its very existence presupposes a civilizational narrative (some countries are ``developed,'' others are ``developing'') that maps onto colonial hierarchy.
The field's institutional infrastructure (World Bank research departments, USAID, bilateral development agencies, academic programs) produces knowledge that tends to frame non-Western economies as deficient versions of Western ones, requiring Western-style reforms to ``catch up.''
This is not to deny that poverty is real or that economic analysis is useful; it is to note that the framing of the problem already contains the colonial assumption of a single developmental path.
\end{example}

\subsection{Epistemic violence and knowledge erasure}

\begin{definition}[Epistemic Violence]
\label{def:epistemic-violence}
\index{epistemic violence}
\emph{Epistemic violence} is the systematic destruction, suppression, or devaluation of a community's knowledge systems through doxonomic means.
It includes:
\begin{enumerate}[nosep]
  \item the destruction of indigenous knowledge repositories (libraries, oral traditions, medicinal practices),
  \item the delegitimation of non-Western epistemologies as ``superstition'' or ``pre-scientific,''
  \item the imposition of colonial languages as the medium of education and governance, rendering indigenous knowledge inaccessible to subsequent generations,
  \item the extraction and appropriation of indigenous knowledge without attribution (biopiracy, appropriation of agricultural techniques, adoption of medicinal compounds without crediting their sources).
\end{enumerate}
\end{definition}

Epistemic violence is the destructive counterpart of doxonomic production: where doxonomic production builds beliefs, epistemic violence destroys them.
The two operate in tandem during colonization: the colonizer's knowledge system is institutionally amplified while the colonized's knowledge system is institutionally suppressed.

\section{Gender as Doxonomic Construction}
\label{sec:ch24b-gender}

\subsection{The gender regime}

\begin{definition}[Gender Regime]
\label{def:gender-regime}
\index{gender regime}
A \emph{gender regime} is a doxonomic system producing and maintaining beliefs about:
\begin{enumerate}[nosep]
  \item natural differences between men and women in capacity, temperament, and social role,
  \item the complementarity or hierarchy of these roles,
  \item the unnaturalness of deviations from prescribed gender behavior,
  \item the naturalness of the sexual division of labor (productive vs.\ reproductive work).
\end{enumerate}
\end{definition}

\subsection{The distinctive features of gender doxonomics}

Gender regimes are sustained through family socialization, religious institutions, educational tracking, media representation, workplace norms, and legal codes.
Their doxonomic cost is distributed across many institutions, making them difficult to account for with the funding-tracing methods of Chapter~\ref{ch:belief-flow-accounting}.
Three features distinguish gender doxonomics from the funding-centered models of earlier chapters:

\begin{enumerate}
  \item \textbf{Reproduction through unpaid labor.}
  Much gender-regime maintenance happens through unpaid labor: parents socializing children, communities enforcing norms, cultures producing representations.
  The ``funding'' is not monetary but temporal and emotional.
  This means that the doxonomic accounting framework must be extended to include non-monetary investment.

  \item \textbf{Embodied belief.}
  Gender beliefs operate heavily through layers 3--5 of the layered ontology (Principle~\ref{prin:layered-ontology}): practical schemas (how one moves, speaks, presents oneself), identity commitments (``being a man/woman''), and affective orientations (what feels natural or disgusting).
  These layers are more resistant to discursive counter-evidence than propositional beliefs.

  \item \textbf{Biological anchoring.}
  Gender regimes anchor themselves to real biological differences (reproductive anatomy, average physical size differences, hormonal variation) while vastly overgeneralizing what those differences imply for social roles.
  The doxonomic operation is not the invention of difference from nothing but the \emph{amplification} of modest biological variation into comprehensive social hierarchy.
\end{enumerate}

\subsection{The gender belief production chain}

\begin{model}[Gender Regime Production]
\label{mod:gender-production}
\index{gender regime!production model}
Gender beliefs are produced through a multi-channel pipeline:
\begin{enumerate}[nosep]
  \item \textbf{Family} (ages 0--5): gendered names, clothing, toys, differential parental behavior (studies show parents talk to daughters more about emotions and to sons more about achievement).
  \item \textbf{Education} (ages 5--22): gendered expectations in classrooms, differential encouragement in STEM, sex-segregated sports, gendered career counseling.
  \item \textbf{Media} (all ages): the Geena Davis Institute found that female characters are underrepresented 2:1 in film, more likely to be shown in domestic or romantic contexts, and less likely to be shown in professional or leadership roles.
  \item \textbf{Religion} (all ages): many religious traditions produce explicit gender-role prescriptions backed by the credibility of divine authority.
  \item \textbf{Workplace} (ages 18--65): gendered job descriptions, differential evaluation standards, mommy-track career penalties, sexual harassment as norm enforcement.
  \item \textbf{Law and policy}: historical coverture laws, ongoing gaps in parental leave policy, differential criminal sentencing.
\end{enumerate}
The combined effect is a total-institution doxonomic environment: the gender regime is encountered in every institutional setting from birth to death.
\end{model}

\subsection{Feminist counter-doxonomics}

Feminist movements constitute the most sustained counter-doxonomic campaign in history.
Over two centuries, feminist organizations have:
\begin{itemize}[nosep]
  \item challenged the naturalization of gender hierarchy through theoretical critique,
  \item built counter-institutions (women's studies programs, feminist media, advocacy organizations),
  \item achieved legal reforms that reduced the gender regime's institutional enforcement (suffrage, equal pay legislation, anti-discrimination law),
  \item shifted norm perceptions measurably: beliefs about women's capabilities have changed dramatically in survey data over 50 years.
\end{itemize}

Yet the gender regime remains partially intact, especially at the embodied and affective levels.
This illustrates a general doxonomic principle: propositional beliefs are the most responsive to counter-narrative, while practical schemas and identity commitments change more slowly.

\section{Intersections and Compound Systems}
\label{sec:ch24b-intersections}

Race, gender, and class regimes do not operate independently.
They intersect, reinforce, and sometimes contradict each other.

\subsection{The interaction structure}

\begin{remark}[Doxonomic Intersection]
\label{prop:intersection}
\index{intersectionality}
When multiple doxonomic regimes (racial, gender, class) operate simultaneously on agent $j$, the combined belief effect is not additive but interactive:
\[
  \belief{j}^{\text{total}} \neq \belief{j}^{\text{race}} + \belief{j}^{\text{gender}} + \belief{j}^{\text{class}}.
\]
Rather, the regimes modulate each other: racial narratives are gendered (stereotypes differ for men and women of the same racialized group), class narratives are racialized (the ``deserving poor'' is often implicitly white), and gender narratives are class-inflected (the ``ideal worker'' presupposes a wife managing domestic labor).
\end{remark}

\begin{example}
\label{ex:intersection}
The narrative ``welfare queen'' (US political discourse, 1976--present) operates at the intersection of race, gender, and class:
\begin{itemize}[nosep]
  \item \emph{Racial}: the image is implicitly (and sometimes explicitly) Black.
  \item \emph{Gender}: the image is female, violating prescriptions about maternal self-sacrifice and sexual restraint.
  \item \emph{Class}: the image is poor, violating the meritocratic narrative of self-sufficiency.
\end{itemize}
No single-axis analysis captures the narrative's doxonomic function.
Its power comes from activating all three regimes simultaneously, producing a figure that concentrates the stigma of racial, gender, and class deviance in a single image.
This compound stigma was deployed to justify welfare reform in the 1990s, a policy change with measurable effects on millions of people, driven in significant part by a doxonomic construction.
\end{example}

\subsection{Sheaf-theoretic intersectionality}

In the sheaf-theoretic framework of Chapter~\ref{ch:contradiction-coherence}, the local belief systems for race, gender, and class may have different gluing obstructions depending on the institutional context.

\begin{example}
\label{ex:sheaf-intersection}
Consider the local belief systems:
\begin{itemize}[nosep]
  \item $U_{\text{workplace}}$: ``women are less competent in leadership'' (gender regime) + ``minorities are less qualified'' (racial regime) + ``pay reflects productivity'' (meritocratic regime).
  \item $U_{\text{home}}$: ``women are natural caregivers'' (gender regime) + ``family is private'' (class regime shields domestic labor from economic analysis).
  \item $U_{\text{political}}$: ``all citizens are equal before the law'' (democratic regime) + ``affirmative action is reverse discrimination'' (meritocratic regime applied selectively).
\end{itemize}
The transition maps between these local systems are inconsistent: the ``natural caregiver'' of $U_{\text{home}}$ contradicts the ``equal citizen'' of $U_{\text{political}}$, but the contradiction is managed by keeping the two domains institutionally separate.
The sheaf obstruction is nonzero ($H^1 \neq 0$), but the social system manages the inconsistency through domain segregation rather than resolution.
\end{example}

\section{Alternative and Complementary Explanations}
\label{sec:ch25-alternatives}

The doxonomic account of racial, colonial, and gender regimes is one explanatory layer among several.
Competing and complementary accounts include:
\begin{itemize}[nosep]
  \item \textbf{Psychological accounts}: in-group/out-group cognition, implicit bias, and social identity theory explain individual-level racial and gender attitudes without reference to institutional funding.
  These mechanisms are real and interact with doxonomic processes: institutions amplify and channel pre-existing psychological tendencies.
  \item \textbf{Cultural-evolutionary accounts}: cultural norms around gender, race, and kinship may have deep evolutionary or historical roots that precede and partially operate independently of any specific institutional investment.
  \item \textbf{Materialist accounts without the doxonomic layer}: Marxist and political-economy explanations of racial and gender hierarchy focus on labor-market segmentation, property relations, and class structure, treating ideology as derivative rather than as an independently analyzable production system.
  \item \textbf{Phenomenological accounts}: lived experience of racial and gender identity involves embodied, affective, and existential dimensions that the doxonomic framework's focus on institutional production does not fully capture.
\end{itemize}
Doxonomics does not replace these accounts.
It adds a specific dimension (the material infrastructure of belief production) that the other accounts tend to leave implicit.
The strongest analysis combines doxonomic institutional tracing with psychological, cultural, and phenomenological insights.

\section{Limitations of the Doxonomic Framework for Identity Systems}
\label{sec:ch24b-limitations}

\subsection{What the framework captures}

The doxonomic framework captures well:
\begin{itemize}[nosep]
  \item the role of funded institutions in producing and circulating racial, colonial, and gender narratives,
  \item the self-fulfilling mechanisms by which these narratives generate confirming evidence,
  \item the material interests served by specific narratives,
  \item the formal dynamics of belief production, competition, and regime change.
\end{itemize}

\subsection{What it risks missing}

\begin{enumerate}
  \item \textbf{Agency and resistance.}
  The doxonomic framework, with its emphasis on institutional production, can understate the agency of those subjected to oppressive regimes.
  People are not blank slates on which doxonomic systems write; they resist, reinterpret, subvert, and transform the narratives directed at them.
  Any adequate doxonomic analysis of race, coloniality, or gender must include the counter-doxonomic activity of the oppressed.

  \item \textbf{Lived experience.}
  The formal models treat belief as a numerical variable $\belief{}(t)$.
  But racial identity, gender identity, and colonial subjectivity are not reducible to belief intensities: they involve lived experience, embodied practice, community belonging, and existential meaning.
  The doxonomic framework provides partial insight, not total comprehension.

  \item \textbf{Non-monetary reproduction.}
  As noted above, gender and racial regimes are reproduced substantially through unpaid labor and institutional routine.
  The funding-tracing methodology that works well for think-tank-produced narratives is less well-suited to regimes maintained through family socialization, peer pressure, and cultural tradition.

  \item \textbf{The positionality problem.}
  The analyst's own position within racial, gender, and class regimes affects what they can see and how they frame the analysis.
  The symmetry obligation (Principle~\ref{prin:symmetry}) is necessary but not sufficient: a white researcher studying racial doxonomics should be aware that their analysis is itself shaped by the racial regime they are analyzing.
\end{enumerate}

These limitations do not invalidate the doxonomic analysis of identity systems; they define its scope.
Doxonomics provides one lens (the institutional-material lens) among several that are needed for a complete understanding.
It should be used alongside, not in place of, phenomenological, ethnographic, and literary approaches to race, coloniality, and gender.

\section{Notes and References}
\label{sec:ch24b-notes}

On race as a social construction with material effects, see \textcite{dubois1903} and \textcite{fanon1961}.
The history of scientific racism is documented in \textcite{gould1981} and Kendi (2016).
Colonial knowledge production is analyzed in \textcite{said1978}.
On coloniality as persisting beyond formal colonialism, see \textcite{quijano2000} and \textcite{mignolo2011}.

On the doxonomic function of education in reproducing class and racial hierarchy, see \textcite{bowles1976} and \textcite{bourdieu1984}.
For epistemic injustice and its relation to social identity, see \textcite{fricker2007} and \textcite{medina2013}.
On epistemic violence, see Spivak (1988).

On gender as a performed and institutionally sustained category, see \textcite{butler1990} and Connell (1987).
On the coloniality of gender, the claim that modern gender categories are themselves a colonial imposition, see \textcite{lugones2010}.
On media representation and gender, see the Geena Davis Institute reports.
For the ``welfare queen'' as a doxonomic construction, see \textcite{gilens1999}.

The intersectionality framework originates with \textcite{crenshaw1989}.
System justification in the context of racial hierarchy is analyzed in \textcite{jost2004}.
On the ``wages of whiteness,'' see \textcite{du-bois2007}.

\section{Exercises}

\begin{exercise}
Construct a doxonomic account for racial hierarchy in a specific national context.
\begin{enumerate}[nosep]
  \item Identify at least four institutional channels that produce or reproduce racial belief.
  \item For each channel, estimate whether the investment is primarily monetary, temporal (unpaid labor), or institutional (routine practices).
  \item Estimate the relative contribution of each channel to $\belief{R}$.
  \item How has the balance among channels changed over the past 50 years?
\end{enumerate}
\end{exercise}

\begin{exercise}
Compare the doxonomic production of racial categories in two colonial contexts (e.g., British India and French West Africa).
\begin{enumerate}[nosep]
  \item What institutional differences led to different racial classification systems?
  \item How do these differences persist in the post-colonial period?
  \item Apply Model~\ref{mod:racial-production} to both cases: which terms differ?
\end{enumerate}
\end{exercise}

\begin{exercise}
Using the model of racial belief production (Model~\ref{mod:racial-production}), argue that residential segregation functions as a self-reinforcing doxonomic feedback loop.
\begin{enumerate}[nosep]
  \item Write $S(t) = h(\belief{R}(t))$ with a specific functional form for $h$.
  \item Show that the coupled system $(\belief{R}, S)$ has two stable equilibria (high-segregation/high-belief and low-segregation/low-belief).
  \item Under what conditions can the system be moved from the high to the low equilibrium?
  \item What policy intervention is implied (intervening on belief, on segregation, or on both)?
\end{enumerate}
\end{exercise}

\begin{exercise}
Analyze the doxonomic infrastructure maintaining gender complementarity beliefs in a religious institutional context of your choice.
\begin{enumerate}[nosep]
  \item Map the production chain (texts, sermons, schools, family practices).
  \item How does this infrastructure differ from the funding-driven models of earlier chapters?
  \item What is the ``return on investment'' when much of the investment is unpaid labor?
  \item How would you measure the doxonomic output of this system using the tools of Chapters~\ref{ch:doxometric-measurement}--\ref{ch:text-media}?
\end{enumerate}
\end{exercise}

\begin{exercise}
Analyze the ``welfare queen'' narrative (Example~\ref{ex:intersection}) as an intersectional doxonomic product.
\begin{enumerate}[nosep]
  \item Trace the production chain: who coined the term, which institutions amplified it, through which media channels?
  \item Apply framing analysis (Chapter~\ref{ch:text-media}) to characterize its causal story, normative valence, subject position, and prescribed action.
  \item What policy outcomes did the narrative enable?
  \item What would a counter-narrative look like, and what institutional infrastructure would be needed to circulate it?
\end{enumerate}
\end{exercise}

\begin{exercise}
Apply the sheaf framework of Chapter~\ref{ch:contradiction-coherence} to the intersection of racial and meritocratic beliefs.
\begin{enumerate}[nosep]
  \item Define local belief systems for at least three institutional contexts.
  \item Identify the transition maps between them.
  \item Compute whether the gluing obstruction $H^1$ is zero (the beliefs are globally consistent) or nonzero (they are inconsistent).
  \item How does the social system manage the inconsistency?
\end{enumerate}
\end{exercise}

\begin{exercise}
Assess the limitations of the doxonomic framework (Section~\ref{sec:ch24b-limitations}) by comparing its analysis of a specific case with an analysis from a different tradition (phenomenological, ethnographic, or literary).
\begin{enumerate}[nosep]
  \item Choose a specific instance of racial, colonial, or gender belief production.
  \item Analyze it using the doxonomic framework (funding, institutions, formal model).
  \item Analyze it using an alternative framework (lived experience, cultural meaning, literary representation).
  \item What does each framework reveal that the other misses?
  \item How could the two be productively combined?
\end{enumerate}
\end{exercise}

\chapter{Austerity, Scarcity, and Fiscal Morality}
\label{ch:austerity}

\epigraph{The national budget must be balanced. The public debt must be reduced. The arrogance of the authorities must be moderated and controlled.}{Attributed to Cicero (apocryphal)}

\noindent
In 2012, a retired Greek schoolteacher named Dimitris Christoulas walked to Syntagma Square in central Athens and shot himself.
His suicide note read: ``I find no other solution than a dignified end before I start searching through garbage for food.''
His pension had been cut by 40\% under the austerity program imposed by Greece's creditors.
He was 77 years old, had worked his entire adult life, and had committed no financial crime.
The narrative that justified his impoverishment---that Greece had ``lived beyond its means'' and must now ``tighten its belt''---is the subject of this chapter.

Austerity narratives---the claim that public spending must be cut, that governments must ``live within their means,'' that debt is inherently dangerous---are among the most successful doxonomic products of the past century.
They have shaped policy in dozens of countries, redirected trillions of dollars from public services to debt repayment, and helped justify a redistribution of risk from capital to labor.
This chapter analyzes their construction, the institutional infrastructure that produces them, the formal dynamics of fiscal belief, and the distributional consequences of treating a policy choice as a moral imperative.

The doxonomic interest is not primarily whether austerity is good or bad economics (though the evidence is relevant) but how a specific fiscal framework---one whose empirical support is contested and whose distributional consequences are highly unequal---achieved the status of common sense across multiple countries within a few years.

\section{The Household Metaphor}
\label{sec:ch25-household}

\subsection{Definition and mechanism}

\begin{definition}[Fiscal Morality Narrative]
\label{def:fiscal-morality}
\index{fiscal morality}
The \emph{fiscal morality narrative} is the claim that sovereign budgets are analogous to household budgets: governments, like families, must not spend more than they earn, must avoid debt, and must ``tighten their belts'' in hard times.
\end{definition}

This analogy is economically misleading---sovereign currency issuers face fundamentally different constraints than households (they cannot go bankrupt in their own currency, they can create money, and their spending creates income for the private sector)---but doxonomically powerful because it maps complex macroeconomics onto intuitive domestic experience.
Every adult has managed a household budget; almost no one has managed a sovereign treasury.
The analogy exploits this asymmetry in lived experience.

\subsection{The moral dimension}

The household metaphor is not merely an analogy; it is a \emph{moral} claim.
Debt is framed as sin, spending as temptation, austerity as virtue.
\textcite{graeber2011} traces this moral loading of debt across five millennia, showing that the equation of debt with moral failing is a culturally constructed norm that serves the interests of creditors.

\begin{proposition}[Debt-Morality Fusion]
\label{prop:debt-morality}
\index{debt morality}
The fiscal morality narrative fuses two distinct questions:
\begin{enumerate}[nosep]
  \item \textbf{Technical question}: what is the optimal level of public debt given growth, interest rates, and multiplier effects?
  \item \textbf{Moral question}: is it ``right'' for governments to borrow?
\end{enumerate}
By fusing these, the narrative transforms a technical debate (where reasonable economists disagree) into a moral imperative (where ``responsible'' people agree).
This fusion makes the narrative resistant to empirical refutation: even if austerity produces bad economic outcomes, it can still be defended as ``the right thing to do.''
\end{proposition}

\subsection{Institutional amplification}

The household metaphor circulates through:
\begin{itemize}[nosep]
  \item \textbf{Political speech}: ``We have maxed out our credit card'' (repeated by politicians across the spectrum in multiple countries).
  \item \textbf{Media framing}: news reports consistently use household-budget language for sovereign finances---``the government's deficit,'' ``the nation's credit card,'' ``living beyond our means.''
  \item \textbf{Think-tank output}: organizations like the Peterson Foundation (US), the Taxpayers' Alliance (UK), and the Initiative for a Responsible Federal Budget produce materials that frame public debt as a crisis requiring immediate belt-tightening.
  \item \textbf{Popular economics}: bestselling books (e.g., \emph{The Coming Generational Storm}, \emph{The Debt Bomb}) translate the narrative for mass audiences.
\end{itemize}

\section{The Doxonomic Genealogy of Austerity}
\label{sec:ch25-genealogy}

\subsection{The classical foundation}

The austerity narrative draws on a long intellectual tradition:
\begin{itemize}[nosep]
  \item \textbf{Ricardo's equivalence} (1817): the proposition that government borrowing is equivalent to taxation because rational agents anticipate future taxes. In its strong form (which Ricardo himself doubted), it implies that fiscal policy is irrelevant.
  \item \textbf{The ``Treasury view''} (1920s--30s): the British Treasury's doctrine that government spending necessarily ``crowds out'' private investment, implying that stimulus cannot increase total output.
  \item \textbf{Monetarism} (1960s--70s): Friedman's argument that fiscal policy is ineffective and that only monetary policy matters, providing intellectual cover for reducing government's fiscal role.
  \item \textbf{Public choice theory} (1970s--80s): Buchanan and Tullock's argument that government spending reflects bureaucratic self-interest rather than public welfare, implying that less spending is generally better.
\end{itemize}

Each intellectual precursor served a specific institutional function: Ricardo's equivalence was developed in the context of debates about war financing; the Treasury view defended the gold standard; monetarism justified reducing the post-war welfare state; public choice theory provided the theoretical foundation for Thatcher/Reagan spending cuts.

\subsection{The institutional build-out}

The modern austerity infrastructure was built during the neoliberal period (1970s--2000s):

\begin{enumerate}
  \item \textbf{Think tanks}: Heritage Foundation's ``Mandate for Leadership'' (1981) included detailed spending-cut proposals for every federal department. The Cato Institute's annual ``Budget Pork'' reports framed all government spending as waste. In the UK, the Institute of Economic Affairs and the Adam Smith Institute played analogous roles.

  \item \textbf{International institutions}: the IMF's ``structural adjustment programs'' (1980s--90s) conditioned loans to developing countries on austerity measures, creating a global institutional norm that debt reduction trumps social spending.
  The ``Washington Consensus'' codified this norm: fiscal discipline was listed as the first of ten policy prescriptions.

  \item \textbf{Fiscal rules and institutions}: balanced-budget amendments, debt brakes (Germany's \emph{Schuldenbremse}), the Maastricht Treaty's 3\%/60\% rules for eurozone members---all encoded the austerity narrative in law, removing fiscal decisions from democratic deliberation.

  \item \textbf{Credit rating agencies}: Moody's, S\&P, and Fitch evaluate sovereign creditworthiness, creating a private-sector enforcement mechanism for fiscal discipline. A credit downgrade raises borrowing costs, punishing governments that reject the austerity narrative.
\end{enumerate}

\section{Academic Credibility Laundering}
\label{sec:ch25-credibility}

\subsection{The Reinhart-Rogoff episode}

The most vivid example of academic credibility laundering in the austerity narrative is the Reinhart-Rogoff affair.

\begin{example}
\label{ex:reinhart-rogoff}
In January 2010, Carmen Reinhart and Kenneth Rogoff published ``Growth in a Time of Debt,'' claiming that countries with public debt exceeding 90\% of GDP experienced dramatically lower growth.
The paper was not peer-reviewed but was presented at the American Economic Association conference and published as an NBER working paper.

\paragraph{The doxonomic chain.}
\begin{enumerate}[nosep]
  \item \emph{Production}: two Harvard economists with high epistemic capital produced a paper with a clear, quotable result (``90\% is the threshold'').
  \item \emph{Amplification}: the Peterson Foundation promoted the finding; EU Commissioner Olli Rehn cited it in official policy documents; Paul Ryan cited it in his US budget proposal; George Osborne referenced it in justifying UK austerity.
  \item \emph{Capture}: within months, ``the 90\% threshold'' became common sense in policy debates---a number that could be cited to end any argument about whether austerity was necessary.
\end{enumerate}

\paragraph{The collapse.}
In 2013, Thomas Herndon, a graduate student at UMass Amherst, discovered that the paper contained a spreadsheet coding error (five countries were accidentally excluded from a key calculation), selective data exclusion, and an unconventional weighting methodology.
Correcting these errors eliminated the sharp threshold: countries above 90\% debt had lower but still positive growth, and the relationship was far from the cliff-edge the original paper suggested.

\paragraph{The doxonomic aftermath.}
The paper had already served its function: austerity policies justified partly by its findings had been implemented across Europe, affecting millions of people.
The correction received far less institutional amplification than the original claim---a classic asymmetry in doxonomic production, where the initial framing captures common sense and corrections arrive too late to undo the policy consequences.
\end{example}

\subsection{The general pattern}

The Reinhart-Rogoff case illustrates a general doxonomic pattern:

\begin{model}[Academic Credibility Laundering]
\label{mod:credibility-laundering}
\index{credibility laundering!academic}
\begin{enumerate}[nosep]
  \item A researcher at a high-credibility institution produces a finding that supports a pre-existing policy preference.
  \item The finding is amplified through think tanks, media, and policy networks before full peer review.
  \item The finding's epistemic capital (derived from the researcher's institutional affiliation) exceeds its evidential warrant.
  \item Policy decisions are made based on the amplified finding.
  \item When the finding is later corrected or qualified, the policy has already been implemented and the narrative has already been captured.
\end{enumerate}
The time lag between amplification and correction is the ``laundering window''---the period during which unvetted research shapes policy through borrowed credibility.
\end{model}

\section{Crisis as Narrative Opportunity}
\label{sec:ch25-crisis}

\subsection{The shock doctrine model}

\begin{model}[Shock Doctrine]
\label{mod:shock-doctrine}
\index{shock doctrine}
Following \textcite{klein2007}, economic crises create windows of doxonomic opportunity.
During a crisis, public attention is high, uncertainty is high, and trust in existing narratives is shaken.
The first coherent narrative to fill the vacuum captures the frame.

Formally, let $H(t)$ be narrative entropy at time $t$.
A crisis causes $H(t) \to H_{\text{max}}$ (maximum uncertainty), followed by rapid entropy reduction as a new dominant frame is installed:
\[
  H(t) = H_{\text{max}} \cdot e^{-\lambda(t - t_{\text{crisis}})} + H_{\text{new}}
\]
The rate $\lambda$ depends on:
\begin{itemize}[nosep]
  \item the \emph{preparedness} of competing narratives (did someone have a pre-packaged explanation ready?),
  \item the \emph{institutional infrastructure} available to amplify the narrative,
  \item the \emph{credibility} of the narrative's proponents.
\end{itemize}
The group that supplies the new frame controls the post-crisis common sense.
\end{model}

\subsection{Narrative preparedness}

The austerity narrative's success after 2008 was not spontaneous; it reflected decades of institutional investment.
When the crisis hit, austerity proponents had:
\begin{itemize}[nosep]
  \item pre-written policy briefs and talking points,
  \item established relationships with journalists and editors,
  \item academic research (including Reinhart-Rogoff) ready for deployment,
  \item institutional credibility from decades of fiscal-conservatism advocacy,
  \item a simple, intuitive story (household metaphor) that required no economic training to understand.
\end{itemize}

Keynesian counter-narratives existed but lacked comparable infrastructure: fewer think tanks, less media presence, no simple metaphor as powerful as ``belt-tightening,'' and an argument (``spend more during a recession'') that contradicts intuitive household logic.

\section{The 2008--2013 Austerity Turn: A Detailed Case Study}
\label{sec:ch25-case}

\subsection{Phase 1: Crisis narrative (2008--2009)}

The initial frame was ``banks caused this.''
Public anger targeted financial institutions.
Narrative entropy was high.
Government intervention (bailouts, stimulus) was initially popular: the American Recovery and Reinvestment Act (February 2009, \$787 billion) passed with broad public support, and similar stimulus packages were enacted in the UK, Germany, China, and elsewhere.

\subsection{Phase 2: Reframing (2009--2010)}

Within 18 months, the dominant frame shifted from ``private-sector recklessness'' to ``public-sector profligacy.''
This reframing was not spontaneous.
It was driven by:
\begin{itemize}[nosep]
  \item \textbf{Think tanks}: the Reinhart-Rogoff paper (January 2010), the Peterson Foundation's ``fiscal responsibility'' campaign, and dozens of policy briefs from Heritage, Cato, and their European counterparts.
  \item \textbf{Financial press}: the Financial Times, Wall Street Journal, and Economist shifted framing from banking crisis to sovereign debt crisis.
  A content analysis would likely show the frequency of ``sovereign debt'' overtaking ``banking crisis'' in headlines by mid-2010.
  \item \textbf{European institutions}: the ECB, the European Commission, and the ``troika'' (ECB/EC/IMF) conditioned bailout assistance on austerity measures in Greece, Ireland, Portugal, and Spain, institutionalizing the reframing in policy.
  \item \textbf{Political actors}: the UK Conservative Party's 2010 campaign centered on ``the mess Labour left,'' reframing a global financial crisis as a domestic fiscal crisis.
  \item \textbf{Credit rating agencies}: S\&P's downgrade of US sovereign debt (August 2011) reinforced the narrative that government spending was the problem.
\end{itemize}

\subsection{Phase 3: Normalization (2010--2013)}

Austerity became common sense across much of Europe and influenced US fiscal policy (the sequester, the debt ceiling battles).
The household metaphor (``belt-tightening'') dominated media framing.
Counter-narratives (Keynesian stimulus, MMT) existed but lacked comparable institutional infrastructure.

\begin{remark}
The austerity turn illustrates several doxonomic mechanisms simultaneously:
the shock doctrine model (crisis opens the window),
narrative infrastructure effects (think tanks had pre-packaged arguments ready),
credibility laundering (an academic paper gave the 90\% threshold scientific authority),
and common-sense capture (within two years, austerity went from a policy choice to ``the only responsible option'').
\end{remark}

\subsection{Phase 4: Partial narrative erosion (2013--2020)}

The austerity consensus began eroding as:
\begin{itemize}[nosep]
  \item the Herndon correction undermined the 90\% threshold,
  \item austerity-adopting countries (Greece, UK, Spain) experienced prolonged recessions while non-austerity countries (US, with partial stimulus; Iceland, with heterodox policy) recovered faster,
  \item the IMF's own chief economist, Olivier Blanchard, published research arguing that fiscal multipliers had been underestimated---austerity was more damaging than predicted,
  \item populist movements (Syriza, Podemos, Sanders, Corbyn) explicitly challenged the austerity narrative.
\end{itemize}

But the institutional infrastructure remained: balanced-budget rules, the Maastricht criteria, credit-rating-agency discipline, and the household metaphor continued to constrain fiscal policy even as the intellectual case for austerity weakened.

\section{A Formal Model of Fiscal Belief}
\label{sec:ch25-model}

\begin{model}[Austerity Belief Dynamics]
\label{mod:austerity-dynamics}
\index{austerity belief}
Let $\belief{A}(t)$ be the strength of austerity belief in the population.
The dynamics are:
\[
  \frac{d\belief{A}}{dt} = \underbrace{\alpha \phi_A (1 - \belief{A})}_{\text{institutional production}} + \underbrace{\beta \cdot \text{Crisis}(t) \cdot (1 - \belief{A})}_{\text{crisis amplification}} - \underbrace{\gamma \cdot \Delta U(t) \cdot \belief{A}}_{\text{unemployment erosion}} - \underbrace{\delta \belief{A}}_{\text{natural decay}}
\]
where:
\begin{itemize}[nosep]
  \item $\phi_A$ is the institutional investment in austerity narratives,
  \item $\text{Crisis}(t)$ is an indicator for economic crisis periods (when the shock doctrine operates),
  \item $\Delta U(t)$ is the change in unemployment (rising unemployment associated with austerity erodes belief),
  \item $\delta$ is the natural decay rate (generational replacement, fading salience).
\end{itemize}
\end{model}

\begin{proposition}[Austerity Belief Persistence]
\label{prop:austerity-persistence}
The steady-state austerity belief (absent crisis) is:
\[
  \belief{A}^* = \frac{\alpha \phi_A}{\alpha \phi_A + \gamma \Delta U^* + \delta}
\]
If $\Delta U^* \approx 0$ (no unemployment change), the steady state depends only on the ratio of institutional production to natural decay: $\belief{A}^* = \alpha \phi_A / (\alpha \phi_A + \delta)$.
High institutional investment produces high equilibrium belief even without a crisis trigger.
\end{proposition}

\section{Distributional Analysis}
\label{sec:ch25-distributional}

\subsection{Who benefits from the austerity narrative?}

The austerity narrative's distributional consequences are highly unequal:

\begin{enumerate}
  \item \textbf{Holders of government bonds}: austerity protects the value of government debt by reducing default risk and inflation risk. Bondholders are disproportionately wealthy individuals and financial institutions.

  \item \textbf{Opponents of public spending}: austerity reduces the size of government, which aligns with the preferences of actors who view government spending as competing with private enterprise (for labor, for legitimacy, for narrative space).

  \item \textbf{Employers}: public-sector wage cuts and workforce reductions weaken the public sector as an alternative employer, increasing labor supply to the private sector and reducing bargaining power.

  \item \textbf{Tax-reduction advocates}: austerity narratives that frame spending as the problem (rather than revenue) implicitly argue for lower taxes, benefiting high-income taxpayers disproportionately.
\end{enumerate}

\begin{proposition}[Austerity as Redistribution]
\label{prop:austerity-redistribution}
\index{austerity!distributional effects}
Austerity programs redistribute income and risk from public-service users (disproportionately lower-income) to bondholders and taxpayers (disproportionately higher-income).
The fiscal morality narrative (Definition~\ref{def:fiscal-morality}) performs the doxonomic function of framing this redistribution as a moral imperative rather than a distributional choice.
\end{proposition}

\subsection{The Greek case}

\begin{example}
\label{ex:greece}
Greece's austerity program (2010--2018) provides a stark illustration:
\begin{itemize}[nosep]
  \item GDP contracted by 26\% (a depression comparable in scale to the 1930s US experience).
  \item Unemployment peaked at 27\% (youth unemployment at 60\%).
  \item Public health spending was cut by 40\%, leading to documented increases in infant mortality, HIV transmission, and mental health crises.
  \item Pension payments were cut by 40--50\%, pushing elderly poverty rates above 20\%.
  \item Meanwhile, the primary beneficiaries of the bailout funds were German and French banks that held Greek government bonds---roughly 90\% of bailout money went to debt servicing rather than the Greek economy.
\end{itemize}
The austerity narrative framed this as ``Greece living beyond its means'' and ``Greeks retiring too early''---a narrative shield (Definition~\ref{def:narrative-shield}) that deflected attention from the structural causes (eurozone design flaws, predatory lending, Northern European trade surpluses) to individual/national moral failings.
\end{example}

\section{Negative Cases and Counter-Narratives}
\label{sec:ch25-negative}

\subsection{Iceland: referendum against austerity}

Iceland rejected bank bailouts and austerity in 2008--2009 through a popular referendum.
Instead of socializing private banking losses, Iceland let the banks fail, imposed capital controls, and maintained social spending.
Its recovery was faster than most austerity-adopting European countries: GDP returned to pre-crisis levels by 2015, unemployment fell below 4\%, and the fiscal position recovered through growth rather than cuts.

The Icelandic case demonstrates that the austerity narrative requires specific institutional conditions to achieve capture.
Iceland's small size, homogeneous population, and strong civic traditions made the counter-narrative (``this was a banking crisis, not a spending crisis'') more credible than in larger, more institutionally fragmented countries.

\subsection{Latin American heterodoxy}

Several Latin American countries in the 2000s explicitly rejected IMF austerity programs: Argentina (after its 2001 default), Ecuador, Bolivia, and to some extent Brazil under Lula.
Their post-rejection economic performance was generally positive, undermining the claim that austerity is necessary for recovery.

These cases produced a regional counter-narrative infrastructure: ALBA (Bolivarian Alliance), the Banco del Sur, and academic networks around Latin American structuralism and dependency theory.
The counter-infrastructure was smaller than the pro-austerity infrastructure but was regionally effective.

\subsection{Modern Monetary Theory}

MMT has emerged as the most theoretically developed counter-narrative to fiscal morality:
\begin{itemize}[nosep]
  \item A sovereign currency issuer cannot ``run out'' of money in its own currency.
  \item The constraint on government spending is not financial (money) but real (inflation, resource availability).
  \item Government deficits are the private sector's surpluses (accounting identity).
  \item The household metaphor is therefore fundamentally misleading: a currency-issuing government is nothing like a household.
\end{itemize}

MMT's doxonomic challenge: despite logical coherence, it struggles against the household metaphor because it lacks an equally intuitive counter-metaphor.
``The government is the issuer of the currency, not a user of it'' is true but requires economic sophistication that ``families can't spend more than they earn'' does not.
The asymmetry in cognitive accessibility---not in evidential quality---explains the household metaphor's dominance.

\section{Media Framing of Fiscal Policy}
\label{sec:ch25-media}

\subsection{Content analysis evidence}

Studies of media framing of fiscal policy consistently find:
\begin{itemize}[nosep]
  \item deficit and debt stories outnumber stories about the benefits of public spending,
  \item the household metaphor is the dominant framing device,
  \item ``expert'' sources quoted on fiscal matters are disproportionately from institutions with pro-austerity positions,
  \item the distributional consequences of austerity are underreported relative to the aggregate fiscal position.
\end{itemize}

\begin{proposition}[Asymmetric Fiscal Framing]
\label{prop:asymmetric-framing}
\index{fiscal framing}
Media coverage of fiscal policy exhibits a systematic asymmetry: costs of government spending are framed as visible and immediate (``the deficit is \$X trillion''), while benefits are framed as diffuse and uncertain (``public investment may promote growth'').
Conversely, costs of spending cuts are framed as diffuse (``some programs may be reduced'') while fiscal benefits are framed as visible and immediate (``the deficit will shrink by \$X billion'').
This asymmetry favors the austerity narrative regardless of the actual evidence.
\end{proposition}

\section{Notes and References}
\label{sec:ch25-notes}

On austerity as ideology, see \textcite{streeck2016} and \textcite{brown2015}.
The shock doctrine is developed in \textcite{klein2007}.
The political economy of debt narratives connects to \textcite{graeber2011}.
For the history of neoliberal fiscal narratives, see \textcite{mirowski2009}, \textcite{harvey2005}, and \textcite{slobodian2018}.

On the Reinhart-Rogoff controversy, see Herndon, Ash, and Pollin (2014).
Blanchard and Leigh (2013) document the underestimation of fiscal multipliers.
On Iceland's alternative path, see Wade and Sigurgeirsd{\'o}ttir (2012).
On the Greek crisis, see Varoufakis (2017).

Modern Monetary Theory is developed in Wray (2015) and Kelton (2020).
On media framing of fiscal policy, see Blyth (2013).
System justification in the context of economic policy is analyzed in \textcite{jost2004}.
For the historical genealogy of fiscal morality, see \textcite{graeber2011} and \textcite{piketty2014}.

\section{Exercises}

\begin{exercise}
Construct a doxonomic genealogy of the phrase ``living beyond our means'' in political discourse.
\begin{enumerate}[nosep]
  \item When did it first appear in budget debates?
  \item What institutional channels popularized it?
  \item How does its frequency in media coverage correlate with economic conditions?
  \item Is the phrase used symmetrically (applied equally to tax cuts that increase deficits and spending that increases deficits)?
\end{enumerate}
\end{exercise}

\begin{exercise}
Using the shock doctrine model (Model~\ref{mod:shock-doctrine}), analyze the 2008 financial crisis and subsequent austerity turn.
\begin{enumerate}[nosep]
  \item Estimate the narrative entropy trajectory $H(t)$ qualitatively (sketch the curve from 2007 to 2013).
  \item Identify who supplied the dominant post-crisis frame and through what institutional channels.
  \item What counter-narratives were available and why did they fail to capture the frame?
  \item Apply the model to a second crisis (e.g., COVID-19). Did the shock doctrine operate similarly?
\end{enumerate}
\end{exercise}

\begin{exercise}
Design a survey to measure belief in the household-budget analogy for government spending.
\begin{enumerate}[nosep]
  \item Write five items that capture different aspects of fiscal morality belief.
  \item Predict which demographic groups will score highest and explain why (consider education, political affiliation, age, income).
  \item Include a treatment condition: half of respondents read a brief explanation of sovereign currency issuance before answering. Does it change responses?
\end{enumerate}
\end{exercise}

\begin{exercise}
Apply the academic credibility laundering model (Model~\ref{mod:credibility-laundering}) to a case other than Reinhart-Rogoff.
\begin{enumerate}[nosep]
  \item Identify the researcher, institution, finding, and amplification channel.
  \item Estimate the ``laundering window'' (time between amplification and correction).
  \item What policy decisions were made during the laundering window?
  \item Were those decisions reversed after the correction? If not, why not?
\end{enumerate}
\end{exercise}

\begin{exercise}
Conduct a content analysis of fiscal-policy coverage in two newspapers with different editorial stances.
\begin{enumerate}[nosep]
  \item Count the frequency of the household metaphor and its variants.
  \item Categorize expert sources by institutional affiliation.
  \item Assess whether costs and benefits of spending are framed symmetrically (Proposition~\ref{prop:asymmetric-framing}).
  \item What doxonomic explanation accounts for the patterns you find?
\end{enumerate}
\end{exercise}

\begin{exercise}
Compare the distributional consequences of austerity in Greece (Example~\ref{ex:greece}) with Iceland's alternative approach.
\begin{enumerate}[nosep]
  \item Tabulate key economic and social indicators for both countries (GDP change, unemployment, poverty, health outcomes) from 2008 to 2018.
  \item Who bore the costs in each case?
  \item Apply the legitimation framework (Chapter~\ref{ch:inequality-legitimacy}): why was the austerity narrative accepted in Greece but rejected in Iceland?
  \item What institutional differences explain the different narrative outcomes?
\end{enumerate}
\end{exercise}

\begin{exercise}[Computational]
Simulate the austerity belief dynamics model (Model~\ref{mod:austerity-dynamics}):
\begin{enumerate}[nosep]
  \item Set $\belief{A}(0) = 0.4$, $\alpha = 0.1$, $\phi_A = 3$, $\beta = 0.5$, $\gamma = 0.3$, $\delta = 0.05$.
  \item Simulate 100 periods with no crisis. Plot $\belief{A}(t)$.
  \item At $t = 30$, introduce a crisis ($\text{Crisis}(t) = 1$ for $t \in [30, 35]$). How does $\belief{A}$ respond?
  \item At $t = 50$, introduce rising unemployment ($\Delta U(t) = 0.5$ for $t \in [50, 70]$) as a consequence of the austerity policies adopted during the crisis. Does this erode austerity belief?
  \item Compare scenarios with different $\phi_A$ (narrative investment): how much investment is needed to sustain austerity belief despite rising unemployment?
\end{enumerate}
\end{exercise}

\input{chapters/ch27-work-discipline}
\input{chapters/ch28-nation-myth}
\input{chapters/ch29-technology-inevitability}
\input{chapters/ch30-crisis-doxonomics}

\begin{interlude}
The applications have shown what doxonomics can see: how institutional investment produces common sense in domain after domain.
But diagnosis without remedy is incomplete.
Part~V asks the forward-looking questions: how do belief regimes break, how can they be challenged, what would a more democratic belief ecology look like, what ethical constraints bind the field, and what remains to be discovered?
These chapters are the most speculative and the most important.
\end{interlude}

% --- Part V: Resistance, Repair, and Reconstruction ---
\part{Resistance, Repair, and Reconstruction}
\label{part:resistance}

\chapter{How Belief Regimes Fracture}
\label{ch:how-belief-regimes-fracture}

\epigraph{There are decades where nothing happens; and there are weeks where decades happen.}{Vladimir Lenin (attr.)}

\noindent
Doxonomic equilibria are stable but not permanent.
The preceding chapters analyzed how belief regimes are constructed, funded, and maintained.
This chapter analyzes how they come apart: the mechanisms by which dominant belief regimes fracture, the dynamics of credibility collapse, the role of contradictions between narrative and lived experience, and the conditions under which fracture produces transformation rather than mere turbulence.

Understanding fracture is essential to doxonomics for two reasons.
First, it completes the theory: a framework that can explain regime maintenance but not regime change is fundamentally incomplete.
Second, it has practical implications: actors seeking to challenge a dominant belief regime need to understand when and how regimes are vulnerable, just as actors seeking to defend a regime need to understand the sources of fragility.

\section{Regime Fracture as Bifurcation}
\label{sec:ch30-bifurcation}

\subsection{The formal model}

\begin{model}[Bifurcation Model of Regime Fracture]
\label{mod:bifurcation}
\index{bifurcation}
Let $\belief{}(t)$ be the strength of the dominant belief and $\sigma(t)$ a stress parameter (accumulating contradictions, declining legitimacy, external shocks).
The dynamics are:
\[
  \frac{d\belief{}}{dt} = \belief{}(\alpha\phi - \delta - \sigma(t)) + \beta\belief{}^2(1 - \belief{})
\]
where $\alpha\phi$ is the institutional production rate, $\delta$ is natural decay, and $\beta$ captures social reinforcement.
\end{model}

\subsection{The critical stress threshold}

\begin{proposition}[Critical Stress for Regime Fracture]
\label{prop:critical-stress}
\index{critical stress}
The high-$\belief{}$ equilibrium exists if and only if $\sigma < \sigma^*$, where:
\[
  \sigma^* = \alpha\phi - \delta + \frac{\beta}{4}.
\]
When $\sigma$ exceeds $\sigma^*$, the high-belief equilibrium disappears through a saddle-node bifurcation.
The system drops to a low-$\belief{}$ state: regime fracture.

The critical stress $\sigma^*$ increases with:
\begin{itemize}[nosep]
  \item institutional investment $\phi$ (more funding makes the regime more resilient),
  \item social reinforcement $\beta$ (stronger peer effects make the regime harder to crack),
  \item conversion efficiency $\alpha$ (more effective institutions are more resilient).
\end{itemize}
It decreases with natural decay $\delta$ (regimes with faster turnover are more fragile).
\end{proposition}

\subsection{The fragility-resilience spectrum}

\begin{example}
\label{ex:fragility-spectrum}
Different belief regimes occupy different positions on the fragility-resilience spectrum:

\begin{center}
\begin{tabular}{p{3.5cm}ccp{4cm}}
\toprule
Regime & $\alpha\phi$ (investment) & $\beta$ (reinforcement) & Assessment \\
\midrule
``Markets are efficient'' & High & Moderate & Fractured partially in 2008; recovered due to high $\phi$ \\[4pt]
``Smoking doesn't cause cancer'' & Moderate & Low & Fractured 1960s--90s as evidence accumulated \\[4pt]
National identity & Very high & Very high & Highly resilient; fractures only in existential crises \\[4pt]
``Human nature is selfish'' & High & High & Very resilient; self-fulfilling properties add extra stability \\[4pt]
Soviet communism & Moderate & High (coerced) & Fractured 1989--91 when coercion was withdrawn \\
\bottomrule
\end{tabular}
\end{center}
\end{example}

\section{Sources of Stress}
\label{sec:ch30-stress}

The stress parameter $\sigma(t)$ aggregates multiple distinct sources of strain on the belief regime.

\subsection{Contradiction accumulation}

\begin{definition}[Doxonomic Contradiction]
\label{def:contradiction}
\index{contradiction!doxonomic}
A \emph{doxonomic contradiction} arises when the belief system predicts $X$ but reality delivers $\neg X$ in a publicly visible way.
Each contradiction adds $\Delta\sigma > 0$ to the stress parameter.
Small contradictions can be absorbed through rationalization, reinterpretation, or blame-shifting.
Large or repeated contradictions accumulate until $\sigma$ approaches $\sigma^*$.
\end{definition}

\begin{example}
\label{ex:contradictions}
Contradictions that accumulated against the ``efficient markets'' regime:
\begin{itemize}[nosep]
  \item Long-Term Capital Management collapse (1998): $\Delta\sigma$ moderate, contained by bailout.
  \item Dot-com crash (2000--01): $\Delta\sigma$ moderate, rationalized as ``irrational exuberance'' (a market failure, but a correctable one).
  \item Enron/WorldCom scandals (2001--02): $\Delta\sigma$ moderate, framed as individual fraud rather than systemic failure.
  \item 2008 financial crisis: $\Delta\sigma$ very large, exceeding the rationalization capacity. Regime fractured partially.
\end{itemize}
Note the pattern: each contradiction was individually absorbed through narrative repair, but the cumulative $\sigma$ brought the system close to $\sigma^*$.
The 2008 crisis pushed it over the threshold, not because it was qualitatively different from previous crises, but because the accumulated stress left no remaining margin.
\end{example}

\subsection{Credibility collapse}

\begin{definition}[Credibility Collapse]
\label{def:credibility-collapse}
\index{credibility collapse}
A \emph{credibility collapse} occurs when the credibility $\credibility{k}$ of a central institution drops rapidly, typically due to a scandal, failure, or exposure of deception.
Because credibility is networked (institutions lend credibility to each other through association), a single collapse can cascade through the institutional network.
\end{definition}

\begin{model}[Credibility Cascade]
\label{mod:credibility-cascade}
\index{credibility cascade}
Let $\credibility{k}(t)$ be the credibility of institution $k$, and let $A_{kj}$ be the association matrix ($A_{kj} > 0$ if institutions $k$ and $j$ are perceived as connected: shared funders, shared personnel, shared narrative space).
When institution $k_0$ suffers a credibility shock:
\[
  \credibility{k}(t+1) = \credibility{k}(t) - \gamma \sum_{j} A_{kj} \cdot \Delta\credibility{j}(t)
\]
where $\Delta\credibility{j}(t)$ is the credibility loss at $j$ in the previous step.
If $\gamma$ is large (strong association effects), a single institutional failure can trigger a credibility cascade that undermines the entire network of institutions supporting the belief regime.
\end{model}

\begin{example}
\label{ex:credibility-cascade}
The Catholic Church sex abuse scandal illustrates a credibility cascade:
\begin{itemize}[nosep]
  \item The initial revelation (Boston Globe, 2002) destroyed the credibility of specific dioceses.
  \item The cascade spread to the institutional Church as a whole (shared identity).
  \item It damaged the credibility of \emph{moral authority} in general (if the Church was complicit, who can be trusted on moral questions?).
  \item It weakened the doxonomic output of religious institutions more broadly (survey data shows declining trust in all organized religion post-2002, not just Catholicism).
\end{itemize}
The cascade followed the association matrix: institutions perceived as closer to the Church (other religious bodies, moral-authority organizations) suffered larger credibility losses than distant institutions.
\end{example}

\subsection{Economic shocks}

Material conditions constrain belief.
When economic conditions deteriorate sharply (unemployment rises, living standards fall, inequality becomes visible), the gap between narrative (``the economy works for everyone'') and experience (``I just lost my job/house/savings'') generates stress.

\begin{proposition}[Material-Narrative Gap]
\label{prop:material-narrative-gap}
\index{material-narrative gap}
The stress contribution from economic conditions is:
\[
  \sigma_{\text{econ}}(t) = \kappa \cdot |Y_{\text{narrative}}(t) - Y_{\text{actual}}(t)|
\]
where $Y_{\text{narrative}}$ is the prosperity level implied by the dominant belief regime and $Y_{\text{actual}}$ is the experienced reality.
Regimes that promise more (``rising tide lifts all boats'') are more vulnerable to economic downturns than regimes with more modest claims (``life is hard but fair'').
\end{proposition}

\subsection{Generational turnover}

Each generation arrives without the formative experiences that established the belief regime for their parents.
The beliefs must be actively reproduced through education and socialization: they are not transmitted automatically.

\begin{proposition}[Generational Decay]
\label{prop:generational-decay}
\index{generational decay}
If the belief regime's institutional reproduction is imperfect (transmission rate $\tau < 1$ per generation), then belief strength decays geometrically across generations:
\[
  \belief{}^{(n)} = \tau^n \cdot \belief{}^{(0)}.
\]
For $\tau = 0.85$ (15\% loss per generation), the regime retains only $0.85^3 = 0.61$ of its original strength after three generations ($\sim$75 years).
This decay is independent of external stress and represents the ``natural lifespan'' of a belief regime in the absence of continuous institutional reinvestment.
\end{proposition}

\subsection{Memetic overload}

The information environment can become so saturated with competing narratives that no single frame achieves dominance.
This raises entropy $H(t)$ and erodes the coherence of the dominant regime even without any specific contradiction.

The digital media environment has accelerated memetic overload: the proliferation of information sources, the algorithmic amplification of controversy, and the collapse of shared informational baselines all contribute to rising entropy.
Some analysts interpret the current political moment as a \emph{permanent memetic overload}: a condition in which no stable common sense can be maintained because the information environment is too fragmented.

\section{Fracture Dynamics}
\label{sec:ch30-dynamics}

\subsection{Slow accumulation, fast collapse}

Regime fracture typically follows a pattern of slow stress accumulation followed by rapid collapse.
The bifurcation model explains why: as $\sigma$ approaches $\sigma^*$, the high-belief equilibrium becomes increasingly fragile (small perturbations produce large swings), but the equilibrium persists until the critical threshold is crossed.
The collapse, when it comes, is rapid and discontinuous.

This pattern (decades of apparent stability followed by sudden, surprising collapse) is familiar from the fall of the Soviet Union, the Arab Spring, the \#MeToo movement, and other regime fractures.
In each case, observers were surprised by the speed of change, but the bifurcation model predicts exactly this pattern: the regime appeared stable because the equilibrium existed, but it was increasingly fragile because $\sigma$ was close to $\sigma^*$.

\subsection{Partial vs.\ complete fracture}

Not all regime fractures are complete.
The 2008 crisis produced a \emph{partial} fracture of the efficient markets regime: the belief was damaged but not destroyed, and institutional investment ($\phi$) proved sufficient to repair much of the damage over the following decade.

\begin{definition}[Partial Fracture]
\label{def:partial-fracture}
\index{partial fracture}
A \emph{partial fracture} occurs when $\sigma$ temporarily exceeds $\sigma^*$ but then falls back below it (because the stress source was temporary), and institutional investment is sufficient to rebuild the high-belief equilibrium before the system fully transitions to the low-belief state.
The regime survives but is weakened: $\belief{}^{\text{post}} < \belief{}^{\text{pre}}$.
\end{definition}

Complete fracture occurs when either the stress is sustained long enough for the system to fully transition, or when the stress destroys the institutional infrastructure ($\phi$ declines), making the high-belief equilibrium unrecoverable.

\subsection{What fills the vacuum?}

Regime fracture creates a narrative vacuum (cf.\ Chapter~\ref{ch:crisis-doxonomics}).
What fills the vacuum determines whether fracture produces progressive transformation, reactionary backlash, or prolonged instability.

\begin{remark}[Post-Fracture Outcomes]
\label{prop:post-fracture}
\index{post-fracture outcomes}
The outcome of regime fracture depends on the relative preparedness of competing narratives:
\begin{enumerate}[nosep]
  \item \textbf{Regime replacement}: a prepared alternative with institutional infrastructure captures the vacuum. Example: the New Deal replacing laissez-faire after 1929.
  \item \textbf{Regime restoration}: the old regime's institutional infrastructure survives the fracture and rebuilds. Example: efficient markets partially restored after 2008.
  \item \textbf{Prolonged instability}: no alternative achieves capture; the system oscillates between competing frames. Example: the post-2016 political landscape in many Western democracies.
  \item \textbf{Authoritarian capture}: a crisis narrative built on fear, identity, and simplification captures the vacuum. Example: Weimar Germany, various 21st-century populist movements.
\end{enumerate}
This typology is descriptive, not predictive: it organizes historical cases but does not (yet) predict which outcome will occur in a given crisis.
\end{remark}

\section{Historical Case Studies}
\label{sec:ch30-cases}

\subsection{The fall of Soviet communism (1989--91)}

The Soviet belief regime fractured through a combination of all five stress sources operating simultaneously:
\begin{itemize}[nosep]
  \item \emph{Contradictions}: the economy failed to deliver the promised prosperity; Soviet citizens could see Western living standards through increasing media exposure.
  \item \emph{Credibility collapse}: Chernobyl (1986) demonstrated that the state would lie about existential threats.
  \item \emph{Economic shock}: declining oil prices and economic stagnation in the 1980s.
  \item \emph{Memetic overload}: glasnost opened the information space to competing narratives.
  \item \emph{Generational turnover}: young Soviets had no memory of WWII, which had been the regime's primary legitimating experience.
\end{itemize}

The critical factor was that Gorbachev's reforms \emph{withdrew coercive enforcement} without building an alternative belief infrastructure.
The social reinforcement term $\beta$ dropped sharply (people could now express dissent), and the institutional investment $\phi$ was not redirected to a new narrative.
The result was complete regime fracture within two years.

\subsection{The \#MeToo fracture (2017--)}

The gender regime (Chapter~\ref{ch:race-coloniality-gender}) experienced a rapid partial fracture beginning in October 2017:
\begin{itemize}[nosep]
  \item \emph{Trigger}: the Harvey Weinstein revelations provided a high-profile credibility shock.
  \item \emph{Cascade}: social media enabled rapid aggregation of individual experiences into a visible pattern, bypassing institutional gatekeepers who had previously suppressed the narrative.
  \item \emph{Norm shift}: the disclosure norm shifted from ``don't talk about it'' to ``speak up,'' enabling thousands of additional disclosures.
  \item \emph{Institutional response}: companies, universities, and governments adopted new policies, institutionalizing the partial fracture.
\end{itemize}

The fracture was partial: behavioral norms in many workplaces shifted, but the deeper gender regime (Chapter~\ref{ch:race-coloniality-gender}) remained substantially intact.
The digital infrastructure that enabled the cascade also enabled counter-narratives (``cancel culture,'' ``due process concerns'') that partially contained the fracture.

\section{Regime Resilience and Repair}
\label{sec:ch30-resilience}

\subsection{Repair mechanisms}

Regimes that survive partial fracture employ several repair mechanisms:

\begin{enumerate}[nosep]
  \item \textbf{Narrative absorption}: incorporate the critique into the regime without changing its core. Example: ``We need \emph{better} regulation, not more government'' absorbs the 2008 critique without challenging the market-primacy narrative.

  \item \textbf{Scapegoating}: attribute the failure to specific bad actors rather than systemic problems. Example: ``A few bad apples on Wall Street'' preserves the system while sacrificing individuals.

  \item \textbf{Amnesia acceleration}: speed the forgetting of the fracture event. Institutional media moves on to new stories; the crisis becomes ``old news.''

  \item \textbf{Credibility restoration}: rebuild institutional credibility through reforms, apologies, or personnel changes that signal change without altering the underlying structure.

  \item \textbf{Counter-narrative suppression}: use institutional power to marginalize or delegitimize the alternative narratives that emerged during the fracture.
\end{enumerate}

\subsection{The resilience paradox}

\begin{remark}[Resilience Paradox]
\label{prop:resilience-paradox}
\index{resilience paradox}
Regime repair that is too successful can increase long-run fragility.
If the regime absorbs a crisis without meaningful reform, the underlying contradictions that produced the stress are not resolved: they are merely papered over.
The stress parameter $\sigma$ resumes its accumulation, and the next crisis finds a regime that is closer to $\sigma^*$ than before, with repair mechanisms that have been partially depleted.
Successful short-run repair can produce a cycle of escalating crises and increasingly brittle repairs, culminating in a fracture that no repair can contain.
\end{remark}

\section{Notes and References}
\label{sec:ch30-notes}

On the dynamics of regime change, see \textcite{gramsci1971} on interregnum and \textcite{polanyi1944} on the double movement.
Threshold and cascade dynamics are in \textcite{granovetter1978} and \textcite{watts2002}.
On economic shocks and ideological change, see \textcite{klein2007}.

Bifurcation theory for social systems draws on Scheffer (2009) and Strogatz (2015).
On the fall of Soviet communism as a belief-system collapse, see Kotkin (2001) and Yurchak (2005).
The \#MeToo movement's dynamics are analyzed in Fileborn and Loney-Howes (2019).

On credibility cascades and institutional trust, see Slovic (1993) and Cvetkovich and L\"{o}fstedt (1999).
Generational effects on belief are studied in Mannheim (1952) and Inglehart (1977).
For memetic overload and information fragmentation, see \textcite{sunstein2017} and Benkler, Faris, and Roberts (2018).

\section{Exercises}

\begin{exercise}
For the bifurcation model (Model~\ref{mod:bifurcation}), compute $\sigma^*$ as a function of $\alpha$, $\phi$, $\delta$, and $\beta$.
\begin{enumerate}[nosep]
  \item What makes a regime more resilient to stress? Interpret each parameter's contribution.
  \item Compare the resilience of the selfishness regime ($\alpha\phi$ high, $\beta$ high) vs.\ a specific scientific consensus ($\alpha\phi$ moderate, $\beta$ low).
  \item How much additional institutional investment $\Delta\phi$ would be needed to survive a stress increase of $\Delta\sigma = 0.5$?
\end{enumerate}
\end{exercise}

\begin{exercise}
Analyze the 2008 financial crisis as a partial fracture of the ``efficient markets'' regime.
\begin{enumerate}[nosep]
  \item Identify the specific contradictions ($\Delta\sigma$ events) that accumulated prior to 2008.
  \item Why was the fracture incomplete? Which repair mechanisms (Section~\ref{sec:ch30-resilience}) were deployed?
  \item Apply the resilience paradox (Proposition~\ref{prop:resilience-paradox}): has the repair made the regime more or less fragile going forward?
  \item What would complete fracture look like, and what stress level would be required?
\end{enumerate}
\end{exercise}

\begin{exercise}
Model a credibility cascade (Model~\ref{mod:credibility-cascade}) using a network of 10 institutions.
\begin{enumerate}[nosep]
  \item Define the association matrix $A_{kj}$ based on shared funding, shared personnel, or shared narrative space.
  \item Shock one institution ($\Delta\credibility{k_0} = -0.5$) and propagate the cascade.
  \item How many institutions lose more than 20\% credibility? Does the cascade reach institutions not directly associated with $k_0$?
  \item What network structure (hub-and-spoke vs.\ distributed) is more vulnerable to cascading credibility collapse?
\end{enumerate}
\end{exercise}

\begin{exercise}
Apply the generational decay model (Proposition~\ref{prop:generational-decay}) to a belief regime of your choice.
\begin{enumerate}[nosep]
  \item Estimate the transmission rate $\tau$ from survey data comparing cohorts.
  \item Compute the regime's ``natural lifespan'' (the number of generations until $\belief{}$ falls below 0.3).
  \item How much institutional investment $\phi$ is needed to offset generational decay?
  \item Compare with the actual investment in the regime's institutional reproduction (education budgets, media, etc.).
\end{enumerate}
\end{exercise}

\begin{exercise}
Analyze the \#MeToo movement as a partial regime fracture.
\begin{enumerate}[nosep]
  \item What was the stress accumulation ($\sigma(t)$) prior to October 2017?
  \item What role did social media play in bypassing institutional gatekeepers?
  \item Was the fracture captured by a prepared alternative, or did it produce prolonged instability?
  \item Apply the partial fracture definition: what aspects of the gender regime fractured, and what aspects survived intact?
\end{enumerate}
\end{exercise}

\begin{exercise}[Computational]
Simulate the bifurcation model of regime fracture:
\begin{enumerate}[nosep]
  \item Set $\alpha\phi = 0.3$, $\delta = 0.1$, $\beta = 0.5$, $\belief{}(0) = 0.8$.
  \item Slowly increase $\sigma$ from 0 to 1 over 500 time steps. Plot $\belief{}(t)$.
  \item Identify the critical $\sigma^*$ at which fracture occurs.
  \item Repeat with $\beta = 0.3$ (weaker social reinforcement). How does $\sigma^*$ change?
  \item Simulate a partial fracture: increase $\sigma$ above $\sigma^*$ for 20 time steps, then reduce it. Does the regime recover?
  \item How long must $\sigma > \sigma^*$ persist for the fracture to become irreversible?
\end{enumerate}
\end{exercise}

\chapter{Counter-Doxonomic Strategies}
\label{ch:counter-doxonomic}

\epigraph{The master's tools will never dismantle the master's house.}{Audre Lorde}

\noindent
If beliefs can be produced, they can also be contested.
The preceding chapters documented how material power shapes belief through institutional investment.
This chapter asks the practical question: what can be done about it?
How can movements with limited resources challenge belief regimes backed by billions of dollars?

The answer is not obvious.
Lorde's epigraph warns against mimicking the oppressor's methods.
But doxonomic analysis suggests that the ``master's tools'' (institutional infrastructure, narrative production, strategic communication) are not inherently tools of domination; they are tools of belief formation, and they can be used by anyone who builds them.
The question is not whether to use institutional infrastructure but \emph{how} to use it with vastly fewer resources.

\section{The Counter-Doxonomic Challenge}
\label{sec:ch31-challenge}

\subsection{The fundamental asymmetry}

Counter-doxonomic movements face a stark resource asymmetry.
The regimes they challenge (selfishness, meritocracy, consumerism, austerity) are supported by institutional budgets in the billions.
Counter-movements typically operate at budgets in the millions.
The ratio is roughly 100:1 to 1000:1.

\begin{proposition}[Resource Asymmetry]
\label{prop:resource-asymmetry}
\index{resource asymmetry}
If the dominant regime invests $F_D$ and the counter-movement invests $F_C$, with $F_C \ll F_D$, then direct competition on the same channels (mass media, advertising, lobbying) is a losing strategy.
The counter-movement's belief share under symmetric competition is:
\[
  \belief{C}^* \approx \frac{F_C}{F_D + F_C} \approx \frac{F_C}{F_D} \ll 1.
\]
At a 100:1 funding ratio, the counter-movement captures roughly 1\% of the belief space through symmetric competition: negligible.
Effective counter-doxonomics requires \emph{asymmetric} strategies that exploit the dominant regime's vulnerabilities rather than competing on its strengths.
\end{proposition}

\subsection{Sources of asymmetric advantage}

Despite the funding gap, counter-movements have several potential advantages:

\begin{enumerate}
  \item \textbf{Truth.}
  If the dominant regime's beliefs are empirically false or significantly distorted (as the preceding chapters have argued for several regimes), the counter-movement has a correspondence advantage: its claims align better with lived experience.
  This advantage grows over time as contradictions accumulate (Chapter~\ref{ch:how-belief-regimes-fracture}).

  \item \textbf{Authenticity.}
  Counter-movements often speak from lived experience rather than institutional script.
  A worker describing precarity is more credible (to other workers) than a think-tank report justifying it.
  Authenticity provides a credibility premium per dollar of investment.

  \item \textbf{Motivation.}
  Participants in counter-movements are often motivated by conviction rather than compensation, providing volunteer labor that reduces the effective cost of narrative production.

  \item \textbf{Network effects.}
  Counter-narratives that resonate can spread virally through social networks at near-zero marginal cost, a channel that was unavailable before digital media and that partially offsets the mass-media funding disadvantage.

  \item \textbf{Regime fragility.}
  Dominant regimes have specific vulnerabilities (contradictions, credibility weaknesses, legitimation costs that grow with inequality) that targeted counter-doxonomic strategies can exploit.
\end{enumerate}

\section{Alternative Institutions}
\label{sec:ch31-alternatives}

\subsection{Building parallel infrastructure}

\begin{definition}[Counter-Doxonomic Institution]
\label{def:counter-institution}
\index{counter-doxonomics}
A \emph{counter-doxonomic institution} is an organization that produces narratives opposing the dominant regime, operating through alternative channels with independent funding.
Examples: labor unions, cooperative enterprises, community media, alternative schools, progressive think tanks, mutual-aid networks.
\end{definition}

The most successful counter-doxonomic movements in history built parallel institutional infrastructure rather than relying solely on critique:

\begin{example}
\label{ex:counter-institutions}
\begin{center}
\begin{tabular}{p{3cm}p{4cm}p{4.5cm}}
\toprule
Movement & Counter-institutions built & Doxonomic output \\
\midrule
Labor movement (1880--1950) & Unions, labor press, workers' education, cooperatives & ``Labor has rights''; ``solidarity, not competition'' \\[4pt]
Civil rights (1950--70) & NAACP, SCLC, Freedom Schools, Black newspapers & ``Segregation is unjust''; ``equality under law'' \\[4pt]
Feminism (1960--) & Women's studies programs, feminist media, shelters, legal organizations & ``Personal is political''; ``equal pay for equal work'' \\[4pt]
Environmental (1970--) & Environmental organizations, green parties, research institutes & ``Pollution has costs''; ``sustainability'' \\
\bottomrule
\end{tabular}
\end{center}
Each movement's lasting impact depended on the institutional infrastructure it built, not just the protests it organized or the critiques it published.
Protest creates attention; institutions sustain narrative over time.
\end{example}

\subsection{Institution design principles}

\begin{model}[Counter-Institution Design]
\label{mod:counter-institution}
\index{counter-institution design}
Effective counter-doxonomic institutions should:
\begin{enumerate}[nosep]
  \item \textbf{Produce, not just critique}: generate positive alternative narratives, not only refutations of the dominant regime.
  \item \textbf{Build credibility}: accumulate epistemic capital through rigorous research, transparent methods, and consistent accuracy.
  \item \textbf{Diversify funding}: avoid dependence on any single funder whose withdrawal could collapse the institution.
  \item \textbf{Invest in persistence}: design for multi-generational operation, not just campaign cycles.
  \item \textbf{Target the production chain}: intervene at the institutional stages of belief production (education, professional norms, policy design), not just at the media stage.
\end{enumerate}
\end{model}

\section{Pedagogical Interventions}
\label{sec:ch31-pedagogy}

\subsection{Critical pedagogy}

\textcite{freire1968} developed ``critical pedagogy'': education designed to reveal the social construction of common sense.
In doxonomic terms, Freire's method attacks the internalization stage of the belief production chain by making the chain itself visible.
When students learn to ask ``who produced this belief, through what institutions, with what funding, and serving whose interests?,'' they are performing doxonomic analysis, even if they have never heard the term.

\subsection{Media literacy as counter-doxonomics}

\begin{model}[Media Literacy Intervention]
\label{mod:media-literacy}
\index{media literacy}
Media literacy education reduces the susceptibility parameter $\alpha$ in the belief dynamics model:
\[
  \alpha^{\text{post-literacy}} = \alpha \cdot (1 - \ell)
\]
where $\ell \in [0, 1]$ is the ``literacy shield'': the fraction of doxonomic influence that the individual can recognize and resist after training.
If $\ell = 0.3$ (a 30\% reduction in susceptibility), the effective doxonomic investment becomes $0.7 \cdot \phi$, equivalent to cutting the dominant regime's budget by 30\% without spending a dollar on counter-narrative production.
\end{model}

This makes media literacy one of the highest-IROI counter-doxonomic investments available: rather than competing dollar-for-dollar in narrative production, it reduces the effectiveness of the opponent's spending.
The analogy is to vaccination rather than treatment: it is cheaper to prevent belief capture than to reverse it.

\subsection{Doxonomic literacy}

Beyond general media literacy, \emph{doxonomic literacy} is the specific ability to:
\begin{enumerate}[nosep]
  \item identify the funding source behind a narrative,
  \item trace the institutional production chain,
  \item recognize the credibility-laundering mechanisms (academic authority, expert panels, ``independent'' research),
  \item distinguish between beliefs held because of evidence and beliefs held because of institutional repetition,
  \item recognize one's own susceptibility to doxonomic influence.
\end{enumerate}

This textbook is, among other things, an exercise in doxonomic literacy.

\section{Narrative Inoculation}
\label{sec:ch31-inoculation}

\subsection{The inoculation model}

\begin{definition}[Narrative Inoculation]
\label{def:inoculation}
\index{narrative inoculation}
\emph{Narrative inoculation} is preemptive exposure to a weakened form of a dominant narrative, accompanied by refutation, so that subsequent full-strength exposure is less effective.
Formally, it reduces the susceptibility parameter $\alpha$ in the SIR model (Model~\ref{mod:sir-belief}).
\end{definition}

The analogy to biological vaccination is precise: a weakened version of the ``pathogen'' (the misleading narrative) is introduced, along with ``antibodies'' (refutation and critical analysis), so that when the full-strength narrative is encountered, the individual has pre-existing resistance \autocite{mcguire1964, vanderlinden2022}.

\subsection{Practical inoculation design}

\begin{model}[Inoculation Program]
\label{mod:inoculation-program}
\index{inoculation program}
An inoculation intervention consists of:
\begin{enumerate}[nosep]
  \item \textbf{Warning}: alert the audience that they will encounter persuasion attempts (``some people will try to convince you that\ldots'').
  \item \textbf{Weakened exposure}: present the core argument of the misleading narrative in a simplified or clearly attributed form.
  \item \textbf{Refutation}: provide the counter-arguments, evidence, and critical analysis that reveal the narrative's weaknesses.
  \item \textbf{Practice}: give the audience opportunities to practice identifying and refuting the narrative in new contexts.
\end{enumerate}
Studies show that inoculation can reduce susceptibility to misinformation by 20--30\% and that the effect persists for months \autocite{vanderlinden2022}.
\end{model}

\begin{example}
\label{ex:inoculation}
An inoculation program against the household-budget metaphor for government spending (Chapter~\ref{ch:austerity}):
\begin{enumerate}[nosep]
  \item \emph{Warning}: ``Politicians often compare the government budget to a household budget. This comparison sounds intuitive but is misleading.''
  \item \emph{Weakened exposure}: ``The argument: `Just like families, governments can't spend more than they earn.' ''
  \item \emph{Refutation}: ``Unlike families, governments that issue their own currency can create money; their spending creates income for the private sector; and they never need to `pay back' debt the way a household does. The analogy confuses a currency user (household) with a currency issuer (sovereign government).''
  \item \emph{Practice}: present three variants of the household metaphor (``maxed out credit card,'' ``belt-tightening,'' ``living beyond our means'') and ask participants to identify and refute each.
\end{enumerate}
\end{example}

\section{Optimal Allocation of Counter-Resources}
\label{sec:ch31-optimal}

\subsection{The optimization problem}

\begin{model}[Counter-Doxonomic Optimization]
\label{mod:counter-optimal}
\index{counter-doxonomics!optimization}
Given a total counter-doxonomic budget $F_C \ll F_D$, the counter-movement solves:
\[
  \max_{\fundingvec_C} \; \Delta\belief{C} = \sum_k \funding{k}^C \cdot \credibility{k}^C \cdot \reach{k}^C \cdot \persistence{k}^C
\]
subject to $\sum_k \funding{k}^C = F_C$.
The optimal allocation equates the marginal IROI across all channels:
\[
  \frac{\partial \Delta\belief{C}}{\partial \funding{k}^C} = \lambda \quad \text{for all active channels } k
\]
where $\lambda$ is the shadow price of the budget constraint.
\end{model}

\subsection{Channel ranking}

The IROI-per-dollar typically ranks channels as follows:

\begin{center}
\begin{tabular}{lcccc}
\toprule
Channel & Credibility & Reach & Persistence & IROI/\$ \\
\midrule
Community organizing & High & Low & Very high & Highest \\
Education/curricula & High & Moderate & Very high & Very high \\
Investigative journalism & High & Moderate & High & High \\
Social media campaigns & Low & Very high & Low & Moderate \\
Mass advertising & Low & Very high & Low & Lowest \\
\bottomrule
\end{tabular}
\end{center}

The optimal strategy for a resource-constrained counter-movement typically concentrates on high-credibility, high-persistence channels (community organizing, education) rather than competing on reach with the dominant regime's mass media.

\subsection{The credibility-reach tradeoff}

\begin{proposition}[Asymmetric Warfare in Belief Space]
\label{prop:asymmetric}
\index{asymmetric warfare}
Counter-doxonomic movements are most effective when they:
\begin{enumerate}[nosep]
  \item target the \emph{weakest link} in the dominant belief production chain (the point where the narrative is most vulnerable to contradiction or exposure),
  \item exploit credibility asymmetries (authentic lived experience vs.\ corporate messaging; independent research vs.\ industry-funded studies),
  \item invest in persistence over reach (deep education that changes schemas over years vs.\ viral campaigns that shift attention for days),
  \item build parallel institutional infrastructure rather than only critiquing existing infrastructure,
  \item time interventions to exploit regime fractures (Chapter~\ref{ch:how-belief-regimes-fracture}) and crisis windows (Chapter~\ref{ch:crisis-doxonomics}).
\end{enumerate}
\end{proposition}

\section{Transparency and Disclosure}
\label{sec:ch31-transparency}

\subsection{Funding transparency as counter-doxonomic tool}

One of the most cost-effective counter-doxonomic strategies is not producing counter-narratives but \emph{revealing the production chain of the dominant narrative}.
If the public knows who funds a think tank, who sponsors a study, or who owns a media outlet, the credibility of the narrative is automatically discounted.

\begin{proposition}[Disclosure Effect]
\label{prop:disclosure}
\index{funding disclosure}
If the funding source of institution $k$ is disclosed, its effective credibility becomes:
\[
  \credibility{k}^{\text{disclosed}} = \credibility{k} \cdot (1 - d \cdot \text{conflict}(k))
\]
where $d \in [0, 1]$ is the disclosure-awareness rate (what fraction of the audience learns the funding information) and $\text{conflict}(k) \in [0, 1]$ measures the perceived conflict of interest.
For a fossil-fuel-funded climate denial institute with $\credibility{k} = 0.6$, $d = 0.5$, and $\text{conflict} = 0.8$:
\[
  \credibility{k}^{\text{disclosed}} = 0.6 \times (1 - 0.4) = 0.36.
\]
Disclosure reduces effective credibility by 40\%, achieved not by producing counter-narrative but by revealing the production chain of the existing one.
\end{proposition}

\subsection{Institutional transparency requirements}

Policy-level counter-doxonomic strategies include:
\begin{itemize}[nosep]
  \item mandatory disclosure of think-tank funding sources,
  \item academic conflict-of-interest disclosure requirements,
  \item media ownership transparency laws,
  \item political advertising source disclosure,
  \item algorithmic transparency requirements for platforms.
\end{itemize}

Each makes the doxonomic production chain visible, reducing its effectiveness by the disclosure effect.

\section{Coalition and Scale}
\label{sec:ch31-coalition}

\subsection{The coordination problem}

Counter-doxonomic movements are often fragmented: environmental groups, labor unions, racial justice organizations, feminist movements, and anti-poverty groups each challenge different aspects of the dominant regime but rarely coordinate their narrative production.

\begin{proposition}[Counter-Doxonomic Coordination Gains]
\label{prop:coordination}
\index{coordination!counter-doxonomic}
If $n$ counter-movements each invest $f$ independently, the total belief shift is $\Delta\belief{} = n \cdot g(f)$ where $g$ is concave (diminishing returns within each channel).
If they coordinate and pool resources $F = nf$ into a unified counter-narrative, the total shift is $\Delta\belief{} = g(nf) > n \cdot g(f)$ by concavity.
Coordination produces more belief shift per dollar than fragmented investment.
\end{proposition}

The practical challenge: coordination requires agreement on a shared counter-narrative, which is difficult when movements have different priorities and constituencies.
The most successful coordinating frameworks (e.g., the Green New Deal, which linked climate, labor, and racial justice) succeed by finding a narrative that encompasses multiple movements without subordinating any.

\section{Notes and References}
\label{sec:ch31-notes}

Critical pedagogy draws on \textcite{freire1968} and \textcite{illich1971}.
On alternative media and networked public spheres, see \textcite{benkler2006}.
Counter-hegemonic strategy is theorized in \textcite{gramsci1971}.
Narrative inoculation is studied in \textcite{mcguire1964} and \textcite{vanderlinden2022}.
For real utopian institution-building, see \textcite{wright2010}.

On the history of counter-doxonomic movements, see the labor movement analysis in Foner (1947), the civil rights movement in Morris (1984), and feminist institution-building in Ferree and Mueller (2004).
Media literacy effectiveness is reviewed in Jeong et al.\ (2012).
Funding transparency and its effects are analyzed in \textcite{mayer2016}.

For the optimization of social-movement resources, see McCarthy and Zald (1977) on resource mobilization theory.
On movement coordination, see Tarrow (2011) and McAdam, Tarrow, and Tilly (2001).

\section{Exercises}

\begin{exercise}
Given a counter-doxonomic budget of \$5 million/year, design an allocation strategy across education, media, community organizing, investigative journalism, and inoculation programs to shift belief in selfish human nature.
\begin{enumerate}[nosep]
  \item Justify your allocation using the IROI framework and the channel ranking table.
  \item What percentage of the budget goes to each channel?
  \item How does your allocation change if the budget doubles to \$10M? Does it scale linearly?
  \item Estimate the expected belief shift $\Delta\belief{C}$ under your allocation.
\end{enumerate}
\end{exercise}

\begin{exercise}
Design a narrative inoculation program for high school students to build resistance to advertising-driven consumerism.
\begin{enumerate}[nosep]
  \item Specify the curriculum (4--6 sessions).
  \item For each session, describe the warning, weakened exposure, refutation, and practice components.
  \item Predict the effect on subsequent consumer behavior.
  \item How would you evaluate the program's effectiveness (study design, outcome measures)?
\end{enumerate}
\end{exercise}

\begin{exercise}
Analyze a historical counter-doxonomic movement (e.g., the labor movement's challenge to laissez-faire).
\begin{enumerate}[nosep]
  \item What counter-institutions did it build?
  \item What was its funding level relative to the dominant regime?
  \item What asymmetric advantages did it exploit?
  \item What doxonomic strategies succeeded, and which failed?
  \item What happened when the movement's institutional infrastructure was defunded or dismantled?
\end{enumerate}
\end{exercise}

\begin{exercise}
Apply the disclosure effect (Proposition~\ref{prop:disclosure}) to a specific doxonomic institution.
\begin{enumerate}[nosep]
  \item Choose an institution whose funding sources are partially opaque.
  \item Estimate its current credibility $\credibility{k}$ and the conflict of interest $\text{conflict}(k)$.
  \item If full disclosure were achieved ($d = 1$), how much would effective credibility decline?
  \item What institutional resistance to disclosure would you expect, and how would you overcome it?
\end{enumerate}
\end{exercise}

\begin{exercise}
Design a coordinated counter-doxonomic campaign linking two or more social movements.
\begin{enumerate}[nosep]
  \item Choose the movements (e.g., labor + environmental, racial justice + economic justice).
  \item Write a shared counter-narrative that encompasses both movements' concerns.
  \item Estimate the coordination gain (Proposition~\ref{prop:coordination}) vs.\ independent operation.
  \item What tensions might arise between the movements, and how would you manage them?
\end{enumerate}
\end{exercise}

\begin{exercise}[Computational]
Simulate the media literacy intervention (Model~\ref{mod:media-literacy}):
\begin{enumerate}[nosep]
  \item Create 1,000 agents exposed to a dominant narrative with $\alpha = 0.1$ and $\phi = 10$.
  \item Without literacy: compute the steady-state belief $\belief{}^*$.
  \item Apply a literacy intervention to 30\% of agents ($\ell = 0.3$). Recompute the steady-state.
  \item How much does the literacy intervention reduce dominant-regime belief compared to spending the same budget on counter-narrative production?
  \item Find the break-even point: at what literacy cost per person does direct counter-narrative production become more cost-effective than literacy training?
\end{enumerate}
\end{exercise}

\chapter{Epistemic Democracy}
\label{ch:epistemic-democracy}

\epigraph{The cure for the ailments of democracy is more democracy.}{John Dewey, \textit{The Public and Its Problems}}

\noindent
If concentrated doxonomic power distorts public belief, the remedy is the democratization of belief production.
But what does that mean in practice?
This chapter explores what epistemic democracy means, whether it is achievable, what institutional designs might support it, and what tradeoffs it entails.

The problem is fundamental: democratic governance presupposes informed citizens, but the information environment is itself shaped by power.
If the ``marketplace of ideas'' is dominated by a few well-funded sellers, the resulting common sense reflects purchasing power rather than evidence or deliberation.
Epistemic democracy is the attempt to design institutions that produce a more pluralistic, evidence-responsive, and democratically accountable belief environment.

\section{The Problem}
\label{sec:ch32-problem}

\subsection{Epistemic concentration and democratic failure}

\begin{definition}[Epistemic Democracy]
\label{def:epistemic-democracy}
\index{epistemic democracy}
An \emph{epistemic democracy} is a social arrangement in which the production, circulation, and evaluation of public belief is not dominated by a narrow set of funders or institutions.
Formally, epistemic democracy requires:
\begin{enumerate}[nosep]
  \item $\ehhi < \bar{h}$: epistemic concentration is below a threshold that would allow a small number of actors to determine common sense.
  \item Narrative pluralism: multiple competing frames for any important public question are institutionally supported and accessible.
  \item Accountability: narrative-producing institutions are transparent about their funding, methods, and interests.
  \item Responsiveness: the belief environment is responsive to evidence, deliberation, and democratic input, not just to funding levels.
\end{enumerate}
\end{definition}

The current situation in most democracies falls far short of this ideal.
The preceding chapters documented epistemic concentration ratios of 100:1 or higher in multiple narrative domains (selfishness, meritocracy, consumerism, austerity).
This concentration is not compatible with the democratic principle that citizens should form their beliefs through free and equal deliberation.

\subsection{Three failure modes}

Current democracies exhibit three related epistemic failures:

\begin{enumerate}
  \item \textbf{Plutocratic capture}: wealthy actors dominate narrative production through think-tank funding, media ownership, and lobbying.
  The beliefs that become common sense disproportionately reflect the interests of funders.

  \item \textbf{Algorithmic distortion}: platform algorithms optimize for engagement rather than truth or democratic deliberation, producing information environments that amplify outrage, simplification, and polarization.

  \item \textbf{Institutional sclerosis}: educational curricula, professional norms, and policy frameworks lag behind evidence, reproducing outdated beliefs through institutional inertia even when the evidence has shifted.
\end{enumerate}

\section{Normative Foundations}
\label{sec:ch32-normative}

\subsection{The deliberative ideal}

\textcite{habermas1962} proposed an ``ideal speech situation'' in which all participants have equal opportunity to make claims, challenge claims, and express needs, free from domination.
In doxonomic terms, this requires:
\[
  \credibility{k} = \credibility{k}(\text{evidence quality, argument strength}) \quad \text{not} \quad \credibility{k}(\text{funding level, institutional prestige}).
\]

The ideal is unrealizable in its pure form (credibility will always be influenced by institutional position), but it provides a benchmark against which actual epistemic arrangements can be evaluated.

\subsection{The agonistic alternative}

\textcite{mouffe2005} argues that the deliberative ideal is not only unrealizable but undesirable: genuine politics involves irreducible conflict between incompatible values, and attempts to reduce politics to rational consensus suppress this conflict rather than resolving it.

In doxonomic terms, the agonistic view suggests that the goal is not consensus but \emph{fair competition}: multiple narrative-producing institutions with roughly equal resources competing to persuade a critically engaged public.
The problem is not that different groups produce different narratives; the problem is that one group's narrative dominates because of \emph{funding} rather than \emph{argument}.

\begin{proposition}[Agonistic Epistemic Democracy]
\label{prop:agonistic}
\index{agonistic democracy}
An agonistic epistemic democracy requires not consensus but \emph{competitive pluralism}: multiple institutionally supported narratives contesting for public adherence, with no single narrative dominating through resource advantage rather than persuasive quality.
The metric is not low entropy (consensus) but \emph{fair entropy}: $H(t)$ is neither so low that one narrative monopolizes (capture) nor so high that no shared framework exists for collective action (fragmentation), and the distribution of narrative resources is roughly proportional to the diversity of genuinely held positions in the population.
\end{proposition}

\section{Institutional Designs}
\label{sec:ch32-designs}

\subsection{Public-interest media}

\begin{model}[Public Media as Epistemic Infrastructure]
\label{mod:public-media}
\index{public media}
Public-interest media (funded by taxation or public endowment, editorially independent, and free from advertising pressure) provides a baseline of information quality that commercial media need not match.

Design principles:
\begin{enumerate}[nosep]
  \item \textbf{Funding independence}: endowment or hypothecated tax, insulated from annual budget negotiations that create political leverage.
  \item \textbf{Editorial autonomy}: governance by independent boards with fixed terms, not by government appointment.
  \item \textbf{Universal access}: free at point of use, available across all platforms and regions.
  \item \textbf{Mandate for diversity}: explicit requirement to represent diverse perspectives, not just the perspectives of the educated metropolitan class.
\end{enumerate}
\end{model}

\begin{example}
\label{ex:public-media-comparison}
Comparison of public media models:

\begin{center}
\begin{tabular}{p{2cm}p{3cm}p{3cm}p{3.5cm}}
\toprule
Country & Funding & Independence & Doxonomic assessment \\
\midrule
UK (BBC) & License fee (\pounds3.7B) & Moderate (Royal Charter, but government sets fee) & Substantial epistemic infrastructure but vulnerable to political pressure \\[4pt]
Germany (ARD/ZDF) & Household levy (\texteuro8.4B) & High (interstate treaty, independent fee commission) & Strong independence; highest per-capita public media investment \\[4pt]
US (PBS/NPR) & Mixed (15\% federal, rest private) & Low (dependent on donations and corporate sponsorship) & Weak epistemic infrastructure; commercial pressures dominate \\[4pt]
Scandinavia & License fee / tax & High & High-trust information environment \\
\bottomrule
\end{tabular}
\end{center}

Countries with stronger public media consistently show higher levels of political knowledge, lower susceptibility to misinformation, and more moderate political polarization, consistent with the hypothesis that public media reduces epistemic concentration.
\end{example}

\subsection{Open curricula}

Educational content developed through transparent, participatory processes with input from educators, researchers, communities, and students:
\begin{itemize}[nosep]
  \item \textbf{Transparent development}: curriculum-writing processes that are publicly documented, with input from diverse stakeholders.
  \item \textbf{Evidence-based revision}: mandatory periodic review of curricula against current evidence (addressing the textbook-evidence gaps documented in Chapter~\ref{ch:human-nature}).
  \item \textbf{Critical thinking}: explicit instruction in identifying narrative sources, evaluating evidence, and recognizing doxonomic influence (the doxonomic literacy of Chapter~\ref{ch:counter-doxonomic}).
\end{itemize}

\subsection{Institutional pluralism}

\begin{definition}[Epistemic Pluralism Requirement]
\label{def:pluralism}
\index{epistemic pluralism}
An \emph{epistemic pluralism requirement} is a regulatory framework ensuring that no single funding source dominates any narrative domain.
Analogous to antitrust regulation in economic markets, it limits the concentration of narrative-production capacity.
Formal condition: $\ehhi_{\text{domain}} < \bar{h}$ for each major narrative domain (economics, health, environment, education, etc.).
\end{definition}

Implementation mechanisms:
\begin{itemize}[nosep]
  \item caps on single-source funding of think tanks (e.g., no more than 20\% of budget from one funder),
  \item media ownership limits (maximum share of national audience),
  \item public matching funds for narrative-producing institutions that meet transparency standards,
  \item antitrust review of mergers that would increase epistemic concentration.
\end{itemize}

\subsection{Transparency requirements}

Mandatory disclosure of funding for all narrative-producing institutions:
\begin{itemize}[nosep]
  \item think tanks and policy institutes: full donor lists above a threshold,
  \item academic research: comprehensive conflict-of-interest disclosure,
  \item media outlets: ownership structure and major advertising clients,
  \item political advertising: source, funding, and targeting criteria,
  \item platform algorithms: public auditing of content-recommendation systems.
\end{itemize}

\section{Mechanism Design for Belief Production}
\label{sec:ch32-mechanism}

\subsection{The design problem}

\begin{model}[Democratic Belief Production]
\label{mod:democratic-belief}
\index{mechanism design}
Treating belief production as a mechanism design problem:
\begin{itemize}[nosep]
  \item agents have private preferences over beliefs and narratives,
  \item a social planner seeks to design institutions that aggregate information and preferences without allowing wealthy agents to dominate,
  \item the mechanism must be incentive-compatible (agents cannot profit by misrepresenting preferences) and individually rational.
\end{itemize}
The impossibility results of \textcite{arrow1963} apply: no perfect mechanism exists.
But approximate mechanisms (public media, deliberative assemblies, citizen juries, open curricula, transparency requirements) can reduce epistemic concentration below the capture threshold.
\end{model}

\subsection{Deliberative mini-publics}

\begin{definition}[Deliberative Mini-Public]
\label{def:mini-public}
\index{deliberative mini-public}
A \emph{deliberative mini-public} is a randomly selected group of citizens who are given time, information, and facilitation to deliberate on a specific policy question.
Examples: citizens' assemblies, consensus conferences, deliberative polls.
\end{definition}

Mini-publics bypass doxonomic concentration by:
\begin{itemize}[nosep]
  \item random selection (not self-selection or wealth-selection),
  \item balanced information provision (curated by neutral facilitators, not by interested parties),
  \item structured deliberation (ensuring that all voices are heard, not just the loudest or most funded),
  \item protected time (participants are compensated for their time, removing the wealth barrier to participation).
\end{itemize}

Evidence from citizens' assemblies (Ireland's constitutional convention on abortion, France's Citizens' Convention on Climate, British Columbia's electoral reform assembly) suggests that deliberative mini-publics produce more nuanced, evidence-responsive, and less polarized recommendations than standard democratic processes.

\subsection{Epistemic public goods}

\begin{proposition}[Belief Quality as Public Good]
\label{prop:epistemic-public-good}
\index{epistemic public good}
High-quality public belief is a public good: it is non-rivalrous (one person's accurate belief does not diminish another's) and non-excludable (everyone benefits from a well-informed electorate, not just the well-informed individuals themselves).
Like all public goods, it is under-provided by market mechanisms because individual consumers of information bear the full cost but capture only a fraction of the benefit.
Public investment in epistemic infrastructure (public media, education, research) is justified on the same grounds as public investment in physical infrastructure (roads, bridges, sanitation): the market alone will not provide the socially optimal level.
\end{proposition}

\section{The Diversity-Coordination Tradeoff}
\label{sec:ch32-tradeoff}

\subsection{The optimal entropy problem}

\begin{model}[Optimal Narrative Entropy]
\label{mod:optimal-entropy}
\index{optimal entropy}
Epistemic democracy faces a fundamental tradeoff between diversity and coordination:
\begin{itemize}[nosep]
  \item Too little diversity ($H$ low): common-sense capture; one narrative dominates, suppressing alternatives.
  \item Too much diversity ($H$ high): fragmentation; no shared framework exists for collective action, mutual understanding, or democratic deliberation.
\end{itemize}
The optimal entropy $H^*$ satisfies:
\[
  H^* = \arg\max_H \; W(H) = \underbrace{\alpha \cdot D(H)}_{\text{diversity benefit}} - \underbrace{\beta \cdot F(H)}_{\text{fragmentation cost}}
\]
where $D(H)$ increases with entropy (more diverse perspectives improves collective decision-making) and $F(H)$ increases with entropy (more fragmentation makes collective action harder).
The optimum balances these: enough diversity to prevent capture, enough coherence to sustain cooperation.
\end{model}

This tradeoff is not merely theoretical.
The current moment in many democracies may involve \emph{too much} entropy (post-truth fragmentation) rather than too little (common-sense capture by a single narrative).
Epistemic democracy requires not just challenging dominant narratives but building \emph{shared epistemic infrastructure} (common reference points, trusted institutions, agreed-upon methods of evaluation) that can sustain democratic deliberation without collapsing into either capture or chaos.

\section{Limitations and Risks}
\label{sec:ch32-limits}

\subsection{Who guards the guardians?}

Any institutional design for epistemic democracy creates new concentrations of power: who selects the facilitators for citizens' assemblies? Who designs the transparency requirements? Who governs the public media?
These questions do not have perfect answers (they are variants of the ancient problem of \emph{quis custodiet ipsos custodes?}), but they can be mitigated through:
\begin{itemize}[nosep]
  \item layered accountability (each institution overseen by a different body),
  \item term limits and rotation for governance positions,
  \item sunset clauses requiring periodic renewal and review,
  \item constitutional protections for editorial independence.
\end{itemize}

\subsection{The risk of epistemic paternalism}

\begin{principle}[Anti-Paternalism Constraint]
\label{prin:anti-paternalism}
\index{epistemic paternalism}
Epistemic democracy aims to democratize belief \emph{production}, not to determine belief \emph{content}.
The goal is an information environment in which citizens can form beliefs through free and informed deliberation, not an environment in which the ``correct'' beliefs are pre-selected by experts.
Any institutional design that empowers a specific group to determine what the public ``should'' believe has replaced plutocratic doxonomic power with technocratic doxonomic power, which is not an improvement.
\end{principle}

The line between providing information and shaping belief is not always clear.
But the principle provides a guide: epistemic-democratic institutions should increase the diversity and quality of available information, not select the conclusions that citizens should draw from it.

\section{Notes and References}
\label{sec:ch32-notes}

On epistemic democracy, see \textcite{dewey1927}, \textcite{habermas1962}, and Landemore (2012).
Public media is analyzed in \textcite{mcchesney1999} and Cushion (2012).
Mechanism design draws on \textcite{arrow1963} and Myerson (1981).
On participatory governance, see \textcite{ostrom1990} and \textcite{wright2010}.

For democratic theory and agonistic pluralism, see \textcite{mouffe2005}.
Deliberative mini-publics are surveyed in Fishkin (2009) and Landemore (2020).
On epistemic public goods, see Anderson (2006) and Kitcher (2011).
Media ownership and epistemic concentration are analyzed in Baker (2007).

For the diversity-coordination tradeoff, see Page (2007) on the diversity prediction theorem and Sunstein (2002) on the risks of group polarization.
On epistemic paternalism, see Goldman (1999).

\section{Exercises}

\begin{exercise}
Design a transparency regime for think-tank funding.
\begin{enumerate}[nosep]
  \item What information should be disclosed (donor identity, amounts, strings attached)?
  \item At what threshold should disclosure be required?
  \item How would you enforce compliance?
  \item Estimate the effect on $\ehhi$ if full transparency were achieved.
  \item What resistance to transparency would you expect, and from whom?
\end{enumerate}
\end{exercise}

\begin{exercise}
Compare the epistemic concentration ($\ehhi$) of a country with strong public media (e.g., Germany) vs.\ one with primarily commercial media (e.g., the US).
\begin{enumerate}[nosep]
  \item Estimate $\ehhi$ for each using available data on media ownership and think-tank funding.
  \item Compare levels of political knowledge, misinformation susceptibility, and polarization.
  \item Does the data support the hypothesis that public media reduces epistemic concentration?
  \item What confounding factors might explain the differences?
\end{enumerate}
\end{exercise}

\begin{exercise}
Formalize the diversity-coordination tradeoff (Model~\ref{mod:optimal-entropy}).
\begin{enumerate}[nosep]
  \item Specify functional forms for $D(H)$ and $F(H)$.
  \item Compute $H^*$ analytically.
  \item Is the current entropy in your country above or below $H^*$?
  \item What institutional changes would move entropy toward the optimum?
\end{enumerate}
\end{exercise}

\begin{exercise}
Design a deliberative mini-public (Definition~\ref{def:mini-public}) to address a doxonomically contested question (e.g., ``Is economic inequality too high?'' or ``Should social media algorithms be regulated?'').
\begin{enumerate}[nosep]
  \item How would you select participants?
  \item What information would you provide, and how would you ensure balance?
  \item How would you structure the deliberation?
  \item How would you connect the mini-public's recommendations to actual policy?
\end{enumerate}
\end{exercise}

\begin{exercise}
Apply the anti-paternalism constraint (Principle~\ref{prin:anti-paternalism}) to a specific epistemic-democratic proposal.
\begin{enumerate}[nosep]
  \item Choose a proposal (e.g., public media reform, curriculum revision, platform regulation).
  \item Identify the point at which the proposal risks crossing from ``democratizing information'' to ``determining beliefs.''
  \item Design institutional safeguards to stay on the right side of the line.
\end{enumerate}
\end{exercise}

\begin{exercise}[Computational]
Simulate the optimal entropy model:
\begin{enumerate}[nosep]
  \item Define $D(H) = a\sqrt{H}$ (diversity benefit, concave) and $F(H) = bH^2$ (fragmentation cost, convex).
  \item Compute $H^*$ analytically for $a = 2$, $b = 0.5$.
  \item Simulate a society with 20 narrative frames and varying institutional arrangements (monopoly, duopoly, public media + commercial, fully fragmented).
  \item Measure $H$ under each arrangement. Which comes closest to $H^*$?
  \item What institutional design parameters (number of public media outlets, funding levels, transparency requirements) most effectively move $H$ toward $H^*$?
\end{enumerate}
\end{exercise}

\input{chapters/ch34-better-belief-ecologies}
\input{chapters/ch35-ethics-of-doxonomics}
\chapter{Toward a Science of Social Self-Understanding}
\label{ch:toward}

\epigraph{The owl of Minerva spreads its wings only with the falling of the dusk.}{G.\ W.\ F.\ Hegel, \textit{Philosophy of Right}}

\noindent
This final chapter surveys what doxonomics can and cannot explain, develops the open problems whose resolution will define the field's future, sketches the research programs that will address them, locates doxonomics within the landscape of neighboring disciplines, and closes with a reflection on what it means to study the social production of belief, knowing that one's own beliefs are, inevitably, socially produced.

\section{What Doxonomics Can Explain}
\label{sec:ch35-can}

The framework developed across the preceding 34 chapters provides tools for answering a specific class of questions:

\begin{enumerate}
  \item \textbf{How specific beliefs achieved dominance through institutional investment.}
  The selfishness regime (Chapter~\ref{ch:human-nature}), meritocracy (Chapter~\ref{ch:meritocracy}), consumerism (Chapter~\ref{ch:consumerism}), and austerity (Chapter~\ref{ch:austerity}) were all analyzed as products of identifiable institutional investment operating through traceable production chains.

  \item \textbf{The financial structure of narrative production.}
  Doxonomic accounting (Chapter~\ref{ch:belief-flow-accounting}) provides tools for mapping the flow of money through narrative-producing institutions and estimating the return on doxonomic investment.

  \item \textbf{Why beliefs persist despite contrary evidence.}
  Funded repetition (the $\alpha\phi$ term in the belief dynamics models), institutional credibility, and self-fulfilling mechanisms all explain how beliefs can survive, and even strengthen, in the face of disconfirming evidence.

  \item \textbf{How beliefs become self-fulfilling through behavioral feedback.}
  The selfishness regime's self-fulfilling prophecy (Theorem~\ref{thm:self-fulfilling}), the segregation-hierarchy loop (Proposition~\ref{prop:segregation-feedback}), and the meritocracy-inequality spiral (Model~\ref{mod:merit-inequality-feedback}) all demonstrate how beliefs can generate the conditions that appear to confirm them.

  \item \textbf{The conditions under which belief regimes fracture.}
  The bifurcation model (Chapter~\ref{ch:how-belief-regimes-fracture}), the credibility cascade mechanism, and the crisis dynamics (Chapter~\ref{ch:crisis-doxonomics}) provide a formal framework for understanding when and how dominant beliefs collapse.

  \item \textbf{The welfare costs of doxonomic externalities.}
  The doxonomic market model (Chapter~\ref{ch:economics-belief}) and the legitimation cost analysis (Chapter~\ref{ch:inequality-legitimacy}) provide tools for estimating the social costs of manufactured belief.
\end{enumerate}

\section{What Doxonomics Cannot Explain}
\label{sec:ch35-cannot}

Equally important are the questions doxonomics does \emph{not} answer:

\begin{enumerate}
  \item \textbf{Why specific individuals resist dominant beliefs.}
  Doxonomics is a social-level theory.
  It predicts population-level belief distributions, not individual belief formation.
  Why some people see through doxonomic production while most do not (the psychology of dissent) is a question for cognitive science, personality psychology, and biography, not for doxonomics.

  \item \textbf{The full content of any belief.}
  Doxonomics studies the \emph{production} of belief, not its \emph{content} or \emph{truth value}.
  Showing that a belief is funded does not show that it is false (the genetic fallacy).
  Showing that a belief is unfunded does not show that it is true.

  \item \textbf{Whether a belief should be held.}
  Doxonomics is descriptive and analytical, not normative.
  It can show that a belief was produced through institutional investment rather than evidence, but the normative conclusion (``therefore the belief should be questioned'') requires an additional premise about the epistemic authority of evidence over institutional repetition, a premise this book endorses but cannot derive from doxonomic analysis alone.

  \item \textbf{The precise mechanisms of individual cognition.}
  Doxonomics treats agents as responsive to narrative exposure, social norms, and institutional credibility, but it does not model the neural or cognitive mechanisms by which narratives are processed.
  It is social, not cognitive.

  \item \textbf{Non-material sources of belief.}
  Religious revelation, aesthetic experience, love, grief, wonder: these shape belief in ways that are not reducible to institutional investment.
  Doxonomics covers the material-institutional dimension of belief production, which is enormous, but it is not the whole story.
\end{enumerate}

\section{Open Problems}
\label{sec:ch35-open}

The field's future will be shaped by the resolution of several foundational problems:

\begin{conjecture}[Open Problems in Doxonomics]
\label{conj:open}
\index{open problems}
\leavevmode

\medskip\noindent\textbf{1.\ The Measurement Problem.}
Can we construct a reliable, comparable, cross-national index of doxonomic concentration?
The epistemic Herfindahl index (Definition~\ref{def:ehhi}) provides a theoretical framework, but practical measurement requires solving the problem of tracing funding through opaque channels, accounting for in-kind contributions, and weighting different narrative channels by their effectiveness.

\medskip\noindent\textbf{2.\ The Attribution Problem.}
How much of a population's belief change is attributable to doxonomic investment vs.\ lived experience, peer influence, generational change, and autonomous reasoning?
The causal methods of Chapter~\ref{ch:causal-inference} provide tools, but the problem of separating these sources remains largely unsolved in practice.

\medskip\noindent\textbf{3.\ The Threshold Problem.}
Is there a universal critical funding level above which common-sense capture becomes inevitable?
Or does the threshold depend on contextual factors (existing beliefs, institutional landscape, cultural resilience) in ways that preclude generalization?
If there is no universal threshold, can we at least identify the \emph{factors} that determine the threshold in specific settings?

\medskip\noindent\textbf{4.\ The Decay Problem.}
What determines the half-life of a doxonomic investment once funding is withdrawn?
The generational decay model (Proposition~\ref{prop:generational-decay}) provides a starting point, but the actual decay rate depends on whether the belief has become self-fulfilling, whether it is embedded in institutional routines, and whether counter-narratives are available.

\medskip\noindent\textbf{5.\ The Self-Fulfillment Problem.}
Can we formally separate beliefs that are true \emph{because} they were funded (self-fulfilling prophecies, where the belief creates the conditions it describes) from beliefs that were funded \emph{because} they were true (legitimate research communication)?
This is the deepest methodological challenge in doxonomics: the funded belief and the true belief may look identical in the data.

\medskip\noindent\textbf{6.\ The Welfare Problem.}
How should we compute the total social cost of a doxonomic externality?
The Pigouvian framework (Chapter~\ref{ch:economics-belief}) provides a structure, but measuring the welfare cost of a belief distortion requires counterfactual reasoning about what policies would have been adopted under different beliefs, a profoundly difficult exercise.

\medskip\noindent\textbf{7.\ The Design Problem.}
What is the optimal institutional architecture for minimizing doxonomic distortion while maintaining social coordination?
Chapter~\ref{ch:better-belief-ecologies} sketches design principles, but the multi-objective optimization problem (Proposition~\ref{prop:design-tradeoffs}) does not have a known solution.

\medskip\noindent\textbf{8.\ The Reflexivity Problem.}
Can doxonomics be applied to itself without infinite regress?
The doxonomic paradox (Remark~\ref{rem:paradox}) suggests that the field must be reflexive, but reflexivity applied recursively (``the doxonomic analysis of the doxonomic analysis of the doxonomic analysis\ldots'') either converges (which would be interesting) or diverges (which would be troubling).

\medskip\noindent\textbf{9.\ The Computational Tractability Problem.}
Can the sheaf-theoretic framework (Chapter~\ref{ch:contradiction-coherence}) and the network models (Chapter~\ref{ch:network-analysis}) be made computationally tractable for real social systems with millions of agents and thousands of institutions?
Current analytical models are illustrative; practical application requires scalable computation.

\medskip\noindent\textbf{10.\ The Historical Accounting Problem.}
Can we retroactively estimate the total doxonomic investment in specific beliefs over centuries?
The archival methods of Chapter~\ref{ch:historical-archival} provide qualitative tools, but quantitative historical accounting faces severe data limitations for periods before the 20th century.
\end{conjecture}

\section{Future Research Programs}
\label{sec:ch35-future}

\subsection{Computational doxonomics}

Large-scale simulation of belief production chains using agent-based models calibrated to real funding data.
The research agenda:
\begin{itemize}[nosep]
  \item build ABMs of specific narrative domains (climate, economics, health) with realistic institutional structure,
  \item calibrate model parameters against historical data (funding flows, survey time series, media content analysis),
  \item use calibrated models to predict the effects of policy interventions (transparency laws, public media expansion, media literacy programs),
  \item validate predictions against natural experiments and policy changes.
\end{itemize}

\subsection{Comparative doxonomics}

Cross-national comparison of narrative infrastructures and their effects on public belief.
The research agenda:
\begin{itemize}[nosep]
  \item construct comparable measures of doxonomic concentration ($\ehhi$) across countries,
  \item compare the institutional infrastructure of belief production (think-tank density, media structure, educational systems) across political-economic regimes,
  \item test whether variation in institutional infrastructure predicts variation in public beliefs, controlling for economic conditions and cultural factors,
  \item identify ``natural experiments'' in institutional change (media deregulation, think-tank founding, curriculum reform) that allow causal estimation.
\end{itemize}

\subsection{Historical doxonomics}

Systematic reconstruction of the doxonomic history of major beliefs.
The research agenda:
\begin{itemize}[nosep]
  \item construct longitudinal datasets of institutional funding for narrative production in specific domains,
  \item digitize and analyze archival materials using the computational methods of Chapter~\ref{ch:text-media},
  \item develop methods for estimating doxonomic investment from indirect evidence (publication volumes, institutional founding dates, personnel flows),
  \item produce definitive doxonomic genealogies of the major belief regimes analyzed in Part IV.
\end{itemize}

\subsection{Applied doxonomics}

The development of practical tools for democratic institutions:
\begin{itemize}[nosep]
  \item doxonomic audit methodologies that can be applied by journalists, NGOs, and government agencies,
  \item real-time monitoring of narrative entropy and epistemic concentration,
  \item media literacy curricula grounded in doxonomic theory,
  \item policy evaluation frameworks that include doxonomic effects.
\end{itemize}

\section{Doxonomics and Its Neighbors}
\label{sec:ch35-neighbors}

Doxonomics draws on and contributes to several established fields:

\begin{center}
\begin{tabular}{p{3.5cm}p{4cm}p{4.5cm}}
\toprule
Field & What doxonomics takes from it & What doxonomics adds \\
\midrule
Political economy & Institutional analysis, power structures & Formal models of belief production \\[4pt]
Sociology of knowledge & Social construction of belief & Quantitative measurement, funding-tracing \\[4pt]
Media studies & Content analysis, framing theory & Institutional economics of narrative production \\[4pt]
Behavioral economics & Cognitive biases, bounded rationality & Institutional sources of bias (not just cognitive) \\[4pt]
Science and technology studies & Social shaping of knowledge & Application to all belief domains, not just science \\[4pt]
Critical theory & Ideology critique, hegemony & Mathematical formalization, empirical testing \\
\bottomrule
\end{tabular}
\end{center}

Doxonomics does not replace any of these fields.
It provides a \emph{connective framework} (a common language of formal models, measurement tools, and causal methods) that allows insights from each field to be integrated into a unified analysis of how power shapes belief.

\section{What Would Success Look Like?}
\label{sec:ch35-success}

If doxonomics succeeds as a field, what changes?

\begin{enumerate}
  \item \textbf{Routine doxonomic auditing}: major policy debates would routinely include an analysis of who funds the competing narratives and through what institutional channels, just as financial auditing is now routine for corporate governance.

  \item \textbf{Doxonomic literacy}: citizens would be trained to ask ``who funded this claim, and why?'' just as they are trained to check sources in academic writing.

  \item \textbf{Institutional transparency}: funding disclosure for narrative-producing institutions would be as standard as financial disclosure for corporations.

  \item \textbf{Better belief ecologies}: policy interventions informed by doxonomic analysis (public media, media literacy, epistemic pluralism requirements) would produce information environments that are more diverse, more evidence-responsive, and more equitable.

  \item \textbf{Self-awareness}: societies would have a vocabulary for discussing the industrial production of common sense, just as they now have vocabularies for discussing economic production, public health, and environmental quality.
\end{enumerate}

None of these outcomes requires that doxonomics achieve the precision of physics or the predictive power of epidemiology.
It requires only that the field provide \emph{useful tools}: tools that make visible what was previously invisible, that measure what was previously unmeasured, and that enable choices where there were previously only defaults.

\section{Final Reflection}
\label{sec:ch35-final}

A society that can measure its economic output, its carbon emissions, and its disease burden, but cannot measure the industrial production of its own common sense, has a blind spot at the center of its self-understanding.

This book has attempted to illuminate that blind spot, to develop the concepts, models, methods, and applications needed to study how material power shapes public belief.
The attempt is preliminary.
The models are stylized.
The empirical base is thin in many areas.
The field does not yet exist as an institutionalized discipline with departments, journals, and professional norms.

But the phenomenon it studies is real.
The funding flows are traceable.
The institutional production chains are visible.
The beliefs they produce have measurable consequences for inequality, democracy, well-being, and collective action.
And the alternative (pretending that common sense is natural, that ideas succeed on their merits, that the marketplace of ideas is free and fair) is not intellectual modesty.
It is a refusal to look at the evidence.

Doxonomics does not promise to eliminate ideology.
It does not promise to free humanity from the influence of power on belief.
It promises something more modest and more useful: the tools to \emph{see} how beliefs are produced, to \emph{measure} what they cost, and to \emph{ask} whether the world might believe otherwise.

That is the beginning of doxonomic literacy.
And doxonomic literacy, like all literacy, is the precondition for freedom: not freedom from all influence, which is impossible, but freedom to \emph{know} what influences you, which is the beginning of choosing for yourself.

\section{Notes and References}
\label{sec:ch35-notes}

On the project of social self-understanding, see \textcite{dewey1927} and \textcite{habermas1962}.
The Hegelian epigraph and its implications for a science that arrives ``at dusk'' (after the phenomena it studies are already in place) connect to Hegel's \emph{Philosophy of Right} (1820) and to Marx's inversion of Hegel (the point is to change the world, not merely to understand it).

On real utopias and institutional imagination, see \textcite{wright2010}.
On the relationship between social science and democracy, see Putnam (2000) and \textcite{dewey1927}.
For the epistemic responsibilities of researchers, see Kitcher (2001) and Douglas (2009).

The open problems listed here are original to this text, though each connects to established research programs in the fields enumerated in Section~\ref{sec:ch35-neighbors}.

\section{Exercises}

\begin{exercise}
Choose one open problem from Conjecture~\ref{conj:open}.
Write a two-page research proposal outlining:
\begin{enumerate}[nosep]
  \item the specific question you would address,
  \item the data sources and methods you would use,
  \item the expected results and how they would advance the field,
  \item the limitations of your proposed approach,
  \item how you would evaluate whether the problem has been ``solved'' or merely ``advanced.''
\end{enumerate}
\end{exercise}

\begin{exercise}
Design a one-semester ``Doxonomic Literacy'' course for undergraduates.
\begin{enumerate}[nosep]
  \item Specify 14 weekly topics (drawing on the structure of this book).
  \item For each week, identify the key concept, the reading, and a practical exercise.
  \item Design the assessment: what should students be able to \emph{do} at the end of the course?
  \item How would you evaluate whether the course has actually improved students' doxonomic literacy (as opposed to merely teaching them the vocabulary)?
\end{enumerate}
\end{exercise}

\begin{exercise}
Apply the full doxonomic toolkit to a belief of your choice.
Produce a mini case study (5--10 pages) covering:
\begin{enumerate}[nosep]
  \item doxonomic genealogy (when and how the belief was produced),
  \item doxonomic account (who funds the belief and through what channels),
  \item network analysis (the institutional structure of production),
  \item causal analysis (evidence that the production causes belief change),
  \item welfare estimation (what the belief costs society),
  \item counter-doxonomic analysis (what counter-narratives exist and how well-funded they are).
\end{enumerate}
\end{exercise}

\begin{exercise}
Write a critical review of this textbook from the perspective of a skeptic.
\begin{enumerate}[nosep]
  \item What are the strongest objections to the doxonomic framework?
  \item Where does the book overclaim relative to its evidence?
  \item What alternative explanations for the phenomena it documents are not adequately addressed?
  \item What would need to be true for doxonomics to be fundamentally wrong?
\end{enumerate}
\end{exercise}

\begin{exercise}
Imagine doxonomics 20 years from now.
\begin{enumerate}[nosep]
  \item Which open problems (Conjecture~\ref{conj:open}) are most likely to have been resolved?
  \item What new problems will have emerged?
  \item What institutional infrastructure will the field need (journals, departments, professional associations)?
  \item What is the greatest risk the field faces (co-optation? irrelevance? technocratic capture?)?
\end{enumerate}
\end{exercise}


% === Back Matter ===
\backmatter

% --- Appendices ---
\appendix

\chapter{Glossary of Doxonomic Terms}
\label{app:glossary}

\begin{description}[style=nextline, leftmargin=3cm]

\item[Anthropological regime] A dominant practical model of human nature specifying motivational theory, social ontology, and capacity claims. See Definition~\ref{def:anthropological-regime}.

\item[Asymptotic capture] The process by which sufficiently asymmetric evidence exposure drives beliefs to certainty regardless of prior. See Conjecture~\ref{conj:asymptotic-capture}.

\item[Belief production chain] The composition of operators from funding through institutional processing, narrative output, circulation, internalization, behavioral consequence, to reproduction. See Definition~\ref{def:bpc-formal}.

\item[Belief stock] The fraction of a population holding a given belief at a given time. See Definition~\ref{def:stock-flow}.

\item[Common sense] A proposition held with high credence and low variance across a population. See Definition~\ref{def:common-sense}.

\item[Common-sense capture] The phase transition by which a funded narrative becomes invisible as a claim, experienced as a feature of reality. See Definition~\ref{def:capture}.

\item[Conversion problem] The foundational problem of doxonomics: how resources are converted into durable public belief.

\item[Counter-doxonomics] The study and practice of resisting, disrupting, or replacing dominant doxonomic regimes. See Chapter~\ref{ch:counter-doxonomic}.

\item[Credibility collapse] Rapid decline in institutional credibility due to scandal, failure, or exposure. See Definition~\ref{def:credibility-collapse}.

\item[Doxa] (Greek) The unquestioned taken-for-granted; beliefs so deeply embedded they are not perceived as beliefs. See Definition~\ref{def:doxa}.

\item[Doxametry] The quantitative measurement branch of doxonomics. See Chapter~\ref{ch:doxometric-measurement}.

\item[Doxonomic equilibrium] A fixed point of the belief production chain where beliefs reproduce the conditions of their own production. See Definition~\ref{def:doxonomic-equilibrium}.

\item[Doxonomic externality] A social cost generated when a subsidized belief reshapes behavior in harmful ways. See Definition~\ref{def:doxonomic-externality}.

\item[Doxonomic game] An experimental economic game administered after exposure to a doxonomic treatment. See Definition~\ref{def:doxonomic-game}.

\item[Doxonomic network] A weighted directed graph of funders, institutions, media, and audiences connected by influence or resource flows. See Definition~\ref{def:doxonomic-network}.

\item[Doxonomic power] The capacity to shape the conditions under which people form beliefs, including the power to set priors, agendas, and epistemic rules. See Definition~\ref{def:doxonomic-power}.

\item[Doxonomic system] A structured ensemble of agents, institutions, narrative spaces, belief spaces, and operators for production and internalization. See Definition~\ref{def:doxonomic-system}.

\item[Doxonomics] The science of how material power shapes public belief. From Greek \textit{doxa} (opinion, belief) and \textit{-nomics} (systematic study).

\item[Epistemic broker] A node with high betweenness centrality in a doxonomic network, bridging otherwise disconnected narrative clusters. See Definition~\ref{def:broker}.

\item[Epistemic capital] Resources convertible to credibility, attention, or authority in belief formation. See Definition~\ref{def:epistemic-capital}.

\item[Epistemic democracy] A social arrangement in which belief production is not dominated by a narrow set of funders or institutions. See Definition~\ref{def:epistemic-democracy}.

\item[Epistemic fragmentation] A condition in which populations inhabit incommensurable belief worlds with no shared epistemic framework. See Definition~\ref{def:fragmentation}.

\item[Epistemic Herfindahl Index (EHI)] A concentration score measuring market share in belief production. See Definition~\ref{def:ehhi}.

\item[Epistemic rent] Excess doxonomic return earned from a monopoly position in legitimate knowledge production. See Definition~\ref{def:epistemic-rent}.

\item[Gluing obstruction] The failure of locally consistent belief systems to combine into a globally consistent whole, measured by sheaf cohomology. See Definition~\ref{def:obstruction}.

\item[Ideological ROI (IROI)] The belief change achieved per unit of doxonomic investment. See Definition~\ref{def:iroi}.

\item[Lock-in] A condition in which a belief is locally asymptotically stable with a large basin of attraction. See Definition~\ref{def:lock-in}.

\item[Narrative entropy] The information-theoretic measure of diversity in the narrative space. See Definition~\ref{def:narrative-entropy}.

\item[Narrative flow] The rate at which narrative content supporting a belief is produced and circulated. See Definition~\ref{def:stock-flow}.

\item[Narrative infrastructure] The durable network of institutions through which a society produces and reproduces its common sense. See Definition~\ref{def:narrative-infrastructure}.

\item[Narrative inoculation] Preemptive exposure to a weakened dominant narrative plus refutation, reducing subsequent susceptibility. See Definition~\ref{def:inoculation}.

\item[Narrative shield] A set of beliefs that deflects attention from structural causes to individual explanations. See Definition~\ref{def:narrative-shield}.

\item[Pluralistic ignorance] A condition in which most individuals privately doubt a proposition but publicly affirm it. See Definition~\ref{def:pluralistic-ignorance}.

\item[Prior-setting power] The capacity to influence the initial belief distribution of a population, typically through education. See Definition~\ref{def:prior-setting}.

\item[Tipping point] A discontinuous jump in equilibrium belief level caused by a small shift in threshold distribution. See Theorem~\ref{thm:tipping}.

\end{description}

\chapter{Mathematical Notation and Basic Models}
\label{app:math-notation}

This appendix collects all mathematical notation and summarizes the core models developed in the text.

\section*{Complete Symbol Table}

See the Notation chapter in the frontmatter for the full symbol table.

\section*{Core Models}

\subsection*{Model 1: First-Order Belief Equation (Model~\ref{mod:first-order})}
\[
  \belief{i}(t) = \sum_{k=1}^{m} \funding{k}(t) \cdot \credibility{k} \cdot \relevance{k}{i} \cdot \reach{k} \cdot \persistence{k} - \counter{i}(t)
\]

\subsection*{Model 2: Stock-and-Flow Dynamics (Definition~\ref{def:stock-flow})}
\[
  \frac{d\belief{i}}{dt} = \alpha \cdot \phi_i(t) \cdot (1 - \belief{i}(t)) - \delta \cdot \belief{i}(t) + \beta \cdot \belief{i}(t)(1 - \belief{i}(t))
\]

\subsection*{Model 3: SIR Belief Diffusion (Model~\ref{mod:sir-belief})}
\begin{align*}
  \frac{dS}{dt} &= -\beta S I - \alpha \phi(t) S \\
  \frac{dI}{dt} &= \beta S I + \alpha \phi(t) S - \gamma I \\
  \frac{dR}{dt} &= \gamma I
\end{align*}
Doxonomic reproduction number: $R_0 = (\beta + \alpha\phi)/\gamma$.

\subsection*{Model 4: Doxonomic Input-Output (Model~\ref{mod:input-output})}
\[
  \bm{n} = (\bm{I} - \bm{A})^{-1}\bm{d}
\]
where $\bm{A}$ is the inter-institutional narrative flow matrix and $\bm{d}$ is the exogenous injection.

\subsection*{Model 5: Granovetter-Schelling Threshold (Model~\ref{mod:granovetter})}
\[
  \mu^* = G(\mu^*)
\]
where $G$ is the CDF of individual adoption thresholds.

\section*{Key Equations}

\begin{longtable}{p{4cm} p{9cm}}
\toprule
\textbf{Equation} & \textbf{Reference} \\
\midrule
\endhead
Optimal allocation & $\credibility{k} \cdot \relevance{k}{i} \cdot c_k'(\funding{k}) \cdot \persistence{k} = \lambda$ (Prop.~\ref{prop:optimal-allocation}) \\
Epistemic HI & $\ehhi = \sum_k s_k^2$ (Def.~\ref{def:ehhi}) \\
Welfare loss & $\int_{\mu^{**}}^{\mu^*} [D'(\mu) - c'(\mu)] \, d\mu$ (Thm.~\ref{thm:welfare-loss}) \\
Polarization & $\pi_A^{(T)} \to 1$, $\pi_B^{(T)} \to 0$ under asymmetric exposure (Thm.~\ref{thm:polarization}) \\
Legitimation & $L(I, \phi_L) = \phi_L / (\phi_L + \kappa(I-I_0)^+)$ (Def.~\ref{def:legitimation}) \\
Fixed point & $\fundingvec^* = \bpc(\fundingvec^*)$ (Def.~\ref{def:doxonomic-equilibrium}) \\
\bottomrule
\end{longtable}

\chapter{Guide to Data Sources}
\label{app:data-sources}

This appendix lists major data sources for doxonomic research, organized by type.

\textbf{Source-quality classification.}
Each source is tagged with an epistemic status to help researchers assess reliability:
\begin{description}[nosep]
  \item[\textsc{primary}] Original data collected through systematic procedures (government filings, surveys, sensor data).
  \item[\textsc{secondary}] Compiled or curated from primary sources (databases, indices, meta-analyses).
  \item[\textsc{commercial}] Proprietary data sold by private firms; methodology may not be fully transparent.
  \item[\textsc{self-report}] Data reported by organizations about themselves (annual reports, transparency reports).
  \item[\textsc{advocacy}] Data compiled by organizations with an explicit political or ideological mission; useful but requires independent verification.
\end{description}
Researchers should triangulate across source types and never rely solely on advocacy or self-report data for causal claims.

\section*{Advertising Data}
\begin{description}[nosep]
  \item[WARC Global Ad Trends] Global and regional advertising expenditure data. \textsc{commercial}
  \item[eMarketer/Insider Intelligence] Digital advertising forecasts and spending data. \textsc{commercial}
  \item[Ad\$pender (Kantar)] Detailed U.S.\ advertising expenditure by company and brand. \textsc{commercial}
\end{description}

\section*{Lobbying and Political Spending}
\begin{description}[nosep]
  \item[OpenSecrets (opensecrets.org)] U.S.\ federal lobbying, campaign finance, and dark money. \textsc{secondary} (compiled from primary filings)
  \item[EU Transparency Register] Lobbying in the European Union. \textsc{primary/self-report} (registration is mandatory but data is self-declared)
  \item[FollowTheMoney (followthemoney.org)] U.S.\ state-level campaign finance. \textsc{secondary}
\end{description}

\section*{Think Tank and Foundation Funding}
\begin{description}[nosep]
  \item[IRS Form 990 Filings] Tax returns for U.S.\ nonprofits; available via ProPublica Nonprofit Explorer. \textsc{primary} (legally required disclosures; incomplete for donor-advised funds)
  \item[InfluenceWatch (influencewatch.org)] Profiles of foundations, think tanks, and advocacy organizations. \textsc{secondary/advocacy} (compiled by the Capital Research Center, a conservative-leaning organization; useful for mapping but interpret editorial framing critically)
  \item[Foundation Directory Online] Comprehensive database of U.S.\ foundations and grants. \textsc{secondary}
  \item[Atlas Network Annual Reports] Self-reported data on the global network of free-market think tanks. \textsc{self-report} (organizationally curated; independent verification recommended)
\end{description}

\section*{Media Ownership}
\begin{description}[nosep]
  \item[Columbia Journalism Review (CJR) Who Owns What] U.S.\ media ownership database.
  \item[Media Ownership Monitor (mom-gmr.org)] Global media concentration data.
  \item[Ofcom Media Ownership Reports] UK media ownership data.
\end{description}

\section*{Survey Data on Beliefs}
\begin{description}[nosep]
  \item[World Values Survey] Cross-national survey of values and beliefs (7 waves, 1981--present).
  \item[General Social Survey (GSS)] U.S.\ attitudes and beliefs since 1972.
  \item[Eurobarometer] EU-wide public opinion surveys.
  \item[Pew Research Center] Global attitudes, media habits, and political beliefs.
  \item[ISSP (International Social Survey Programme)] Cross-national surveys on specific topics.
\end{description}

\section*{Academic Funding}
\begin{description}[nosep]
  \item[NSF Award Search] U.S.\ National Science Foundation grants.
  \item[NIH Reporter] U.S.\ National Institutes of Health funding.
  \item[OECD Science and Technology Indicators] International R\&D spending.
  \item[UnKoch My Campus] Database of Koch-funded academic programs.
\end{description}

\section*{Text and Media Corpora}
\begin{description}[nosep]
  \item[Google Ngrams] Word frequency in books, 1500--2019.
  \item[GDELT Project] Global media monitoring and event database.
  \item[Media Cloud] Open-source media analysis platform.
  \item[Common Crawl] Web archive for large-scale text analysis.
  \item[LexisNexis / Factiva] Commercial news archives.
\end{description}

\section*{Platform Transparency}
\begin{description}[nosep]
  \item[Meta Ad Library] Political and social advertising on Facebook/Instagram.
  \item[Google Transparency Report] Political ad spending on Google platforms.
  \item[Twitter/X Transparency Center] Data on content moderation and advertising.
\end{description}

\chapter{Sample Doxometric Survey Instruments}
\label{app:survey}

\textbf{Status}: The instruments below are \textbf{pilot items} intended to illustrate the structure of doxometric measurement.
They have not been validated through factor analysis, invariance testing, or external criterion studies.
Before use in research, they require: (a)~exploratory factor analysis to assess dimensionality, (b)~confirmatory factor analysis with holdout samples, (c)~measurement invariance testing across demographic groups, (d)~social desirability bias assessment, and (e)~test-retest reliability estimation.
Where validated instruments exist in the literature (e.g., Kluegel and Smith's stratification beliefs scale, Jost and colleagues' system justification measures), researchers should prefer those.

\section*{Instrument 1: Belief in Selfish Human Nature}

Respondents rate agreement on a 7-point Likert scale (1 = Strongly Disagree, 7 = Strongly Agree).

\begin{enumerate}[nosep]
  \item Most people are primarily motivated by self-interest.
  \item If given the chance, most people would take advantage of others.
  \item Cooperation only happens when people are forced or incentivized.
  \item People who help others usually expect something in return.
  \item Competition is more natural to humans than cooperation.
  \item[\textbf{R6.}] Most people are willing to help a stranger even when there is no reward. \textit{(reverse-coded)}
  \item[\textbf{R7.}] Humans naturally form cooperative communities without being forced. \textit{(reverse-coded)}
  \item[\textbf{R8.}] People generally try to be fair, even when no one is watching. \textit{(reverse-coded)}
\end{enumerate}

\textbf{Scoring}: Reverse-code items R6--R8 before summing. Range 8--56. Higher scores indicate stronger belief in selfish human nature.
Internal consistency (Cronbach's $\alpha$) should be assessed empirically; a threshold of $\alpha > 0.70$ is minimally acceptable.

\textbf{Note}: This instrument conflates beliefs about motivation, norm compliance, and institutional necessity.
Exploratory factor analysis may reveal a multi-factor structure (e.g., ``cynicism about others'' vs.\ ``necessity of enforcement'').
Researchers should test dimensionality before treating it as a single scale.

\section*{Instrument 2: Meritocratic Belief Scale}

\begin{enumerate}[nosep]
  \item People who work hard almost always get ahead.
  \item Success in life is mostly determined by effort and talent.
  \item The wealthy generally deserve their wealth.
  \item Everyone has a roughly equal opportunity to succeed.
  \item[\textbf{R5.}] Success depends more on connections and family background than on individual effort. \textit{(reverse-coded)}
  \item[\textbf{R6.}] The economic system is set up in a way that gives some people unfair advantages. \textit{(reverse-coded)}
\end{enumerate}

\textbf{Scoring}: Reverse-code R5--R6 before summing. Range 6--42. Higher scores indicate stronger meritocratic belief.

\textbf{Note}: This scale folds together beliefs about effort, deservingness, equal opportunity, and poverty attribution.
These may be distinct constructs; pilot testing should assess whether a single-factor model fits the data.

\section*{Instrument 3: Behavioral Vignette (Norm Misperception)}

\begin{quote}
\textit{Imagine you find a wallet containing \$200 and the owner's ID. No one is watching.}
\begin{enumerate}[nosep]
  \item \textit{What would you do?} (Return it / Keep the money / Other)
  \item \textit{What do you think most people would do?} (Return it / Keep the money / Other)
  \item \textit{What percentage of people do you think would return it?} (0--100\%)
\end{enumerate}
\end{quote}

The gap between own behavior (item 1) and perceived norm (items 2--3) measures the \emph{norm misperception gap} (Definition~\ref{def:perceived-norm}).
A large gap (most people say they would return it, but predict others would keep it) indicates doxonomic installation of selfishness as perceived norm.

\section*{General Methodological Notes}

For all instruments:
\begin{itemize}[nosep]
  \item Reverse-code counter-directional items (marked R) before summing.
  \item Report Cronbach's $\alpha$ and factor loadings.
  \item Assess social desirability bias (e.g., include a short social desirability scale and partial out its effects).
  \item For cross-national use, test measurement invariance (configural, metric, scalar) across country groups.
  \item Compare stated-belief scores with behavioral measures to assess divergence (Proposition~\ref{prop:divergence}).
  \item Consider supplementing Likert items with list experiments or endorsement experiments to mitigate response bias on sensitive beliefs.
\end{itemize}

\chapter{Network and Accounting Templates}
\label{app:templates}

This appendix provides templates for constructing doxonomic networks and accounting tables.

\section*{Template 1: Donor-Institution Adjacency Matrix}

\begin{center}
\begin{tabular}{l|cccc}
\toprule
 & Institution 1 & Institution 2 & Institution 3 & $\cdots$ \\
\midrule
Donor A & \$\_\_\_\_ & \$\_\_\_\_ & \$\_\_\_\_ & \\
Donor B & \$\_\_\_\_ & \$\_\_\_\_ & \$\_\_\_\_ & \\
Donor C & \$\_\_\_\_ & \$\_\_\_\_ & \$\_\_\_\_ & \\
$\vdots$ & & & & \\
\bottomrule
\end{tabular}
\end{center}

\textbf{Instructions}: Fill in annual grant amounts from each donor to each institution. Data sources: IRS 990 filings, foundation annual reports, InfluenceWatch.

\section*{Template 2: Doxonomic Input-Output Table}

\begin{center}
\begin{tabular}{l|cccc|c}
\toprule
\textbf{To $\downarrow$ / From $\to$} & Think tanks & Media & Universities & Gov't & Final output \\
\midrule
Think tanks & --- & $a_{12}$ & $a_{13}$ & $a_{14}$ & $d_1$ \\
Media & $a_{21}$ & --- & $a_{23}$ & $a_{24}$ & $d_2$ \\
Universities & $a_{31}$ & $a_{32}$ & --- & $a_{34}$ & $d_3$ \\
Gov't & $a_{41}$ & $a_{42}$ & $a_{43}$ & --- & $d_4$ \\
\bottomrule
\end{tabular}
\end{center}

\textbf{Instructions}: $a_{ij}$ is the fraction of institution $j$'s narrative output used as input by institution $i$. Estimate from citation patterns, sourcing data, and content analysis. Compute the Leontief inverse $(\bm{I} - \bm{A})^{-1}$ to find multiplier effects.

\section*{Template 3: Belief-Flow Account}

\begin{center}
\begin{tabular}{lllrr}
\toprule
\textbf{Source} & \textbf{Institution} & \textbf{Channel} & \textbf{Amount (\$)} & \textbf{Throughput ($\Theta$)} \\
\midrule
 & & & & \\
 & & & & \\
 & & & & \\
\midrule
\textbf{Total} & & & \$\_\_\_\_ & \_\_\_\_ \\
\bottomrule
\end{tabular}
\end{center}

\textbf{Instructions}: For a specific belief $p$, trace all identifiable funding flows from source through institution to channel. Estimate throughput as $\Theta_k = \funding{k} \cdot \credibility{k} \cdot \reach{k}$.

\chapter{Suggested Datasets, Archives, and Corpora}
\label{app:datasets}

This appendix lists datasets suitable for doxonomic research.

\begin{longtable}{p{4cm} p{9.5cm}}
\toprule
\textbf{Dataset} & \textbf{Description} \\
\midrule
\endhead

World Values Survey & Cross-national survey of beliefs and values in 100+ countries (1981--present). Core resource for doxometric measurement. \\

General Social Survey (GSS) & U.S.\ attitudes, behaviors, and beliefs since 1972. Annual or biennial. \\

Eurobarometer & EU-wide public opinion surveys on trust, institutions, and social values. \\

OpenSecrets Bulk Data & U.S.\ federal lobbying, PAC contributions, and outside spending. Machine-readable. \\

ProPublica Nonprofit Explorer & Searchable database of IRS 990 filings for U.S.\ nonprofits. \\

GDELT Project & Real-time global media monitoring: events, tone, themes, and locations. \\

Media Cloud & Open-source platform for analyzing media ecosystems and narrative flows. \\

Google Ngrams & Word and phrase frequency in the Google Books corpus (1500--2019). \\

Google Trends & Search interest over time by region. Proxy for public attention. \\

Common Crawl & Petabyte-scale web archive for large-scale text analysis. \\

Facebook Ad Library & Archive of political and social issue advertisements on Meta platforms. \\

OECD Education at a Glance & Cross-national data on education spending, enrollment, and curricula. \\

UNESCO Institute for Statistics & Global data on media, education, and cultural spending. \\

Pew Global Attitudes & Cross-national surveys on democracy, economics, technology, and social values. \\

ISSP & International Social Survey Programme: annual thematic cross-national surveys. \\

ICPSR & Inter-university Consortium for Political and Social Research: massive archive of social science datasets. \\

HathiTrust Digital Library & 17+ million digitized volumes for computational text analysis. \\

Internet Archive & Web and media archives including the Wayback Machine. \\

InfluenceWatch & Profiles and funding data for advocacy organizations, unions, and think tanks. \\

Corporate Europe Observatory & Data on EU corporate lobbying and revolving doors. \\

\bottomrule
\end{longtable}

\chapter{Replication Exercises}
\label{app:replication}

These exercises guide the reader through replicating key doxonomic analyses using real data.

\section*{Exercise 1: Epistemic Herfindahl Index for Cable / Broadcast News}

\textbf{Goal}: Compute $\ehhi$ for a national television-news market. The US cable-news worked example is given below; the same methodology applies to UK broadcast news (BBC, Sky, ITV, GB News), Brazilian open TV (Globo, Record, SBT, Band), Indian news TV (Republic, Times Now, NDTV, Aaj Tak, India Today), or any other partitioned national market.

\textbf{Data}: Nielsen audience-share data (available from Pew Research Center's annual ``State of the News Media'' reports for the US; BARB for the UK; Kantar IBOPE Media for Brazil; BARC for India; ATM-style audience services in EU member states).

\textbf{Steps}:
\begin{enumerate}[nosep]
  \item Obtain average primetime viewership for the major channels in the chosen market (e.g., in the US: CNN, Fox News, MSNBC, and other cable news channels; comparable channel lists for the UK, Brazil, India, etc.).
  \item Compute market shares $s_k = \text{viewership}_k / \text{total viewership}$.
  \item Compute $\ehhi = \sum_k s_k^2$.
  \item Interpret: is the cable / broadcast news market concentrated by doxonomic standards ($\ehhi > 0.25$)?
  \item Repeat for 2005, 2010, 2015, 2020. Plot the trend.
\end{enumerate}

\section*{Exercise 2: Donor Network of Free-Market Think Tanks}

\textbf{Goal}: Construct a donor-institution network for a national free-market think-tank ecosystem. The US case is given below; the same methodology applies to the UK (IEA, Adam Smith Institute, Centre for Policy Studies, Institute for Free Trade, TaxPayers' Alliance), Brazil (Instituto Millenium, Instituto Mises Brasil, Instituto Liberal), India (CCS, Liberty Institute, Takshashila), and other Atlas-Network national ecosystems.

\textbf{Data}: IRS 990 filings via ProPublica Nonprofit Explorer; InfluenceWatch profiles. (UK Charity Commission filings; Brazilian and Indian equivalent registries where available.)

\textbf{Steps}:
\begin{enumerate}[nosep]
  \item Select 10 major think tanks for the chosen market (in the US: Heritage Foundation, Cato Institute, AEI, Hoover, Manhattan Institute, etc.; substitute the national-equivalent list for non-US replications).
  \item For each, identify the top 5 foundation donors from 990 filings.
  \item Construct a bipartite network (donors $\times$ institutions) with edge weights = grant amounts.
  \item Compute eigenvector centrality for all nodes.
  \item Identify narrative clusters using modularity maximization.
\end{enumerate}

\section*{Exercise 3: Semantic Drift of ``Freedom''}

\textbf{Goal}: Measure how the word ``freedom'' has changed meaning in American political discourse.

\textbf{Data}: Google Ngrams or a corpus of presidential speeches.

\textbf{Steps}:
\begin{enumerate}[nosep]
  \item Train word embeddings (Word2Vec or GloVe) on decade-specific subcorpora.
  \item Extract the embedding vector for ``freedom'' in each decade.
  \item Compute cosine distance between consecutive decades.
  \item Identify the decade with maximum drift.
  \item Examine nearest neighbors of ``freedom'' in each period to interpret the drift.
\end{enumerate}

\section*{Exercise 4: Meritocratic Belief Across Countries}

\textbf{Goal}: Compare meritocratic belief across 20 countries.

\textbf{Data}: World Values Survey (items on effort, success, fairness).

\textbf{Steps}:
\begin{enumerate}[nosep]
  \item Select relevant items from WVS Wave 7.
  \item Construct a meritocratic belief index (factor analysis or simple sum).
  \item Plot country-level means against Gini coefficients.
  \item Test the hypothesis: does higher inequality correlate with stronger meritocratic belief (legitimation function)?
\end{enumerate}

\section*{Exercise 5: Narrative Entropy of Climate Discourse}

\textbf{Goal}: Measure narrative diversity in media coverage of climate change over time.

\textbf{Data}: Media Cloud or GDELT.

\textbf{Steps}:
\begin{enumerate}[nosep]
  \item Collect climate-related articles from 2000--2024.
  \item Apply LDA topic modeling with $K = 10$ topics.
  \item Compute narrative entropy $H(t) = -\sum_k p_k(t) \log p_k(t)$ for each year.
  \item Plot $H(t)$ over time. Identify periods of low entropy (narrative capture).
  \item Correlate entropy changes with known doxonomic events (e.g., fossil fuel campaign spending).
\end{enumerate}

\chapter{Recommended Reading}
\label{app:reading}

Organized by topic, with brief annotations.

\section*{Political Economy}

\begin{itemize}[nosep]
  \item \textcite{marx1846} --- The foundational text on ideology as material production.
  \item \textcite{polanyi1944} --- How markets were politically constructed, not naturally emergent.
  \item \textcite{north1990} --- Institutions as the rules of the game shaping economic behavior.
  \item \textcite{acemoglu2012} --- How political institutions determine economic outcomes.
  \item \textcite{piketty2014} --- The dynamics of wealth concentration and its legitimation.
  \item \textcite{stiglitz2012} --- How inequality is maintained through institutional design.
  \item \textcite{streeck2016} --- The contradictions of contemporary capitalism.
  \item \textcite{harvey2005} --- The institutional history of neoliberalism.
\end{itemize}

\section*{Sociology of Knowledge and Symbolic Power}

\begin{itemize}[nosep]
  \item \textcite{bourdieu1984} --- The social production of taste, distinction, and doxa.
  \item \textcite{bourdieu1991} --- Language as an instrument of symbolic power.
  \item \textcite{gramsci1971} --- Hegemony and the production of consent through civil society.
  \item \textcite{foucault1971} --- The rules governing what counts as knowledge.
  \item \textcite{lukes2005} --- Three dimensions of power, including the power to shape desires.
  \item \textcite{scott1998} --- How states impose legibility on complex social realities.
  \item \textcite{eyal2019} --- The sociology of expertise and its social authority.
\end{itemize}

\section*{Propaganda, Media, and Public Opinion}

\begin{itemize}[nosep]
  \item \textcite{lippmann1922} --- The manufacture of public opinion through media.
  \item \textcite{dewey1927} --- The public, communication, and democratic life.
  \item \textcite{bernays1928} --- The engineering of consent as a professional practice.
  \item \textcite{adorno1944} --- The culture industry and mass deception.
  \item \textcite{herman1988} --- The propaganda model of media: five structural filters.
  \item \textcite{mcchesney1999} --- Corporate media and the erosion of democratic communication.
  \item \textcite{bagdikian2004} --- Media ownership concentration and its effects.
  \item \textcite{zuboff2019} --- Surveillance capitalism and the commodification of behavior.
  \item \textcite{sunstein2017} --- Echo chambers and the fragmentation of public discourse.
\end{itemize}

\section*{Behavioral Science and Cooperation}

\begin{itemize}[nosep]
  \item \textcite{kahneman2011} --- Cognitive biases and dual-process theory.
  \item \textcite{fischbacher2001} --- Experimental evidence for conditional cooperation.
  \item \textcite{henrich2001} --- Cross-cultural evidence against homo economicus.
  \item \textcite{bowles2011} --- Cooperation as an evolved human capacity.
  \item \textcite{ostrom1990} --- Self-governance of commons without external coercion.
  \item \textcite{gintis2005} --- Moral sentiments and material interests in economic behavior.
  \item \textcite{nowak2006} --- Evolutionary mechanisms for cooperation.
\end{itemize}

\section*{Formal Methods}

\begin{itemize}[nosep]
  \item \textcite{jackson2008} --- Network theory for social and economic systems.
  \item \textcite{schelling1978} --- Micromotives, macrobehavior, and threshold effects.
  \item \textcite{granovetter1978} --- Threshold models of collective behavior.
  \item \textcite{pearl2009} --- Structural causal models and DAGs.
  \item \textcite{angrist2009} --- Causal inference for empirical research.
  \item \textcite{blei2003} --- Latent Dirichlet Allocation for text analysis.
  \item \textcite{shannon1948} --- Information theory: entropy and communication.
\end{itemize}

\section*{Case Studies and Applied Doxonomics}

\begin{itemize}[nosep]
  \item \textcite{oreskes2010} --- How doubt was manufactured about climate science and tobacco.
  \item \textcite{mayer2016} --- The dark-money network behind American conservatism.
  \item \textcite{mirowski2009} --- The institutional history of neoliberal thought.
  \item \textcite{klein2007} --- Crisis as an opportunity for doxonomic capture.
  \item \textcite{mudge2018} --- How left parties adopted neoliberal frameworks.
  \item \textcite{brown2015} --- Neoliberalism and the undoing of democratic subjectivity.
  \item \textcite{graeber2011} --- The doxonomic history of debt and moral obligation.
\end{itemize}

\section*{Alternatives and Futures}

\begin{itemize}[nosep]
  \item \textcite{wright2010} --- Real utopias: institutional designs for a more just society.
  \item \textcite{benkler2006} --- Networked information economy and peer production.
  \item \textcite{sen1999} --- Development as the expansion of capabilities.
  \item \textcite{freire1968} --- Critical pedagogy and the education of political consciousness.
  \item \textcite{srnicek2015} --- Post-work futures and counter-hegemonic strategy.
  \item \textcite{mason2015} --- Postcapitalism and the information economy.
\end{itemize}


% --- Bibliography ---
\printbibliography[heading=bibintoc,title={References}]

% --- Index ---
\printindex

\end{document}
